\documentclass[11pt]{amsart}
\usepackage[a4paper,vmargin=27mm,hmargin=22.45mm]{geometry}
\usepackage[T1]{fontenc}
\usepackage{lmodern,microtype}
\usepackage{amsmath,amssymb,amsthm,mathtools,mathrsfs}
\usepackage{booktabs,longtable,array,enumitem}
\usepackage{float}
\usepackage{xr-hyper}
\usepackage[hidelinks]{hyperref}
\newtheorem{theorem}{Theorem}[section]
\newtheorem{proposition}[theorem]{Proposition}
\newtheorem{lemma}[theorem]{Lemma}
\newtheorem{corollary}[theorem]{Corollary}
\newtheorem{claim}[theorem]{Claim}
\theoremstyle{definition}
\newtheorem{definition}{Definition}[section]
\newtheorem{example}[theorem]{Example}
\theoremstyle{remark}
\newtheorem{remark}[theorem]{Remark}
\newcommand{\R}{\mathbb R}
\newcommand{\Z}{\mathbb Z}
\newcommand{\N}{\mathbb{Z}_{\ge0}}
\newcommand{\C}{\mathbb C}
\newcommand{\eps}{\varepsilon}
\newcommand{\norm}[1]{\lVert #1\rVert}
\newcommand{\abs}[1]{\lvert #1\rvert}
\newcommand{\ip}[2]{\langle #1,#2\rangle}
\DeclareMathOperator{\Sym}{Sym}
\DeclareMathOperator{\vol}{vol}
\DeclareMathOperator{\rank}{rank}
\DeclareMathOperator{\diag}{diag}
\DeclareMathOperator{\dist}{dist}
\DeclareMathOperator{\supp}{supp}
\DeclareMathOperator{\tr}{tr}
\DeclareMathOperator{\SO}{SO}
\newcommand{\smin}{s_{\min}}
\newcommand{\HS}{\mathrm{HS}}
\newcommand{\op}{\mathrm{op}}
\newcommand{\Id}{I}
\newcommand{\ind}{\mathbf1}
\numberwithin{equation}{section}
\allowdisplaybreaks[1]
\setlength{\emergencystretch}{3em}
\setcounter{tocdepth}{2}
\raggedbottom
\hypersetup{pdftitle={Complete affine jets and sharp integrability in three variables. II. Degrees two through thirteen},pdfauthor={}}
\title[Complete affine jets II]{Complete affine jets and sharp integrability in three variables\\II. Degrees $2\le k\le13$}
\author{}
\date{28 September 2026}
\externaldocument[ext-supp:]{Complete_Affine_Jets_Supplement}


\makeatletter
\newcommand{\equationalias}[1]{\begingroup\protected@edef\@currentlabel{\p@equation\theequation}\protected@edef\@currentHref{equation.\theHequation}\label{#1}\endgroup}
\makeatother

\newcommand{\partIref}[1]{\hyperref[I:#1]{I.\ref*{I:#1}}}
\newcommand{\partIeqref}[1]{\textup{(\partIref{#1})}}
\newcommand{\partIIref}[1]{\hyperref[II:#1]{II.\ref*{II:#1}}}
\newcommand{\partIIeqref}[1]{\textup{(\partIIref{#1})}}

\externaldocument[I:]{Complete_Affine_Jets_I}

\begin{document}
\begin{abstract}
We complete the sharp integrability theorem for polynomial oscillatory integrals in three variables in degrees $2\le k\le13$. The sharp range is $p>\binom{k+3}{4}+3$. We use the analytic reduction, geometric assignment and coefficient coordinates of the companion paper. The remaining arguments consist of cubic adaptive localization, finite regular sampling estimates, and surface estimates in degrees four through nine. All coefficient counts retain the complete nonconstant polynomial coefficient space.
\end{abstract}
\maketitle\enlargethispage{3pt}
\tableofcontents


\section{Introduction}\label{root:introduction}

\subsection{Tarry's problem and the main theorem}\label{intro:main}

For an integer $k\ge2$, let $V_k=\R[x,y,z]_{\le k}/\R$ be the vector space of real polynomials of total degree at most $k$, modulo constants. We equip $V_k$ with Lebesgue measure in the nonconstant monomial coefficients. Its dimension and the coefficient character used below are
\[
N_k=\binom{k+3}{3}-1,\qquad \chi_k=\binom{k+3}{4}.
\]
For the representative with constant coefficient zero, define
\[
 \mathcal I_k(F)=\int_{[0,1]^3}e^{iF(x,y,z)}\,dx\,dy\,dz.
\]
The absolute value of this integral is independent of the representative.

\begin{theorem}\label{root:main}\label{low:4:pf:assembly}\label{low:5:qpf:assembly}\label{low:5:qga:actual-upper}
For every integer $2\le k\le13$ and every $0<p<\infty$,

\noindent\textup{(i)} \[ \mathcal I_k\in L^p(V_k)\quad\Longleftrightarrow\quad p>\chi_k+3. \]

\noindent\textup{(ii)} For every fixed $a\in C_c^\infty(\R^3)$, the integral $F\mapsto\int a(x)e^{iF(x)}\,dx$ also belongs to $L^p(V_k)$ when $p>\chi_k+3$, with constants depending on the fixed degree, the support, and finitely many smooth seminorms of $a$.
\end{theorem}


\subsection{Earlier work and the three-dimensional problem}\label{intro:earlier}

In one variable, Arkhipov, Karatsuba, and Chubarikov proved that the convergence exponent is $k(k+1)/2+1$ for every $k\ge2$, using estimates for polynomial oscillatory integrals in terms of their $H$-function \cite[Theorems 1 and 5]{AKC80}. Basu, Guo, Zhang, and Zorin-Kranich proved that the corresponding convergence exponent in two variables is $\binom{k+2}{3}+2$ for every $k\ge2$ \cite[Theorem 1.2]{BGZZK}. Mockenhaupt resolved the quadratic case in every dimension at least two in 1996 \cite{Moc96}; see also \cite[Section 1.3]{BGZZK}. The ternary quadratic upper bound is recorded in Proposition~\ref{root:quadratic-upper}. Paper I \cite{PaperI} proves the ternary theorem for $k\ge14$ and supplies the shared estimates used below.

The planar stationary-set argument uses determinant factorizations whose role becomes more difficult in higher dimension; see \cite[Remark 5.11]{BGZZK}. The present proof instead organizes coefficient measure through complete affine jets and local geometric configurations. Its coefficient estimates keep every nonconstant coefficient, including those lost if one differentiates the phase and then discards its kernel.

The binary cubic example in Paper I, Section~\partIref{pl:section}, explains this organization. It recovers the known planar threshold $12$ and introduces the coupling between derivatives of different orders.

\subsection{Proof strategy and organization}\label{intro:strategy}

An affine map $u\mapsto c+Bu$ gives the polynomial $F(c+Bu)-F(c)$, whose coefficients form the complete jet. Changing $c$ mixes its homogeneous parts, whereas a linear substitution preserves each degree. The determinant of this coefficient substitution is $|\det B|^{\chi_k}$. We first choose a family of physical affine maps independently of $F$. When integrating the contribution of a label $j$, we use the coefficient coordinates of its fixed map $T_j$.

We bound the original integral by a weighted sum over the selected labels. Labels are assigned either to regular pieces, where polynomial sampling gives coefficient-volume bounds, or to pieces near a surface, where the implicit graph and its transverse derivatives give those bounds. The assignment is chosen once for all labels of a phase. The coefficient set for a fixed label includes every allowed choice of sample points; a determinant bound for one choice must therefore be combined with a bound for the number of choices.

The uniform regular estimate of Paper I is sufficient for $k\ge14$. Its formal exponent has insufficient decay at $k=13$; Section~\ref{reg:combined} gives the margin. Degrees four through thirteen therefore use the finite regular estimates below, while degree three uses adaptive localization. The surface estimate of Paper I applies for $k\ge10$; degrees four through nine use the surface estimates in this paper. Every coefficient count retains the complete coefficient space and the original fixed frame.

For degrees four through thirteen, Proposition~\ref{an:unified-incidence} combines the coefficient estimates on an interval immediately above $p=\chi_k+3$. In degree four it retains a middle-side loss, whose summability is checked in Section~\ref{completion:physical-II}. The smallest physical side gives the limiting condition $p>\chi_k+3$. The pointwise bound in Paper I extends the smooth upper bounds to larger $p$, and its smooth-to-cube reduction gives the cube upper bound. Paper I also proves the matching lower bound and divergence at the endpoint.

Section~\ref{root:quadratic} records the quadratic result, and Section~\ref{low:degree3} gives the cubic adaptive argument. Section~\ref{reg:finite-section} proves the finite regular estimates. Section~\ref{surf:finite-section} treats surfaces in degrees seven through nine. Sections~\ref{low:degree45} and~\ref{low:degree6} treat quartic and quintic surfaces, and sextic surfaces, respectively. Section~\ref{an:assembly} combines the coefficient estimates and completes the proof. The common supplement contains the finite certificate data.

References prefixed I refer to Paper I \cite{PaperI}. We use its complete jets and fixed physical families (Section~\partIref{an:section}), augmented stopping condition and geometric assignments (Section~\partIref{geo:section}), regular reference frames and sampling nets (Section~\partIref{reg:geometry-section}), and coefficient coordinates (Section~\partIref{atlas:section}), with the hypotheses and degree ranges stated there. All coefficient sets and assignments below use the conventions of Definition~\partIref{def:assignment}.

\subsection*{Notation}

For nonnegative quantities, $A\lesssim B$ means $A\le CB$ for some constant $C>0$. Subscripts indicate parameter dependence: $A\lesssim_\theta B$ means $A\le C_\theta B$. We write $A\sim B$ when $A\lesssim B$ and $B\lesssim A$, with the same convention for subscripts. The symbols $\ll$ and $\gg$ denote sufficiently small and large ratios, respectively.




\section{The quadratic case}\label{root:quadratic}

We use Mockenhaupt's quadratic theorem \cite{Moc96} with the following coefficient and amplitude conventions.

For the analytic estimates in Sections~\ref{root:quadratic}, \ref{low:degree3} and~\ref{an:assembly}, the degree is fixed before scale parameters are chosen. Constants may depend on the fixed degree, prescribed finite derivative orders, fixed amplitude support and indicated seminorms, and displayed exponent parameters. They are independent of the phase coefficients, frequency scales, physical side ratios and long geometric scales unless a dependence is stated.

\begin{proposition}\label{root:quadratic-upper}

\noindent\textup{(i)} For every $a\in C_c^\infty(\R^3)$ and every $p>8$, the function $F\mapsto\int_{\R^3}a(x)e^{iF(x)}\,dx$ belongs to $L^p(V_2)$. On a fixed compact support, its norm is bounded by a constant times finitely many derivatives of $a$.

\noindent\textup{(ii)} The same conclusion holds for the amplitude $\mathbf1_{[0,1]^3}$.
\end{proposition}

\begin{proof}
Put $d=3$, $Q=(x^2,xy,xz,y^2,yz,z^2)$, and $E^Q_{[0,1]^3}g(\xi,\zeta)=\int_{[0,1]^3}g(x)e^{i(\xi\cdot x+\zeta\cdot Q(x))}\,dx$ for $(\xi,\zeta)\in\R^3\times\R^6$. Mockenhaupt's theorem \cite{Moc96}, in the bounded-amplitude form recorded in \cite[Section 1.3]{BGZZK}, gives
\[
\|E^Q_{[0,1]^3}g\|_{L^p(\mathbb R^{3+6})}\lesssim_p \|g\|_{L^\infty([0,1]^3)},\qquad p>2d+2=8.
\]
The three linear frequencies and six quadratic frequencies are precisely the independent coordinates of $V_2$. Changing from the cited Fourier convention to $e^{iF}$ is a fixed invertible scaling of these coordinates. The cube case follows by taking $g=1$. For a smooth amplitude on any fixed compact support, a fixed translation and dilation place that support in $[0,1]^3$. Pullback is an invertible linear change on complete nonconstant coefficients with constant Jacobian; the constant phase has modulus one. Applying the estimate to the transformed amplitude proves the assertion, with its $L^\infty$ norm bounded by its zeroth smooth seminorm on the fixed support.
\end{proof}


\section{Cubic adaptive localization}\label{low:degree3}

\subsection{Adaptive normalization and its hypotheses}\label{cubic:normalization}

Work in $V_3$, with its coefficient measure; in this section $P$ denotes the cubic phase. Write
\[
\mathcal E a(P)=\int_{\R^3}e^{iP(x)}a(x)\,dx,\qquad a\in C_c^\infty(\R^3).
\]
A representative of $P$ with zero constant coefficient is used when taking its coefficient norm. Set $\mathcal A_R=\{P\in V_3:R\le\|P\|\le2R\}$ for $R\ge2$. Fix bounded physical neighborhoods $U\Subset X$ with $\supp a\Subset U$. All centers and all physical carriers below lie in a fixed enlargement of $X$.

For $B\in\operatorname{GL}^+(3)$ and $c\in X$, write
\begin{equation}\label{low:3:eq:jet}
P(c+Bz)-P(c)=\ell[z]+\frac12Q[z,z]+\frac16C[z,z,z]=:F(z),\qquad v=\det B.
\end{equation}
Here $Q$ and $C$ are symmetric derivative tensors, not ordinary monomial coefficient arrays. Here and below, all tensor norms are full Hilbert--Schmidt norms unless otherwise specified. Define
\begin{equation}\label{low:3:eq:moment}
(G_C)_{ij}=\sum_{k,l=1}^3 C_{ikl}C_{jkl},\qquad \mathscr M(Q,C)=2Q^2+3G_C,\qquad \kappa(C)=\min_{\abs h=1}\norm{\iota_hC}_{\mathrm{HS}}.
\end{equation}
The cubic is essential precisely when $\kappa(C)>0$. For a matrix $N$, the notation $L_NC$ denotes the sum of insertion of $N$ in the three tensor slots. The complete affine derivative in translation and linear directions is
\begin{equation}\label{low:3:eq:affine-derivative}
\mathcal D_F(h,N)=\bigl(Qh+N^T\ell,\ \iota_hC+N^TQ+QN,\ L_NC\bigr).
\end{equation}
In particular, $\mathcal D_F(h,N)=0$ is equivalent to $(h+Nx)\cdot\nabla F$ being constant. The degree-one equation is retained throughout the rank-degenerate argument.

The coefficient map $P\mapsto F$ in \eqref{low:3:eq:jet} has determinant $v^{15}$ by Lemma~\partIref{an:character}, with degreewise exponents $1,4,10$. Consequently,
\begin{equation}\label{low:3:eq:one-window}
\bigl|\{P:\abs\ell\leq H,\ \norm Q_{\mathrm{HS}}+\norm C_{\mathrm{HS}}\leq C_*\}\bigr|\lesssim_{C_*}H^3v^{-15},\qquad H\geq1,
\end{equation}
for every fixed frame and fixed $C_*$.

\paragraph{Cubic stabilizers and the small-product condition}

An essential cubic may have a nilpotent linear stabilizer. For example, with $(x,y,z)=(x_1,x_2,x_3)$,
\[
P_*(x,y,z)=xz^2-\frac12y^2z,\qquad (y\partial_x+z\partial_y)P_*=0.
\]
The ordered moment condition in the selected frame excludes this kernel below.

A second example concerns small cubic tensors and a large linear coefficient. For $q,t,c\ne0$, let
\begin{equation}\label{low:3:eq:small-counterexample}
F_{q,t,c}=qxy+\frac t2z^2+\frac c2xz^2+\frac{tq}{c}y.
\end{equation}
Then
\begin{equation}\label{low:3:eq:small-counterexample-field}
\left(-\frac q c\partial_z+z\partial_y\right)F_{q,t,c}=0,
\end{equation}
and its tensors satisfy
\begin{equation}\label{low:3:eq:small-counterexample-moment}
Q=\begin{pmatrix}0&q&0\\q&0&0\\0&0&t\end{pmatrix},\quad G_C=\diag(c^2,0,2c^2),\quad \mathscr M=\diag(2q^2+3c^2,2q^2,2t^2+6c^2).
\end{equation}
The moment can approach the identity while the cubic tends to zero and the linear coefficient tends to infinity. For prescribed $m_1>m_2\geq m_3>2(m_1-m_2)$, setting
\[
q^2=\frac{m_2}{2},\qquad c^2=\frac{m_1-m_2}{3},\qquad t^2=\frac{m_3-2(m_1-m_2)}2
\]
gives $\mathscr M=\diag(m_1,m_2,m_3)$. This example falls outside the small-product condition $\norm C_{\mathrm{HS}}\abs\ell\ll1$ used below. Localization on the larger ellipsoids makes the part with large linear term nonstationary before it is counted.

The phase retains every cubic term and mixed quadratic coefficient. The two kernel calculations below establish nonvanishing on compact parameter sets for the later count of selected frames in the partition of unity in \eqref{low:3:eq:positive-resolution}.

\paragraph{The cubic conclusion}

\begin{theorem}\label{low:3:thm:main}\label{low:cubic}

\noindent\textup{(i)} For $a\in C_c^\infty(\R^3)$ and $p>18$, one has $\mathcal E a\in L^p(V_{3})$.

\noindent\textup{(ii)} More precisely, for each $18<p<19$ there is $\epsilon_p>0$ such that
\begin{equation}\label{low:3:eq:annular-main}
\int_{\{R\leq\norm P\leq2R\}}\abs{\mathcal E a(P)}^p\,dP\lesssim_{p,a} R^{-\epsilon_p},\qquad R\geq2.
\end{equation}

\noindent\textup{(iii)} The same global integrability conclusion holds for the amplitude $\mathbf 1_{[0,1]^3}$.
\end{theorem}

\par\smallskip
\noindent\textit{Proof stages for Theorem~\ref{low:3:thm:main}.}\par\nobreak
\noindent\textup{Step 1.} Minimizing frames and adaptive local comparison.\par\nobreak
\noindent\textup{Step 2.} Kernel exclusions and quantitative differential bounds.\par\nobreak
\noindent\textup{Step 3.} Frame counts and the adaptive partition.\par\nobreak
\noindent\textup{Step 4.} Two classes, tails and annular summation.\par\nobreak
\smallskip

Fix $R\geq2$ and $K\geq K_0$, and put
\begin{equation}\label{low:3:eq:parameters}
L_R=\log(C_XR),\qquad \lambda=\frac{K^{-1}}{C_0L_R^2},\qquad \eta=K^{-1/3},\qquad L=K^{1/4}.
\end{equation}
The fixed constants $C_X,C_0,K_0$ are chosen large enough for the estimates below. The relation $K=R^\beta$ is imposed only in the final summation.

At a center $c$, set $Q_0=D^2P(c)$ and $C^{\mathrm{raw}}=D^3P$. On positive-determinant frames consider
\begin{equation}\label{low:3:eq:functional}
\mathcal F_{\lambda,c}(B)=\norm B_{\mathrm{HS}}^2+\norm{Q_0[B,B]}_{\mathrm{HS}}^2+\norm{C^{\mathrm{raw}}[B,B,B]}_{\mathrm{HS}}^2-2\log\det B+\lambda\norm{\log(BB^T)}_{\mathrm{HS}}^2.
\end{equation}

\begin{lemma}[Existence, identity, and bounded jets]\label{low:3:lem:regularizer}

\noindent\textup{(i)} For $\norm P\leq2R$, \eqref{low:3:eq:functional} has a unique minimizing ellipsoid, represented by $B(P,c)$ up to right orthogonal multiplication. It depends smoothly on $c$ and on the coefficients.

\noindent\textup{(ii)} Its normalized tensors satisfy
\begin{equation}\label{low:3:eq:regularized-identity}
B^TB+\mathscr M(Q,C)+2\lambda\log(B^TB)=I,
\end{equation}

\noindent\textup{(iii)} \begin{equation}\label{low:3:eq:frame-bounds} c_XR^{-1}\leq s_{\min} B\leq\norm B_{\mathrm{op}}\leq1,\qquad \norm Q_{\mathrm{HS}}+\norm C_{\mathrm{HS}}\leq C_*, \end{equation} and
\begin{equation}\label{low:3:eq:balance-error}
\norm{\mathscr M(Q,C)+B^TB-I}_{\mathrm{op}}\lesssim \lambda L_R.
\end{equation}

\noindent\textup{(iv)} In an ordered right singular basis, $B^TB=E^2=\diag(d_1^2,d_2^2,d_3^2)$ with $d_1\leq d_2\leq d_3$, and
\begin{equation}\label{low:3:eq:ordered-moment}
\mathscr M=\diag(m_1,m_2,m_3),\quad m_i=1-d_i^2-2\lambda\log d_i^2,\quad \mathscr M+E^2\succeq I,
\end{equation}
\begin{equation}\label{low:3:eq:gap}
m_i-m_j\geq2\lambda(1-r_{ij}^2),\qquad r_{ij}=d_i/d_j,\quad i<j.
\end{equation}
\end{lemma}

\begin{proof}
The coercivity, convexity, and smooth-dependence argument of Lemma~\partIref{an:minimizer} applies to \eqref{low:3:eq:functional}: it uses $\lambda>0$ and $\lambda L_R\lesssim 1$, both satisfied here. Its first variation gives \eqref{low:3:eq:regularized-identity}, and its singular-vector test gives \eqref{low:3:eq:frame-bounds}. In particular $\|\log(B^TB)\|\lesssim L_R$, which proves \eqref{low:3:eq:balance-error}. The matrix-logarithm Hessian used in that argument is Lemma~\partIref{an:log-hessian}, with formula \partIeqref{an:log-hessian-equation}.

In the ordered singular basis, \eqref{low:3:eq:regularized-identity} gives the displayed $m_i$. Since $d_i\le1$, the logarithmic term is negative semidefinite, so $\mathscr M+E^2\succeq I$. Finally,
\[
m_i-m_j=d_j^2-d_i^2-2\lambda\log r_{ij}^2\ge2\lambda(1-r_{ij}^2),
\]
which proves \eqref{low:3:eq:gap}.
\end{proof}

The next estimate retains the normalized cubic norm in the local comparison argument.

\begin{lemma}[Cubic-dependent local comparison]\label{low:3:lem:larger-radius}
There are dimensional $c_{\mathrm{cmp}},C>0$ with the following property.

\noindent\textup{(i)} Let $B_0=B(P,c)$ and let $C_{c}$ be its normalized cubic tensor. If
\begin{equation}\label{low:3:eq:larger-radius}
\abs{B_0^{-1}(c'-c)}\leq\frac{c_{\mathrm{cmp}}\sqrt\lambda}{\eta+\norm{C_c}_{\mathrm{HS}}},
\end{equation}
then the minimizing frames at $c$ and $c'$ are relatively comparable by a factor two. More precisely, after a right orthogonal adjustment,
\begin{equation}\label{low:3:eq:larger-radius-close}
\norm{B_0^{-1}B(P,c')-I}\leq C\frac{\norm{C_c}_{\mathrm{HS}}}{\sqrt\lambda}\abs{B_0^{-1}(c'-c)}.
\end{equation}

\noindent\textup{(ii)} Define
\begin{equation}\label{low:3:eq:adaptive-radius}
\rho(P,c)=\frac{c_*\sqrt\lambda}{\eta+\norm{C_c}_{\mathrm{HS}}}.
\end{equation}
For sufficiently small fixed $c_*$, the ellipsoids $c+\rho(P,c)B(P,c)\mathbb B$ have uniform local comparison, including fixed dilations used in their covering construction. On intersecting such dilates, the products $\rho(P,c)B(P,c)$ are relatively comparable by a fixed constant. Also
\begin{equation}\label{low:3:eq:rho-bounds}
c_{\mathrm{cmp}}\sqrt\lambda\leq\rho(P,c)\leq C\sqrt\lambda/\eta<1.
\end{equation}
\end{lemma}

\begin{proof}
At a minimizing frame let $\mathcal H$ be the Hessian in symmetric frame directions and write $L_HQ=HQ+QH$. The squared quadratic norm and \partIeqref{an:log-hessian-equation} give
\begin{equation}\label{low:3:eq:mixed-hessian}
\mathcal H(H,H)\geq4\norm{L_HQ}_{\mathrm{HS}}^2+8\lambda\norm H_{\mathrm{HS}}^2.
\end{equation}
A physical center variation $Bh$ changes the higher jets at fixed $B$ only by $\dot Q=\iota_hC$. Its mixed derivative with the frame variation is exactly $4\ip{\iota_hC}{L_HQ}$. Choose smooth representing frames satisfying $\dot B=BS$ with $S$ symmetric, as in the proof of Lemma~\partIref{an:comparison}. Differentiating criticality and testing by $S$ gives
\[
\mathcal H(S,S)=-4\ip{\iota_hC}{L_SQ},\qquad \norm S\leq C\lambda^{-1/2}\norm C_{\mathrm{HS}}\abs h.
\]
The quadratic term in \eqref{low:3:eq:mixed-hessian} is used in Cauchy--Schwarz; a crude inverse-Hessian estimate would lose another factor $\lambda^{-1/2}$.

Follow $c(t)=c+tB_0u$ and set $A(t)=B_0^{-1}B(P,c(t))$ with this choice of representing frame. Put $q(t)=\max(\norm{A(t)},\norm{A(t)^{-1}})$. Then $\abs{h(t)}\leq q(t)\abs u$ and $\norm{C_{c(t)}}\leq\norm{C_c}q(t)^3$, because the raw cubic is independent of the center. Therefore the upper Dini derivative obeys
\[
\frac{d}{dt}\log q(t)\leq C\lambda^{-1/2}\norm{C_c}\abs u\,q(t)^4,\qquad q(0)=1.
\]
Integration gives $q(t)\leq(1-4C\lambda^{-1/2}\norm{C_c}\abs u\,t)^{-1/4}$ while the denominator is positive. Under \eqref{low:3:eq:larger-radius} it is at most two. Integrating $\dot A=AS$ proves \eqref{low:3:eq:larger-radius-close}. If $C_c=0$, the raw cubic vanishes and the minimizing ellipsoid is independent of $c$, which also proves the assertion.

Comparable frames give comparable normalized cubic norms, in both directions, by tensor pullback. Thus $\eta+\norm{C_{c'}}$ and $\eta+\norm{C_c}$ are comparable under \eqref{low:3:eq:larger-radius}. For an intersection point of two fixed dilates, compare each center to that point; choosing $c_*$ smaller than the fixed dilation constants permits both applications. This proves the assertion about products of frames and radii. The bounds \eqref{low:3:eq:rho-bounds} follow from \eqref{low:3:eq:frame-bounds} and \eqref{low:3:eq:parameters}.
\end{proof}

\subsection{Affine kernels and quantitative reference frames}\label{cubic:frames}\label{low:3:sec:quantitative}

Both cases use the following tensor calculation, including all lower-degree equations.

\begin{lemma}[A rank-two triangular linear part]\label{low:3:lem:rank-two}
Let
\[
N=\begin{pmatrix}0&\alpha&\beta\\0&0&\gamma\\0&0&0\end{pmatrix},\qquad \alpha\gamma\ne0,
\]
and suppose $L_NC=0$ and $\iota_hC+N^TQ+QN=0$.

\noindent\textup{(i)} Then $Q_{11}=Q_{12}=0$.

\noindent\textup{(ii)} \begin{equation}\label{low:3:eq:rank-two-gap} \mathscr M(Q,C)_{33}-\mathscr M(Q,C)_{11}\geq\norm C_{\mathrm{HS}}^2. \end{equation}
\end{lemma}

\begin{proof}
Write $f_3(x)=C[x,x,x]/6$. The homogeneous invariant cubics are precisely
\begin{equation}\label{low:3:eq:invariant-cubic}
f_3=s z^3+t\left(xz^2-\frac{\alpha}{2\gamma}y^2z-\frac{\beta}{\gamma}yz^2\right).
\end{equation}
Expand $f_3=a_3x^3+x^2f_1(y,z)+xf_2(y,z)+f_0(y,z)$ in $(\alpha y+\beta z)\partial_x+\gamma z\partial_y$. The $x^2$ coefficient first gives $a_3=0$ and $\partial_yf_1=0$. Writing $f_1=a_1z$, the $x$ coefficient gives $f_2=-(\alpha a_1/\gamma)y^2-(2\beta a_1/\gamma)yz+tz^2$. The coefficient independent of $x$, evaluated at $z=0$, then forces $a_1=0$. Integrating its remaining $y$ derivative gives \eqref{low:3:eq:invariant-cubic}, with a free coefficient $s$ of $z^3$. This coefficient computation also works when $s$ or $t$ is zero.

Let $f_2(x)=\tfrac12x^TQx$. In $(Nx)\cdot\nabla f_2+h\cdot\nabla f_3=0$, the $xy$ coefficient is $\alpha Q_{11}$, so $Q_{11}=0$. The $xz$ and $y^2$ coefficients are, respectively,
\[
\gamma Q_{12}+2t h_3,\qquad \alpha Q_{12}-\frac{\alpha t}{2\gamma}h_3.
\]
They give $Q_{12}=0$ and $t h_3=0$. Thus
\[
(Q^2)_{33}-(Q^2)_{11}=Q_{23}^2+Q_{33}^2.
\]
Direct differentiation of \eqref{low:3:eq:invariant-cubic} gives
\begin{align*}
(G_C)_{33}-(G_C)_{11}&=36s^2+t^2\left(4+\frac{\alpha^2}{\gamma^2}+\frac{8\beta^2}{\gamma^2}\right),\\
\norm C_{\mathrm{HS}}^2&=36s^2+t^2\left(12+\frac{3\alpha^2}{\gamma^2}+\frac{12\beta^2}{\gamma^2}\right).
\end{align*}
Therefore $3((G_C)_{33}-(G_C)_{11})-\norm C_{\mathrm{HS}}^2=72s^2+12\beta^2t^2/\gamma^2\geq0$, which proves \eqref{low:3:eq:rank-two-gap}.
\end{proof}

We now include shape boundaries. Use parameters
\begin{equation}\label{low:3:eq:compact-shapes}
0\leq d_1\leq d_2\leq d_3\leq1,\quad 0\leq r_{12},r_{23}\leq1,\quad d_1=r_{12}d_2,\quad d_2=r_{23}d_3,\quad r_{13}=r_{12}r_{23}.
\end{equation}
The ratios are retained as separate parameters when some lengths are zero. Let $E_{ij}$ be the standard matrix unit. For $u=(u_{12},u_{13},u_{23})$ set
\begin{equation}\label{low:3:eq:shape-fields}
N_r(u)=\sum_{i<j}u_{ij}(E_{ij}-r_{ij}^2E_{ji}),\qquad \Omega_r(u)=\sum_{i<j}r_{ij}u_{ij}(E_{ij}-E_{ji}),\qquad E=\diag(d_i).
\end{equation}
For positive lengths, $N_r=E^{-1}\Omega_rE$. In these coordinates define the augmented derivative
\begin{equation}\label{low:3:eq:J}
\mathcal J_{E,r,F}(h,u)=\bigl(\mathcal D_F(h,N_r(u)),\ Eh,\ \Omega_r(u)\bigr).
\end{equation}
All entries are polynomial in the compact shape parameters and in the jets.

\begin{lemma}[Non-small cubic tensor]\label{low:3:lem:non-small}
Suppose $\norm{\mathscr M(Q,C)+E^2-I}_{\mathrm{op}}\leq\epsilon<1$ and $\norm C_{\mathrm{HS}}^2>2\epsilon$. Then \eqref{low:3:eq:J} is injective, without any assumption on the contraction rank of $C$.
\end{lemma}

\begin{proof}
Let $(h,u)$ belong to the kernel. The equation $\Omega_r(u)=0$ makes $N_r(u)=N$ strictly upper triangular and gives $EN=0$. Also $Eh=0$. If $N=0$, then $Qh=0$ and $\iota_hC=0$, whence $h^T\mathscr M h=0$. The balance condition and $Eh=0$ instead give $h^T\mathscr M h\geq(1-\epsilon)\abs h^2$. Thus $h=0$.

If $\rank N=2$, then $d_1=0$ and Lemma~\ref{low:3:lem:rank-two} gives $\mathscr M_{33}-\mathscr M_{11}\geq\norm C_{\mathrm{HS}}^2$. Balance gives $\mathscr M_{33}-\mathscr M_{11}\leq-d_3^2+2\epsilon\leq2\epsilon$, a contradiction.

It remains to treat $\rank N=1$. Write $N=s\,a\otimes\nu$ with $\abs a=\abs\nu=1$ and $s\ne0$, identifying the covector $\nu$ with its Euclidean vector. Nilpotence gives $a\perp\nu$, and $EN=0$ gives $Ea=0$. The cubic equation implies $\iota_aC=0$. Put $q=Qa$. Evaluating the quadratic equation twice on $a$ and then once on $a$ gives $a^TQa=0$: more directly, contraction once with $a$ yields $s\nu\,(a^TQa)=0$, since $\iota_aC=0$ and $a\cdot\nu=0$. Hence $q\perp a$. The linear equation gives $Qh=-s\nu(a\cdot\ell)$, so $q\cdot h=a\cdot Qh=0$.

If $q=0$, then $a^T\mathscr M a=0$, contradicting balance and $Ea=0$. Write $h=h_a a+k$ with $k\perp a$. If $k=0$, the quadratic equation is $\nu\otimes q+q\otimes\nu=0$, again forcing $q=0$. Thus $k\ne0$. Put $e=k/\abs k$ and choose a unit $w\perp a,e$. Since $q\perp a,k$, one has $q=q_w w$. In the orthonormal basis $(a,e,w)$,
\[
Q=\begin{pmatrix}0&0&q_w\\0&q_e&q_{ew}\\q_w&q_{ew}&q_{ww}\end{pmatrix}.
\]
The tensor $C$ has no $a$ slot. Evaluating $\iota_hC=-s(\nu\otimes q+q\otimes\nu)$ on $(e,e)$ gives $C_{eee}=0$. Its only possibly nonzero components, up to symmetry, are $A=C_{eew}$, $B=C_{eww}$ and $S=C_{www}$. Therefore
\begin{align*}
\norm C_{\mathrm{HS}}^2&=3A^2+3B^2+S^2,\\
w^T\mathscr M w-a^T\mathscr M a&=2(q_{ew}^2+q_{ww}^2)+3(A^2+2B^2+S^2)\geq\norm C_{\mathrm{HS}}^2.
\end{align*}
Since $Ea=0$, balance bounds the same difference by $-\abs{Ew}^2+2\epsilon\leq2\epsilon$. This is the required contradiction.
\end{proof}

\begin{lemma}[Small cubic tensor and small linear product]\label{low:3:lem:small-product}
There is a dimensional $c_0>0$ such that the following holds. Suppose
\begin{equation}\label{low:3:eq:small-conditions}
\norm C_{\mathrm{HS}}\leq c_0,\qquad \norm C_{\mathrm{HS}}\abs\ell\leq c_0,\qquad \mathscr M+E^2\succeq I,
\end{equation}
and $\mathscr M$ is diagonal with
\begin{equation}\label{low:3:eq:small-gap}
\mathscr M_{ii}-\mathscr M_{jj}\geq g(1-r_{ij}^2),\qquad i<j,\qquad g>0.
\end{equation}
Then \eqref{low:3:eq:J} is injective on the full compactified shape set \eqref{low:3:eq:compact-shapes}.
\end{lemma}

\begin{proof}
For a kernel vector, $N=N_r(u)$ is strictly upper triangular, $EN=0$ and $Eh=0$, as above. The case $N=0$ follows from stopping, exactly as in Lemma~\ref{low:3:lem:non-small}. If $\rank N=2$, its $12$ and $23$ entries are nonzero. Thus $r_{12}=r_{23}=r_{13}=0$. Lemma~\ref{low:3:lem:rank-two} gives $\mathscr M_{33}\geq\mathscr M_{11}$, whereas \eqref{low:3:eq:small-gap} gives $\mathscr M_{11}-\mathscr M_{33}\geq g$. This is impossible, including when $C=0$.

For $\rank N=1$, use $N=s a\otimes\nu$, $Ea=0$ and $\iota_aC=0$. Stopping gives $\abs{Qa}\geq1/\sqrt2$. Also $Eh=0$ implies
\[
2\abs{Qh}^2+3\norm{\iota_hC}_{\mathrm{HS}}^2\geq\abs h^2,
\]
so $\abs{Qh}\gtrsim\abs h$ if $c_0$ is small. The linear equation yields $\abs h\lesssim \abs s\abs\ell$. The quadratic equation gives
\[
\abs s\norm{\nu\otimes Qa+Qa\otimes\nu}_{\mathrm{HS}}=\norm{\iota_hC}_{\mathrm{HS}}\leq\norm C_{\mathrm{HS}}\abs h.
\]
For any vectors $v,w$, $\norm{v\otimes w+w\otimes v}_{\mathrm{HS}}^2=2\abs v^2\abs w^2+2(v\cdot w)^2$. Hence the left side is $\gtrsim\abs s$. Combining the inequalities gives $1\lesssim \norm C_{\mathrm{HS}}\abs\ell$, contrary to \eqref{low:3:eq:small-conditions} for sufficiently small $c_0$.
\end{proof}


The compact semialgebraic {\L}ojasiewicz inequality \cite[Theorem 2.2]{low:lojasiewicz} gives $|g|^N\lesssim |f|$ for continuous semialgebraic $f,g$ on a compact semialgebraic set when $\{f=0\}\subseteq\{g=0\}$. Only a fixed finite exponent is needed. The following applications check this inclusion by explicit kernel exclusion.

\begin{proposition}[Polynomial dependence on the thresholds]\label{low:3:prop:quantitative}
All jets in this proposition satisfy $\norm Q_{\mathrm{HS}}+\norm C_{\mathrm{HS}}\leq C_*$. There are $c>0$ and finite positive exponents $A_2,A_3$ such that:

\smallskip\noindent
(a) If $H\geq1$, $0<\eta_0\leq1$, $\abs\ell\leq H$, $\norm C_{\mathrm{HS}}\geq\eta_0$ and
\[
\norm{\mathscr M+E^2-I}_{\mathrm{op}}\leq\eta_0^2/8,
\]
then $s_{\min}\mathcal J\geq c(\eta_0/H)^{A_2}$.

\smallskip\noindent
(b) If $H\geq1$, $0<g\leq1$, $\abs\ell\leq H$, $\norm C_{\mathrm{HS}}\leq c_0/H$, $\mathscr M$ is diagonal, $\mathscr M+E^2\succeq I$, and \eqref{low:3:eq:small-gap} holds, then $s_{\min}\mathcal J\geq c(g/H)^{A_3}$.
\end{proposition}

\begin{proof}
For (a), put $t=H^{-1}$ and $e=t\ell$. Include $t,\eta_0\in[0,1]$, $\abs e\leq1$, bounded $Q,C$, the compact shape parameters \eqref{low:3:eq:compact-shapes}, and the displayed polynomial balance inequality in one compact semialgebraic set. The entries of $t\mathcal J$ are polynomial in these variables. On this compact set, $\{s_{\min}(t\mathcal J)=0\}\subseteq\{t\eta_0=0\}$ by Lemma~\ref{low:3:lem:non-small}. Apply the inequality with $g_0=t\eta_0$, then divide by $t$ for $t>0$ and weaken the result to the stated bound.

For (b), use the same compactification with $t=H^{-1}$, $e=t\ell$, $g\in[0,1]$, and $\norm C\leq c_0t$. The diagonal, stopping and gap conditions are polynomial inequalities. Lemma~\ref{low:3:lem:small-product} implies that $\{s_{\min}(t\mathcal J)=0\}\subseteq\{gt=0\}$ on this compact set. Applying the compact inequality proves (b). In particular, no logarithmic defining condition is used in this application. The polynomial gap condition is a relaxation of the exact logarithmic moment identity.
\end{proof}

\begin{lemma}[Pointwise comparison with an exact frame]\label{low:3:lem:reference-transfer}
Let $B=OE$ and $\widehat B=\widehat O\widehat E$ be ordered singular-frame representatives of two ellipsoids at the same center. Suppose $G=B^{-1}\widehat B$ satisfies $\norm G,\norm{G^{-1}}\leq A_*$, where $A_*$ is fixed.

\noindent\textup{(i)} If the augmented physical derivative at $\widehat B$ has lower bound $\sigma$ in its normalized translation--rotation coordinates, then the derivative at $B$ has lower bound $\gtrsim_{A_*}\sigma$.

\noindent\textup{(ii)} The same assertion holds for the map recording only the quadratic and cubic jets and the physical orientation. All bounds are independent of the singular values and their ratios.
\end{lemma}

\begin{proof}
Use the same physical center variation $v$ and physical skew matrix $W$ at both frames. Their normalized tangent parameters are
\[
h=B^{-1}v,\quad M=B^{-1}WB,\qquad \widehat h=G^{-1}h,\quad \widehat M=G^{-1}MG.
\]
Thus $\abs h+\norm M$ and $\abs{\widehat h}+\norm{\widehat M}$ are comparable with constants depending only on $A_*$. For an ordered diagonal frame, $M=N_r(u)$ and $\abs u\leq\norm M_{\mathrm{HS}}\leq\sqrt2\abs u$. The same comparison holds at $\widehat B$.

Along these physical variations, $\widehat B(t)=B(t)G$ with $G$ fixed, and hence $\widehat F(t)=F(t)\circ G$ modulo constants. The complete coefficient derivatives are related by the bounded invertible linear pullback $T_GH=[H\circ G]$. For the higher-jet version this pullback is block diagonal on degrees two and three. The physical center output is $v$ for both frames, while the norm of the physical orientation derivative is $\norm W$ for both. These observations prove the lower-bound transfer.
\end{proof}

\begin{corollary}[Differential bounds at frames comparable to the minimizing frame]\label{low:3:cor:eligible-differential}
Consider frames from the fixed family relatively comparable by a fixed factor to the regularized minimizing frame at their own centers, and with bounded higher jets. For all sufficiently large $K$, uniformly over their physical side lengths:

\smallskip\noindent
(i) If $\norm C_{\mathrm{HS}}\geq\eta=K^{-1/3}$ and $\abs\ell\leq K$, the complete augmented derivative has lower bound $\gtrsim K^{-A}$ for a fixed finite $A$.

\smallskip\noindent
(ii) If $\norm C_{\mathrm{HS}}\leq\eta$ and $\abs\ell\leq L=K^{1/4}$, the complete augmented derivative has lower bound $\gtrsim(\lambda/L)^A$ for a fixed finite $A$.
\end{corollary}

\begin{proof}
For (i), the exact reference has cubic norm at least a fixed multiple of $\eta$ and linear norm at most a fixed multiple of $K$. By \eqref{low:3:eq:balance-error}, its balance defect is at most $C K^{-1}/L_R=o(\eta^2)$. Proposition~\ref{low:3:prop:quantitative}(a), with the fixed comparison factors included in the thresholds, applies. For (ii), the exact reference has $\norm{\widehat C}\leq C\eta$ and $\abs{\widehat\ell}\leq CL$. Since $\eta L=K^{-1/12}\to0$, it satisfies the small-product conditions with $H=CL$. Equations \eqref{low:3:eq:ordered-moment} and \eqref{low:3:eq:gap} give the stopping and gap hypotheses with $g=2\lambda$, or with $g=\lambda$ after weakening. Apply Proposition~\ref{low:3:prop:quantitative}(b) and transfer pointwise. No rounding error is amplified by an inverse side ratio.
\end{proof}

\subsection{Frame counts and amplitude decomposition}\label{cubic:assembly}\label{low:3:sec:counting}\label{low:3:sec:resolution}\label{low:3:sec:assembly}

Fix positive side lengths $E=\diag(d_1,d_2,d_3)$ and identify a frame with $(c,O)$, its physical frame being $OE$. The normalized tangent is
\[
\dot c=OEh,\qquad \dot O=O\Omega_r(u).
\]
Write $\Lambda_{c,B}P=P(c+B\,\cdot)-P(c)$ for the complete normalized jet. The complete augmented frame map is
\begin{equation}\label{low:3:eq:state-map}
\Psi_{P,E}(c,O)=\bigl(\Lambda_{c,OE}P,c,O\bigr).
\end{equation}
Its derivative, after orthogonal changes in the physical output coordinates, is \eqref{low:3:eq:J}. The normalized volume measure is
\begin{equation}\label{low:3:eq:haar}
d\mu_E=(\det E)^{-1}\prod_{i<j}r_{ij}^{-1}\,dc\,dO.
\end{equation}
Here $dO$ is Haar measure on $\operatorname{SO}(3)$ with a fixed normalization. The six normalized tangent fields have volume comparable to one for this measure.

\begin{lemma}[Uniform local frame neighborhoods]\label{low:3:lem:state-neighborhoods}

\noindent\textup{(i)} The normalized frame group has small neighborhoods $U_s$, $0<s<s_0$, with $\mu_E(U_s)\asymp s^6$ and uniform product and inverse containment.

\noindent\textup{(ii)} On a neighborhood of radius $s$ about a frame with $\abs\ell\leq H$ and bounded higher jets, the complete jet remains $\lesssim H$, and its higher jets remain bounded.

\noindent\textup{(iii)} On the same neighborhood, the normalized differential of \eqref{low:3:eq:state-map} varies by $\lesssim Hs$. All constants are independent of the side ratios.
\end{lemma}

\begin{proof}
The group is $\R^3\rtimes(E^{-1}\operatorname{SO}(3)E)$, with right translations represented by $h$ and rotations by the three matrices $N_r$. Their entries are uniformly bounded. Put $N_{ij}=E_{ij}-r_{ij}^2E_{ji}$. Writing $a=r_{12}$ and $b=r_{23}$, their brackets satisfy
\[
[N_{12},N_{23}]=N_{13},\quad [N_{12},N_{13}]=-a^2N_{23},\quad [N_{13},N_{23}]=-b^2N_{12}.
\]
Their action on the translation coordinates also has bounded coefficients. Ordered exponential coordinates therefore have uniformly bounded derivatives and an invertible differential at zero. Equivalently, projection of the relative linear matrix onto its three strict upper entries has differential the identity at zero and uniform higher-derivative bounds. This gives uniform local charts, volume comparison, and product and inverse containment, after shrinking $s_0$.

In such a neighborhood, the relative affine map has bounded translation and a linear matrix uniformly close to $I$. The complete polynomial pullback has bounded derivatives of every fixed order. In degree at most three the quadratic and cubic coefficients depend only on the old higher coefficients, while the linear coefficient may also contain the old linear coefficient. The bounds and the differential variation follow directly from the complete affine Taylor formula. The physical center and orthogonal matrix outputs have uniform derivative bounds as well.
\end{proof}

\begin{lemma}[Regular-state count]\label{low:3:lem:count}
Let $\mathcal S$ be a set of frames of one fine shape $E$, with physical centers in a fixed bounded region, separated by at least $a_0\leq1$ in the normalized local frame metric. Suppose their higher jets are bounded, $\abs\ell\leq H$, $H\geq1$, and $s_{\min} d\Psi_{P,E}\geq\sigma$, $0<\sigma\leq1$.

\noindent\textup{(i)} Then
\begin{equation}\label{low:3:eq:regular-count}
\#\mathcal S\lesssim a_0^{-6}H^9\sigma^{-12}.
\end{equation}

\noindent\textup{(ii)} The bound \eqref{low:3:eq:regular-count} also holds if the differential lower bound is assumed for the map recording only $(Q,C,O)$.
\end{lemma}

\begin{proof}
Take disjoint normalized neighborhoods of radius $s=c a_0\sigma/H$. By Lemma~\ref{low:3:lem:state-neighborhoods}, their union has bounded higher jets, $|\ell|\le CH$, differential lower bound $\sigma/2$, and normalized measure $\gtrsim\#\mathcal S\,s^6$.

Cauchy--Binet supplies a six-coordinate projection with normalized Jacobian $\gtrsim\sigma^6$. On finitely many bounded rational charts of $\operatorname{SO}(3)$ the projected frame maps are rational of uniformly bounded degree. Lemma~\ref{ext-supp:pl:rational-volume} therefore gives normalized measure $\lesssim H^3\sigma^{-6}$ for the union: only the three linear coordinates have range $O(H)$, while the physical outputs and higher jets have bounded ranges. Equivalently, apply its change-of-variables proof in each chart, inserting the density of \eqref{low:3:eq:haar} in both the domain measure and the normalized Jacobian. These factors cancel. Only nonsingular solutions of the projected polynomial systems are counted, so singular components of a fiber do not affect the bound. Thus
\[
\#\mathcal S\lesssim H^3\sigma^{-6}s^{-6}\lesssim a_0^{-6}H^9\sigma^{-12}.
\]
Adjoining the linear-jet rows cannot decrease the smallest singular value. Thus a differential lower bound for the map recording $(Q,C,O)$ implies the one used above, with the same $\sigma$.
\end{proof}

We use the following convention for the fixed family. Singular values use a multiplicative $(1+\vartheta)$ mesh. For each fine shape $E$, the orientations are a separated net at normalized scale $c\vartheta$. At one fixed support-radius bin $r$, centers for orientation $O$ lie on a lattice of mesh $c\vartheta r\,OE\mathbb Z^3$. There are $\lesssim\vartheta^{-3}$ fine length triples in one coarse dyadic shape and at most $T$ radius bins. Distinct frames in one fine shape and radius bin are separated at scale $c\vartheta r$; frames with different orientation have the required separation in their linear part, and frames with the same orientation have it in their translation part.

For any such family, let $\psi_j$ be the fixed cutoff functions, let $\chi_j$ be the indicators of the selected frames satisfying the small-linear-term restriction, and set
\begin{equation}\label{low:3:eq:weights}
\omega_j=\frac{C_a}{v_j}\int\abs a\psi_j,\qquad M_D=\sum_{j\in\mathcal J_D}\omega_j\chi_j,\qquad Z_D=\sum_{j\in\mathcal J_D}\omega_jv_j^{-15}.
\end{equation}
We take $C_a$ large enough to dominate the reconstructed pieces. The cutoff functions give $\omega_j\lesssim_a 1$ since their radii are at most one. The weights and the family are fixed before coefficient integration.

\begin{proposition}[Weighted consequence of a differential bound]\label{low:3:prop:weighted}
Restrict the selected frames at a coarse shape to any coefficient-dependent subclass for which the hypotheses of Lemma~\ref{low:3:lem:count} hold, with the same $H,\sigma$ and support radii at least $r_{\min}$. Put
\begin{equation}\label{low:3:eq:count-polynomial}
\mathcal P=C_a T\vartheta^{-9}r_{\min}^{-6}H^9\sigma^{-12}.
\end{equation}
The corresponding subincidence $M_D^{\mathrm{sub}}$ satisfies
\begin{equation}\label{low:3:eq:weighted-moment}
\begin{aligned}
\textup{(i)}\quad &M_D^{\mathrm{sub}}(P)\leq\mathcal P,\\
\textup{(ii)}\quad &\int_{\mathcal A_R}(M_D^{\mathrm{sub}})^p\,dP\lesssim \mathcal P^{p-1}H^3 Z_D,\qquad p\geq1.
\end{aligned}
\end{equation}
\end{proposition}

\begin{proof}
Apply \eqref{low:3:eq:regular-count} in each fine shape and radius bin, use $a_0=c\vartheta r_{\min}$, and sum over $\lesssim T\vartheta^{-3}$ classes. Multiplication by the upper weight bound gives the first assertion. For the second, use $(M_D^{\mathrm{sub}})^p\leq\mathcal P^{p-1}M_D^{\mathrm{sub}}$ and integrate the individual selectors with \eqref{low:3:eq:one-window}. The selectors can be arbitrary further measurable restrictions of the actual low-linear windows; discarding these restrictions only increases the first moment.
\end{proof}

We construct a positive partition on the adaptive carriers \eqref{low:3:eq:adaptive-radius}, before any linear-frequency restriction. Fix the global precision
\begin{equation}\label{low:3:eq:precision}
\vartheta=c_1\lambda K^{-2}.
\end{equation}
This precision is independent of the physical shape.

\begin{proposition}[An adaptive partition of unity]\label{low:3:prop:resolution}
For $\norm P\leq2R$ there is a finite coefficient-independent family $(c_j,B_j,r_j,\psi_j)$ and a selected subcollection $\widetilde{\mathcal I}(P)$ with the following properties.

\noindent\textup{(i)} For each fixed index $j$, the set $\{P:\|P\|\le2R,\ j\in\widetilde{\mathcal I}(P)\}$ is Borel measurable in the cubic phase coefficients. The $r_j$ belong to at most $T\lesssim (1+\log K)$ dyadic bins, with
\[
c\sqrt\lambda\leq r_j\leq C\sqrt\lambda/\eta.
\]

\noindent\textup{(ii)} Every selected $B_j$ is relatively comparable to $B(P,c_j)$, and $r_j\asymp\rho(P,c_j)$. The fixed cutoff functions are supported in $c_j+C r_jB_j\mathbb B$.

\noindent\textup{(iii)} There are smooth pieces $\widetilde a_{P,j}$ satisfying
\begin{equation}\label{low:3:eq:positive-resolution}
a=\sum_{j\in\widetilde{\mathcal I}(P)}\widetilde a_{P,j},\qquad
\sum_j\abs{\widetilde a_{P,j}}=\abs a,\qquad
\abs{\widetilde a_{P,j}}\leq C\abs a\psi_j,
\end{equation}

\noindent\textup{(iv)} \begin{equation}\label{low:3:eq:adaptive-mass} \sum_{j\in\widetilde{\mathcal I}(P)}r_j^3v_j\lesssim 1,\qquad \norm{\widetilde a_{P,j}(c_j+r_jB_j\cdot)}_{C^N}\lesssim_N \norm a_{C^N}. \end{equation} The constants in the last display are independent of $R,K,\vartheta$ and the side ratios, for each fixed $N$. The supports of the selected cutoff functions have bounded overlap.
\end{proposition}

\begin{proof}
Use the fine length and orientation construction of Lemma~\partIref{an:catalogue} at precision $\vartheta$. For each fine diagonal $E$ its normalized angular widths are $c\vartheta d_i/d_j$, giving
\begin{equation}\label{low:3:eq:orientation-count}
\lesssim\vartheta^{-3}\prod_{i<j}\frac{d_j}{d_i}
\end{equation}
orientations. At each radius bin $r$ use centers on $c\vartheta rB_j\mathbb Z^3$ in a fixed enlargement of $X$, and attach fixed bumps with a small plateau and a larger support.

Fix $\delta_{\rm sel}=C_{\rm sel}\vartheta<1/2$ and a fixed radius comparison factor $C_{\rm rad}\ge2$, large enough to include adjacent radius bins. Retain for selection the indices satisfying
\[
\begin{gathered}
(1-\delta_{\rm sel})^2|B(P,c_j)^Tu|^2\le|B_j^Tu|^2\le(1+\delta_{\rm sel})^2|B(P,c_j)^Tu|^2\qquad(u\in\mathbb R^3),\\
C_{\rm rad}^{-1}\rho(P,c_j)\le r_j\le C_{\rm rad}\rho(P,c_j).
\end{gathered}
\]
For fixed $j$ these tests define closed sets in the phase coefficient space: the minimizing quadratic form and the radius~\eqref{low:3:eq:adaptive-radius} depend continuously on $P$. No continuous choice of singular vectors is needed. To cover a point, approximate its exact frame and radius, then choose a nearest lattice center. Its displacement is $O(\vartheta\rho)$ in the parent coordinates. Equation~\eqref{low:3:eq:larger-radius-close} changes the minimizing ellipsoid by $O(\vartheta)$, and tensor pullback changes $\eta+\|C\|$ by the same relative order. Choose $C_{\rm sel}$ to accommodate these comparisons. The radius comparison and adjacent bins then give an index passing both tests near every point of $U$.

Apply the disjoint-core selection from Proposition~\partIref{an:resolution}, in a fixed ordering, to the supporting matrices $r_jB_j$. For $x\in U$, define
\[
S(P,x)=\sum_{j\in\widetilde{\mathcal I}(P)}\psi_j(x),\qquad \widetilde a_{P,j}(x)=\mathbf1_{\{j\in\widetilde{\mathcal I}(P)\}}\,a(x)\psi_j(x)/S(P,x),
\]
and put $\widetilde a_{P,j}(x)=0$ outside $U$. The same finite selection argument makes membership Borel in $P$ and these functions Borel in $(P,x)$. Lemma~\ref{low:3:lem:larger-radius} supplies the required relative comparison of both frames and radii on intersecting fixed dilations. This gives bounded overlap, \eqref{low:3:eq:positive-resolution}, and \eqref{low:3:eq:adaptive-mass}; the derivative bounds use $r_j\|B_j\|\le1$ and are independent of $\vartheta$ and the side ratios. Finally, \eqref{low:3:eq:rho-bounds} gives a radius ratio $O(\eta^{-1})$, hence $T\lesssim (1+\log K)$ bins.
\end{proof}

Let $D=(\delta_1,\delta_2,\delta_3)$ be a coarse dyadic shape and put $v_D=\delta_1\delta_2\delta_3$. Define
\begin{equation}\label{low:3:eq:W}
W_p(D)=\delta_1^{p-18}\delta_2^{p-16}\delta_3^{p-14}.
\end{equation}
Write $\widetilde{\mathcal J}_D$ for all labels of the fixed adaptive family in shape $D$. For this family use tildes on $\widetilde\omega_j$, $\widetilde Z_D$ and all reconstructed quantities, with the same mass-weight formula \eqref{low:3:eq:weights}.

\begin{lemma}[Weighted mass bound]\label{low:3:lem:resource}
For every $p>18$,
\begin{equation}\label{low:3:eq:resource}
\sum_{j\in\widetilde{\mathcal J}_D}\widetilde\omega_jv_j^{p-15}\lesssim_{p,a} T\vartheta^{-9}W_p(D).
\end{equation}
This estimate concerns the full fixed family, before selection.
\end{lemma}

\begin{proof}
The first-moment lattice calculation in the proof of Lemma~\partIref{an:resource}, applied in each radius bin, gives $\sum_{c_j}\widetilde\omega_j\lesssim_a \vartheta^{-3}v^{-1}$ independently of the radius. Multiplying by $v^{p-15}$ and summing the orientations, fine length triples and $T$ radius bins gives
\[
\lesssim_{p,a}T\vartheta^{-9}v_D^{p-16}\prod_{i<j}\frac{\delta_j}{\delta_i}=T\vartheta^{-9}\delta_1^{p-18}\delta_2^{p-16}\delta_3^{p-14}.
\]
This is \eqref{low:3:eq:resource}.
\end{proof}

On the adaptive selected resolution, put $C_j=D^3P[B_j,B_j,B_j]$ and $\ell_j=B_j^T\nabla P(c_j)$, as in \eqref{low:3:eq:jet}. Divide frames at their own normalized cubic tensor norm. Define
\begin{align}\label{low:3:eq:new-selectors}
\widetilde\chi_j^{>}(P)&=\mathbf 1_{\{j\in\widetilde{\mathcal I}(P)\}}\mathbf 1_{\{\norm{C_j}>\eta\}}\mathbf 1_{\{\abs{\ell_j}\leq K\}},\\
\widetilde\chi_j^{\leq}(P)&=\mathbf 1_{\{j\in\widetilde{\mathcal I}(P)\}}\mathbf 1_{\{\norm{C_j}\leq\eta\}}\mathbf 1_{\{\abs{\ell_j}\leq L\}}.\nonumber
\end{align}
Let $\widetilde I_j(P)=\int e^{iP(x)}\widetilde a_{P,j}(x)\,dx$ for selected frames, and set it to zero otherwise. Put
\begin{equation}\label{low:3:eq:new-F}
\widetilde F_D^\bullet=\sum_{j\in\widetilde{\mathcal J}_D}\widetilde I_j\,\widetilde\chi_j^\bullet,\qquad
\widetilde M_D^\bullet=\sum_{j\in\widetilde{\mathcal J}_D}\widetilde\omega_j\widetilde\chi_j^\bullet,\qquad \bullet\in\{>,\leq\}.
\end{equation}
Set $\widetilde M_D=\widetilde M_D^>+\widetilde M_D^{\leq}$.

\begin{theorem}[Uniform weighted bounds for the grouped frames]\label{low:3:thm:two-regimes}
For every $p\geq1$ there are finite $A_p,B_p$ such that, for each $\bullet\in\{>,\leq\}$,

\noindent\textup{(i)} \begin{equation}\label{low:3:eq:two-moment} \int_{\mathcal A_R}(\widetilde M_D^\bullet)^p\,dP\lesssim_{p,a} K^{A_p}L_R^{B_p}\,K^3\widetilde Z_D, \end{equation} uniformly over all physical shapes.

\noindent\textup{(ii)} The same bound, with adjusted constants, holds for $\widetilde M_D$. In particular, there is no loss $\delta_i^{-b_i}$ in either case.
\end{theorem}

\begin{proof}
Use Corollary~\ref{low:3:cor:eligible-differential}(i) for the first class and (ii) for the second. Both differential lower bounds are inverse powers of $K,L_R$ by \eqref{low:3:eq:parameters}. All support radii are $\gtrsim\sqrt\lambda$, and $\vartheta^{-1}$ is a fixed power of $K,L_R$. Proposition~\ref{low:3:prop:weighted}, using $H=K$, which bounds both linear-jet ranges, in the count and in the first moment, proves \eqref{low:3:eq:two-moment}. The count itself still uses the correct differential bound established with the smaller linear range in the second class. The inequality $(u+v)^p\leq2^{p-1}(u^p+v^p)$ proves the combined statement.
\end{proof}

\begin{lemma}[Large linear terms after grouping]\label{low:3:lem:tail}
For every positive integer $N$, the remaining high-linear terms in the adaptive resolution satisfy the pointwise bound
\begin{equation}\label{low:3:eq:tail}
\abs{\mathcal E a(P)-\sum_D(\widetilde F_D^>(P)+\widetilde F_D^{\leq}(P))}\lesssim_{N,a} \left(\frac{K^{1/12}}{L_R}\right)^{-N},\qquad \norm P\leq2R.
\end{equation}
The constant depends on finitely many derivatives of the fixed amplitude, but not on $R,K$ or the side ratios.
\end{lemma}

\begin{proof}
Under $x=c_j+r_jB_jz$, the scaled amplitude satisfies \eqref{low:3:eq:adaptive-mass} and the higher phase jets are bounded. The pointwise estimate in the proof of Lemma~\partIref{an:linear-tail}, with $\rho=r_j$, gives
\begin{equation}\label{low:3:eq:single-tail}
\abs{\widetilde I_j(P)}\lesssim_{N,a} r_j^3v_j(1+r_j\abs{\ell_j})^{-N}.
\end{equation}
In the class $\|C_j\|>\eta$, $r_jK\gtrsim\sqrt\lambda K\asymp K^{1/2}/L_R$. In the other class, bounded frame comparison and \eqref{low:3:eq:adaptive-radius} give $r_j\gtrsim\sqrt\lambda/\eta$, hence
\begin{equation}\label{low:3:eq:cutoff-compatibility}
r_jL\gtrsim\frac{\sqrt\lambda L}{\eta}\asymp\frac{K^{1/12}}{L_R}.
\end{equation}
Sum \eqref{low:3:eq:single-tail} over the omitted parents using the exact resolution and $\sum_jr_j^3v_j\lesssim 1$. This proves \eqref{low:3:eq:tail}.
\end{proof}

\begin{proposition}[Small linear term estimate]\label{low:3:prop:low-assembly}
For $18<p<19$, both classes of parent integrals satisfy
\begin{equation}\label{low:3:eq:coherent-small}
\sum_D\norm{\widetilde F_D^\bullet}_{L^p(\mathcal A_R)}\lesssim_{p,a} K^{A_p}L_R^{B_p}R^{-(p-18)/(3p)},\qquad \bullet\in\{>,\leq\},
\end{equation}
provided $K=R^\beta$ with fixed $0<\beta<2/3$ and $R$ is sufficiently large. The finite exponents $A_p,B_p$ are independent of $\beta$.
\end{proposition}

\begin{proof}
For a selected low-linear piece, \eqref{low:3:eq:positive-resolution} gives $\abs{\widetilde I_j}\leq v_j\widetilde\omega_j$. Hence
\[
\abs{\widetilde F_D^\bullet}\lesssim v_D\widetilde M_D^\bullet.
\]
Combining Theorem~\ref{low:3:thm:two-regimes} with Lemma~\ref{low:3:lem:resource}, and using $v_D\asymp v_j$ on the coarse shape, yields
\begin{equation}\label{low:3:eq:shape-integral}
\int_{\mathcal A_R}\abs{\widetilde F_D^\bullet}^p\,dP\lesssim_{p,a} K^{A_p}L_R^{B_p}W_p(D).
\end{equation}

At a selected low-linear frame the exact inverse Taylor map, the bounded physical center and the bounded higher jets give
\begin{equation}\label{low:3:eq:shortest-side}
R\leq\norm P\lesssim \left(K\delta_1^{-1}+\delta_1^{-2}+\delta_1^{-3}\right).
\end{equation}
If $K=R^\beta$ with $\beta<2/3$, this implies $\delta_1\leq C R^{-1/3}$ for sufficiently large $R$. The quadratic term and the linear term would otherwise be $o(R)$ after taking a sufficiently large fixed multiple of $R^{-1/3}$ as a lower bound for $\delta_1$; the cubic term is made less than $R/2$ by that fixed multiple.

Since $p>18$, each exponent in \eqref{low:3:eq:W} is positive. Dropping the ordering restrictions in the last two length sums gives
\begin{align}\label{low:3:eq:shape-sum}
\sum_{\substack{D\text{ dyadic}\\\delta_1\leq C R^{-1/3}}}W_p(D)^{1/p}
&\lesssim_p\sum_{\delta_1\leq C R^{-1/3}}\delta_1^{(p-18)/p}\sum_{\delta_2\leq1}\delta_2^{(p-16)/p}\sum_{\delta_3\leq1}\delta_3^{(p-14)/p}\\
&\lesssim_p R^{-(p-18)/(3p)}.\nonumber
\end{align}
Apply Minkowski to \eqref{low:3:eq:shape-integral} and use \eqref{low:3:eq:shape-sum}. This proves \eqref{low:3:eq:coherent-small}.
\end{proof}

\paragraph{Completion of the upper bound}

\begin{proof}[Proof of Theorem~\ref{low:3:thm:main}]
Fix $18<p<19$. Proposition~\ref{low:3:prop:low-assembly} and the two-term triangle inequality give a low-linear $L^p(\mathcal A_R)$ bound
\[
\lesssim_{p,a}K^{A_p}L_R^{B_p}R^{-\alpha_p},\qquad \alpha_p=\frac{p-18}{3p}>0.
\]
Choose $0<\beta<2/3$ with $\beta A_p<\alpha_p/4$ and put $K=R^\beta$. Absorbing the logarithmic factor gives the low-linear bound $\lesssim_{p,a}R^{-\alpha_p/2}$.

The coefficient annulus has measure $\lesssim R^{19}$. Lemma~\ref{low:3:lem:tail} gives its tail norm
\[
\lesssim_{N,a}R^{19/p}\left(\frac{R^{\beta/12}}{L_R}\right)^{-N}.
\]
After choosing $\beta$, take an integer $N$ with $\beta N/12>19/p+\alpha_p/2$. Then $R^{19/p}(R^{\beta/12}/L_R)^{-N}=O(R^{-\alpha_p/2})$. The bounded coefficient region is controlled by $\abs{\mathcal E a}\leq\norm a_1$. Thus \eqref{low:3:eq:annular-main} holds with $\epsilon_p=p\alpha_p/2=(p-18)/6$, and dyadic annular summation gives global integrability. For $p\geq19$, choose $q\in(18,19)$ with $q<p$ and use $\abs{\mathcal E a}^p\leq\norm a_1^{p-q}\abs{\mathcal E a}^q$.

The cube conclusion follows from Lemma~\partIref{an:cube-reduction}, applied with degree three and coefficient character $\chi_3=15$. Its hypothesis $p>\chi_3$ follows from $p>18$. The smooth-amplitude estimate just proved has the required finite-order amplitude dependence, because every step uses only finitely many fixed amplitude derivatives.
\end{proof}


\section{Regular coefficient estimates in degrees four through thirteen}\label{reg:finite-section}\label{low:appendix}\label{low:section}\label{reg:physical-specializations}

We use the fixed physical frames, augmented stopping condition and regular assignment of Paper I, Section~\partIref{reg:geometry-section}. The finite sampling certificates below give complete coefficient-volume bounds at each original fixed frame, including the union over the allowed sample nets. Section~\ref{reg:combined} lists the resulting exponents. These estimates enter Proposition~\ref{an:unified-incidence}; the quartic cutoff and its side loss are retained there.

For the coefficient estimates from this section onward, $P\ge2$ is the conditioning parameter. The degree and prescribed finite derivative orders are fixed before scale parameters are chosen. For nonnegative $X,Y$, the notation $X\lesssim_P Y$ means $X\le P^CY$, and $X\gtrsim_PY$ means $X\ge P^{-C}Y$; $X\asymp_PY$ means both inequalities. An $O_P(Y)$ bound has absolute value or norm at most $P^CY$. Positive constants $C_P,c_P$ satisfy $C_P+c_P^{-1}\le P^C$, with the smallness or largeness specified in the corresponding argument. The exponents may depend on the fixed degree, chosen finite certificates, prescribed derivative orders, fixed amplitude support and indicated seminorms, and displayed exponent parameters. They are independent of the phase coefficients and varying geometric scales unless a dependence is stated. The scalar and matrix pivot bounds are stated below.

For $k=4,5,6$, use the physical family, positive reconstruction and mass weights of Paper I, Section~\partIref{an:section}. Lemmas~\partIref{an:character} and~\partIref{an:resource} give the physical exponents $p-\chi_k-3$, $p-\chi_k-1$ and $p-\chi_k+1$ before any incidence loss; their numerical values are in Table~\ref{app:table-physical}. In degree four, put $q=(1+\eta)(p-1)$, with $38<p<39$ and $\eta>0$ chosen so that $q<38$. The regular cutoff $H\le P^C/d_2$ gives the factor $d_2^{-(q-36)}$. Section~\ref{completion:physical-II} checks the physical sum with this factor.

In the physical estimates write $K=K_{\mathrm{lin}}$ and apply Lemma~\partIref{an:linear-tail} with physical amplitude derivatives, choosing its integration-by-parts order after the incidence exponents in $K,L_R$.

\subsection{Quartic regular estimates and the physical cutoff}\label{reg:quartic}\label{reg:finite4}\label{low:4:cf:framework}\label{low:4:rc:section}\label{low:4:cf:allocation}\label{low:4:cf:spine_geometry}\label{low:4:cf:counting}\label{low:4:rd:section}\label{low:4:rr:f}\label{low:4:rr:a}\label{low:4:rr:central}

Use the degree-four physical family of Section~\partIref{an:section}, with $K=K_{\mathrm{lin}}$, and the augmented states, generators, and stop~\partIeqref{reg:stop} of Section~\partIref{reg:geometry-section}. The complete windows, reciprocal stopping constants, and physical comparisons are bounded by powers of $P$; $\Lambda$ is the relative-frame packing multiplicity.

Apply the coefficient-independent RQ binning of Lemma~\partIref{an:rq}. In each bin the fixed representatives $T_j^0=D_\ell g_j$ retain the complete admissibility and the actual singular lengths $d_i$; the unordered diagonal entries of $D_\ell$ do not replace them. Recompute the rank and velocity tests in this presentation. Lemma~\partIref{reg:regular-preparation} supplies their polynomial conditioning and logarithmic multiplicity.

\paragraph{The augmented orbit differential and the rank-four boundary}

Here and below, coefficient norms use a fixed basis. Write $\mathcal V=V_4$ and retain the closed augmented window
\begin{equation}\label{low:4:rc:stop}
\|\widehat J\|\le P,\qquad \|[u\cdot\nabla J]\|^2+|Bu|^2\ge\kappa^2|u|^2\quad(u\in\mathbb R^3),\qquad P\ge2,\quad\kappa\ge P^{-1}.
\end{equation}
For a scalar pivot $q$, conditioning below includes $|q|\ge P^{-C}$ and the stated upper bound. For a matrix pivot $L$, it includes the stated bounds $s_{\min}(L)\ge P^{-C}$ and $\|L\|_{\mathrm{op}}\le P^C$ on the specified columns and coefficient space. Use the augmented orbit matrix $R=\mathcal R$ of Definition~\partIref{def:augmented-state}. Lemma~\partIref{geo:normal-pivot} gives the normal-column lower bound and $\sigma_4(R)\ge P^{-C}$, allowing singular physical matrices in the closed window. Theorem~\partIref{geo:rank-four-low} and Corollary~\partIref{geo:rank-four-entry-low} give the stratum~\partIeqref{low:4:rc:rankfour-form} and the small-fifth-pivot neighborhood~\partIeqref{low:4:rc:surface-family}, with its phase-only stop. Its actual critical equation is $J_x=0$.

Use the priority allocation of Lemma~\partIref{reg:regular-preparation}. A quantitative full-rank component has count $P^C\Lambda$. Write $L_i$ for $R$ with column $i$ omitted; $L_d,L_e,L_f$ omit $D,T_y,T_z$. After assigning the full-rank, small-fifth-pivot and good-$L_d$ regions, the remaining region is
\begin{equation}\label{low:4:cf:envelope}
\sigma_5(R)\ge\kappa,\qquad \sigma_6(R)<\varepsilon,\qquad \sigma_{\min}(L_d)<\mu.
\end{equation}
For a unit right singular vector $w$ corresponding to $\sigma_6(R)$, the common deleted-column estimate is
\begin{equation}\label{low:4:cf:delete}
\sigma_{\min}(L_i)\ge\kappa|w_i|.
\end{equation}
Together with the normal-column bound, this exhausts this remaining region by the $T_z$ and $T_y$ charts. All thresholds have the polynomial margins of Lemma~\partIref{low:4:cf:polynomial}.

Retain the following numerical margins for their later uses. Choose the central first-shear threshold $a_0$ first, then the $T_z$ threshold $\nu$, the threshold $\mu$ for $\sigma_{\min}(L_d)$, and $\varepsilon$ last. If $\delta_f$ and $\delta_A$ are the permitted fitted-$D$ cutoffs in the terminal and first-shear branches, require
\[
 \mu/\kappa\le\nu\delta_f/4,\qquad \varepsilon\le\kappa\nu^2\delta_f/4,\qquad \mu/\kappa\le a_0\delta_A/4.
\]
The fitted terminal-chart error is at most $\varepsilon/(\kappa\nu^2)$; choose $\varepsilon$ small enough for the first-shear chart as well. Assign $|w_{T_z}|\ge\nu$ first. On its complement, let $v_{T_y,A},v_{T_y,C}$ be the fitted $A,C$ components in the $T_y$ chart, as in the proof of Lemma~\partIref{reg:regular-preparation}. Assign $|v_{T_y,A}|\ge a_0$ to the first-shear chart and its complement, where $|v_{T_y,C}|\ge P^{-C}$, to the central chart. For each label the normalized representative is $g_j\circ r_j$, with $r_j$ in the selected transverse coordinate slice and $\|r_j-\mathrm{id}\|\le P^{-C_1}$ as in that lemma. Its deterministic ties and outer margins give exactly one assigned piece for every original admissible label; coefficient integration remains in the original frame.

Lemma~\partIref{reg:regular-preparation} gives the common level-set curves, their small right comparisons and their outer margins. In the $T_y$ chart the transverse slice is $h(a,b,c,d,0,f)$.

For a regular algebraic parametrization $s(t)$ of a $T_y$ section curve, its five fixed-output equations solve the tangent as
\begin{equation}\label{low:4:cf:clock}
s^{-1}s'=\alpha(t)\big(V_A A+V_C C+V_xT_x+V_DD+T_y+V_zT_z\big),
\qquad \alpha=e'-d f'.
\end{equation}
Orient the tangent so that $d\tau=\alpha\,dt>0$ in \eqref{low:4:cf:clock}; its body $T_y$ component is then one. Lemma~\partIref{reg:regular-preparation} applies to this transverse minor and intrinsic parameter.

Apply Lemma~\partIref{reg:regular-preparation} with the intrinsic parameter~\eqref{low:4:cf:clock}, preparing the frame-change formulas before amplitude division and retaining the final band in the equations. Its identity for $V_{\rm other}-v$ transfers the fitted tests with error $O_P(H^{-1})$; pieces with $H\le P^C$ enter the bounded term. The complete $34$-column sampling map has volume bound $P^C\Delta_0^{-1}\prod_{i=1}^{m-1}N_i$ by the same fixed-net argument, where $\Delta_0$ is a nonzero complete minor, $N_i$ are the free original-coordinate net sizes, and the anchor supplies the first block.

For a diagonal dilation $D_*=\diag(s_x,s_y,s_z)$ and selected output exponent sums $(r_x,r_y,r_z)$, the character calculation in Section~\partIref{reg:sampling} gives
\begin{equation}\label{low:4:cf:weight}
s_x^{35-r_x}s_y^{35-r_y}s_z^{35-r_z}.
\end{equation}
The row certificates in Section~\ref{ext-supp:low:4:rr:certificates} retain their triangular gradings and are compared on normalized compact sets. Use the slice-covering and one-sided sampling bounds of Lemma~\partIref{reg:regular-preparation}, including the factors $1+\sup|d|$ on an $e$ level and $1+\sup|a|$ on a $b$ level.

Ordered dilations $s_x\ge c s_y\ge c s_z\ge c$ absorb bounded upper comparisons in the sampling boxes. The nonordered $f$ normalization retains its one-parameter anchor cost in Section~\ref{low:4:rr:f}. The swapped central branch uses its non-upper change of frame only for Gram comparison, with nets still in the original upper-affine coordinates. Lemma~\partIref{reg:middle-cutoff}, applied to the original prepared sections and their retained physical comparisons, gives $H\le P^C/d_2$.

\paragraph{The regular chart omitting the second shear}

We retain the fixed upper-affine presentation of Lemma~\partIref{an:rq} and the augmented jet of Definition~\partIref{def:augmented-state}. In this regular-chart calculation, $d(g)$ denotes the fourth coordinate of $g\in\mathcal U$, not the ambient dimension. Write $A=y\partial_x$, $C=z\partial_x$, $D=z\partial_y$, and $T_i=\partial_i$, and put
\[
L_d=(A\widehat J,C\widehat J,T_x\widehat J,T_y\widehat J,T_z\widehat J),\qquad v=-L_d^+D\widehat J.
\]
We consider the part of the admissibility condition on which $\sigma_{\min}(L_d)\ge P^{-C}$. The number and degrees of the defining equations and inequalities are fixed.

Use the common sections and accumulated displacements of Section~\partIref{reg:geometry-section}. The model flows below give determinant bounds on short portions of these sections.

\begin{theorem}[The complete regular $d$ chart]\label{low:4:rd:theorem}

\noindent\textup{(i)} For the regular assignment of Lemma~\partIref{reg:regular-preparation}, length-$H$ pieces have count $P^C\Lambda H$ and the multiplicity stated there. At a fixed frame, the quartic coefficient set has volume at most
\[
P^C H^{-36}\bigl(1+\log d_1^{-1}\bigr)^C.
\]

\noindent\textup{(ii)} Consequently, for every fixed nonnegative sequence $b_j$, $38<p<39$ and sufficiently small $\eta>0$, the fixed-shape integral $\int(\sum_j b_jv_j\mathbf1_{\mathcal A_j^d(F)})^p\,dF$, where $\mathcal A_j^d$ is the regular $d$ admissibility condition, is
\[
\lesssim_p P^{C_p}\Lambda^{p-1}\bigl(1+\log d_1^{-1}\bigr)^{C_p}d_2^{-((1+\eta)(p-1)-36)}\sum_j b_j^pv_j^{p-35}.
\]

\noindent\textup{(iii)} The existing physical family sum is convergent for this bound.
\end{theorem}

\paragraph{Sections and the three cases}

Use the level-set curves of Lemma~\partIref{reg:regular-preparation}, whose corrections lie in $\mathcal U_d=\{d=0\}$ and preserve the regular parameter $t=d$. Lemma~\partIref{reg:d-sections} gives the fiber comparison and bounded count over unit intervals. Use the body rates and $\tau=\xi-\beta\beta_y$ of Section~\partIref{reg:d-channels}. Its stopped inequality assigns every sufficiently long prepared section to the significant-first-shear, significant-terminal or significant-mixed case.

Prepare the augmented jet, body rates and frame-change expressions together by Lemma~\partIref{reg:regular-preparation}. The length-$H$ pieces have count $P^C\Lambda H$ and
\begin{equation}\label{low:4:rd:preparation}
\|\widehat J^{(r)}\|\le P^C H^{-r},\qquad |u^{(r)}|\le P^C\sup_I|u|\,H^{-r}
\end{equation}
Retain the lemma's treatment of the final small-amplitude band. Equation~\partIeqref{reg:d-section-rates} gives the $O(P^C/H)$ least-squares error, so the prescribed inner and outer margins apply to this same assignment.

\paragraph{Sampling minors}

Use the sampling matrix~\partIeqref{reg:sample-entry} with $k=4$, ordering monomials by total degree and then lexicographically. The seven row sets are in supplementary Section~\ref{ext-supp:low:4:rd:rows}; their nodes are $t=0,\ldots,m-1$. Write $l_t,q_t,p_t,m_t,r_t$ for the models $L_t,Q_t,P_t,M_t,R_t$ of Section~\partIref{reg:sampling}.

For each certificate, the selected coefficients of the sampled polynomials give a $34\times34$ linear map on the nonconstant monomial coefficients. Its entries are those of~\partIeqref{reg:sample-entry}. Polynomial expansion and exact rational elimination give the displayed determinant magnitudes; the elimination procedure is in the proof of supplementary Lemma~\ref{ext-supp:low:4:rr:minor_lemma}. ``At unit scale'' means before the diagonal dilation~\partIeqref{reg:minor-scaling}. Only absolute determinants are used; their signs depend on row ordering.

For each certificate let $(R_x,R_y,R_z)$ be the sum of its selected output exponents. The identities are
\begin{center}
\begin{tabular}{cclc}
Certificate & States & Determinant at unit scale & $(R_x,R_y,R_z)$\\\hline
$L_4$ & $l_0,\ldots,l_3$ & $531441/1024$ & $(1,31,73)$\\
$Q_3$ & $q_0,\ldots,q_2$ & $-1925/867623536069632$ & $(6,20,53)$\\
$P_4$ & $p_0,\ldots,p_3$ & $-18600435/64$ & $(35,4,31)$\\
$M_5$ & $m_0,\ldots,m_4$ & $15850845241344$ & $(4,11,59)$\\
$L_5$ & $l_0,\ldots,l_4$ & $-9795520512$ & $(11,6,88)$\\
$R_4$ & $r_0,\ldots,r_3$ & $139314069504$ & $(8,11,50)$\\
$Q_4$ & $q_0,\ldots,q_3$ & $-14105/248832$ & $(11,7,54)$
\end{tabular}
\end{center}
The scale factor is~\partIeqref{reg:minor-scaling}.

The $L_4,L_5$ row sets have total-degree multiplicities $(3,6,10,15)$, matching the input multiplicities. Lemma~\partIref{reg:lem:transfer} therefore makes their determinants independent of arbitrary separate sample translations.

By the degree-three ordered-exponential comparison in Lemma~\partIref{reg:regular-preparation}, a $\delta$ perturbation of bounded normalized rates changes the sample tuple by at most $P^C\delta$. Choosing $\delta$ smaller than the polynomial conditioning margin of the displayed nonzero minors preserves their lower bounds, with sampling fractions and dominance thresholds given by powers of $P$.

Apply the complete Gram comparison in Lemma~\partIref{reg:regular-preparation}. It replaces the bounded coefficient block at the corrected first sample by the identity coefficient block at the original fixed frame. Only the other $m-1$ samples require fixed original-coordinate nets. Since $d$ is a homomorphism and the section correction lies in $\mathcal U_d$, the relative $d$ levels are fixed multiples of $H$ in this coefficient frame.

Use the original-coordinate box and fixed-$d$ net in~\partIeqref{reg:d-box}, with the RQ bin and coefficient anchor fixed. A certificate with $m$ states and determinant $\Delta$ therefore gives coefficient volume at most $P^C\Delta^{-1}(H^3U^3W^3)^{m-1}$.

\paragraph{The three cases}

Apply the dilations and short-interval comparisons of Section~\partIref{reg:d-channels} with the quartic minors above. In the first-shear case $U\asymp_PH$ and $W=\max(1,H\sup|\gamma|)$; in the terminal case $W\asymp_PH$ and $U=\max(1,H\sup|\alpha|)$. In the mixed case use the exact rates in~\partIeqref{reg:d-gauge}, with $u=H\sup|\alpha+\beta'|$, $w=H\sup|\gamma-\beta_y'|$, $U=\max(1,u)$ and $W=\max(1,w)$. The return to original coordinates proved there gives~\partIeqref{reg:d-box} also when $u<1$; in particular its original $b$ width is $H$ in that range.

The determinant lower bounds and the products with the $m-1$ original-coordinate nets are as follows, with powers of $P$ suppressed. Bounded and large ranges use the overlapping conditioned thresholds of Section~\partIref{reg:d-channels}.
\begin{center}\small
\begin{tabular}{llll}
\toprule
Case & Certificate & Determinant & Coefficient volume\\
\midrule
First shear, bounded $W$ & $L_4$ & $H^{72}$ & $H^{-54}$\\
First shear, large $W$ & $Q_3$ & $H^{73}W^{26}$ & $H^{-61}W^{-20}$\\
Terminal, bounded $U$ & $P_4$ & $H^{66}$ & $H^{-48}$\\
Terminal, large $U$ & $Q_3$ & $H^{44}U^{29}W^{26}$ & $H^{-38}U^{-23}W^{-20}$\\
Mixed, bounded $u,w$ & $M_5$ & $H^{55}$ & $H^{-43}$\\
Mixed, large $u$, bounded $w$ & $L_5$ & $H^{53}U^{24}$ & $H^{-41}U^{-12}$\\
Mixed, large $w$, bounded $uw$ & $R_4$ & $H^{51}w^9$ & $H^{-42}$\\
Mixed, large $w,uw$, $u\ge1$ & $Q_4$ & $H^{52}u^{24}w^{33}$ & $H^{-43}u^{-15}w^{-24}$\\
Mixed, large $w,uw$, $u<1$ & $Q_4$ & $H^{52}u^{24}w^{33}$ & $H^{-43}(uw)^{-24}$\\
\bottomrule
\end{tabular}
\end{center}

The first-shear bounds sum in $W$. In the terminal case $W\asymp_PH$, so the large-$U$ row is at most $P^CH^{-58}U^{-23}$ and sums in $U$. In the mixed case the bounded-$w$ rows sum in $U$. For large $w$ and bounded $uw$, the bound is independent of the small $u$, and the $O_P(1+\log H)$ possible $w$ bands cost only a logarithm. For large $w,uw$, the $u\ge1$ row has two convergent displacement sums; for $u<1$, sum first in the product $uw$ and then in $w\lesssim_PH$. No lower cutoff for the instantaneous first-shear rate is required.

Use the threshold order in Section~\partIref{reg:d-channels}: first the cubic fraction and large-product cutoff, then a shorter parabolic fraction and terminal-displacement cutoff, then the linear and bounded-corner fractions and the minimum long $H$. Its overlapping ranges cover every $u,w\ge0$. The table consequently gives the summed bound $P^CH^{-36}(1+\log d_1^{-1})^C$ required in Theorem~\ref{low:4:rd:theorem}.

Proposition~\partIref{an:common-incidence}, with $q=(1+\eta)(p-1)$ and Lemma~\partIref{reg:middle-cutoff}, gives the weight $v_j^{p-35}$ and loss $d_2^{-(q-36)}$. For $38<p<39$, choose $1-\eta(p-1)>0$; the physical summation is~\partIeqref{an:physical-positivity} with $k=4$ and $\gamma=q-36$. This proves Theorem~\ref{low:4:rd:theorem}.
\begin{theorem}[Residual regular coefficient-volume bound]\label{low:4:rr:tail}

\noindent\textup{(i)} For every original index assigned to the region \eqref{low:4:cf:envelope}, the common preparation of Lemma~\partIref{reg:regular-preparation} has count at most $P^C\Lambda H$ on a length-$H$ piece.

\noindent\textup{(ii)} In the coefficient coordinates $J=\mathcal L_{T_j}F$ of its fixed label $j$, with the representative convention of Lemma~\partIref{an:rq}, the coefficient set $E_j(H)$ that this label belongs to such a piece satisfies
\[
\vol E_j(H)\le P^C H^{-36}(1+\log P+\log d_1^{-1})^C,
\qquad 1\le H\le P^C d_2^{-1}.
\]
At most $P^C(1+\log P+\log d_1^{-1})^C$ pieces occur at each dyadic length. The construction preserves the original labels and all their weights.
\end{theorem}

\paragraph{Step 1. \texorpdfstring{The small-$D$ chart omitting $T_z$}{The small-D chart omitting Tz}}
Use the residual chart omitting $T_z$ from Section~\partIref{reg:gauges}, with $f$ as parameter and $\chi=\zeta$. Its bounded frame change, exact body equations and significant-rate alternative apply here. Retain the scales
\[
u=H^2\sup|\delta|,\qquad v=H\sup|\chi|,\qquad U=\max(1,u),\qquad r=H\sup|\alpha|.
\]
The original-coordinate box and sampling net are \partIeqref{reg:f-box}. Choose the comparison fractions after the quartic determinant margins below.
The following certificates will be used:
\begin{center}\small
\begin{tabular}{llll}
\toprule
Model & States & Output sums & Determinant after dilation\\
\midrule
$A+T_z$ &5&$(0,70,15)$&$15850845241344\,H^{55}$\\
$C+T_z$ &4&$(4,35,31)$&$\frac{18600435}{64}H^{35}v^{31}$\\
Cubic chain &5&$(8,48,8)$&$-4173258660249600\,H^{54}U^{14}$\\
Cubic chain &4&$(6,29,29)$&$\frac{416912}{19683}H^6r^{29}u^{35}$\\
\bottomrule
\end{tabular}
\end{center}
The respective dilations are $\diag(H,1,H)$, $\diag(Hv,U,H)$, $\diag(HU,U,H)$, and $\diag(ru,u,H)$. The mixed five-state minor depends only on the $a_i,f_i$ coordinates: grade inputs by $\deg_x+\deg_y$ and selected rows by $y$ degree. Its square diagonal blocks of sizes $4,8,9,8,5$ ignore every $b_i,c_i,d_i,e_i$. The parabolic minor at $a_i=0$ is independent of separate $d_i,e_i$, by $y$-degree blocks of sizes $14,10,6,3,1$. These observations give uniform stability in the indicated bounded additional coordinates.

When $u\lesssim H$, remove by right multiplication the bounded $d$ coordinate. In the $C$-dominated boxes, put $U=\max(1,u)\lesssim v\lesssim H$. At prescribed $f$ levels the remaining widths are $(a,b,c,e)=(H,v,Hv,U)$, giving net cardinality per additional sample
\begin{equation}\label{low:4:rr:fwitness}
N=CH^2v^2U.
\end{equation}
The common bounded $D$-anchor comparison contributes one additional grid factor. Factor a bounded $E\in\mathcal U$ as $E=\exp(sD)R$, $R\in\{d=0\}$. Conjugation by $R$ preserves these boxes and the prescribed $f$ levels, up to fixed enlargement. A mesh in $s$ of width $cU/(H^2v)$ makes the normalized error $O((Hv)^{-1})$ under $\diag(Hv,U,H)$. Conjugation through the bounded normalized sample frames and return to original coordinates then gives bounded right errors. Thus the grid cardinality is
\begin{equation}\label{low:4:rr:dcost}
Q=CH^2v/U.
\end{equation}
This is one shared factor, not one factor per additional sample. It handles the nonordered middle dilation length. When the dilation lengths are ordered, use the full fixed-$f$ box with no additional anchor net.

The displacement table follows. All inequalities allow fixed overlap constants. By Lemma~\partIref{low:4:cf:polynomial}, these constants and the sampling fractions below cost at most a power of $P$. First choose a cubic freezing fraction, then a large cubic-dominance threshold. For the bounded complementary ratio choose a separate smaller parabolic fraction, then a large-$v$ cutoff. The remaining bounded corner is treated by the mixed certificate. 

If $\abs\alpha\gtrsim1$ and $u,v$ are bounded, the five-sample bound is $H^{-55}(H^2)^4H^2=H^{-45}$. If $u\gg\max(1,v)$ and $1\lesssim U\asymp u\lesssim H$, the cubic certificate gives
\[
H^{-54}U^{-14}(H^2U^3)^4H^2=H^{-44}U^{-2}.
\]
For the same dominance with $U\gtrsim H$, four cubic samples in ordered scales give
\[
H^{-35}U^{-35}(HU^4)^3=H^{-32}U^{-23},
\]
whose dyadic $U$ sum is $O(H^{-55})$.

If $\abs\alpha\gtrsim1$, $v$ is large, and $u\lesssim v$, both mixed and parabolic certificates apply after a fixed short fraction. Indeed, on $h=\theta H$ the scales are $u_h=\theta^2u$, $v_h=\theta v$, and the normalized $A$ effect is $\lesssim\theta u/v+1/(\theta v)$. Choose $\theta$ first and then the large-$v$ threshold to make this small. The normalized $d,e$ coordinates stay bounded and disappear from the diagonal grading blocks at $A=0$, so continuity is uniform in them. With \eqref{low:4:rr:fwitness}--\eqref{low:4:rr:dcost} the certificates give two bounds for the same coefficient set:
\[
H^{-55}N^4Q=H^{-45}v^9U^3,\qquad
H^{-35}v^{-31}N^3Q=H^{-27}v^{-24}U^2.
\]
Sum dyadic $U\lesssim v$, obtaining the minimum of $H^{-45}v^{12}$ and $H^{-27}v^{-22}$. Their crossover is $v=H^{9/17}$. Both geometric sums therefore give
\begin{equation}\label{low:4:rr:ftail}
\sum_v\min(H^{-45}v^{12},H^{-27}v^{-22})\lesssim H^{-657/17}.
\end{equation}

If instead $\abs\chi\gtrsim1$, then $v\asymp H$. For $u\lesssim H$, the parabolic estimate above sums to $H^{-49}$. For $u\gtrsim H$ and $\abs\alpha u\lesssim v$, the ordered dilation $\diag(Hv,u,H)$ and full fixed-$f$ cost $Hv^2u^2$ give
\[
H^{-35}v^{-31}(Hv^2u^2)^3=H^{-32}v^{-25}u^6.
\]
Since $v\asymp H$ and $u\lesssim H^2$, its dyadic sum is $O(H^{-45})$. Finally, if $u\gtrsim H$ and $\abs\alpha u\gg v$, put $r=H\sup\abs\alpha$; then $ru\gtrsim H^2$. The ordered cubic dilation has fixed-$f$ cost $r^3u^4/H^2$, and
\[
H^{-6}r^{-29}u^{-35}(r^3u^4/H^2)^3
=H^{-12}r^{-20}u^{-23}\lesssim H^{-52}u^{-3}.
\]
Summing first $r$ and then $u\gtrsim H$ gives $O(H^{-55})$. All small $\alpha,\delta,\chi$ rates are retained and relatively prepared where needed. These alternatives are exhaustive and prove \eqref{low:4:rr:ftail}, with at most the physical logarithmic preparation multiplicity. Since $657/17>38$, this chart needs no physical power loss for $38<p<39$.

\paragraph{Step 2. \texorpdfstring{The significant-$A$ part of the $T_y$ chart}{The significant-A part of the Ty chart}}
Here $T_y$ and $A$ kernel components are significant and $D,T_z$ components are small. Equation~\eqref{low:4:cf:delete} supplies a good chart omitting $A$. Its homomorphic coordinate is $a=t$. Write
\[
V=A+\beta C+\tau T_x+\delta D+\eta T_y+\gamma T_z.
\]
Use the bounded change of frame and coordinate box from the significant-first-shear calculation in Section~\partIref{reg:gauges}. After choosing the terminal cutoff and separating short intervals, $|\eta+\beta\gamma-\tau'|\ge P^{-C}$ on the retained long pieces. Put $u=H\sup\abs{\delta-\beta'}$, $w=H\sup\abs\gamma$, $U=\max(1,u)$, $W=\max(1,w)$. The original-coordinate sampling nets below include the bounded comparisons specified there.
If $UW\lesssim H$, four parabolic states under $\diag(H^2,H,W)$ have determinant
\[
-624127031377920\,H^{66}.
\]
Its output exponent sums are $(4,31,35)$, so no $W$ power remains. The determinant is independent of separate $z$ translations, by a triangular $z$-degree grading. The normalized $D$ correction is made small on a fixed fraction. The fixed-$a$ sampling net cardinality is $H^4U^2W$, so the actual fixed-upper-group count is
\[
H^{-66}(H^4U^2W)^3=H^{-54}U^6W^3.
\]
Summing $U\lesssim H/W$ gives $O(H^{-48}W^{-3})$ and then $O(H^{-48})$.

In the complementary case $u,w$ before taking maxima with one are bounded below. Three cubic samples, using dilation $\diag(Huw,uw,w)$ and output sums $(6,21,51)$, give
\[
\frac{1925}{1301435304104448}H^{29}u^{43}w^{27}.
\]
The fixed-$a$ cost is $H^2u^4w^3$. Two additional samples give
\[
H^{-29}u^{-43}w^{-27}(H^2u^4w^3)^2
=H^{-25}u^{-35}w^{-21}
=H^{-46}(uw/H)^{-21}u^{-14}.
\]
The displayed negative powers give convergent dyadic sums over the displacement parameters. At a boundary where a power is zero, the number of permitted dyadic values is bounded by a power of $1+\log P+\log d_1^{-1}$. The small-product argument extends to any fixed large boundary ratio by shortening the sampling interval by a fixed factor; the cubic argument applies beyond that ratio. These estimates imply the stated weaker tail $H^{-38}$, with the permitted physical logarithms. Both dilations are ordered up to fixed factors in their stated ranges, so the original-coordinate nets already include every required frame choice.

\paragraph{Step 3. Four cases in the central \texorpdfstring{$C,T_y$}{C,Ty} chart}
Use the prepared central normal form and four overlapping cases of Section~\partIref{reg:gauges}. Its significant rate is $\beta$, and its equations and original-coordinate box are \partIeqref{reg:central-equations}--\partIeqref{reg:central-box}. Retain
\[
U=\max(1,H\sup|\alpha|),\qquad V_0=\max(1,H\sup|\delta|),\qquad W=\max(1,H\sup|\gamma|),
\]
and write $M=\max(U,W)$. Choose the comparison fractions and thresholds in the order specified there, after fixing the quartic determinant margins. The cases below use the two ratios $UV_0/H$ and $V_0W/H$.

\paragraph{Both products bounded}
Suppose $V_0M\lesssim H$. The boxes become
\[
(a,b,c,d,e,f)=O(U,H,HM,V_0,H,W).
\]
Under $\diag(HM,H,M)$ the normalized $D$ effect becomes small on the prescribed short fraction. The $b,e$ entries follow nonzero linear progressions up to small normalized errors. The four-state mixed $C+T_y$ certificate is
\begin{equation}\label{low:4:rr:centralmixed}
\Delta=-61917364224\,H^{51},\qquad (r_x,r_y,r_z)=(4,15,66).
\end{equation}
At $d_i=0$, grade inputs and outputs by $\deg_x+\deg_z$. The blocks have sizes $4,8,9,8,5$ on both sides. Their diagonals depend only on $b_i,e_i$; separate $a_i,c_i,f_i$ only lower the grade. Thus \eqref{low:4:rr:centralmixed} is independent of all those additional coordinates, and is uniformly stable under the small normalized $d_i$ perturbation. Formula~\eqref{low:4:cf:weight} shows exact cancellation of the $M$ power.

Prescribe the relative $e$ levels. Equation~\partIeqref{reg:central-equations} has $e'=1+d\gamma$ bounded away from zero after the chosen fraction. The fixed-$e$ five-coordinate volume is $H^2UV_0WM$; covering its intersection with relative affine neighborhoods requires the additional factor $V_0$ proved in Section~\ref{low:4:cf:counting}. The net cardinality per additional sample is therefore
\[
N_e=CH^2UV_0^2WM.
\]
Four samples and three additional samples give
\[
H^{-51}N_e^3\lesssim H^{-45}U^3V_0^6W^3M^3\lesssim H^{-36},
\]
because $U,W\le M$, $V_0M\lesssim H$ and $M\lesssim H$. Dyadic displacement sums add at most powers of $\log H$. This includes the cancellation $\alpha+\beta\gamma=0$; no significance of that combination is assumed.

\paragraph{\texorpdfstring{$UV_0$ large and $V_0W$ bounded}{UV large and VW bounded}}
On the chosen fraction, $UV_0/H$ is large and $V_0W/H$ small. Normalize by $\diag(UH,H,H/V_0)$. The $A,D,T_y$ rates are significant and relatively frozen; the $C$ and $T_z$ effects have respective sizes $H/(UV_0)$ and $V_0W/H$. The model flow is
\[
(x+t y+t^2z/2+t^2/2,\ y+t z+t,\ z).
\]
Its four-state certificate has output sums $(11,11,46)$ and determinant
\[
\Delta=-2125764\,H^{37}U^{24}V_0^{11}.
\]
All normalization lengths are bounded below, and they are ordered up to fixed factors. At prescribed homomorphic $d$ levels the five widths $(a,b,c,e,f)$ can be enlarged to $(U,UV_0,UH,H,H/V_0)$, of product $H^3U^3$. Hence
\[
\vol E\lesssim H^{-37}U^{-24}V_0^{-11}(H^3U^3)^3
=H^{-28}U^{-15}V_0^{-11}\lesssim H^{-39}.
\]
The last inequality uses $UV_0\gtrsim H$ and $U\ge1$. The model field retains the $T_y$ term.

\paragraph{Both products large}
Normalize by $\diag(UV_0W,V_0W,W)$. The relatively frozen $A,D,T_z$ rates form a cubic chain. The other two rates have normalized sizes $H/(UV_0)$ and $H/(V_0W)$ and are small by the thresholds. Three states give
\[
\Delta=\frac{1925}{1301435304104448}U^{29}V_0^{43}W^{27}.
\]
At prescribed $d$ levels the five widths are $(U,UV_0,UV_0W,V_0W,W)$, of product $U^3V_0^3W^3$. Two additional samples yield
\[
U^{-23}V_0^{-37}W^{-21}\lesssim H^{-37},
\]
since $U,W\gtrsim H/V_0$ and $V_0\lesssim H$ imply $U^{23}V_0^{37}W^{21}\gtrsim H^{44}V_0^{-7}\gtrsim H^{37}$.

\paragraph{\texorpdfstring{$V_0W$ large and $UV_0$ bounded: the coordinate interchange}{VW large and UV bounded: the coordinate interchange}}\label{low:4:rr:swap}
Put $R_0=UV_0/H$. Apply the complete coordinate-interchange comparison \partIeqref{reg:swapped-field}--\partIeqref{reg:swapped-error}, retaining its exact $u_*$, $X=\max(1,u_*W)$, lower-shear and diagonal terms, and ordered $Q/R$ comparison fractions. That comparison permits sampling at original relative $b$ levels and retains the original box
\[
(a,b,c,d,e,f)=O_P(U,H,HW,V_0,V_0W,W).
\]
Its bounded-product and large-product cases give the two quartic estimates below.

The model is precisely the significant-mixed $d$ system treated by the finite certificates below. If $u_*W$ is bounded, a further fixed fraction gives its four-state mixed determinant $cH^{51}W^9$. If $u_*W$ is large, the four-state cubic determinant is $cH^{52}u_*^{24}W^{33}$. The row sets $\mathcal M_d,\mathcal C_d$ in Section~\ref{ext-supp:low:4:rr:certificates} of the companion supplement give absolute constants $139314069504$ and $14105/248832$, respectively. The first needs significant normalized $T_x$ drift, which holds because it is $1/\beta$; the cubic certificate does not need that drift to remain significant after dilation. Relative preparation of the exact $u_*$ rate and $\gamma$ gives uniform stability. Equation~\partIeqref{reg:swapped-error} transfers it to the actual full system.

The determinant is transferred back by full Gram comparability as in Section~\ref{low:4:cf:counting}. The fixed permutation and common $S$ at the anchor have coefficient determinants of absolute value one; each separate right $S$ and its inverse is bounded. We then net the original upper-affine boxes just displayed. 

At prescribed $b$ levels, the five-coordinate volume in $(a,c,d,e,f)$ is $HU V_0^2W^3$. Covering this slice by relative affine neighborhoods requires the additional factor $U$, so one additional sample costs
\[
N_b=CHU^2V_0^2W^3.
\]
Three additional samples give respectively
\[
H^{-51}W^{-9}N_b^3\lesssim H^{-48}(UV_0)^6=H^{-42}R_0^6,
\]
\[
H^{-52}u_*^{-24}W^{-33}N_b^3
\lesssim H^{-49}(UV_0)^6(u_*W)^{-24}
=H^{-43}R_0^6(u_*W)^{-24}.
\]
All displacement sums converge or cost only logarithms. Together the four central cases prove $\vol E_j(H)\lesssim_P H^{-36}(1+\log P+\log d_1^{-1})^C$ for every assigned central label.

\begin{proof}
\par\smallskip\noindent\textit{Step 4. Assembly at the original label.}\par\nobreak\smallskip

The exhaustive priority assignment is Section~\ref{low:4:cf:allocation}; its level-set curves and their preparations are Lemma~\partIref{reg:regular-preparation}. Section~\ref{low:4:rr:f} gives the stronger tail $H^{-657/17}$, Section~\ref{low:4:rr:a} gives a tail stronger than $H^{-38}$, and the four exhaustive cases of Section~\ref{low:4:rr:central} give respectively $H^{-36},H^{-39},H^{-37},H^{-42}$, up to powers of the physical logarithm. There are only polynomially many choices of chart, sign, output row, output-grid value and sampling fraction by Lemma~\partIref{low:4:cf:polynomial}. Lemma~\partIref{reg:regular-preparation} applies after the fixed normalization at the fixed frame to each selected tuple. Thus summing all choices already counted gives the displayed volume. Lemma~\partIref{reg:middle-cutoff} gives the length cutoff. 
\end{proof}

\subsection{Quintic and sextic sampling bounds}\label{reg:five-six}\label{reg:finite5}\label{low:5:qcf:framework}\label{low:5:qcf:allocation}\label{low:5:qcf:spine_geometry}\label{low:5:qcf:counting}\label{low:5:qrc:section}\label{low:5:rd5:section}\label{low:5:qrr:residual}\label{reg:finite6}\label{low:6:s36cf:framework}\label{low:6:s6r:section}

For $k=5,6$, use the coefficient spaces and physical family of Section~\partIref{an:section} and the regular construction of Section~\partIref{reg:geometry-section}. Lemma~\partIref{an:rq} gives $T_j=D_\ell g_jA_j$, with $1\le|\det A_j|<8$, preserving the coefficient factor $v_j^{p-\chi_k}$ up to fixed factors. Perform the rank and velocity tests in this presentation, retaining the singular values, complete windows, augmented stop~\partIeqref{reg:stop}, and packing multiplicity $\Lambda$.

Here and below, an algebraic or semialgebraic family has fixed format when its number of variables, defining polynomials, and polynomial degrees are bounded in terms of the fixed degree and the stated auxiliary data. For semialgebraic paths we use \cite[Theorems 5.22 and 5.23]{BPR}; the uniform-format and variation bounds are recorded in Lemma~\partIref{geo:paths}. Lemma~\partIref{geo:conditioning} supplies polynomial dependence on $P$ on the closed stopped windows with their prescribed outer slack. On the regular region, Lemma~\partIref{reg:regular-preparation} gives the assignment, intrinsic preparation, bounded full-rank count and anchored count~\partIeqref{reg:fixed-net-bound}. Lemma~\partIref{reg:middle-cutoff} gives $H\lesssim_Pd_2^{-1}$. The exact parameters, moving frames, original-coordinate nets and overlapping scale cases are those of Section~\partIref{reg:gauges}; their sampling fractions are chosen after the finite row sets below. In particular all derived rates are prepared before amplitude division.

Use the sampling matrix, diagonal scaling and original-coordinate sample-net cardinalities of Section~\partIref{reg:sampling}. Each application retains its selected rows and verifies the grading hypothesis of Lemma~\partIref{reg:lem:transfer}.

\paragraph{Quintic regular pieces}

The physical incidence estimate for $k=5$ is proved in Proposition~\ref{an:unified-incidence}. Here we establish its regular coefficient input.

For the quintic rank-four specialization, retain the conditioned augmented window
\begin{equation}\label{low:5:qrc:stop}
\|\widehat J\|\le P,\qquad \|[u\cdot\nabla J]\|^2+|Bu|^2\ge\kappa^2|u|^2\quad(u\in\mathbb R^3),\qquad P\ge2,\quad\kappa\ge P^{-1}.
\end{equation}
Here $R_{\widehat J}$ is the six-column augmented orbit matrix of Section~\partIref{geo:augmented-models}; $B$ may be singular. Lemma~\partIref{geo:normal-pivot} supplies the normal-column pivot and the lower bound for the fourth singular value.

\paragraph{Classification of exact rank four}

Apply Theorem~\partIref{geo:rank-four-low} and Corollary~\partIref{geo:rank-four-entry-low} with $k=5$ to \eqref{low:5:qrc:stop}. The rank-four phases have the form \partIeqref{low:4:rc:rankfour-form}; for every prescribed defect $P^{-L}$, a sufficiently small fifth singular value places the phase in the square-quadric neighborhood \partIeqref{low:4:rc:surface-family} and gives $\|B\|_{\mathrm{op}}\le\kappa/2$ and the phase-only stop $\|[u\cdot\nabla J]\|\ge(\sqrt3/2)\kappa|u|$.

The critical-surface construction applies to the actual quartic polynomial $J_x$ in this neighborhood.

\paragraph{The quintic regular chart omitting the second shear}

Use the regular $d$ chart of Section~\partIref{reg:d-channels}, with $\sigma_{\min}(L_d)\ge P^{-C}$. 

\begin{theorem}[The complete quintic regular $d$ chart]\label{low:5:rd5:theorem}

\noindent\textup{(i)} Under the same regular assignment, now with $k=5$, the fixed-anchor coefficient set has volume at most
\[
P^C H^{-89}\bigl(1+\log d_1^{-1}\bigr)^C.
\]

\noindent\textup{(ii)} Consequently, for every fixed nonnegative sequence $b_j$, $73<p<74$, and sufficiently small $\eta>0$, the fixed-shape integral $\int(\sum_j b_jv_j\mathbf1_{\mathcal A_j^{d}(F)})^p\,dF$ is
\[
\lesssim_p P^{C_p}\Lambda^{p-1}\bigl(1+\log d_1^{-1}\bigr)^{C_p}\sum_j b_j^pv_j^{p-70}.
\]
Here $\mathcal A_j^d$ is the regular $d$ admissibility condition.

\noindent\textup{(iii)} For the physical choice $b_j=\omega_j$, its family shape sum converges. The conclusion concerns the regular $d$ allocation only.
\end{theorem}

Use the level-set curves of Lemma~\partIref{reg:regular-preparation}, parametrized by $t=d$. Their transverse corrections lie in $\mathcal U_d$, so the parameter is unchanged. Lemma~\partIref{reg:d-sections} supplies the fiber comparison, while the actual sampling rates are those of these level-set curves. Section~\partIref{reg:d-channels} gives the exhaustive significant-$\alpha$, significant-$\gamma$ and significant-$\tau$ allocation, with $\tau=\xi-\beta\beta_y$. All exact change of frame rates are included in the common preparation~\partIeqref{reg:prepared-derivatives}.

\paragraph{Sampling minors}

Use~\partIeqref{reg:sample-entry} with $k=5$, ordering monomials by total degree and then lexicographically. The seven $55$-row sets are in supplementary Section~\ref{ext-supp:low:5:rd5:rows}, at nodes $t=0,\ldots,m-1$. The models are those of Section~\partIref{reg:sampling}.

For each certificate let $(R_x,R_y,R_z)$ be the sum of its selected output exponents. The identities are
\begin{center}
\begin{tabular}{cclc}
Certificate & States & $|\det|$ at unit scale & $(R_x,R_y,R_z)$\\\hline
$L_4$ & $l_0,\ldots,l_3$ & $3^{20}/2^{20}$ & $(4,64,142)$\\
$Q_3$ & $q_0,\ldots,q_2$ & $37/(2^{19}3^{44})$ & $(14,38,112)$\\
$P_4$ & $p_0,\ldots,p_3$ & $2^3 3^{24}5^2 7$ & $(70,25,30)$\\
$M_5$ & $m_0,\ldots,m_4$ & $2^{59}3^{23}5$ & $(5,29,111)$\\
$L_5$ & $l_0,\ldots,l_4$ & $2^{25}3^{28}$ & $(24,15,171)$\\
$R_4$ & $r_0,\ldots,r_3$ & $2^{27}3^{19}5^3 59$ & $(16,24,103)$\\
$Q_4$ & $q_0,\ldots,q_3$ & $163\cdot1907/(2^{16}3^{16})$ & $(17,18,120)$
\end{tabular}
\end{center}
The scale of each minor is given by~\partIeqref{low:56:weight} in degree five.

The $L_4,L_5$ row sets have total-degree multiplicities $(3,6,10,15,21)$, matching the input multiplicities. Lemma~\partIref{reg:lem:transfer} therefore makes their determinants independent of arbitrary separate sample translations.

The model comparisons of Section~\partIref{reg:d-channels} apply to these row sets, with sampling fractions chosen for their nonzero determinant margins. The additional samples use the box and fixed-$d$ net~\partIeqref{reg:d-box}; Lemma~\partIref{reg:regular-preparation} absorbs the bounded section curve comparison into the identity-anchor block.

\paragraph{The three cases}

Use the three exhaustive cases and their exact rates and scales from Section~\partIref{reg:d-channels}. With the original net $N_d\lesssim_PH^3U^3W^3$ of \partIeqref{reg:d-box}, including $u<1$, the determinant bounds and \partIeqref{reg:fixed-net-bound} give the following estimates, with powers of $P$ suppressed:
\begin{center}\small
\begin{tabular}{llll}
\toprule
Case & Certificate & Determinant & Coefficient volume\\
\midrule
First shear, bounded $W$ & $L_4$ & $H^{138}$ & $H^{-120}$\\
First shear, large $W$ & $Q_3$ & $H^{144}W^{46}$ & $H^{-132}W^{-40}$\\
Terminal, bounded $U$ & $P_4$ & $H^{130}$ & $H^{-112}$\\
Terminal, large $U$ & $Q_3$ & $H^{88}U^{56}W^{46}$ & $H^{-122}U^{-50}$\\
Mixed, bounded $u,w$ & $M_5$ & $H^{106}$ & $H^{-94}$\\
Mixed, large $u$, bounded $w$ & $L_5$ & $H^{101}U^{46}$ & $H^{-89}U^{-34}$\\
Mixed, large $w$, bounded $uw$ & $R_4$ & $H^{100}w^{13}$ & $H^{-91}w^{-4}$\\
Mixed, large $w,uw$, $u\ge1$ & $Q_4$ & $H^{105}u^{53}w^{55}$ & $H^{-96}u^{-44}w^{-46}$\\
Mixed, large $w,uw$, $u<1$ & $Q_4$ & $H^{105}u^{53}w^{55}$ & $H^{-96}(uw)^{-53}w^7$\\
\bottomrule
\end{tabular}
\end{center}

The spectator exponent $R_x=70$ removes the first coordinate scale in the bounded-$U$ parabolic case. The first-shear and terminal rows sum in $W$ and $U$, respectively. In the mixed parabolic row the bound is independent of the small $u$, so there is no small-$u$ amplitude union. In the last row, sum first in $uw$; then $w\lesssim_PH$ gives $\sum_w w^7\lesssim_PH^7$, leaving $H^{-89}$. The remaining mixed rows have convergent displacement sums. The overlapping thresholds of Section~\partIref{reg:d-channels}, chosen after the finite certificates, therefore give $P^CH^{-89}(1+\log d_1^{-1})^C$ throughout the $d$ chart.

For $73<p<74$, choose $q=(1+\eta)(p-1)<89$. Proposition~\partIref{an:common-incidence} gives the claimed incidence bound, and~\partIeqref{an:physical-positivity} with $k=5$, $\gamma=0$ gives its physical summation.

The complete ordered row lists for the seven certificates are printed in the certificate supplement, Section~\ref{ext-supp:low:5:rd5:rows}. Together with~\partIeqref{low:5:rd5:matrix_entry}, they specify each minor without an unspecified pivot choice; the table above uses their absolute values.

\paragraph{Residual regular components for ternary quintics}
The residual charts use the same assignment in $V_5$.

\begin{theorem}[Quintic residual regular tail]\label{low:5:qrr:tail}

\noindent\textup{(i)} Every original index assigned to a residual regular component has a common preparation with count at most $P^C\Lambda H$ and cutoff $1\le H\le P^Cd_2^{-1}$.

\noindent\textup{(ii)} If $E_j(H)$ is the coefficient set that this label belongs to a prepared component of length $H$, then
\[
\operatorname{vol}_{55} E_j(H)\le P^C H^{-73}\bigl(1+\log P+\log d_1^{-1}\bigr)^C.
\]
The volume is measured in the original fixed-frame coefficient coordinates.
\end{theorem}

\paragraph{Step 1. Finite certificates and their uniform use}
Use the complete quintic basis above and the scaling rule~\partIeqref{low:56:weight}; in this table its output sums are denoted by $(r_x,r_y,r_z)$.

The nine minors used below are specified by the complete row sets in Section~\ref{ext-supp:supp:quintic-residual-rows} of the certificate supplement. They use integer sample times $t=0,\ldots,m-1$ in the models $M_f,P_f,Q,P_a,M_e,L_e,R$ of Section~\partIref{reg:sampling}; the certificate named $M_d$ uses $R$.

Here $C_f,C_e,C_d$ refer to different row sets of the same model $Q$. Their exact constants and output sums are
\[
\begin{array}{c|c|c|r}
\text{certificate}&m&(r_x,r_y,r_z)&|\det|\\\hline
M_f&4&(26,114,14)&5546527766851092480\\
P_f&4&(24,70,32)&10983370863150\\
C_f&4&(17,50,56)&330176000/6561\\
P_a&4&(24,32,70)&772886014477410984630681600\\
M_e&4&(26,14,114)&5546527766851092480\\
L_e&4&(24,24,84)&17433922005\\
C_e&3&(23,27,107)&487/1161683025219092623645605888\\
M_d&4&(16,24,103)&1150471189139816448000\\
C_d&4&(17,18,120)&310841/2821109907456.
\end{array}
\]
Exact elimination of~\partIeqref{reg:sample-entry} gives these rational determinants. The executable checks recompute all nine determinants and their output-exponent sums.

Lemma~\partIref{reg:lem:transfer} applies with the following grades:
\[
\begin{array}{c|c|c}
\text{certificate and condition}&\text{grade}&\text{entries lowering the grade}\\\hline
M_f&\deg_x+\deg_y&b,c,d,e\\
P_f,\ a_t=0&\deg_y&d,e\\
P_a&\deg_z&f\\
M_e,\ d_t=0&\deg_x+\deg_z&a,c,f
\end{array}
\]
The input and selected-row multiplicities agree: they are $(5,10,12,12,10,6)$ for the mixed grades and $(20,15,10,6,3,1)$ for a coordinate grade. Hence the determinants are independent of the listed entries at each sample. The comparisons of Section~\partIref{reg:gauges} then transfer these complete minors to the prepared curves, with sampling fractions chosen after their determinant margins.

\paragraph{Step 2. The chart omitting $T_z$}

Use the $T_z$-omitted chart in Section~\partIref{reg:gauges}, writing $\chi$ for its rate $\zeta$. Thus $u=H^2\sup|\delta|\lesssim_PH^2$, $v=H\sup|\chi|\lesssim_PH$ and $B_0=\max(H,u)$. The original box and net cardinality per additional sample are~\partIeqref{reg:f-box}.

If $\alpha$ is significant, the $M_f$ certificate and the grade check above give determinant at least $P^{-C}H^{100}$ under the comparison of Section~\partIref{reg:gauges}. If $u\le P^CH$, then $B_0\le P^CH$ and $N_f\le P^CH^5$, giving volume $O(P^CH^{-85})$. If $u\ge cP^{-C}H$, then $B_0\le P^Cu$ and $N_f\le P^CHu^4$, giving
\[
P^CH^{-100}(Hu^4)^3=P^CH^{-97}u^{12}\le P^CH^{-73}.
\]
The two ranges overlap only within a fixed multiplicative interval. Together they cover the significant-$\alpha$ branch without a smallness assumption on the other four affine coordinates.

Now assume $\chi$ is significant, so $v\asymp_PH$. For $u\lesssim_PH$, the $P_f$ comparison of Section~\partIref{reg:gauges} and the grade check above give determinant $P^{-C}H^{84}v^{46}$. Its original net is $N_f\le P^CH^5$, so this branch gives $H^{-84}v^{-46}(H^5)^3\lesssim_PH^{-115}$.

For $u\gtrsim_PH$ and $|\alpha|u\lesssim_Pv$, the ordered parabolic comparison gives the same determinant $H^{84}v^{46}$ and original fixed-$f$ net at most $Hv^2u^2$. Therefore
\[
H^{-84}v^{-46}(Hv^2u^2)^3=H^{-81}v^{-40}u^6\lesssim_PH^{-109}.
\]
In the remaining case $u\gtrsim_PH$ and $|\alpha|u\gg_Pv$, put $r=H\sup|\alpha|$. The cubic comparison of Section~\partIref{reg:gauges} gives $ru\gtrsim_PH^2$, $r\lesssim_PH$, and original fixed-$f$ net $r^3u^4/H^2$. The $C_f$ certificate has determinant $H^{14}r^{53}u^{73}$. Thus
\[
H^{-14}r^{-53}u^{-73}(r^3u^4/H^2)^3=H^{-20}r^{-44}u^{-61}\lesssim_PH^{-108}u^{-17}\lesssim_PH^{-125}.
\]
Use the overlapping thresholds of Section~\partIref{reg:gauges}. Dyadic sums over $u,v,r$ add at most a power of the physical logarithm, proving the asserted $H^{-73}$ tail for the entire $T_z$-omitted chart.

\paragraph{Step 3. The significant-$A$ part of the $T_y$ chart}

Use the significant-$A$ chart of Section~\partIref{reg:gauges}, with $u=H\sup|\delta-\beta'|$, $w=H\sup|\gamma|$, $U=\max(1,u)$ and $W=\max(1,w)$. Its original fixed-$a$ nets are $H^4U^2W$ for $UW\lesssim_PH$ and $H^2u^4w^3$ for $uw\gg_PH$.

If $UW\lesssim_PH$, the $P_a$ minor under $\operatorname{diag}(H^2,H,W)$ gives $H^{130}$. Use the grade check above and the original fixed-$a$ net $H^4U^2W$ from Section~\partIref{reg:gauges}. Four samples therefore give
\[
H^{-130}(H^4U^2W)^3=H^{-118}U^6W^3\lesssim_PH^{-112}W^{-3}.
\]
If $uw\gg_PH$, the scales $u,w$ are bounded below before taking maxima with one. The dilation $\operatorname{diag}(Huw,uw,w)$ gives the $C_e$ cubic determinant $H^{47}u^{90}w^{53}$. The fixed-$a$ net costs $H^2u^4w^3$, and three samples give
\[
H^{-47}u^{-90}w^{-53}(H^2u^4w^3)^2=H^{-43}u^{-82}w^{-47}=H^{-90}(uw/H)^{-47}u^{-35}.
\]
The comparisons of Section~\partIref{reg:gauges} cover the transition between these ranges. Both bounds are stronger than required for the significant-$A$ part.

\paragraph{Step 4. The central $C,T_y$ chart}

Use the central chart and its four overlapping cases from Section~\partIref{reg:gauges}, with the rates and scales there; write $M=\max(U,W)$. Its equations and original box are~\partIeqref{reg:central-equations}--\partIeqref{reg:central-box}. Each coefficient-volume bound below is multiplied by the cardinalities of the original-coordinate sampling nets established there.

If both products are bounded, $\operatorname{diag}(HM,H,M)$ gives the $M_e$ determinant $H^{100}$, with cancellation of the $M$ power. The grade check above and the central comparison of Section~\partIref{reg:gauges} apply. The original net is $N_e\lesssim_PH^2UV_0^2WM$. Thus
\[
H^{-100}N_e^3\lesssim_PH^{-94}U^3V_0^6W^3M^3\le P^CH^{-94}(V_0M)^6M^3\lesssim_PH^{-85}.
\]
This includes the possible cancellation between $\alpha$ and $\beta\gamma$.

If $UV_0/H$ is large and $V_0W/H$ bounded, the central comparison with model $A+D+T_y$ gives the $L_e$ determinant $H^{78}U^{46}V_0^{14}$ and original fixed-$d$ net $H^3U^3$. Consequently
\[
H^{-78}U^{-46}V_0^{-14}(H^3U^3)^3=H^{-69}U^{-37}V_0^{-14}\lesssim_PH^{-83}U^{-23}.
\]
The additional $T_y$ term is part of this exact model field and has not been discarded.

If both products are large, the central cubic comparison gives the $C_e$ determinant $U^{47}V_0^{90}W^{53}$ and original fixed-$d$ net $U^3V_0^3W^3$. Three samples give
\[
U^{-41}V_0^{-84}W^{-47}\lesssim_PH^{-88}V_0^4\lesssim_PH^{-84}.
\]
Here we used $U,W\gtrsim_PH/V_0$ and $V_0\le P^CH$.

When $V_0W/H$ is large and $UV_0/H$ is bounded, apply the complete comparison after coordinate interchange~\partIeqref{reg:swapped-field}--\partIeqref{reg:swapped-error}, with its $u_*$ and $X=\max(1,u_*W)$. It controls the lower-shear and diagonal errors on locally rescaled fractions and prescribes original $b$ levels. The original net is $N_b\lesssim_PHU^2V_0^2W^3$.

If $u_*W$ is bounded, the $M_d$ determinant is $H^{100}W^{13}$; if $u_*W$ is large, the $C_d$ determinant is $H^{105}u_*^{53}W^{55}$. Use the ordered comparison fractions from Section~\partIref{reg:gauges} and the same original net $N_b$. Four samples give respectively
\[
H^{-100}W^{-13}N_b^3\lesssim_PH^{-97}(UV_0)^6W^{-4}\lesssim_PH^{-91},
\]
\[
H^{-105}u_*^{-53}W^{-55}N_b^3\lesssim_PH^{-102}(UV_0)^6(u_*W)^{-53}W^7\lesssim_PH^{-89}(u_*W)^{-53}.
\]
The pointwise Gram comparison transfers these bounds through the fixed coordinate permutation, the common determinant-one anchor change of frame, and separate bounded right frame changes. The union still runs over the original upper-affine boxes. 

\begin{proof}[Completion of Theorem~\ref{low:5:qrr:tail}]
\par\smallskip\noindent\textit{Step 5. Common assignment and summation.}\par\nobreak\smallskip

The common allocation first selects the small-$D$ $T_z$-omitted chart, or the significant-$A$ part of the $T_y$ chart, or its central $C,T_y$ part. The preceding cases exhaust each assigned chart. Their weakest power is $H^{-73}$, in the enlarged-box significant-$\alpha$ branch of the $T_z$ chart. All displacement dyadic sums cost at most physical logarithms; convergent sums may only improve that bound. The common section curve assignment and preparation give count $P^C\Lambda H$, a logarithmic number of pieces per scale, and cutoff $H\le P^Cd_2^{-1}$. Bounded frame changes give the pointwise Gram comparisons. The anchored Gram union bound uses the sampling boxes in the original catalogue-anchor coordinates. Each bound includes the stated sampling-net cardinalities. Summing over the finite row and chart choices and the permitted dyadic parameter ranges gives the stated estimate.
\end{proof}

\paragraph{Regular components for sextics}

Use $V_6$ with the common augmented admissibility condition~\partIeqref{reg:stop}.

\begin{theorem}[Sextic regular count]\label{low:6:s6r:theorem}
Consider the part of the admissibility condition where the fifth singular value of
\[
R=(A\widehat J,C\widehat J,T_x\widehat J,D\widehat J,T_y\widehat J,T_z\widehat J)
\]
is bounded below by $P^{-C}$.

\noindent\textup{(i)} There is a coefficient-measurable assignment to pieces with bounded counts and prepared one-dimensional pieces of dyadic length $H\ge1$, with count at most $P^C\Lambda H$.

\noindent\textup{(ii)} For every fixed assigned index, its coefficient set at length $H$ satisfies
\begin{equation}\label{low:6:s6r:tail}
\operatorname{vol}_{83}E_j(H)\le P^C H^{-148}(1+\log P+\log d_1^{-1})^C.
\end{equation}
The good chart omitting $D$ has the stronger exponent $164$. The estimates use all $83$ original coefficients and the original-coordinate sampling boxes.

\noindent\textup{(iii)} Consequently, for $129<p<130$ and sufficiently small $\eta>0$, the regular incidence sum converges over all dyadic $H\ge1$; no upper cutoff on $H$ is needed for this sum.
\end{theorem}

\par\smallskip\noindent\textit{Step 1. Complete minors and grading multiplicities.}\par\nobreak\smallskip

Apply Lemma~\partIref{reg:regular-preparation} and the chart comparisons of Section~\partIref{reg:gauges}. All counts below use~\partIeqref{reg:fixed-net-bound}, with the actual anchor supplying the first sample. The theorem concerns only the region with quantitative fifth singular value; the surface normalization above is used separately in the surface argument.

\paragraph{Sixteen sampling minors}\label{low:6:s6r:certificates}

Use~\partIeqref{low:5:rd5:matrix_entry} with all input and output monomials of degrees one through six, ordered by total degree and then lexicographically. Each certificate has $83$ rows, listed in supplementary Section~\ref{ext-supp:low:6:s6r:rows}, at sample times $t=0,\ldots,m-1$. The lower-case states $l,q,p,m,r,m^f,p^f,p^a,m^e,l^e$ denote the corresponding models $L,Q,P,M,R,M_f,P_f,P_a,M_e,L_e$ of Section~\partIref{reg:sampling}.

If $(R_x,R_y,R_z)$ is the sum of selected output exponents, the identities are as follows. Only absolute determinant values are needed.
\begin{center}
\begin{tabular}{ccll}
Certificate & $m$ & $(R_x,R_y,R_z)$ & $|\det|$\\\hline
$L_4$ & $4$ & $(10,118,250)$ & $3^{40}/2^{36}$\\
$Q_3$ & $3$ & $(27,73,192)$ & $11/(2^{29}3^{82})$\\
$P_4$ & $4$ & $(126,46,55)$ & $2^3 3^{38}5^2\cdot7\cdot11$\\
$M_5$ & $5$ & $(17,48,204)$ & $2^{93}3^{39}5$\\
$L_5$ & $5$ & $(46,30,302)$ & $2^{43}3^{48}$\\
$R_4$ & $4$ & $(31,46,181)$ & $2^{46}3^{38}5^4$\\
$Q_4$ & $4$ & $(35,35,205)$ & $7\cdot53\cdot101/(2^{33}3^{24})$\\
$M_f$ & $4$ & $(50,202,27)$ & $2^{60}3^{26}5$\\
$P_f$ & $4$ & $(46,126,55)$ & $2^3 3^{38}5^2\cdot7\cdot11$\\
$C_f$ & $4$ & $(34,89,96)$ & $2^{15}5^3\cdot7\cdot11\cdot13/3^{24}$\\
$P_a$ & $4$ & $(46,55,126)$ & $2^{83}3^{38}5^2\cdot7\cdot11$\\
$M_e$ & $4$ & $(50,27,202)$ & $2^{60}3^{26}5$\\
$L_e$ & $4$ & $(46,46,143)$ & $2\cdot3^{41}$\\
$C_e$ & $3$ & $(46,49,183)$ & $11\cdot23\cdot241/(2^{37}3^{75})$\\
$M_d$ & $4$ & $(30,47,181)$ & $2^{46}3^{37}5^4$\\
$C_d$ & $4$ & $(35,35,205)$ & $7\cdot53\cdot101/(2^{33}3^{24})$
\end{tabular}
\end{center}
The first seven use $l,q,p,m,l,r,q$, respectively. The next nine use $m^f,p^f,q,p^a,m^e,l^e,q,r,q$, respectively. The $R_4$ and $M_d$ row sets are different, and their terminal-scale exponents differ by one. They are not interchanged below.

Their scaling factor is~\partIeqref{low:56:weight} in degree six.

For $L_4,L_5$, matching total-degree multiplicities $(3,6,10,15,21,28)$ and Lemma~\partIref{reg:lem:transfer} give independence of separate sample translations. For $M_f,P_f,P_a,M_e$, use the grades and coordinate restrictions in the quintic table above. The sextic selected rows have the input multiplicities $(6,12,15,16,15,12,7)$ for the mixed grades and $(27,21,15,10,6,3,1)$ for a coordinate grade, so the same triangular argument applies.

\paragraph{Step 2. The good chart omitting \texorpdfstring{$D$}{D}}\label{low:6:s6r:dchart}

Use the cases, exact rates, scales and comparisons of Section~\partIref{reg:d-channels}, with the original net $N_d\lesssim_PH^3U^3W^3$ from \partIeqref{reg:d-box}. The determinant certificates and \partIeqref{reg:fixed-net-bound} give, up to powers of $P$,
\begin{center}\small
\begin{tabular}{llll}
\toprule
Case & Certificate & Determinant & Coefficient volume\\
\midrule
First shear, bounded $W$ & $L_4$ & $H^{240}$ & $H^{-222}$\\
First shear, large $W$ & $Q_3$ & $H^{251}W^{86}$ & $H^{-239}W^{-80}$\\
Terminal, bounded $U$ & $P_4$ & $H^{231}$ & $H^{-213}$\\
Terminal, large $U$ & $Q_3$ & $H^{152}U^{99}W^{86}$ & $H^{-226}U^{-93}$\\
Mixed, bounded $u,w$ & $M_5$ & $H^{187}$ & $H^{-175}$\\
Mixed, large $u$, bounded $w$ & $L_5$ & $H^{176}U^{80}$ & $H^{-164}U^{-68}$\\
Mixed, large $w$, bounded $uw$ & $R_4$ & $H^{175}w^{25}$ & $H^{-166}w^{-16}$\\
Mixed, large $w,uw$, $u\ge1$ & $Q_4$ & $H^{182}u^{91}w^{103}$ & $H^{-173}u^{-82}w^{-94}$\\
Mixed, large $w,uw$, $u<1$ & $Q_4$ & $H^{182}u^{91}w^{103}$ & $H^{-173}(uw)^{-91}w^{-3}$\\
\bottomrule
\end{tabular}
\end{center}

The scale of the unchanged $x$ coordinate in $P_4$ has exponent $\chi_6-R_x=126-126=0$, and its frozen $x$-$y$ shear is bounded. Relative preparation and the fixed comparison fractions control the normalized perturbations. The original net and cubic comparison remain valid for $u<1$. In the last row sum $uw$ first; both remaining powers are favorable. The bounded-product row has no small-$u$ amplitude union. With the overlapping thresholds of Section~\partIref{reg:d-channels}, these rows give the $d$-chart exponent $164$.

\paragraph{Step 3. The residual chart omitting \texorpdfstring{$T_z$}{Tz}}\label{low:6:s6r:fchart}

Use the $T_z$-omitted chart of Section~\partIref{reg:gauges}, writing $\chi=\zeta$. Retain $u=H^2\sup|\delta|\lesssim_PH^2$, $v=H\sup|\chi|\lesssim_PH$, $B_0=\max(H,u)$, and the original box and net \partIeqref{reg:f-box}. For significant $\alpha$, use $M_f$ and the grade check above, including the range $H\lesssim_Pu\lesssim_PH^2$ in Table~\ref{low:6:regular-products}.

For significant $\chi$, one has $v\asymp_PH$. Use $P_f$ when $u\lesssim_PH$, or when $u\gtrsim_PH$ and $|\alpha|u\lesssim_Pv$. Otherwise use $C_f$, with $r=H\sup|\alpha|$, $ru\gtrsim_PH^2$ and $r\lesssim_PH$. The comparisons, original nets and overlapping thresholds are those of Section~\partIref{reg:gauges}; the sextic determinant and net products are listed in Table~\ref{low:6:regular-products}.

\paragraph{The significant first-shear part of the \texorpdfstring{$T_y$}{Ty} chart}

Use the significant-$A$ chart and scales $u,w,U,W$ of Section~\partIref{reg:gauges}. The certificate is $P_a$ for $UW\lesssim_PH$ and $C_e$ for $uw\gg_PH$. The grade check above, original nets and overlapping comparisons give the corresponding rows of Table~\ref{low:6:regular-products}.

\paragraph{The central \texorpdfstring{$C,T_y$}{C,Ty} chart and coordinate interchange}\label{low:6:s6r:central}

Use \partIeqref{reg:central-equations}--\partIeqref{reg:central-box} and the four overlapping cases of Section~\partIref{reg:gauges}. Retain $U,V_0,W$ and put $M=\max(U,W)$. The bounded-product case uses $M_e$, with cancellation of the $M$ power and the grade check above. The case $UV_0/H$ large and $V_0W/H$ bounded uses $L_e$ with the complete model $A+D+T_y$. When both products are large, use $C_e$. Their original sampling nets and sextic determinants are retained in Table~\ref{low:6:regular-products}.

For $V_0W/H$ large and $UV_0/H$ bounded, use the complete comparison \partIeqref{reg:swapped-field}--\partIeqref{reg:swapped-error}, including its original $b$-derivative lower bound, lower-shear and diagonal terms, and ordered comparison fractions. Retain its exact $u_*$ and $X=\max(1,u_*W)$. Bounded $u_*W$ uses $M_d$; large $u_*W$ uses $C_d$. The original fixed-$b$ net is $HU^2V_0^2W^3$, as recorded in the same table.

\paragraph{Regular summation}

The determinant and net products for these residual cases are collected below. The sample count $m$ is that of the certificate table above; each bound is $P^C\Delta^{-1}N^{m-1}$ after the indicated scale restrictions.
{\small
\begin{longtable}{lcccc}
\caption{Sextic residual sampling products. The last column is the uniform exponent in $H^{-b}$.}\label{low:6:regular-products}\\
Case & Certificate & $\Delta$ & $N$ & $b$\\\hline
\endfirsthead
Case & Certificate & $\Delta$ & $N$ & $b$\\\hline
\endhead
$f:\alpha,\ u\lesssim_PH$ & $M_f$ & $H^{175}$ & $H^5$ & $160$\\
$f:\alpha,\ H\lesssim_Pu\lesssim_PH^2$ & $M_f$ & $H^{175}$ & $Hu^4$ & $148$\\
$f:\chi,\ u\lesssim_PH$ & $P_f$ & $H^{151}v^{80}$ & $H^5$ & $216$\\
$f:\chi$, ordered parabolic & $P_f$ & $H^{151}v^{80}$ & $Hv^2u^2$ & $210$\\
$f:\chi$, ordered cubic & $C_f$ & $H^{30}r^{92}u^{129}$ & $r^3u^4/H^2$ & $236$\\
$A:\ UW\lesssim_PH$ & $P_a$ & $H^{231}$ & $H^4U^2W$ & $213$\\
$A:\ uw\gtrsim_PH$ & $C_e$ & $H^{80}u^{157}w^{100}$ & $H^2u^4w^3$ & $170$\\
Central, both products bounded & $M_e$ & $H^{175}$ & $H^2UV_0^2WM$ & $160$\\
Central, first product large & $L_e$ & $H^{143}U^{80}V_0^{17}$ & $H^3U^3$ & $151$\\
Central, both products large & $C_e$ & $U^{80}V_0^{157}W^{100}$ & $U^3V_0^3W^3$ & $151$\\
Swapped, bounded $u_*W$ & $M_d$ & $H^{175}W^{24}$ & $HU^2V_0^2W^3$ & $166$\\
Swapped, large $u_*W$ & $C_d$ & $H^{182}u_*^{91}W^{103}$ & $HU^2V_0^2W^3$ & $173$\\
\end{longtable}
}

\par\smallskip\noindent\textit{Step 4. Summation of the residual products.}\par\nobreak\smallskip

For the cubic $f$ row, $ru\gtrsim_PH^2$ and $u\gtrsim_PH$ give $H^{-36}r^{-83}u^{-117}\lesssim_PH^{-236}$. In the two $A$ rows use $UW\lesssim_PH$ and $uw\gtrsim_PH$, respectively; the latter product is $H^{-170}(uw/H)^{-94}u^{-55}$. For bounded central products use $U,W\le M$, $V_0M\lesssim_PH$ and $M\lesssim_PH$. For the first-large row use $UV_0\gtrsim_PH$; for the both-large row use $U,W\gtrsim_PH/V_0$ and $V_0\lesssim_PH$, giving $H^{-168}V_0^{17}\lesssim_PH^{-151}$. In the swapped rows use $UV_0\lesssim_PH$ and retain the factors $W^{-15}$ and $(u_*W)^{-91}W^{-3}$, respectively. These inequalities include all intermediate ratios. All original scales lie between fixed powers of the physical cutoff, so their dyadic sums cost only the stated physical logarithms.

The allocation exhausts the regular region. The $d$-chart exponent is $164$, and the residual minimum is $148$ from the large-$u$ first-shear row of Table~\ref{low:6:regular-products}, proving~\eqref{low:6:s6r:tail}. For $129<p<130$, choose $q=(1+\eta)(p-1)<148$ and apply Proposition~\partIref{an:common-incidence}; the physical summation is~\partIeqref{an:physical-positivity} with $k=6$, $\gamma=0$.

\subsection{Sampling bounds in degrees seven through thirteen}\label{reg:seven-thirteen}\label{low:78:h3r:section}\label{low:9:root9:regular}\label{low:ten-thirteen}

\paragraph{Regular coefficient counts in degrees seven and eight}

Lemma~\partIref{reg:regular-preparation} and the equations for the frame change of Section~\partIref{reg:geometric-regimes} supply the regular assignment by the fifth singular value and original-coordinate sampling boxes. We construct complete certificates in dimensions $119$ and $164$, verify every scale region, and give a degree-uniform mixed-core formula. The assignment where the fifth singular value is small is treated separately.

\paragraph{Finite regular estimates}

Use the upper-affine group~\partIeqref{geo:group}, its generators, and the augmented translation nondegeneracy condition. 

The parameter in the chart omitting $T_y$ is $de-d\,df$. In particular, the residual $T_z$ chart retains $u=H^2\sup|\delta|$ up to size $O_P(H^2)$. The swapped central field remains exactly
\begin{equation}\label{low:78:h3r:swap}
\widetilde V=\beta D+(r'-r^2\alpha)A+T_x+\gamma T_z+\alpha x\partial_y+r\alpha(-x\partial_x+y\partial_y),\qquad r=\delta/\beta.
\end{equation}
The lower shear and both diagonal terms are included in the normalized matrix comparison.

\begin{proposition}\label{low:78:h3r:regular-count}
On a bounded conditioned coefficient window satisfying augmented translation nondegeneracy, restrict to the region where the fifth singular value of the augmented upper-affine orbit is at least $P^{-C}$.

\noindent\textup{(i)} The regular assignment has pieces of dyadic length $H\ge1$ with count $O_P(\Lambda H)$.

\noindent\textup{(ii)} At each fixed original assigned index, their coefficient sets satisfy
\begin{equation}\label{low:78:h3r:regular-tails}
\begin{aligned}
\operatorname{vol}_{119}E_j(H)&\lesssim P^C H^{-243}(1+\log P+\log d_1^{-1})^C &&(k=7),\\
\operatorname{vol}_{164}E_j(H)&\lesssim P^C H^{-382}(1+\log P+\log d_1^{-1})^C &&(k=8).
\end{aligned}
\end{equation}

\noindent\textup{(iii)} Consequently this part of the incidence sum converges for $p$ in an interval immediately above $213$ when $k=7$, and immediately above $333$ when $k=8$.
\end{proposition}

\par\smallskip
\noindent\textit{Proof stages for Proposition~\ref{low:78:h3r:regular-count}.}\par\nobreak
\noindent\textup{Step 1.} Transfer of the degree-independent geometry.\par\nobreak
\noindent\textup{Step 2.} Exact minors and the vertex procedure.\par\nobreak
\noindent\textup{Step 3.} Original-scale sampling costs and exact tables.\par\nobreak
\noindent\textup{Step 4.} Coefficient-volume and incidence assembly.\par\nobreak
\smallskip

\paragraph{Transfer of the geometric allocation}\label{low:78:h3r:geometry}

Apply Lemma~\partIref{reg:regular-preparation} and the frame changes of Section~\partIref{reg:geometric-regimes} with $k=7,8$. They retain the normal columns $A,C,T_x$, the six group coordinates, transverse rank five, exact parameters and all accumulated small-rate terms. The complete swapped equation is~\eqref{low:78:h3r:swap}. Choose its comparison fractions and significant-rate margins after the degree-specific determinant bounds. The remaining inputs are the complete minors and sampling net cardinalities below.

\paragraph{Finite sampling certificates}\label{low:78:h3r:certificates}\label{reg:finite-verification}

For degrees seven through thirteen, construct the complete matrix~\partIeqref{reg:sample-entry} from the models of Section~\partIref{reg:sampling}. Order input monomials by total degree $n$, then by $i=0,\ldots,n$ and $j=0,\ldots,n-i$, with exponent $(i,j,n-i-j)$. Each exact row list specifies $N_k$ distinct rows, its sample count $m$, and its output sums $R=(R_x,R_y,R_z)$. The required grading multiplicities are checked against the complete input basis; Lemma~\partIref{reg:lem:transfer} then shows that the corresponding off-diagonal entries do not change the determinant.

All sample denominators have prime factors only $2,3$. Exact elimination modulo the prime $1000033$ gives the recorded nonzero residues, so the rational minors are nonzero. The row lists and direct multinomial implementation are specified in supplementary Sections~\ref{ext-supp:low:78:h3r:row-appendix} and~\ref{ext-supp:low:verification-method}--\ref{ext-supp:low:row-masks}.

For diagonal scales $s_i=H^{a_i(\theta)}$ and original free-net bound $P^CH^{n(\theta)}$, equations~\partIeqref{reg:minor-scaling} and~\partIeqref{reg:fixed-net-bound} give the complete coefficient exponent
\begin{equation}\label{low:finite-tail}
b_k(\theta)=\sum_{i=1}^3(\chi_k-R_i)a_i(\theta)-(m-1)n(\theta).
\end{equation}
\equationalias{low:9:root9:regular-tail}
Each scale domain in Table~\partIref{reg:all-regimes} is a nonempty bounded rational polytope. Enumerate its vertices by solving all full-size boundary systems over $\mathbb Q$, discarding singular systems and points violating a defining inequality. Minimize the affine exponent $b_k$ at these vertices. The saved data retain every vertex and value.

For $k=7,8$, the sample counts are in Table~\ref{low:78:h3r:domains} and the complete rows are specified in supplementary Section~\ref{ext-supp:low:78:h3r:row-appendix}.

\begin{proposition}[Exact complete regular minors]\label{low:78:h3r:exact-certificates}

\noindent\textup{(i)} For every one of the $22$ cases in Table~\ref{low:78:h3r:domains} and each $k\in\{7,8\}$, the rows specified in the certificate supplement, Section~\ref{ext-supp:low:78:h3r:row-appendix}, form a complete $N_k$-row minor of \partIeqref{low:78:h3r:entry}. Their output exponent sums and nonzero determinant residues modulo $1000033$, up to row-order sign, are exactly those of Table~\ref{low:78:h3r:residues}. They have the grading multiplicities specified below.

\noindent\textup{(ii)} With the diagonal scales and net cardinalities per additional sample in Table~\partIref{reg:all-costs}, the minima of the resulting exact affine exponents are recorded in Table~\ref{low:78:h3r:ledger}; their smallest values are $243$ and $382$.
\end{proposition}

The exact procedure above gives the output sums and nonzero residues in Table~\ref{low:78:h3r:residues}; supplementary Section~\ref{ext-supp:supp:seven-eight:verification} identifies the saved-row checks.

The independence assertions retain their original grades. For $L$, the selected row multiplicities agree with the input multiplicities in total degree, so arbitrary sample translations lie off the diagonal blocks. For $M_f$, $P_f$, $P_a$, and $M_e$, the corresponding grades are $\nu_x+\nu_y$, $\nu_y$, $\nu_z$, and $\nu_x+\nu_z$. In each case the saved row multiplicities exactly equal the input multiplicities. Consequently the same square diagonal blocks remove the same additional coordinates as in the sextic proof. Perturbative comparisons are made only after the original normalization has bounded every remaining coordinate.

\paragraph{Counts at fixed frames in each scale range}\label{low:78:h3r:scales}

Use the anchored bound~\partIeqref{reg:fixed-net-bound}, with original-coordinate nets for the $m-1$ free samples.

The cases, parameters and logarithmic regions are those of Table~\partIref{reg:all-regimes}, and their diagonal exponents $a=(a_x,a_y,a_z)$ and original covering exponents $n$ are in Table~\partIref{reg:all-costs}. The finite-degree sample counts are listed below. For the last swapped row, $T=\max(1,u_*W)$ and $N\lesssim_PH^3W^3$ follow from $HU^2V_0^2W^3$ and $UV_0\lesssim_PH$; enlarging to $1\le T,W\lesssim_PH$ weakens the count.
\begin{table}[ht]
\centering
\caption{Finite sample counts for the cases of Table~\partIref{reg:all-regimes}.}\label{low:78:h3r:domains}
\begin{tabular}{rc@{\qquad}rc}
\toprule
ID & Model/$m$ & ID & Model/$m$\\
\midrule
1 & $L/4$ & 12 & $P_f/4$\\
2 & $Q/3$ & 13 & $P_f/4$\\
3 & $P/4$ & 14 & $Q/4$\\
4 & $Q/3$ & 15 & $P_a/4$\\
5 & $M/5$ & 16 & $Q/3$\\
6 & $L/5$ & 17 & $M_e/4$\\
7 & $R/4$ & 18 & $M_e/4$\\
8 & $Q/4$ & 19 & $L_e/4$\\
9 & $Q/4$ & 20 & $Q/3$\\
10 & $M_f/4$ & 21 & $R/4$\\
11 & $M_f/4$ & 22 & $Q/4$\\
\bottomrule
\end{tabular}
\end{table}


{\small
\begin{longtable}{r c r c r}
\caption{Output exponent sums and nonzero determinant residues, up to sign, modulo $1000033$. The residues certify the fixed sample minors before diagonal conjugation.}\label{low:78:h3r:residues}\\
ID & $R$ for $k=7$ & Residue & $R$ for $k=8$ & Residue\\\hline
\endfirsthead
ID & $R$ for $k=7$ & Residue & $R$ for $k=8$ & Residue\\\hline
\endhead
1 & $(20,200,410)$ & 923056 & $(36,318,636)$ & 901154\\
2 & $(51,123,305)$ & 201270 & $(82,201,472)$ & 133186\\
3 & $(210,80,91)$ & 638780 & $(330,130,142)$ & 998542\\
4 & $(81,116,228)$ & 62019 & $(135,180,356)$ & 833545\\
5 & $(32,84,336)$ & 344894 & $(51,136,524)$ & 660493\\
6 & $(80,54,496)$ & 855435 & $(130,90,770)$ & 286685\\
7 & $(82,80,220)$ & 470786 & $(129,130,346)$ & 229701\\
8 & $(64,62,318)$ & 787571 & $(95,109,511)$ & 252651\\
9 & $(97,52,255)$ & 729946 & $(168,86,375)$ & 42270\\
10 & $(86,334,44)$ & 383432 & $(141,519,75)$ & 20247\\
11 & $(86,334,44)$ & 383432 & $(141,519,75)$ & 20247\\
12 & $(80,210,91)$ & 361253 & $(130,330,142)$ & 1491\\
13 & $(80,210,91)$ & 361253 & $(130,330,142)$ & 1491\\
14 & $(78,117,166)$ & 386368 & $(130,184,256)$ & 819200\\
15 & $(80,91,210)$ & 430444 & $(130,142,330)$ & 990511\\
16 & $(45,148,273)$ & 945236 & $(75,236,424)$ & 683387\\
17 & $(86,44,334)$ & 616601 & $(141,75,519)$ & 20247\\
18 & $(86,44,334)$ & 616601 & $(141,75,519)$ & 20247\\
19 & $(80,88,213)$ & 373054 & $(130,139,334)$ & 780344\\
20 & $(69,93,328)$ & 397179 & $(112,150,511)$ & 116464\\
21 & $(82,80,220)$ & 470786 & $(129,130,346)$ & 229701\\
22 & $(81,59,274)$ & 767455 & $(132,95,430)$ & 488863\\
\end{longtable}
}

Applying~\eqref{low:finite-tail} and the rational-vertex procedure above gives Table~\ref{low:78:h3r:ledger}.

\begin{longtable}{r rr@{\qquad}r rr}
\caption{Minimum tail exponents on the scale regions of Table~\partIref{reg:all-regimes}, with the finite sample counts above.}\label{low:78:h3r:ledger}\\
\toprule
ID & $k=7$ & $k=8$ & ID & $k=7$ & $k=8$\\\midrule
1 & 372 & 582 & 12 & 364 & 573\\
2 & 393 & 613 & 13 & 358 & 567\\
3 & 361 & 570 & 14 & 386 & 605\\
4 & 416 & 652 & 15 & 361 & 570\\
5 & 292 & 461 & 16 & 319 & 500\\
6 & 274 & 428 & 17 & 275 & 429\\
7 & 249 & 392 & 18 & 275 & 429\\
8 & 285 & 447 & 19 & 243 & 382\\
9 & 262 & 397 & 20 & 252 & 392\\
10 & 275 & 429 & 21 & 249 & 392\\
11 & 263 & 417 & 22 & 271 & 424\\
\bottomrule
\end{longtable}

\begin{proof}[Completion of the proof of Proposition~\ref{low:78:h3r:regular-count}]
The smallest entries of Table~\ref{low:78:h3r:ledger} are $243$ and $382$. The geometric priority allocation exhausts the regular region defined by the fifth singular value by the degree-stable argument in Section~\ref{low:78:h3r:geometry}; its bounded pieces have bounded count. Every remaining logarithmic scale lies between fixed powers of the physical cutoff, as in the original preparation. Summing those scales costs only the logarithms displayed in \eqref{low:78:h3r:regular-tails}, and favorable powers can only improve the estimate. This proves the two coefficient-volume bounds.

The regular incidence conclusion follows from Proposition~\partIref{an:common-incidence}, choosing $p>\chi_k+3$ and $\eta>0$ with $q=(1+\eta)(p-1)<b_k$, where $b_7=243$ and $b_8=382$.
\end{proof}


\paragraph{Sampling bounds in degree nine}
\begin{lemma}[Block-triangular degree lift]
For fixed affine sample states, let $S_{k-1}$ be a sampling minor on $V_{k-1}$ and let $H_k$ be a sampling minor on homogeneous degree-$k$ polynomials, using degree-$k$ output rows of the linear parts of the same samples. Their combined full minor on $V_k$ equals $\det S_{k-1}\det H_k$.
\end{lemma}
\begin{proof}
Order the columns by $V_{k-1}$ and homogeneous degree $k$, and retain the old rows before the degree-$k$ rows. Affine substitution cannot increase degree, so the full matrix has the exact form
\[
\begin{pmatrix}S_{k-1}&B_*\\0&H_k\end{pmatrix}.
\]
The lower-right block is independent of translations.
\end{proof}

For $k=9$, append $55$ homogeneous rows to the saved $164$ rows in each degree-eight case. Formula~\eqref{low:9:root9:regular-tail}, with $\chi_9=495$, gives the exponent $b(\theta)=b_9(\theta)$. Adding rows at existing sample states does not increase the net cardinality. The same $22$ scale domains give the following exact minima.
\begin{center}
\begin{tabular}{c|rrrrrrrrrrr}
Case &1&2&3&4&5&6&7&8&9&10&11\\\hline
Minimum &867&867&711&811&611&637&533&637&576&570&558\\[3pt]
Regime &12&13&14&15&16&17&18&19&20&21&22\\\hline
Minimum &714&708&746&711&626&570&570&526&518&533&614
\end{tabular}
\end{center}

Every degree-nine certificate has $219$ distinct rows. The common exact procedure recomputes their nonzero residues, output sums, and vertex minima. The checked grades are total degree for $L$, $x+y$ for $M_f$, $y$ for $P_f$, $z$ for $P_a$, and $x+z$ for $M_e$. The complete rows, scale inequalities, and residues are in \texttt{regular\_degree\_9\_lifted\_certificates.json} and its verification file.

For the weakest case, $R=(182,251,835)$ and $b=307\theta_1+551\theta_2+211\theta_3$ on its prescribed central-cubic polytope. Its minimum is $518$, attained at $(1,0,1)$; its determinant residue is $531495$. Thus the regular coefficient set has tail $P^CH^{-518}$, with the common logarithmic multiplicity already described. 

\paragraph{Degrees ten through thirteen}
Use the twenty-two cases of Tables~\partIref{reg:all-regimes} and~\ref{low:78:h3r:domains}, with the affine equations and original-coordinate sampling boxes proved in Paper I, Sections~\partIref{reg:geometric-regimes} and~\partIref{reg:sampling}. Only the coefficient dimension changes; the six group variables, parameters and dimensions of the sample nets remain fixed. Since extending a lower-degree row set can leave insufficient weighted capacity, we select new complete minors in each degree.

\begin{proposition}[The four remaining finite regular estimates]\label{low:regular-ten-thirteen}
For $k=10,11,12,13$, let $E_j^{\mathrm{reg}}(H)$ denote the normalized coefficient set that an original fixed index is assigned to a regular piece of length comparable to $H\ge1$ in the regular assignment.

\noindent\textup{(i)} Then
\[
\vol E_j^{\mathrm{reg}}(H)\le P^{C_k}(1+\log P+\log d_1^{-1})^{C_k}H^{-b_k},
\qquad (b_{10},b_{11},b_{12},b_{13})=(819,1139,1544,2046).
\]

\noindent\textup{(ii)} In each degree, $b_k>\chi_k+2$; the measure is the full $N_k$-dimensional coefficient measure.
\end{proposition}

\begin{proof}
Apply the common finite verification procedure to the $N_k$ rows specified in supplementary Section~\ref{ext-supp:low:row-masks}, with the models and sample counts of Table~\ref{low:78:h3r:domains}. Lemma~\partIref{reg:regular-preparation} supplies the anchored Gram comparison, and Sections~\ref{low:6:s6r:section} and~\ref{low:78:h3r:geometry} supply the prepared samples and original slice nets.

The elimination of the additional coordinates is also exact. For model $L$, order the inputs and selected rows by total degree. For $M_f,P_f,P_a,M_e$, use respectively the grades $i+j,j,\ell,i+\ell$. In each case the fixed saved row set has the same grade multiplicities as the complete input basis. Terms containing the corresponding additional variables lower that grade. The square diagonal blocks therefore give the same determinant for arbitrary values of the additional coordinates. The remaining normalized parameters lie in the compact comparison windows supplied by the geometric construction. Continuity of these finitely many nonzero complete minors then permits the comparison fractions and dominance thresholds to be chosen after the minors, with constants depending on the fixed degree.

The rational-vertex calculation for~\eqref{low:finite-tail} gives the following minima on the scale polytopes of Table~\partIref{reg:all-regimes}.
\begin{center}\small
\begin{tabular}{r|rrrr@{\qquad}r|rrrr}
\toprule
Case&10&11&12&13&Case&10&11&12&13\\\midrule
1&1242&1725&2334&3090&12&1230&1711&2319&3073\\
2&1303&1809&2450&3239&13&1224&1705&2313&3067\\
3&1227&1708&2316&3070&14&1299&1807&2446&3245\\
4&1391&1932&2614&3465&15&1227&1708&2316&3070\\
5&988&1373&1857&2463&16&1067&1484&2009&2658\\
6&912&1266&1713&2267&17&922&1290&1740&2311\\
7&834&1154&1565&2067&18&922&1290&1740&2311\\
8&951&1319&1786&2364&19&819&1139&1544&2046\\
9&853&1186&1602&2129&20&832&1154&1558&2064\\
10&922&1290&1740&2311&21&834&1154&1565&2067\\
11&910&1278&1728&2299&22&907&1256&1700&2253\\\bottomrule
\end{tabular}
\end{center}

Every saved complete minor has nonzero residue modulo $1000033$ under the common exact procedure. The finite normalized comparisons and fixed-degree algebraic preparation give polynomial dependence on the coefficient windows, independently of the physical side ratios. Taking the smallest table entry proves the assertion.
\end{proof}

\subsection{The combined regular bounds}\label{reg:combined}\label{reg:finite-guide}

The formal exponent $\beta_k=\frac{1105}{1026}\chi_k-2k-153$ used in Paper I, Theorem~\partIref{reg:regular-tail}, satisfies
\[
 \beta_k-(\chi_k+2)=\frac{79}{1026}\chi_k-2k-155.
\]
This formal margin is $-20963/513$ at $k=13$ and $131/513$ at $k=14$. The uniform calculation therefore does not supply the required regular bound at degree thirteen. The finite estimates above use the same coefficient-volume argument with their own rows and sample-net cardinalities; degree four also uses the physical cutoff. Their weakest regular exponents are
\[
\begin{array}{c|rrrrrrrrrr}
k&4&5&6&7&8&9&10&11&12&13\\\hline
b_k&36&73&148&243&382&518&819&1139&1544&2046
\end{array}
\]
The degree-nine bound uses the proved lifted certificate. Section~\ref{low:degree3} gives the cubic adaptive argument. The finite-degree admissibility, normalization and rank models also apply to the surface branches, completed in Sections~\ref{low:degree45}--\ref{low:degree6} and~\ref{surf:finite-section}.


\section{Surface coefficient estimates in degrees seven through nine}\label{surf:finite-section}

Each coefficient count retains its stated model-box restrictions, good-area fraction, seed-net cardinality and original fixed coefficient frame. We use the coefficient coordinates and ribbon calculations of Paper I, Sections~\partIref{atlas:section} and~\partIref{surf:section}. Section~\ref{surf:finite-completion} sums the coordinate estimates and converts their aspects to the assigned geometric aspect. The resulting complete surface bounds enter Proposition~\ref{an:unified-incidence}.

\subsection{Geometry and estimates in degrees seven through nine}\label{surf:common-geometry}\label{low:78:s78pf:section}\label{low:78:s78cf:framework}\label{low:9:root9:interfaces}\label{low:degree78}\label{low:78:s37g:section}\label{low:78:s7c:section}\label{low:degree9}

\paragraph{Degrees seven and eight}

For $a\in C_c^\infty(\mathbb R^3)$, write $E_a(F)=\int_{\mathbb R^3}a(x)e^{iF(x)}\,dx$.

\begin{theorem}\label{low:78:s78:main}
For $k=7,8$, every $a\in C_c^\infty(\mathbb R^3)$ satisfies $E_a\in L^p(V_k)$ whenever $p>\chi_k+3$.
\end{theorem}

The proof of Theorem~\ref{low:78:s78:main} is completed in Section~\ref{completion:physical-II}, using the coefficient estimates below.

\par\smallskip
\noindent\textit{Proof stages for Theorems~\ref{low:78:s78:main} and~\ref{low:9:root9:main}.}\par\nobreak
\noindent\textup{Step 1.} Reduction and common geometric assignment.\par\nobreak
\noindent\textup{Step 2.} Parent atlases and null-coordinate tails.\par\nobreak
\noindent\textup{Step 3.} Gauss tails and surface exhaustion.\par\nobreak
\noindent\textup{Step 4.} Incidence and physical completion.\par\nobreak
\smallskip

Use the analytic construction of Section~\partIref{an:section}, the augmented stop~\partIeqref{reg:stop}, and the fixed RQ presentation of Lemma~\partIref{an:rq}. For $k=7,8$, it remains to establish~\partIeqref{an:physical-incidence} with $\gamma=0$ for arbitrary fixed nonnegative weights in the range $\chi_k+3<p<\chi_k+3+1/16$. Table~\ref{app:table-seven-eight} records the numerical dimensions and ranges. All coefficient integrations use the original fixed frames. The weighted mass bound and high-linear remainder are Lemmas~\partIref{an:resource} and~\partIref{an:linear-tail}.

\begin{lemma}[Polynomial dependence on the window parameter]\label{low:78:s78cf:polynomial}

\noindent\textup{(i)} For the fixed-dimensional, fixed-degree constructions below, all quantitative constants, reciprocal pivot thresholds, cardinalities of auxiliary bounded grids, and losses from fixed sampling fractions are bounded by $P^C$.

\noindent\textup{(ii)} The common preparation multiplicity is at most $P^C(1+\log P+\log d_1^{-1})^C$.
\end{lemma}
\begin{proof}
Apply Lemmas~\partIref{geo:conditioning}--\partIref{geo:paths} to the fixed-format complete phase and physical windows, with closed inner windows, outer slack and conditioned denominator pivots. The regular preparation, maximum parameter length and short normalized comparison after coordinate interchange are those of Lemma~\partIref{reg:regular-preparation}. Only finitely many derivative lists, local inversions and grids are added, so their constants remain polynomial in $P$.
\end{proof}
\paragraph{Exhaustion of the stopped labels}

We make explicit the passage from the rank-four classification to the actual sheets used in the coefficient counts. The augmented translation nondegeneracy condition, all parent coefficient conditions, and physical dyadic shape comparisons remain part of the admissibility condition. 

\begin{lemma}[Quantitative entry and the normal multiplicities]\label{low:78:s78:entry}
For $k=7,8$,

\noindent\textup{(i)} every stopped label is assigned either to the quantitative region defined by the fifth singular value of Proposition~\ref{low:78:h3r:regular-count}, or, after a bounded polynomial-in-$P$ right normalization, to one of the following actual sheets with a conditioned normal pivot:
\[
\begin{array}{c|c|c}
\text{normal order }m&\text{actual equation}&\text{pivot}\\\hline
2&F_x|_c=0&F_{xx}|_c\\
3&F_{xx}|_c=0&F_{xxx}|_c\\
4\ (k=8)&F_{xxx}|_c=0&F_{xxxx}|_c.
\end{array}
\]
The selected sheet is near a conditioned exact rank-four graph.

\noindent\textup{(ii)} All finite choices and ties can be made measurably before the assignment.

\noindent\textup{(iii)} Including the normalization multiplicity, the original-label count is at most $P^C\Lambda$ in every relative-frame neighborhood whose matrix, inverse and translation bounds are at most $P^C$, by \partIeqref{an:local-packing}.
\end{lemma}
\begin{proof}
Apply Lemma~\partIref{geo:rank-entry} with $k=7,8$, retaining the full parent coefficient window. In the quadratic-invariant branch, Theorem~\partIref{geo:classification} permits profile degree at most three and four respectively; in the cubic-invariant branch it permits degree at most two. Its persistent-root cover therefore gives exactly $m=2,3$, and additionally $m=4$ when $k=8$, with the displayed equations and pivots. The same common lemma provides finite Borel root selection, all prescribed graph derivatives and the polynomial packing loss. No higher multiplicity is omitted: profile degree five requires degree at least ten, and a cubic invariant with cubic profile requires degree at least nine.
\end{proof}

\begin{lemma}[Null and nonnull alternatives]\label{low:78:s78:chart-split}

\noindent\textup{(i)} The sheets in Lemma~\ref{low:78:s78:entry} have a finite, exhaustive assignment to the Gauss charts of Section~\ref{low:78:s78gauss:section} and the real-null charts of Sections~\ref{low:78:s37g:section} and~\ref{low:78:s7c:section}.

\noindent\textup{(ii)} This assignment retains polynomial conditioning as the cubic-invariant parameter tends to zero.
\end{lemma}
\begin{proof}
The exact graph Hessian has conditioned nonzero determinant $\Delta$; choose the full coefficient neighborhood so that the actual Hessian remains conditioned. Apply Lemmas~\partIref{geo:chart-split} and~\partIref{geo:two-chart-shear}. Above the inverse-polynomial threshold for $|r|=|c_{yy}|$, the Gauss coordinate $p=c_y$ has conditioned inverse and permits both signs of $\Delta$; the exact model is quadratic. Below that threshold the bounded simple null root has $|B|=|c_{yz}+dc_{yy}|\ge P^{-C}$, and subtracting it gives quadratic graph part $ry^2/2+Byz$. The profile order changes the normal equation and pivot, not this tangential split.

At the exact null models, $u=Yd=0$ and $Yv=Zv=0$. The local model-box and collar proofs use only these identities, conditioned pivots, and fixed-degree polynomial pullbacks. They do not divide by the cubic-invariant parameter or by $v$. Thus the same proofs apply at zero rate and at the intersection of the quadratic- and cubic-invariant families. Choose the local model-box radius first and its defect afterward, as in the common-cover construction. The full parent admissibility is pulled back on that box, including for the derivative-defined sheets with $m=3,4$.
\end{proof}


\paragraph{Application of the common geometry in degree seven}

On the ordinary degree-seven branch use the actual graph $F_x(c,y,z)=0$ and selected null root in \partIeqref{geo:null-equation}. Its pivots $\lambda=F_{xx}|_c$ and $B=c_{yz}+dc_{yy}$ satisfy $P^{-C}\le|\lambda|,|B|\le P^C$. Retain all parent coefficient and physical restrictions of Definition~\partIref{def:admissibility}. The augmented stop and small-fifth-pivot reduction give the model neighborhood of Lemmas~\partIref{geo:model-box}--\partIref{geo:core}. Choose $m_0,\tau_0$ and the model-box radius before the coefficient defect, as in Section~\partIref{geo:common-assignment}. These lemmas give the unbanded core and \partIeqref{geo:small-rates}, including zero rate; higher-normal sheets use the same full parent restrictions.

Apply Theorem~\partIref{geo:common-allocation}, with the collars, caps, endpoint ranges and moderate ribbons of Lemmas~\partIref{geo:local-flow}--\partIref{geo:moderate}. The common Borel assignment has count $P^C\Lambda LW$ and inverse-area bound \partIeqref{geo:inverse-area-weight}, with the original anchors and weights. For a moderate ribbon of coordinate lengths $w,H$ and residual bound $h$, the assigned lengths remain $L=\max(w,H,hH)$ and $W=wH/L$; coordinate and assigned aspects are distinct. The graph-window and coefficient calculations follow.

\begin{lemma}[The parabolic collar supplies the ordinary graph window]\label{low:78:s37g:graph-window}
For the degree-seven critical branch, consider the parabolic collar with the admissibility and preparation conditions of Lemmas~\partIref{geo:collar} and~\partIref{geo:collar-jets}. There is a common positive-area prepared subset on which the following assertions hold.

\noindent\textup{(i)} The full $119$-coefficient map has Jacobian of absolute value at least $P^{-C}$.

\noindent\textup{(ii)} The joint $27$-dimensional graph-coordinate window has volume at most $P^C$.

\noindent\textup{(iii)} The seed has the ranges used in the fixed-width collar table.
\end{lemma}
\begin{proof}
Prepare the algebraic seed data in its flux parameter on a positive proportion of its interval of length $L$. In the actual normalized frame at an evaluation point, the centered tangential germ has the form
\[
Y=A(s)+D_*(s)t+\tfrac12V_*(s)t^2,\qquad Z=B_*(s)+t,
\]
where $A(0)=B_*(0)=0$. The seed equations in the proof of Lemma~\partIref{geo:collar}, the fixed-frame ranges of Lemma~\partIref{geo:collar-endpoints}, and finite derivative preparation give, through every needed fixed order,
\[
\begin{aligned}
|A_1|&\le P^C,& |A_j|&\le P^CL^{2-j}\quad(j\ge2),\\
|(B_*)_j|+|(D_*)_j|&\le P^CL^{1-j}\quad(j\ge1),& |(V_*)_j|&\le P^CL^{-j}.
\end{aligned}
\]
Also $D_*(0)$ is bounded and is small as a null-direction defect. At a point with original transverse parameter $t_0$, these statements follow by substituting $t=t_0+\tau$ and subtracting its actual null slope $d_0$ from the horizontal coordinate. After shifting the longitudinal parameter origin to the evaluation point, the centered functions are
\[
\begin{aligned}
A(s)&=a(s)-a(0)+t_0(D(s)-D(0))+\tfrac12t_0^2(V(s)-V(0))-d_0(b(s)-b(0)),\\
B_*(s)&=b(s)-b(0),\\
D_*(s)&=D(s)+t_0V(s)-d_0,\qquad V_*(s)=V(s).
\end{aligned}
\]
Since $|t_0|\le P^C$ and $|d_0|\le P^CL$, the displayed ranges follow. The bounded actual coframe gives the stronger first-derivative bound for $A_1$.

Every finite seed jet in this list is bounded by a power of $P$ when $L\ge1$. The normalized phase and its normal pivot bound every finite graph derivative appearing in the $27$ graph coordinates, by repeated implicit differentiation of $F_x(c,y,z)=0$. Composition with the prepared seed therefore bounds all $27$ outputs by $P^C$. The joint graph window has volume at most $P^C$. No small interval for $c_{zzz}$ or $v$ is imposed, and the points where $v$ vanishes need no separate treatment.

At the seed, write $\nu=z_s$ as in Lemma~\partIref{geo:collar}. The linear tangential matrix is $\left(\begin{smallmatrix}1&0\\\nu&1\end{smallmatrix}\right)$. At any retained point after recentering its first column is a small perturbation of $(1,\nu)$ and its second column a small perturbation of $(0,1)$, by \partIeqref{geo:small-rates} and the coframe computation. If the actual normalized graph has quadratic part $Q(y,z)=(r/2)y^2+Byz$, then $Q(s,\nu s+t)=(r/2+B\nu)s^2+Bst$ at the exact matrix. The quadratic-grid determinant is therefore a conditioned nonzero multiple of $B^{121}$ because its factor $g_7$ satisfies $g_7(0)\ne0$. At the actual nearby matrix, the same determinant remains bounded away from zero by a power of $P$. The weighted triangular proof of the ordinary atlas shows that higher graph and seed jets do not enter its determinant. This supplies the conditioned ordinary atlas on the actual positive-area certificate.

In the original algebraic seed parameter $\zeta$, the flux derivative is $\partial_s=\alpha(\zeta)^{-1}\partial_\zeta$ with nonvanishing algebraic $\alpha$. Every prescribed derivative remains algebraic in $(\zeta,t)$, so the flux version of Lemma~\partIref{geo:preparation} supplies simultaneous good points in the actual $s$ interval. Apply it to the bounded normal scalars, all four third-normal outputs, and the critical-value gradient. Rounding the seven seed jets with the stated parabolic meshes preserves the ordinary collar widths by the finite anisotropic chain rule, and the bounded rounded jets retain the graph-window argument.
\end{proof}

\begin{lemma}[Degree-seven ordinary transverse-collar count]\label{low:78:s37g:high}
The narrow and high-residual transverse collars of Lemma~\partIref{geo:collar-endpoints}, and the projection caps of Lemma~\partIref{geo:caps}, satisfy the complete fixed-frame coefficient bound $P^CL^{-232}$ in the degree-seven critical branch.
\end{lemma}
\begin{proof}
Lemma~\ref{low:78:s37g:graph-window} supplies the conditioned full-coefficient atlas, bounded graph-coordinate ranges, and actual seed preparation. Lemma~\partIref{geo:collar-endpoints} supplies the point factor $P^CL^2$. For a projection cap the same endpoint argument applies with a horizontal seed and a bounded initial characteristic comparison. 

For completeness, the ordinary degree-seven collar has raw scalar power $L^{-269}$, seven-jet seed cost $L^{34}$, and point factor $L^2$. The coefficient, seed and point factors give $L^{-269}L^{34}L^2=L^{-233}$. Conditioning on endpoint shear adds at most one further factor $L$, giving $L^{-269}L^{34}L^2L=L^{-232}$. These coefficient ingredients are independent of the geometric replacement and retain the same fixed-frame fiber as in the stated table.
\end{proof}

\paragraph{Geometric application}

The actual algebraic cells and the seed grids determine the long-scale certificates; the model supplies the local admissibility box.

\paragraph{Degree nine}

\begin{theorem}\label{low:9:root9:main}
Every $a\in C_c^\infty(\mathbb R^3)$ satisfies $E_a\in L^p(V_9)$ whenever $p>498$.
\end{theorem}

The proof of Theorem~\ref{low:9:root9:main} is completed in Section~\ref{completion:physical-II}, using the coefficient estimates below.

Use the analytic construction in degree nine. Lemmas~\partIref{an:resource} and~\partIref{an:linear-tail} supply the weighted mass bound and remainder. The regular assignment, parameters and $22$ displacement cases are those of Lemma~\partIref{reg:regular-preparation} and Section~\partIref{reg:geometric-regimes}; the complete degree-nine minors and the cardinalities of their additional-sample nets were established in Section~\ref{low:9:root9:regular}.

\paragraph{Approximation by rank-four models and surface assignment}
The classification is $F=f(q)+\ell(y-\gamma z^2/2)+nz$, with a nonsingular graph invariant of degree two or three and the relation $r\gamma=0$. A quadratic invariant permits $\deg f\le4$ in degree nine; an invariant of degree three permits $\deg f\le3$. Thus the only required normal orders are $m=2,3,4$.

One does not assume a simple root of $f'$. For a real nonlinear profile, some derivative $f^{(m-1)}$ has a real simple root: a highest nonconstant derivative with degree one suffices. Compact root neighborhoods in the conditioned profile family and semialgebraic growth bound the root location and inverse pivot by $P^C$. The exact root is translated to the origin, and the implicit function theorem gives the selected sheet
\[
\partial_x^{m-1}F|_c=0,\qquad \kappa=\partial_x^mF|_c,\qquad |\kappa|\ge P^{-C}.
\]
The coefficient defect is chosen after these bounds. All finite graph derivatives and normalized coefficients of the original phase remain controlled. The large-first-curvature branch uses the Gauss chart; there the exact model is quadratic because $r\gamma=0$. In the complementary branch the nonsingular Hessian supplies a bounded simple real null slope and a mixed pivot $B$. This includes the cubic invariant and the intersection $\gamma=0$, without division by $\gamma$ or a rate.

Apply Lemmas~\partIref{geo:model-box}--\partIref{geo:core} and Theorem~\partIref{geo:common-allocation} for $m=2,3,4$. At each model, the derivative-defined sheet is the same invariant graph and satisfies $Yd=Yv=Zv=0$; the conditioned normal pivot selects its actual perturbation. Fix the local box and derivative lists before the complete degree-nine coefficient defect, retaining all parent restrictions. The resulting unbanded assignment has contained prepared certificates, count $O_P(\Lambda S)$, assigned width at least a conditioned constant, aspect $K\lesssim_PS$, and the stated conditioning and physical logarithmic multiplicity.

Lemmas~\partIref{geo:original-subgroup-pivots}, \partIref{surf:haar}, and~\partIref{surf:null-density} give the fixed-frame null incidence with four normalization columns of determinant $2\kappa^4B$ (or three of determinant $\kappa^3$). Every other one of the $219$ coefficient directions is retained. The bounded original subgroup fiber follows from its pivot, bounded variation of the complete coefficient vector on fixed-format paths, and the step-three upper-unipotent group law. Wide ribbon point cost is $P^CH^2$; vertical ribbon point cost is $P^CL^2K_0$; collar point cost is $P^CL^2$, or conservatively $P^CL^3$ if the endpoint shear is separately conditioned. Gauss incidence instead has cost $P^CK$ as proved below.

\subsection{Null coefficient bounds by normal order}\label{surf:finite-null-counts}\label{low:78:s37i:section}\label{low:78:s38o:section}\label{low:78:s38i:section}\label{low:78:d8q:section}\label{low:9:u9:formulas}\label{low:9:h9d:section}\label{low:9:u9:collar}

\paragraph{Hypotheses and coefficient maps in degrees seven through nine}

Fix $k\ge7$, a normal order $2\le m\le\lfloor k/2\rfloor$ supplied by the assignment, and $D=k-m+2$. Use the normal restrictions of Definition~\partIref{def:normal-restrictions} and geometric certificates of Definition~\partIref{def:geometric-certificate}, with a fixed positive proportion of common good points and conditioned $\kappa_m,B$. Apply the simultaneous scalar, curvature and seed preparation of Section~\partIref{surf:null-counts}, with its scalar bounds and no denominator from a lower normal scalar; write $\phi=\phi_0$.

Use the index sets of Section~\partIref{atlas:restriction}, with $\mathcal I_d=I_d$, $\mathcal T_d=T_d$, and $\mathcal E_D=E_D$, and the full $N_k$-dimensional parent measure. Apply the normalization and fixed-seed coefficient integration in Section~\partIref{surf:null-counts}. The four deleted graph rows $(0,0),(1,0),(0,1),(0,2)$ are present because $T_{D-2}\subset E_D,E_D^{\rm w}$. In Lemma~\partIref{surf:null-density}, the four-column determinant is $2\kappa_m^4B$ with $N_k-4$ slice directions; the three-column alternative has determinant $\kappa_m^3$ and retains a separate integration over the actual endpoint shear interval. Use the point factor of Definition~\partIref{def:point-factor}.

The unbanded cover uses the complete bounded graph-rate interval, including zero rate. In the four cases $(k,m)=(7,2),(7,3),(8,3),(8,4)$, the bounds also have a factor $\nu$ when the certificate supplies the graph-rate window: the original lower graph coordinates kept in the respective count have joint volume $P^C\nu$, with $c_{zzz}=-2Bv$ retained. All applications to the final theorem take $\nu=1$. Containment, full admissibility, good area, and preparation are geometric hypotheses, not consequences of a nonsingular coefficient determinant.

\begin{proposition}[Finite-degree parent atlases]\label{surf:finite-null-atlases}\label{low:78:s37i:atlas}\label{low:78:s38o:atlas}\label{low:78:s38i:atlas}\label{low:78:d8q:atlas}\label{low:9:h9d:atlases}

\noindent\textup{(i)} For $k=7,8,9$, use the coefficient output lists of Lemma~\partIref{surf:parent-lift}, including the constant critical-value row of the derivative phase and all lower parent triangles when $m\ge3$. Only the constant of $\phi_0$ is omitted. The block dimensions are
\[
\begin{array}{c|c|c|c|c}
(k,m)&D&\text{ordinary or derivative blocks}&\text{parent blocks}&N_k\\\hline
(7,2)&7&63+25+27+4&0&119\\
(7,3)&6&49+16+18+1&35&119\\
(8,2)&8&80+36+38+9+1&0&164\\
(8,3)&7&64+25+27+4&44&164\\
(8,4)&6&49+16+18+1&36+44&164\\
(9,2)&9&99+49+51+16+4&0&219\\
(9,3)&8&81+36+38+9+1&54&219\\
(9,4)&7&64+25+27+4&45+54&219.
\end{array}
\]
The weighted lists replace $I_d,E_D$ by $I_d^{\rm w},E_D^{\rm w}$ and leave the parent triangles unchanged.

\noindent\textup{(ii)} For an external tangential germ with invertible linear part $M$ and transformed quadratic graph $Ay^2+byz+Cz^2$, their absolute Jacobians, up to fixed positive coordinate factors, are
\[
|\det M|^{\chi_k}|\kappa_m(0)|^{-|E_D|}|G_D(A,b,C)|,
\qquad |\det M|^{\chi_k}|\kappa_m(0)|^{-|E_D|}|A|^{t_D},
\]
where $G_D=b^{t_D}g_D(AC/b^2)$ is interpreted polynomially at $b=0$, $g_D(0)\ne0$, and $(t_6,t_7,t_8,t_9)=(73,121,189,282)$. These identities do not require bounded $M$ or $M^{-1}$ and are independent of higher graph and external seed jets. Thus they apply to the actual cubic graphs as well.

\noindent\textup{(iii)} These maps give local coefficient coordinates on the ordinary locus $G_D(A,b,C)\ne0$ and the weighted locus $A\ne0$, respectively.
\end{proposition}
\begin{proof}
Lemma~\partIref{surf:strong-atlas} applies since $D\ge6$. For $m\ge3$, the critical equation for $G$ is $\partial_x^{m-1}F|_c=0$ with second-normal pivot $\kappa_m$. The block factorization in Lemma~\partIref{surf:parent-lift} supplies every parent block and the determinant identity, with moving-graph corrections in lower off-diagonal blocks. On the stated nonzero-Jacobian loci, Proposition~\partIref{low:78:smult:atlas} gives local coefficient coordinates. The dimensions follow by counting the stated sets. By~\partIeqref{surf:parent-jacobian}, the parent triangles contribute $\chi_k-\chi_D$ to the character. Equation~\partIeqref{surf:tD} gives the four quadratic degrees. 
\end{proof}

\par\smallskip\noindent\textit{Step 1. Complete parent outputs and normal orders.}\par\nobreak\smallskip

Both ribbon parametrizations below have $|\det M|=1$; collar linear parts are conditioned.

The full normal blocks in Proposition~\ref{surf:finite-null-atlases} are retained even when their values vanish at the reference model. For example, with $Q=yz$ and $q=x-Q$, the degree-seven variation $(q^4-Q^4)/4$ at $q^3$ is detected by the fourth-normal output, whereas the degree-eight family $G=q^2+\epsilon q^4$ satisfies
\[
\left.(G,G_x,G_{xx},G_{xxx},G_{xxxx})\right|_{x=Q}=(0,0,2,0,24\epsilon).
\]
Thus the first four listed outputs do not determine the fourth normal derivative. Lower normal scalars may vanish; all denominators are the selected nonzero pivots.

For the ordinary degree-eight branch, apply Lemmas~\partIref{geo:model-box}--\partIref{geo:core} with $k=8,m=2$ and then Theorem~\partIref{geo:common-allocation}. The quadratic-invariant profile may have degree four and the cubic-invariant profile degree two; both are preserved by the same affine model flow. The pivots remain $\lambda,2B$, and full parent polynomial pullback changes only the conditioning constants. Thus the preceding common geometry applies.

For the degree-eight quartic-normal case, write $G=F_{xx}$, $\psi_2=\phi_2$, $\psi_1=\phi_1$, $\kappa=\phi_4$, and $\mu=\phi_5$. Its output list is
\begin{equation}\label{low:78:d8q:outputs}
(\psi_2)_{\mathcal I_6},\quad \kappa_{\mathcal I_3},\quad c_{\mathcal E_6},\quad \mu(0),\quad (\psi_1)_{\mathcal T_7},\quad \phi_{\mathcal T_8\setminus\{0\}}.
\end{equation}

In this quartic-normal branch the lower scalars need not vanish and remain parent outputs. The fifth-normal value is also needed: at $F_0=q^4$, where $Q=yz$ and $q=x-Q$, the degree-eight variation
\begin{equation}\label{low:78:d8q:hidden}
R_*=\frac{q^5+Q^5-5xQ^4}{20}
\end{equation}
satisfies
\[
\begin{aligned}
(R_*)_{xx}&=q^3,& (R_*)_x|_{x=Q}&=-Q^4/4,\\
R_*|_{x=Q}&=-Q^5/5,& (R_*)_{xxxxx}|_{x=Q}&=6.
\end{aligned}
\]
The parent restrictions are invisible to their order-seven and order-eight triangles, so the fifth-normal value is retained in~\eqref{low:78:d8q:outputs}.

The exact classification gives profile degree at most four on a quadratic invariant and at most three on a cubic invariant. Thus the necessary normal orders in degree nine are $m=2,3,4$. On the selected simple sheet write
\[
 \partial_x^{m-1}F(c(y,z),y,z)=0,\qquad
 \kappa_m=\partial_x^mF|_c,\qquad |\kappa_m|\ge P^{-C}.
\]
The same compact root-neighborhood argument used for quantitative entry in degrees seven and eight applies to profiles of these bounded degrees: select a root of a derivative having a conditioned next derivative, translate by a polynomially bounded amount, and apply the implicit function theorem while keeping the full parent coefficient defect. The coefficient dimension changes to $219$ but no new profile multiplicity occurs beyond the displayed three orders.

The cubic-normal profile is $F_\gamma=q_\gamma^3$. Its cubic-normal graph is
\[
c(y,z)=yz-\gamma z^3/3,\qquad F_{xxx}|_c=6,
 \qquad d=\gamma z,\quad Yd=0,\quad Zd=\gamma.
\]
The identities $Yd=Y(Zd)=Z(Zd)=0$ and the commuting affine fields preserving the full parent phase verify the model hypotheses of Lemmas~\partIref{geo:model-box}--\partIref{geo:core} with $k=9,m=3$. Fixing the local box and widths before the complete coefficient defect gives the unbanded geometry already obtained in Section~\ref{surf:common-geometry}. Only the normal pivot and scalar atlas change. The nonnull quadratic-graph chart uses the Gauss argument.

The same requirement persists on cubic graphs in degree nine. With $Q_0=yz$, $h=\gamma z^3/3$ and $Q_\gamma=Q_0-h$, the degree-nine variation $R_*=((x-Q_0)^5+Q_0^5)/5$ at $F_0=(x-Q_\gamma)^3$ has graph restrictions
\[
\bigl(R_*,(R_*)_x,(R_*)_{xx},(R_*)_{xxx},(R_*)_{xxxx},(R_*)_{xxxxx}\bigr)|_{Q_\gamma}
=\bigl((Q_0^5-h^5)/5,h^4,-4h^3,12h^2,-24h,24\bigr).
\]
The parent triangle misses the first entry, and the next four entries have excluded indices $(0,12),(0,9),(0,6),(0,3)$ in $I_8,E_8,I_5,I_2$. Only the retained fifth-normal value detects this variation. All higher-normal blocks in the other branches are likewise kept.

The finite slope estimates for $D=6,7,8,9$ are proved in Section~\partIref{atlas:slopes}. For $(k,m)=(7,3),(8,4)$, use its degree-six atlas, with $|E_6|=18$, $t_6=73$ and inverse-determinant cost $H^3$.

\paragraph{Step 2. Wide ribbons}

For $k=7,8,9$, the complete coefficient maps in Proposition~\ref{surf:finite-null-atlases} have seed order at most $k$, the required complete lower graph jet, and extra ordinary graph rows with horizontal index at least two. The calculation in the proof of Proposition~\partIref{surf:wide} gives
\[
\begin{array}{c|c|c|c}
(k,m)&\text{scalar product}&\text{graph product}&\mu_D\\\hline
(7,2)&w^{-213}H^{-269}&w^{-18}H^{-30}&5\\
(7,3)&w^{-220}H^{-248}&w^{-6}H^{-12}&3\\
(8,2)&w^{-343}H^{-415}&w^{-40}H^{-60}&5\\
(8,3)&w^{-352}H^{-388}&w^{-18}H^{-30}&5\\
(8,4)&w^{-331}H^{-367}&w^{-6}H^{-12}&3\\
(9,2)&w^{-524}H^{-614}&w^{-75}H^{-105}&8\\
(9,3)&w^{-534}H^{-579}&w^{-40}H^{-60}&5\\
(9,4)&w^{-507}H^{-552}&w^{-18}H^{-30}&5.
\end{array}
\]
Sum the scalar indices and $(i-2,j)$ over the indicated graph rows. For the scalar and graph products $V_{\mathrm{sc}},V_{\mathrm{gr}}$ in each table row, $\operatorname{vol}E_j\lesssim_P V_{\mathrm{sc}}V_{\mathrm{gr}}H^{2k-1}H^{\mu_D}H^2$. The last three factors are the seed count, inverse determinant and point factor, respectively, giving the wide bounds in Proposition~\ref{surf:finite-null-tails}. The proof of Proposition~\partIref{surf:wide} preserves the widths under rounding and slope shifts; the joint graph-window hypothesis gives the additional factor in the specified cases.

\paragraph{Step 3. Vertically long ribbons}

For $k=7,8,9$, the determinant dichotomy in the proof of Proposition~\partIref{surf:vertical} applies. Before the $k$ purely longitudinal amplitude factors, the ordinary products are
\[
\begin{array}{c|c|c}
(k,m)&\text{scalar product}&\text{extra graph product}\\\hline
(7,2)&L^{-269}W^{-220}&L^{20}\\
(7,3)&L^{-248}W^{-227}&L^{-8}W^{-2}\\
(8,2)&L^{-415}W^{-351}&L^{-40}W^{-20}\\
(8,3)&L^{-388}W^{-360}&L^{-20}W^{-8}\\
(8,4)&L^{-367}W^{-339}&L^{-8}W^{-2}\\
(9,2)&L^{-614}W^{-533}&L^{-70}W^{-40}\\
(9,3)&L^{-579}W^{-543}&L^{-40}W^{-20}\\
(9,4)&L^{-552}W^{-516}&L^{-20}W^{-8}.
\end{array}
\]
Except for the first row, the graph product is obtained from $\partial_\tau^2(c\circ\Psi)=r\circ\Psi$ on $E_D\setminus T_{D-2}$. For $D=6$ these are $(2,3),(3,2),(3,3)$, giving the needed $L^{-8}W^{-2}$ gain. In the first row the six extra graph indices have longitudinal sum $20$; bounded original graph jets and the seed ranges give interval product at most $P^CL^{20}$. Restoring the complete lower graph jet is triangular with determinant one in every case. The graph window and curvature intervals are separate coordinates. With $V_{\mathrm{sc}}$ and $V_{\mathrm{gr}}$ equal to the scalar and graph products in this table row, $\operatorname{vol}E_j\lesssim_P V_{\mathrm{sc}}V_{\mathrm{gr}}L^kL^{2k-1}K_0^kL^2K_0$.

In the weighted chart, let $a$ denote the purely longitudinal amplitude power remaining after determinant absorption, and let $J$ be the required seed order. The complete data are
\[
\begin{array}{c|c|r|r|r}
(k,m)&\text{scalar product}&\text{graph shear power}&a&J\\\hline
(7,2)&L^{-367}W^{-122}&36&14&14\\
(7,3)&L^{-314}W^{-161}&15&0&12\\
(8,2)&L^{-566}W^{-200}&70&16&16\\
(8,3)&L^{-493}W^{-255}&36&8&14\\
(8,4)&L^{-433}W^{-273}&15&0&12\\
(9,2)&L^{-837}W^{-310}&120&18&18\\
(9,3)&L^{-738}W^{-384}&70&9&16\\
(9,4)&L^{-657}W^{-411}&36&9&14.
\end{array}
\]
Here $J=2D\ge k$. The nonlinear substitutions $(i,j)\mapsto(i-b,j+1)$ for $b\ge2$ stay in the weighted graph set, so the same triangular removal leaves precisely the displayed linear-shear powers. Each is absorbed by $|A|^{-t_D}$. For $(7,3)$ and $(8,4)$, the sextic power $73$ also absorbs the seven or eight purely longitudinal parent factors, since $73\ge15+8$; in the other rows retain the full amplitude. With $V_{\mathrm{sc}}$ the scalar product in the indicated row, $\operatorname{vol}E_j\lesssim_P V_{\mathrm{sc}}L^aL^{2J-1}K_0^JL^2K_0$. This proves the weighted entries in Proposition~\ref{surf:finite-null-tails}. The anisotropic comparison uses the prepared auxiliary jets of Proposition~\partIref{surf:vertical}.

\paragraph{Step 4. Parabolic collars}

In finite degree, the raw scalar long exponents and seed costs are
\[
\begin{array}{c|r|r}
(k,m)&\text{raw exponent}&\text{seed exponent}\\\hline
(7,2)&269&34\\
(7,3)&248&30\\
(8,2)&415&39\\
(8,3)&388&35\\
(8,4)&367&27\\
(9,2)&614&44\\
(9,3)&579&44\\
(9,4)&552&44.
\end{array}
\]
The raw exponent is $U-k$. Lemma~\partIref{surf:collar-net} gives $L^{34},L^{39}$ in the ordinary degree-seven and degree-eight cases, and $L^{44}$ for all degree-nine atlases. For the latter, including the mixed row $(9,9)$, prepare scalar order $18$ and retain unrounded seed order $19$ during comparison. Apply the determinant and fiber calculation of Proposition~\partIref{surf:collar-tail}: subtract the seed exponent and point-and-shear cost $3$, or $2$ for the four-dimensional fiber. The improved remaining seed counts are proved next.

\emph{The degree-seven cubic-normal refinement.} Consider a contained fixed-width parabolic collar with bounded coframe, parameterized by
\[
\Psi(s,t)=\left(a(s)+D(s)t+\tfrac12V(s)t^2,b(s)+t\right),\qquad |s|<L,\quad |t|<P^C.
\]
At a normalized endpoint its local parametrization is
\[
Y=A(s)+D_*(s)t+\tfrac12V_*(s)t^2,\qquad Z=B_*(s)+t,
\]
where $A(0)=B_*(0)=0$, while $D_*(0)$ and $V_*(0)$ retain their actual values. As in Lemma~\partIref{geo:collar-jets}, the linear part is a small conditioned perturbation of $\left(\begin{smallmatrix}1&0\\m&1\end{smallmatrix}\right)$, $|m|\le1$. Assume the standard prepared ranges
\[
\begin{gathered}
|A_1|\le P^C,\qquad |A_j|\le P^CL^{2-j}\ (j\ge2),\\
|(B_*)_j|\le P^CL^{1-j},\qquad |(D_*)_j|\le P^CL^{1-j},\qquad |(V_*)_j|\le P^CL^{-j}.
\end{gathered}
\]
with bounded $(D_*)_0$, and retain all fixed extra true jets needed for derivative comparison. The original point factor is at most $P^CL^2$. If the endpoint shear is separately conditioned, its interval has length at most $P^CL$.

\begin{lemma}[A $30$-power seed net]\label{low:78:s37i:collar-net}

\noindent\textup{(i)} The $119$ outputs for $(k,m)=(7,3)$ in Proposition~\ref{low:78:s37i:atlas} require a fixed relative seed net with at most $P^CL^{30}$ elements.

\noindent\textup{(ii)} The net preserves bounded-scalar widths $L^{-i}$ and the parent widths $L^{-i}$ for $j\ge1$ and $L^{1-i}$ for $j=0$, $i\ge1$.
\end{lemma}
\begin{proof}
The first $84$ outputs involve at most six longitudinal derivatives. For the first six jets use the prescribed meshes
\[
\begin{array}{c|c|c}
\text{parameter}&\text{mesh}&\text{total number of choices}\\\hline
A_1&c_PL^{-1}&P^CL\\
A_j,\ 2\le j\le6&c_PL^{-j}&P^CL^{10}\\
(B_*)_j,\ 1\le j\le6&c_PL^{-j-1}&P^CL^{12}\\
(D_*)_j,\ 1\le j\le6&c_PL^{-j}&P^CL^6\\
(D_*)_0,\ (V_*)_j,\ 0\le j\le6&c_P,\ c_PL^{-j}&P^C.
\end{array}
\]
Their product is $P^CL^{29}$. No constant position is rounded. The extra parent triangle has exactly one row with longitudinal index seven, namely $\partial_s^7\phi$. At $t=0$ this row depends on $A_7$ and $(B_*)_7$, but not on $(D_*)_7$ or $(V_*)_7$. With the phase and all lower seed derivatives fixed, the change caused by the two highest derivatives is
\[
\Delta\bigl[(\partial_s^7(\phi\circ\Psi))(0)\bigr]=\phi_y(0)\Delta A_7+\phi_z(0)\Delta(B_*)_7.
\]
The coefficients $\phi_y(0),\phi_z(0)$ are bounded by complete parent admissibility and the conditioned endpoint linear part. The allowed width of this row is $P^CL^{-6}$. Accordingly use mesh $c_PL^{-6}$ for each of $A_7$ and $(B_*)_7$. Their ranges are respectively $P^CL^{-5}$ and $P^CL^{-6}$, so their additional costs are $P^CL$ and $P^C$. The whole net costs $P^CL^{30}$.

For the width comparison, first round the jets through order six, retaining the unrounded higher jets during this comparison. The actual inverse-coordinate fields have the form
\[
\partial_y=J^{-1}(\partial_s-m\partial_t),\qquad \partial_z=\partial_t-(D_*+V_*t)\partial_y,
\]
where $J=\det D\Psi$ and $m(s)=B_*'(s)$. The first field's coefficients have longitudinal derivatives $O_P(L^{-i})$, and those of the second have bounds $O_P(L^{1-i})$. The finer $(B_*)_j$ meshes compensate the latter bound. The induced interpolation vector field therefore has coefficients with derivatives $O(P^Cc_PL^{-i})$ through every fixed order needed here. It vanishes at the centered endpoint. The finite chain rule consequently preserves each mixed bounded-scalar width and each parent width, including the purely longitudinal seventh row: any transverse derivative applied to a purely longitudinal parent jet reaches a term with at least one longitudinal derivative, reducing its extra factor $L$. The unrounded higher jets need no net because the compared $84$ outputs and the mixed parent rows do not depend on them.

For the highest order in the preceding comparison, the longitudinal coefficient of the interpolation field is explicitly $J^{-1}\{\delta A+\delta D\,t+\delta V\,t^2/2-(D_*+V_*t)\delta B_*\}$. At $t=0$ its seventh derivative has no $\delta A_7$ or $\delta B_7$ term because the unrounded seventh jets have been retained. The product terms have size $L^{1-j}L^{-(7-j)-1}=L^{-7}$ for $1\le j\le6$; the $j=7$ term uses the zero position error. Derivatives of $J^{-1}$ retain the prepared powers because the exact flux Jacobian and its derivatives are prepared jointly; under interpolation, a term $\delta D_0(B_*)_{i+1}$ has the allowed size $O_P(c_PL^{-i})$. This verifies the seventh-order comparison without an isolated loss from a high horizontal seed jet.

After that comparison, round $A_7$ and $(B_*)_7$ alone. No lower row changes; the only affected output is the purely longitudinal seventh parent row. The preceding exact linear dependence bounds its change by $P^Cc_PL^{-6}$. No seventh ruling jet appears in any retained row. Thus the net preserves all the required widths.
\end{proof}

For $(k,m)=(8,3)$ the derivative phase block has $120$ outputs and uses seven longitudinal jets. The lower net costs $L^{34}$. The parent triangle adds just one row of longitudinal order eight, $\partial_s^8\phi$, of allowed width $L^{-7}$. With the phase and all lower seed derivatives fixed,
\[
\Delta\bigl[(\partial_s^8(\phi\circ\Psi))(0)\bigr]=\phi_y(0)\Delta A_8+\phi_z(0)\Delta(B_*)_8.
\]
The two coefficients are bounded, and no eighth ruling jet enters this row. The ranges of $A_8,(B_*)_8$ are $L^{-6},L^{-7}$, so mesh $c_PL^{-7}$ adds only $L$. Apply the same two-stage comparison as in Lemma~\ref{low:78:s37i:collar-net}, with the lower jets rounded while unrounded higher jets are retained, followed by these two top jets. This preserves all mixed rows and the purely longitudinal eighth row and gives the stated $L^{35}$ net.

\emph{The degree-eight quartic-normal refinement.} Let a contained fixed-width parabolic collar be
\[
\Psi(s,t)=\left(a(s)+D(s)t+\tfrac12V(s)t^2,b(s)+t\right),\qquad a'=1+Db',\quad D'=u+Vb'.
\]
Its long scale is $L\ge1$, its transverse width is a fixed power of $P$, and its coframe is conditioned. The model defect is chosen after that fixed width, so the tangential Jacobian $1+u(s)t+V'(s)t^2/2$ stays between two positive conditioning constants. All full parent admissibility and fixed-frame bounds are retained.

At an arbitrary good endpoint with transverse parameter $t_0$, recenter $t=t_0+\tau$ and subtract its actual null slope $d_0$. The parameter functions become, up to centered position constants,
\[
A=a+Dt_0+Vt_0^2/2-d_0b,\quad B_*=b,\quad D_*=D+Vt_0-d_0,\quad V_*=V.
\]
In general $D_*(0)\ne0$ and $A'(0)\ne1$. Their exact relation is
\begin{equation}\label{low:78:d8q:off-seed}
A'-D_*B_*'=J_0(s;t_0),\qquad J_0=1+u(s)t_0+\tfrac12V'(s)t_0^2.
\end{equation}
The seed-grid construction below uses this identity at every recentered endpoint.

\begin{lemma}[Longitudinal reparametrization and coupled rounding]\label{low:78:d8q:coupled-net}

\noindent\textup{(i)} The full $164$-output atlas on these collars admits a fixed relative seed net of size at most $P^CL^{27}$.

\noindent\textup{(ii)} The net retains all evaluation points after recentering and preserves the prepared bounded-scalar widths $L^{-i}$ and the parent widths $L^{-i}$ for mixed rows and $L^{1-i}$ for purely longitudinal rows. 
\end{lemma}
\begin{proof}
For the fixed endpoint parameter $t_0$, change only the longitudinal parameter by
\[
\sigma(s)=\int_{s_0}^sJ_0(q;t_0)\,dq.
\]
The derivative is uniformly positive, so the new interval has length comparable to $L$ and the area comparison is unchanged up to conditioning factors. Prepare $u,V'$ and their fixed derivative lists together with the original scalar and seed jets. At the common good points, $\partial_s^jJ_0=O_P(L^{-j})$. Repeated use of $\partial_\sigma=J_0^{-1}\partial_s$ then preserves every stated scalar width and every prepared seed range. The inverse derivatives have the same bounds by the finite reciprocal formula. In an original algebraic seed parameter, $\partial_\sigma$ is differentiation divided by a nonzero algebraic scalar. Thus its fixed derivative lists have the same uniform finite format, although the primitive $\sigma$ itself need not be algebraic.

In this parameter, \eqref{low:78:d8q:off-seed} becomes the identity
\begin{equation}\label{low:78:d8q:flux-identity}
A_\sigma=1+D_* (B_*)_\sigma.
\end{equation}
The recentered germ is still quadratic in $\tau$. Write $B_j=\partial_\sigma^jB_*(0)$, $D_j=\partial_\sigma^jD_*(0)$, and $V_j=\partial_\sigma^jV_*(0)$. The endpoint positions are zero. The prepared ranges are
\[
|B_j|+|D_j|\lesssim_PL^{1-j}\ (j\ge1),\qquad |V_j|\lesssim_PL^{-j},\qquad |D_0|\lesssim_P1.
\]
Here $D_0$ is the actual small transverse slope error; it is not set to zero. The horizontal jets are determined by
\begin{equation}\label{low:78:d8q:A-jets}
A_1=1+D_0B_1,\qquad A_j=\sum_{r=0}^{j-1}\binom{j-1}{r}D_r B_{j-r}\quad(j\ge2).
\end{equation}
Use the following fixed grids, with meshes multiplied by a sufficiently small $c_P>0$:
\[
\begin{array}{c|c|c}
\text{parameter}&\text{mesh}&\text{total number of choices}\\\hline
D_0&L^{-1}&P^CL\\
B_j,D_j,\ 1\le j\le6&L^{-j-1}&P^CL^{24}\\
B_7,D_7&L^{-7}&P^CL^2\\
B_8&L^{-7}&P^C\\
V_j,\ 0\le j\le6&L^{-j}&P^C.
\end{array}
\]
Their product is $P^CL^{27}$. Define the rounded horizontal jets from the rounded $B,D$ jets by \eqref{low:78:d8q:A-jets}; they are not independently gridded. These rounded Taylor data can be realized by polynomial $B_*,D_*,V_*$ and $A(\sigma)=\sigma+\int_0^\sigma D_*(q)(B_*)'(q)\,dq$. The resulting parametrization has linear determinant one. The small range of $D_0$ keeps its linear part in the conditioned neighborhood required by the sextic grid determinant.

For $j\le7$, every change in the right side of \eqref{low:78:d8q:A-jets} is $O_P(c_PL^{-j})$. The term with $r=0$ uses $\delta D_0=O(c_PL^{-1})$ and $B_j=O_P(L^{1-j})$; the other terms use one factor with mesh $L^{-r-1}$ and the other with range $L^{1-(j-r)}$, or the symmetric bound. For $j=8$, the only less restrictive terms involve $D_7B_1$, $D_1B_7$, or $D_0B_8$; each changes by at most $O_P(c_PL^{-7})$. Thus
\[
|\delta A_j|\lesssim_Pc_PL^{-j}\ (1\le j\le7),\qquad |\delta A_8|\lesssim_Pc_PL^{-7}.
\]
Products of two rounding errors satisfy the same estimates. These estimates include the freedom after recentering in $D_0$.

For the first comparison, round the jets through long order six and $D_0$, while retaining the unrounded values of $A_7,A_8,B_7,B_8,D_7$ and all additional higher unrounded jets used in preparation. The temporary germs need not satisfy the coupling identity. The changed jets have errors no larger than the collar meshes: $A_j$ has mesh $L^{-j}$, $B_j$ has the finer mesh $L^{-j-1}$, $D_j$ has at most mesh $L^{-j}$, and $V_j$ has mesh $L^{-j}$. The inverse-coordinate vector fields are
\[
\partial_y=J^{-1}(\partial_\sigma-m\partial_\tau),\qquad \partial_z=\partial_\tau-(D_*+V_*\tau)\partial_y,\qquad m=(B_*)_\sigma.
\]
The first field has longitudinal coefficient derivatives $O_P(L^{-i})$; the second has the bounds $O_P(L^{1-i})$. The extra power in the $B_j$ meshes compensates the latter. Interpolate through these temporary germs. The induced vector field vanishes at the centered endpoint, and its coefficients have derivatives $O_P(c_PL^{-i})$ through every fixed order required by the output list. In particular, the comparison applies through longitudinal order eight: the higher jet errors are zero, and the retained unrounded higher derivatives satisfy the prepared bounds. The finite chain rule therefore preserves all retained bounded-scalar widths and all mixed and purely longitudinal parent widths, including the new high rows. Only finite jets are used; no global inverse ribbon is differentiated.

Next round $B_7,D_7,B_8$ and reset $A_7,A_8$ to the coupled values in \eqref{low:78:d8q:A-jets}. Only the bounded-scalar row $(\psi_1)_{7,0}$ and the parent rows $\phi_{7,0}$, $\phi_{7,1}$, and $\phi_{8,0}$ can change in this second step. Their dependence on the highest new seed jets is linear, with coefficients bounded by the first two coordinate derivatives of the scalar and the conditioned endpoint linear part. Those derivatives are bounded by full parent admissibility and the normal pivot. The allowed errors are $L^{-7}$ for the bounded-scalar seventh row and the mixed parent row, and $L^{-7}$ for the purely longitudinal eighth row. The displayed errors for $A_7,B_7,D_7,A_8,B_8$ satisfy these requirements. In the purely longitudinal seventh parent row the allowed width is even larger. No $D_8$ or $V_7$ occurs: a $D_j$ term carries one transverse derivative and a $V_j$ term carries two, while all parent indices have total order at most eight. The derivative phase square has long index at most six. This proves preservation of every output in \eqref{low:78:d8q:outputs}.

Finally, the reparametrization selects the finite Taylor data listed in the table and the horizontal jets determined by \eqref{low:78:d8q:A-jets}. These data have the fixed ranges above and are rounded into grids independent of the endpoint phase. After a grid element is fixed, the coefficient output map is fixed. All coefficient integration therefore remains at the original fixed frame. The unchanged physical good set is counted with its original area; the longitudinal change has bounded positive Jacobian.
\end{proof}

For this last refinement the raw product is $L^{-367}$: the $\psi_2,\kappa,\psi_1,\phi$ blocks contribute $147,24,84,112$, respectively, while $\mu(0)$ is bounded. The linear part of the rounded germ has determinant one and stays in the conditioned sextic chart, and the graph coordinates retain their bounded window. The raw, seed and point factors give $L^{-367}L^{27}L^2=L^{-338}$ for a valid intrinsic $H_0$ map. With a separate endpoint-shear interval and the $H_3$ fiber, they give $L^{-367}L^{27}L^2L=L^{-337}$. This verifies the quartic row of the collar table.

\paragraph{Step 5. Coefficient bounds and assigned scales}

\begin{proposition}[Finite-degree null-certificate bounds]\label{surf:finite-null-tails}\label{low:78:h3n:wide}\label{low:78:h3n:vertical}\label{low:78:h3n:collar}\label{low:78:s37i:wide}\label{low:78:s37i:vertical}\label{low:78:s37i:collar}\label{low:78:s38o:wide}\label{low:78:s38o:vertical}\label{low:78:s38o:collar}\label{low:78:s38i:wide}\label{low:78:s38i:vertical}\label{low:78:s38i:collar}\label{low:78:d8q:wide}\label{low:78:d8q:vertical}\label{low:78:d8q:collar}\label{low:9:h9d:wide}\label{low:9:h9d:vertical}\label{low:9:u9:collar-tails}
Assume the common full-parent certificate hypotheses above. Wide ribbons have the representation, relative seed ranges, and original point factor in Proposition~\partIref{surf:wide}; vertical ribbons have those in Proposition~\partIref{surf:vertical}. Collars have the prepared germ of Lemma~\partIref{surf:collar-net} and linear part in the stated small neighborhood of the triangular null model, with the additional hypotheses of the two refined nets when used, original point factor $P^CL^2$, and alternative endpoint shear interval $P^CL$.

\noindent\textup{(i)} Then the full $N_k$-dimensional coefficient set has the following bounds, with a factor $P^C$ understood:
\begin{equation}\label{surf:finite-wide-collar}
\begin{array}{c|c|c}
(k,m)&\text{wide ribbon}&\text{fixed-width collar}\\\hline
(7,2)&w^{-231}H^{-279}&L^{-232}\\
(7,3)&w^{-226}H^{-242}&L^{-215}\\
(8,2)&w^{-383}H^{-453}&L^{-373}\\
(8,3)&w^{-370}H^{-396}&L^{-350}\\
(8,4)&w^{-337}H^{-359}&L^{-337}\\
(9,2)&w^{-599}H^{-692}&L^{-567}\\
(9,3)&w^{-574}H^{-615}&L^{-532}\\
(9,4)&w^{-525}H^{-558}&L^{-505}.
\end{array}
\end{equation}
\equationalias{low:78:h3n:wide-tail}

\noindent\textup{(ii)} The vertical bound is the sum of the two chart contributions ($S=LW$, $K_0=L/W$):
\begin{equation}\label{surf:finite-vertical}
\begin{array}{c|c|c}
(k,m)&\text{grid chart}&\text{weighted chart}\\\hline
(7,2)&S^{-447/2}K_0^{9/2}&S^{-223}K_0^{-86}\\
(7,3)&S^{-463/2}K_0^{11/2}&S^{-225}K_0^{-51}\\
(8,2)&S^{-801/2}K_0^{-41/2}&S^{-717/2}K_0^{-283/2}\\
(8,3)&S^{-751/2}K_0^{3/2}&S^{-711/2}K_0^{-171/2}\\
(8,4)&S^{-691/2}K_0^{9/2}&S^{-681/2}K_0^{-109/2}\\
(9,2)&S^{-1229/2}K_0^{-63/2}&S^{-546}K_0^{-217}\\
(9,3)&S^{-577}K_0^{-4}&S^{-540}K_0^{-139}\\
(9,4)&S^{-534}&S^{-515}K_0^{-89}.
\end{array}
\end{equation}
\equationalias{low:78:h3n:vertical-tail}\equationalias{low:78:d8q:vertical-bound}

\noindent\textup{(iii)} The additional joint graph-window hypothesis multiplies the four specified cases by $\nu$; otherwise all bounds are unbanded.
\end{proposition}
\begin{proof}
Multiply the factors in the wide, vertical, and collar tables above. For the vertical conversion use $L=S^{1/2}K_0^{1/2}$ and $W=S^{1/2}K_0^{-1/2}$. The two vertical charts exhaust the slopes, so their estimates add. The finite atlas and fixed-frame disintegration established at the start of the subsection therefore give the stated coefficient volumes on the actual good endpoint sets.
\end{proof}

For $k=7$, the two vertical bounds imply $P^CS^{-218}K_0^{-1}$ and $P^CS^{-225}K_0^{-1}$, respectively. For $(k,m)=(8,3),(8,4)$ they imply $P^CS^{-355}K_0^{-1}$ and $P^CS^{-340}K_0^{-1}$. These comparisons use only $K_0\le S$. The wide cubic-normal bounds in degrees seven and eight similarly imply $P^CS^{-225}K^{-1}$ and $P^CS^{-369}K^{-1}$ for $K=w/H$. The degree-nine quartic grid term may be replaced by $S^{-533}K_0^{-1}$ before summing coordinate aspects.
\paragraph{Ordinary degree-seven consequence}

\begin{corollary}[Ordinary common coefficient estimate]\label{low:78:s7c:ordinary-tail}
The ordinary small-fifth-pivot branch has a common unbanded assignment with count $O_P(\Lambda S)$ and fixed-frame tails
\[
\operatorname{vol}_{119} E_j(S,K)\lesssim_P S^{-435/2}K^{-1/2},\qquad K=L/W\ge1.
\]
This statement includes zero and nonzero rates together. The estimate concerns the critical branch.
\end{corollary}

\begin{proof}
Use the fixed-frame $H_0$ incidence and ordinary degree-seven coefficient atlases. The vertical moderate count, with its rate row integrated over the full bounded interval, is $P^C S^{-218}K_0^{-1}$, where $K_0$ is the coordinate aspect. It also covers fixed-width vertically long horizontal ribbons. Discard $K_0^{-1}$ since $K_0\ge1$. Since $K\le S$ by $W\ge1$, we have $S^{-218}\le S^{-435/2}K^{-1/2}$. The wide count $P^Cw^{-231}H^{-279}\le P^CS^{-231}$ is stronger under the same conversion. So is the collar count $P^CL^{-232}$, whose area is comparable to $L$ up to conditioning factors. Bounded-area caps cost a conditioning constant; unbounded horizontal caps use the collar count. These estimates include the fixed-frame fibers and seed grids. The graph-rate row uses its full bounded interval at each endpoint, without a lower rate bound or rate-dependent construction cost. Finally, $435/2$ exceeds $\chi_7+2=212$.
\end{proof}


\begin{corollary}[Finite-degree assigned pieces]\label{low:78:s38o:common}
On the common unbanded cover, ordinary degree-eight pieces have count $O_P(\Lambda S)$ and fixed-frame tail $P^CS^{-358}K^{-1/2}$. Quartic-normal degree-eight pieces, under the stated common-cover hypotheses, have
\begin{equation}\label{low:78:d8q:assigned-tail}
\operatorname{vol}_{164}E_j(S,K,\nu)\lesssim_P\nu S^{-673/2}K^{-1/2}.
\end{equation}
Here $S$ and $K$ are the assigned geometric area and aspect; the unbanded case has $\nu=1$.
\end{corollary}
\begin{proof}
The ordinary vertical bound is at most $P^CS^{-717/2}$ with summable coordinate-aspect decay; its wide and collar bounds have area exponents $383$ and $373$. In the quartic branch the corresponding area exponents are $340,337,337$. Sum first over coordinate aspects; the wide products have a negative power of the shorter coordinate, and fixed-width collar area is comparable to length. Narrow horizontal enlargements use the vertical bound, bounded caps contribute a conditioning factor, and unbounded horizontal caps use their collars. Finally $K\lesssim_PS$ gives $S^{-a}\lesssim_PS^{-a+1/2}K^{-1/2}$. This proves both assertions. The combination with the other normal and Gauss branches is carried out in Section~\ref{surf:finite-completion}.
\end{proof}

\subsection{Gauss coefficient counts}\label{surf:finite-gauss-counts}\label{low:78:s78gauss:section}\label{low:9:u9:gauss}

\paragraph{The specializations in degrees seven, eight and nine}

For $k=7,8,9$, all relevant normal orders are $2\le m\le\lfloor k/2\rfloor$. The selected graph and pivot are
\begin{equation}\label{low:78:s78gauss:normal-graph}
 \partial_x^{m-1}F(c(y,z),y,z)=0,\qquad \kappa_m=\partial_x^mF|_c\ne0.
\end{equation}
In the curvature chart, $\kappa_m$, $r_0=c_{yy}$, $s_0=\det(D^2c)/r_0$, and their inverses are conditioned. Choosing the complete coefficient defect after this curvature threshold restricts the reference models to
\[
F_0=f(x-Q(y,z))+\ell_y y+\ell_z z,\qquad f^{(m-1)}(0)=0,\qquad |f^{(m)}(0)|\ge P^{-C},
\]
with $Q$ a nonsingular quadratic form. The invariant of degree three has vanishing second derivative in the first tangential direction and belongs to the null chart. Both Hessian signs are allowed: for example, the elliptic graph of $x-(y^2+z^2)/2$ is covered without a real null direction. Near a hyperbolic null long direction, use the grid alternative in Lemma~\partIref{surf:gauss-atlas}.

All hypotheses of the common Gauss estimates hold on these actual full-parent branches. Theorem~\partIref{geo:gauss} gives the contained strips, buffer, count, simultaneous finite derivative preparation, and candidate area $P^CL^2$, with bounded scales absorbed into $P^C$. Lemma~\partIref{surf:null-density}, specifically \partIeqref{surf:gauss-density}, supplies the transverse determinant $|\kappa_m|^4|r_0|^2$ and the identity block on all other $N_k-4$ coefficients. Lemma~\partIref{surf:haar} applies on the original $H=\ker(a,f)$ fiber, using the fixed-frame pivot; there is no additional endpoint shear factor. Lemma~\partIref{surf:gauss-atlas} then supplies the complete grid or weighted map, with graph jets at the fixed physical origin. For the scalar jets, fix the rounded map from strip parameters to $(p,z)$ and compose it with $(p,z)\mapsto(Y_F(p,z),z)$, where $c_y(Y_F(p,z),z)=p$. The external map is fixed during coefficient differentiation, whereas $Y_F$ varies with $F$. Thus all lower parent directions and all moving-graph corrections are included.

Table~\ref{app:table-gauss-rows} records the row counts and fixed-normal powers given by Lemmas~\partIref{surf:normal-filtration} and~\partIref{surf:gauss-atlas}. Replacing block $m-1$ by graph rows contributes $(-1/\kappa_m)^{(k-2m+3)^2}$ at the model. In degree nine the regular Gauss-point branch bounds $(10-m)(9-m)$ are $56,42,30$ for $m=2,3,4$, respectively. These are precisely the normal-block and root-count formulas proved in Paper I, Section~\partIref{surf:gauss-counts}.

\begin{proposition}[Finite Gauss coefficient-volume bounds]\label{low:9:u9:gauss-tails}

\noindent\textup{(i)} Under the actual-strip and fixed-frame hypotheses just verified, the complete coefficient-volume bounds are $P^CS^{-\beta}K^{-\gamma_{\mathrm g}}$ in the grid chart and $P^CS^{-\beta}K^{-\gamma_{\mathrm w}}$ in the weighted chart, with
\begin{equation}\label{low:78:s78gauss:tails}
\begin{array}{c|c|r|r|r}
k&m&\beta&\gamma_{\mathrm g}&\gamma_{\mathrm w}\\\hline
7&2&437/2&33/2&195/2\\
7&3&569/2&33/2&245/2\\
8&2&343&23&151\\
8&3&440&23&187\\
8&4&481&23&203\\
9&2&1033/2&61/2&447/2\\
9&3&1301/2&61/2&545/2\\
9&4&1433/2&61/2&595/2.
\end{array}
\end{equation}

\noindent\textup{(ii)} The weakest area exponents exceed $\chi_k+2$, namely $437/2>212$, $343>332$, and $1033/2>497$, and all aspect exponents are positive.
\end{proposition}
\begin{proof}
Apply Proposition~\partIref{surf:gauss-tail} with $\mathsf a(k)=\chi_k+\tau_k$, $d_g=k-2m+2$, $G=\mathsf a(k)-\mathsf s(d_g)$, and $\Delta=\tau_k-\mathsf t(d_g)$. Its raw scalar pairs are $(a_{\mathrm g},b_{\mathrm g})=(G-k,G-k(k+1))$ and $(a_{\mathrm w},b_{\mathrm w})=(G+\Delta-2k,G-\Delta-k^2)$, respectively; Table~\ref{app:table-gauss-pairs} records their numerical values. The seed-net cardinalities are $P^{C_k}K^k$ for the grid and $P^{C_k}K^{2k}$ for the weighted list. In both cases the candidate area divided by the prepared area is at most $P^{C_k}K$. Thus, for the grid and weighted coefficient sets $E_j^{\mathrm g},E_j^{\mathrm w}$,
\[
\operatorname{vol}E_j^{\mathrm g}\lesssim_P L^{-a_{\mathrm g}}W^{-b_{\mathrm g}}K^{k+1},\qquad \operatorname{vol}E_j^{\mathrm w}\lesssim_P L^{-a_{\mathrm w}}W^{-b_{\mathrm w}}K^{2k+1}.
\]
Therefore $\beta=(a+b)/2$, $\gamma_{\mathrm g}=(a_{\mathrm g}-b_{\mathrm g})/2-k-1$, and $\gamma_{\mathrm w}=(a_{\mathrm w}-b_{\mathrm w})/2-2k-1$, giving \eqref{low:78:s78gauss:tails}. The reference directions have the conditioning cost included in Lemma~\partIref{surf:gauss-atlas}.
\end{proof}

\subsection{Exhaustion, exponent bounds and incidence}\label{surf:finite-completion}\label{low:78:s78:exhaustion}\label{low:78:s78:closing}\label{low:9:root9:completion}

\paragraph{Summation of the finite-degree coordinate tails}\label{surf:finite-tail-summation}
On a common full-parent cover in degrees seven through nine, apply Proposition~\ref{surf:finite-null-tails} before converting to the assigned geometric aspect. On moderate pieces the assigned area is comparable to the ordinary coordinate area, up to powers of $P$. For vertical ribbons, use the summable $K_0^{-1}$ consequences in degrees seven and eight; the degree-nine aspect powers are negative except for the grid term $S^{-534}$, whose allowed logarithmic multiplicity is retained. For wide ribbons,
\[
w^{-c}H^{-d}\lesssim_PS^{-c}H^{-(d-c)},\qquad S\asymp_PwH,
\]
and every row of~\eqref{surf:finite-wide-collar} has $d>c$, so the shorter-scale sum converges. Fixed-width collars have $S\asymp_PL$; bounded caps cost a power of $P$, and unbounded horizontal caps use their collars. Thus the weakest null area powers are
\[
\begin{array}{c|ccc}
k&7&8&9\\\hline
\text{area exponent}&215&337&505
\end{array}
\]
All graph-rate factors here use $\nu=1$. The assigned width gives $K\lesssim_PS$, independently of the ordinary coordinate aspect $K_0$, and hence
\begin{equation}\label{surf:finite-aspect-conversion}
S^{-a}\lesssim_PS^{-a+1/2}K^{-1/2}.
\end{equation}
The conditioning and permitted physical logarithms are unchanged.

\begin{proposition}[Complete common surface estimate]\label{low:78:s78:surface}

\noindent\textup{(i)} Every small-fifth-pivot label in degree $k=7$ or $8$ admits an assignment to pieces with $L\ge W\ge1$, $S=LW$, $K=L/W$, count at most $P^C\Lambda S$, and at most $P^C(1+\log P+\log d_1^{-1})^C$ pieces at each dyadic scale.

\noindent\textup{(ii)} At each original assigned label $j$, the entire assigned-scale coefficient set satisfies
\[
\operatorname{vol}_{N_k}E_j(S,K)\le P^C(1+\log P+\log d_1^{-1})^C S^{-\beta_k}K^{-1/2},\qquad \beta_7=429/2,\quad\beta_8=673/2.
\]
There is no rate index or rate-dependent multiplicity in this statement.
\end{proposition}
\begin{proof}
Apply Lemmas~\ref{low:78:s78:entry} and~\ref{low:78:s78:chart-split}, retaining all original labels. The full-parent local comparison supplies the buffered unbanded core for each null sheet. Theorem~\partIref{geo:common-allocation} assigns it to direct moderate corridors, narrow horizontal ribbons, parabolic collars, and projection caps. Its proof depends only on the null geometry and the finite degree of the full parent phase, so it applies to every sheet just listed. The degree-specific scalar maps are supplied by the ordinary, cubic-normal and quartic-normal sections. Their good-point preparation can be done jointly because it consists of finitely many derivative lists of fixed format.

The preceding coordinate-tail calculation gives the stated null bounds, with the cubic-normal collar decisive in degree seven and the quartic-normal wide and collar bounds decisive in degree eight. The Gauss charts have stronger area exponents $437/2$ and $343$, with stronger aspect decay, by~\eqref{low:78:s78gauss:tails}. Their count and fixed-frame incidence follow from Theorem~\partIref{geo:gauss} and Lemma~\partIref{surf:null-density}. The finite normal, tangential, and branch choices cost only the prescribed conditioning and logarithms; the individual bounds already include their certificate nets. This gives the estimate for the entire assigned coefficient set.
\end{proof}


\paragraph{Common surface tails and completion of the degree-nine sum}
The rank-four approximation and actual cover of Section~\ref{low:9:root9:interfaces} supply the full-parent hypotheses of Proposition~\ref{surf:finite-null-tails}. The common coordinate sum above gives $P^CS^{-505}$, with the allowed logarithmic multiplicity. Equation~\eqref{surf:finite-aspect-conversion} therefore yields
\begin{equation}\label{low:9:root9:common-surface}
\vol E_j(S,K)\lesssim P^CS^{-1009/2}K^{-1/2},\qquad \frac{1009}{2}>497.
\end{equation}
The degree-nine Gauss exponents in Proposition~\ref{low:9:u9:gauss-tails} are larger, with stronger negative aspect powers, so the same bound covers all surfaces. The coefficient measure in \eqref{low:9:root9:common-surface} is the complete normalized $219$-dimensional parent measure. Every original physical anchor contributes its exact factor $v_j^{-495}$. 

These are the surface inputs for $k=7,8,9$ in Table~\ref{an:incidence-inputs}; Section~\ref{an:assembly} combines them with the regular bounds and completes the proof.


\section{Quartic and quintic surface estimates}\label{low:degree45}\label{low:degree4}

The first two subsections concern $k=4$ and use the physical family, augmented rank models, regular assignment and coefficient estimates of Section~\ref{reg:finite4}, with the same notation and conditioning order. They conclude with Theorem~\ref{low:4:sf:surface-completion}. The last two subsections concern $k=5$ and conclude with Theorem~\ref{low:5:qsc:surface-count}. These surface bounds enter Proposition~\ref{an:unified-incidence}, which retains the quartic regular cutoff $H\le P^C/d_2$ and its middle-side loss.

\subsection{Quartic coordinates and characteristic transport}\label{quartic:normalization}\label{low:4:sf:normalization}\label{low:4:sf:sec:packing}\label{low:4:sf:preparation}\label{quartic:coordinates}\label{low:4:sf:sec:gauss}\label{low:4:sf:null}
\par\smallskip\noindent\textit{Step 1. Normalized frames and projected packing.}\par\nobreak\smallskip

Use the square-quadric family of Corollary~\partIref{geo:rank-four-entry-low} and the nested admissibility windows of Definition~\partIref{def:admissibility}, in the complete nonconstant coefficient space. The outer window retains the transported physical matrix, center and stopping inequalities.

For ordinary square-quadric surfaces in degrees four through six, set $G=F_x$ and use $\mathcal U_0=\{e=f=0\}$. Its right action fixes the tangential center, and
\begin{equation}\label{low:4:sf:eq:covariance}
(F\circ g)_x=G\circ g.
\end{equation}
\equationalias{low:6:qsf:eq:covariance}
At the model $J=a(x-Q)^2+\lambda y+\mu z$, the normal root has derivative $2a$ and tangent Hessian $-2aD^2Q$. Lemma~\partIref{geo:rank-entry} therefore supplies a map $k\in\mathcal U_0$ such that $h=g\circ k$ is normalized and the affine matrices of $k,k^{-1}$ are bounded by $P^C$ throughout a sufficiently small complete coefficient neighborhood. Apply Lemma~\partIref{geo:two-chart-shear} to its conditioned tangent Hessian. A normalized frame $h=gk$ satisfies
\begin{equation}\label{low:4:sf:eq:normalized}
\widetilde G(0)=\widetilde G_y(0)=\widetilde G_z(0)=0,\qquad \widetilde G=G\circ h,
\end{equation}
\equationalias{low:6:qsf:eq:normalized}
and either chart I, $\widetilde G_{yz}(0)=0$ with $\widetilde G_{yy}(0)\ne0$, or chart II, $\widetilde G_{zz}(0)=0$ with $\widetilde G_{yz}(0)\ne0$. Chart II occurs only for an indefinite Hessian. The coefficient vectors $[F\circ h]$ and $\widetilde G$ have norm at most $P^C$. Also $|\lambda|=|\widetilde G_x(0)|\ge P^{-C}$, the tangent Hessian and its inverse have norm at most $P^C$, and the absolute value of the selected scalar pivot is at least $P^{-C}$.

Lemma~\partIref{atlas:normal-fibers} bounds the number of normalized frames at a fixed $(e,f)$ by $3(k-1)$. Its transverse Jacobian in $(c,a,b,d)$ is
\begin{align}
\mathrm{I}:&\quad |\lambda|^3|\widetilde G_{yy}(0)|,\label{low:4:sf:eq:jacI}\\
\mathrm{II}:&\quad 2|\lambda|^3|\widetilde G_{yz}(0)|.\label{low:4:sf:eq:jacII}
\end{align}
Table~\ref{app:table-normal-frames} records the specializations of the bound $3(k-1)$ and the coefficient slice dimension $N_k-4$ for $k=4,5,6$. All these normalizations use the actual equation $F_x=0$.

Each active frame lies in one of at most $3(k-1)$ sets $h\mathcal B_0$, with $\mathcal B_0$ a fixed bounded neighborhood in $\mathcal U_0$. Left multiplication preserves the fiber measure $da\,db\,dc\,dd$. Consequently, for $k=4,5,6$,
\begin{equation}\label{low:4:sf:eq:fibervol}
\sup_{(e,f)}\vol_4\{(a,b,c,d):[F\circ g]\text{ is in the specified rank-four neighborhood}\}\le P^C.
\end{equation}
\equationalias{low:5:qsf:eq:fibervol}\equationalias{low:6:qsf:eq:fibervol}

Let $\mathcal N_F^+$ be the outer set of normalized frames from Definition~\partIref{def:admissibility}, retaining the bounds on $[F\circ h]$, its physical data, the inverse pivots, and each permitted normal-root branch. Put $E_F^+=\{(e(h),f(h)):h\in\mathcal N_F^+\}$. The preceding coefficient and pivot bounds give normalized solutions in relative coordinates on an inverse-power neighborhood. Choose its radius below the fixed normalized inner-to-outer margins, so the solutions remain in $\mathcal N_F^+$. The fixed family satisfies \partIeqref{an:local-packing}, with parameter $\Lambda$, after the transfer in Lemma~\partIref{an:rq}. Proposition~\partIref{atlas:projected-packing} therefore gives
\begin{equation}\label{low:4:sf:eq:area}
N(F)\le P^C\Lambda|E_F^+|.
\end{equation}
The same convention for $E_F^+$ and $\Lambda$ is used in the later ordinary-surface applications.

Implicit roots retain a specified regular branch. The number of connected components, isolated fibers and monotonicity intervals is consequently bounded by a constant depending only on this format.

Apply Lemma~\partIref{geo:good-points-strips} to the selected regular branches and retained inverse pivots. Its common contained strips and simultaneous derivative bounds will be used for the following scalar lists.

If a critical value $\phi$ has bounded gradient on such a strip, apply the derivative lemma to its two first derivatives in strip coordinates. At good points this gives
\begin{equation}\label{low:4:sf:critical-widths}
 |\partial_u^i\partial_v^j\phi|\lesssim_P L^{-i}W^{1-j}\ (j\ge1),\qquad |\partial_u^i\phi|\lesssim_P L^{1-i}\ (i\ge1).
\end{equation}
\equationalias{low:5:qsf:critical-widths}\equationalias{low:6:qsf:critical-widths}
At a fixed evaluation point, round the $j$th centerline derivative in steps $c_PW/L^j$. Its range $O_P(L^{1-j})$ costs $O_P(K)$ choices, where $K=L/W$. This gives $K^8$ for orders through eight and $K^4$ for orders through four. The evaluation point is integrated. The coordinate change only shifts the transverse variable; its mixed derivative rows cannot acquire a purely longitudinal derivative without a transverse factor. The chain rule therefore preserves \eqref{low:4:sf:critical-widths}. We include every ordinary derivative through order eight in the good-point set, including rows that occur only after rounding.

Coefficient-volume estimates below use Lemma~\partIref{geo:coefficient-map}, with the conditional coordinates retained in the incidence map.

\paragraph{Step 2. Nonnull surfaces and intrinsic null coordinates}
Use the ordinary ($m=2$) Gauss frame~\partIeqref{geo:gauss-frame} on a selected actual branch $x=c_F(y,z)$ of $F_x=0$. For the ordinary surface applications write $\lambda=F_{xx}|_{\rm graph}$, $r=c_{yy}$, $s=\det(D^2c)/r$, $p=c_y$, $Y=Y_F(p,z)$, $q_0=c_z(Y,z)$ and $d=-c_{yz}/r$. The pivots $\lambda,r,s$ and their inverses are $P$-bounded. The complete normalized $J=F\circ h_{p,z}$ lies near the stated compact square-quadric family, with $J_x=J_{xy}=J_{xz}=J_{xyz}=0$. Its absolute shear $d$ need not be bounded.

\paragraph{Intrinsic coordinates and an assignment}
Define the actual Gauss critical value
\[
\Phi_F(p,z)=F(c_F(Y_F(p,z),z),Y_F(p,z),z).
\]
Write $\phi(y,z)=F(c_F(y,z),y,z)$. Since $Y_p=1/r$ and $Y_z=d$, the normalized coefficient window gives
\[
\Phi_p=\phi_y/r,\qquad \Phi_z=\phi_z+d\phi_y,
\qquad |\nabla\Phi|\lesssim_P1.
\]
The map $\pi_{\mathrm G}(g)=(a,f)$ is a homomorphism with normal kernel
\[
H=\ker\pi_{\mathrm G},\qquad \mathfrak h=\operatorname{span}(C,T_x,D,T_y).
\]
For $g=(a,b,c_0,d_0,e,f)$, the graph transforms by $c_{F\circ g}(y,z)=c_F(y+d_0z+e,z+f)-ay-bz-c_0$. Its $y$ slope is the original slope minus $a$, giving the covariance
\begin{equation}\label{low:4:sf:eq:gausscov}
\Phi_{F\circ g}(p,z)=\Phi_F(p+a,z+f)\pmod{\R}.
\end{equation}

At fixed $(p,z)$, apply the root count in the proof of Theorem~\partIref{geo:gauss} with $m=2$ and $G=F_x$. There are at most $(k-1)(k-2)$ regular solutions of $G=G_y+pG_x=0$. The remaining coordinates are $a=p$, $f=z$, $e=Y$, $c_0=c_F(Y,z)$, $d=-c_{yz}/r$ and $b=c_z+dp$. Retain all regular roots in the fixed-format admissibility condition.

Apply Theorem~\partIref{geo:gauss} with $m=2$, the present normal and Gauss pivots, and the preceding degree-specific root count. Its frame differential~\partIeqref{geo:gauss-body} and the bounded normalized third graph derivatives control relative frames independently of the absolute shear. The same intrinsic-square packing and common strip assignment give count $O_P(\Lambda LW)$, contained strips and centerline bounds $|\zeta^{(j)}|\lesssim_PL^{1-j}$. The inverse-polynomial lower width used here is absorbed by a fixed conditioning rescaling of the common local units.

The common theorem handles all regular root branches, including artificial cuts. Apply Lemma~\partIref{geo:good-points-strips} to the selected normalized fields $\Phi_p,\Phi_z$, with their fixed-format switching sets retained. The good endpoint and fixed frame belong to the same polynomial but need not carry the same root label. The covariance~\eqref{low:4:sf:eq:gausscov} reanchors the Gauss plane by translation. The parent coefficient and physical restrictions remain in the common Borel admissibility condition.

\paragraph{Step 3. Joint coefficient and frame integration}
For the ordinary Gauss normalization, Lemma~\partIref{surf:null-density} gives
\[
dF=C_k|\lambda|^4|r|^2\,dJ_{N_k-4}\,dh_H,
\]
where $C_k>0$ is the fixed Taylor-coordinate factor and the slice dimensions are listed in Section~\ref{quartic:normalization}. The four rows $(J_x,J_{xy},J_{xz},J_{xyz})$ on $(T_x,T_y,C,D)$ have diagonal $(\lambda,-\lambda r,\lambda,-\lambda r)$. Thus the original Gauss anchor has a conditioned $H$ pivot. Bounded right changes preserve it, since $H$ is normal. A full intrinsic slice map with Jacobian at least $P^{-C}$ has bounded inverse multiplicity by Lemma~\partIref{geo:coefficient-map}; conditioning on its outputs leaves the $O_P(1)$ Haar integral of Lemma~\partIref{surf:haar}.
\paragraph{Restriction determinants}
We record the finite-dimensional computation used in all the surface coefficient maps. Put $Q(u,v)=Au^2+Buv+Cv^2$. After the four slice rows and six normal rows have been removed, the $24$ input columns restrict to
\[
 \Pi_{\le4}(u,v)/\mathbb R,\quad Q\Pi_3^{\rm hom},\quad Q^2\Pi_2^{\rm hom},\quad Q^3\Pi_1^{\rm hom},\quad Q^4.
\]
Their ordinary-degree blocks have dimensions $2,3,4,5,4,3,2,1$. The weighted row set is
\[
 \mathcal I_{\rm w}=\{(i,j):0<i+2j\le8\},
\]
and the rectangular row set is
\[
 \mathcal I_{\rm g}=\{(i,j):0\le i,j\le4,\ (i,j)\ne(0,0)\}.
\]
Both sets have $24$ elements. Apply the restriction calculation in the proof of Lemma~\partIref{surf:normal-filtration} with degree four, omitting the constant row and column. Its ordinary-degree blocks have determinants $1$ through degree four, followed by $A^4,A^6,A^6,A^4$ in the weighted chart and $B^4,B^6,B^6,B^4$ in the grid chart at $A=0$ or $C=0$. Thus the determinants are $A^{20}$ and $B^{20}$, up to row-ordering sign. Higher graph terms raise ordinary degree. The effects of a prescribed nonlinear curve and of varying the critical point are accounted for below.

\paragraph{Step 4. The thirty intrinsic outputs}
Take the six normal coefficients
\[
x^2,\ xy^2,\ xz^2,\ x^2y,\ x^2z,\ x^3
\]
and 24 Gauss critical Taylor coefficients. The derivative at fixed Gauss coordinates is
\begin{equation}\label{low:4:sf:eq:gaussvariation}
D_F\Phi_F[R]
=R|_{\rm graph}+
\frac{\phi_y}{r}\,\partial_y\bigl((R_x/\lambda)|_{\rm graph}\bigr),
\end{equation}
evaluated in Gauss coordinates. To derive the formula, differentiate $F_x(c(y,z),y,z)=0$ at fixed $(y,z)$, obtaining $\delta c=-(R_x/\lambda)|_{\rm graph}$. Holding $p=c_y$ fixed gives $\delta Y=-\delta c_y/r$. The envelope identity gives $\delta\phi=R|_{\rm graph}$, so $\delta\Phi=R|_{\rm graph}+\phi_y\delta Y$, which is the displayed expression. 

At a normalized rank-four model, $\lambda,r$ are constant and $\phi_y=\ell_y$. After removing the four slice constraints and six normal rows, the 24 input directions are the fourteen nonconstant binary monomials through degree four, followed by $x\Pi_3^{\rm hom},x^2\Pi_2^{\rm hom},x^3\Pi_1^{\rm hom},x^4$. If restriction of a column has tangential degree $m$, the correction in \eqref{low:4:sf:eq:gaussvariation} has degree $m-3$. The restriction matrix is block diagonal at the model and the correction lies strictly above its ordinary-degree blocks. It does not change the determinant, even when $\ell_y\ne0$.

In bounded intrinsic tangent/transverse coordinates let the effective quadratic form be $A_0u^2+B_0uv+C_0v^2$. Fix a sufficiently small cutoff $\delta_{\rm q}=P^{-C}$, depending only on the current conditioning bounds. If $|A_0|\ge\delta_{\rm q}$, use the 24 indices $0<i+2j\le8$; the determinant is a bounded nonzero factor times $A_0^{20}$. If $|A_0|<\delta_{\rm q}$, nonsingularity gives $|B_0|\ge P^{-C}$; use all $0\le i,j\le4$ except $(0,0)$, whose determinant at $A_0=0$ is a bounded nonzero factor times $B_0^{20}$. Choose $\delta_{\rm q}$ small enough that this grid determinant remains bounded away from zero on the latter branch. Smooth algebraic dependence of the finite implicit jets and compact continuity extend these lower bounds to an actual model neighborhood chosen after $\delta_{\rm q}$.

A fixed nonlinear higher curve jet preserves the first weighted row space triangularly. The second bidegree row space is not invariant under arbitrary nonlinear substitution. For it, use the actual bounds $\zeta^{(j)}=O_P(L^{1-j})$: higher derivatives are small for $L$ above a fixed cutoff, so compact continuity applies. Bounded $L$ has bounded count.

\begin{theorem}[Common unbounded-shear surface count]\label{low:4:sf:thm:gauss}

\noindent\textup{(i)} In the stated nonnull chart, all actual common components have coefficient-volume bounds at the fixed frames
\[
O_P((LW)^{-38}(L/W)^{-15})
\quad\text{or}\quad
O_P((LW)^{-38}(L/W)^{-3}).
\]
In particular the summed tail is $O_P((LW)^{-38})$, without a bound on tangential shear or on the graph slopes in the original coordinates.

\noindent\textup{(ii)} Common weighted summation holds for $1\le p<39$, with the fixed-frame factor $v_j^{p-35}$.
\end{theorem}
\begin{proof}
Use~\eqref{low:4:sf:critical-widths} on good Gauss area $\gtrsim_PLW$, preparing all ordinary derivatives through order eight. The centerline rounding in Section~\ref{quartic:normalization} costs $K^8$ for the weighted rows and $K^4$ for the grid rows, where $K=L/W$.

The products of the 24 critical widths are $L^{-62}W^{-14}$ and $L^{-46}W^{-30}$. All six normal outputs may retain unit-sized windows. The possible evaluation points lie in an anchor disk of area $O_P(L^2)$, while each phase supplies good area $\gtrsim_PLW$; the candidate area divided by the good area is at most $P^CK$. Thus
\[
L^{-62}W^{-14}K^9=(LW)^{-38}K^{-15},
\qquad
L^{-46}W^{-30}K^5=(LW)^{-38}K^{-3}.
\]
Each piece has at most $P^C\Lambda LW$ assigned original labels. Summable component weights and dyadic summation work whenever $(1+\epsilon)(p-1)<38$. The physical variables remain fixed restrictions throughout. At the original fixed frame $T_j$, Lemma~\partIref{an:character} gives $dF=v_j^{-35}dJ$. Combined with the weight $v_j^p$, this gives $v_j^{p-35}$ in Proposition~\partIref{an:common-incidence}. Using the fixed representative $T_j^0$ changes this factor by only a fixed constant, by Lemma~\partIref{an:rq}.
\end{proof}

In the second normalization chart the first curvature may vanish. We retain the chosen null line as a two-sheet branch label. This branch label is continued along the component.

\paragraph{A uniform coefficient map through zero first curvature}
\begin{theorem}[Intrinsic null-jet coordinates]\label{low:4:sf:thm:nulljets}
On the null slice $J_x=J_{xy}=J_{xz}=J_{xzz}=0$ at the origin, let $x=c(y,z)$ be the actual branch $J_x=0$, with $\lambda=J_{xx}(0)\ne0$, $r=c_{yy}(0)$ and $B=c_{yz}(0)\ne0$. Use the selected null root with $d(0)=0$ and the fields of Definition~\partIref{def:null-coordinates}, and write $\phi=J|_c$.

\noindent\textup{(i)} The six normal coefficient outputs $x^2,xy^2,xyz,x^2y,x^2z,x^3$, together with the 24 ordered jets
\[
\frac{(Y^iZ^j\phi)(0)}{i!j!},\qquad 0\le i,j\le4,\quad (i,j)\ne(0,0),
\]
form a uniformly nonsingular 30-coordinate map in a sufficiently small fixed neighborhood of the compact null model family.

\noindent\textup{(ii)} Its local inverse bounds do not require division by $r$.
\end{theorem}
\begin{proof}
The critical-graph and null-root implicit functions use denominators only $\lambda$ and $2B$. Thus every finite jet above depends smoothly and algebraically on the original coefficients, uniformly through $r=0$.

At $J_0=a(x-Q)^2+\ell_y y+\ell_z z$, with $Q=r y^2/2+B yz$, the fields are $Y=\partial_y$ and $Z=\partial_z$. After eliminating the four slice equations and six normal outputs, use the same 24 input monomials as in Section~\ref{low:4:sf:sec:gauss}. For a remaining perturbation $R$ whose restriction to $x=Q$ has tangential degree $m$,
\[
\delta\phi=R|_Q,\qquad
\delta c=-(R_x/\lambda)|_Q,\qquad
\delta d=-(\delta c)_{zz}/(2B).
\]
For $j\ge1$, differentiation of every occurrence of $Z$ gives exactly
\[
\delta(Y^iZ^j\phi)(0)=
\partial_y^i\partial_z^j(R|_Q)(0)
+\ell_y\partial_y^i\partial_z^{j-1}(\delta d)(0).
\]
All other operator-variation terms vanish because $\ell$ is affine. The correction enters only degree $m-3$. It is therefore strictly above the ordinary-degree restriction blocks. In consistent Taylor and input coordinates their determinant is exactly $B^{20}$, for arbitrary bounded $r$ and $\ell_y$. The six normal rows are identity rows. Compact continuity proves the actual nearby assertion. The required neighborhood radius is $P^{-C}$ because all finite jets are algebraic with denominators bounded below by powers of $P^{-1}$.
\end{proof}

The anisotropic bounds for the ordered jets in Theorem~\ref{low:4:sf:thm:nulljets} follow from the geometric argument below, with the selected null branch retained as above.

\begin{lemma}[The full null coefficient--point incidence]\label{low:4:sf:null-incidence}
On the slice $J_x=J_{xy}=J_{xz}=J_{xzz}=0$, let $h=g(a,b,c,d,e,f)$ be the actual normalized endpoint frame and put $\eta=e-df$.

\noindent\textup{(i)} Up to the fixed Taylor-coordinate convention,
\begin{equation}\label{low:4:sf:null-density}
dF\,de\,df=C|\lambda^4B|\,dJ_{30}\,da\,db\,dc\,dd\,de\,df
 =C|\lambda^4B|\,dJ_{30}\,dh_{H_0}\,d\eta\,df,
\end{equation}
where $H_0=\{e=f=0\}$ has generators $A,C,T_x,D$.

\noindent\textup{(ii)} For a full $30$-output slice map with Jacobian of absolute value at least $P^{-C}$ and the fixed-format algebraic branches of Lemma~\partIref{geo:coefficient-map} that depends only on the normalized endpoint phase $J$, fixing its outputs and $(\eta,f)$ leaves a conditional $H_0$ integral bounded by $P^C$ under the fixed frame window.

\noindent\textup{(iii)} An ordinary-coordinate output map may instead hold $d,e,f$ fixed in the coefficient map and use only $H_3=\{d=e=f=0\}$.
\end{lemma}
\begin{proof}
Apply Lemma~\partIref{surf:null-density} with $k=4$ and $m=2$; its factor $2|\lambda|^4|B|$ gives the displayed density after absorbing the fixed factor into $C$. The fixed-frame $H_0$ pivot is supplied by Lemma~\partIref{geo:original-subgroup-pivots}, including $r=0$. Fixing the full $30$ outputs leaves boundedly many endpoint phases by Lemma~\partIref{geo:coefficient-map}. The factorization and Haar estimate in Lemma~\partIref{surf:null-density} then give the $H_0$ conclusion, or its $H_3$ alternative when the endpoint shear is external. The substitution $e=\eta+df$ is made in the full joint incidence identity.
\end{proof}

\paragraph{An exact quartic identity controls the accumulated connection}
In null normalization write $r=c_{yy}$, $s=c_{yz}+dr$ with $|s|\gtrsim_P1$, and impose $c_{zz}+2dc_{yz}+d^2r=0$. With
\[
\theta=dy-d\,dz,\qquad dd=u\theta+v\,dz,
\]
implicit differentiation gives
\[
(d\partial_y+\partial_z)s=-2su+rv,
\qquad
d(\sqrt{|s|}\,\theta)
=-\frac{\operatorname{sgn}(s)rv}{2\sqrt{|s|}}\,\theta\wedge dz.
\]
The exact $r=0$ chart therefore has a closed coframe. The accumulated term $rv$ at small nonzero $r$ can now be controlled by a new degree-four identity. In the following statements, $B$ denotes the normalized mixed curvature $s$ at the moving point.

\begin{lemma}[Quartic connection identity]\label{low:4:sf:lem:connection}
For a quartic $K$ let $[m]K$ denote its coefficient of the monomial $m$, and define
\[
\mathcal L_B(K)=[yz^3]K-\frac B3[xz^2]K,
\qquad
\Delta(J)=2[y^2z^2]J-\frac B3[xyz]J.
\]

\noindent\textup{(i)} Then $\mathcal L_B((BA+T_z)J)=0$ for every quartic $J$, and $\mathcal L_B(DJ)=\Delta(J)$. At the compact null model, $\Delta=(8/3)aB^2$, so it is uniformly nonzero nearby.

\noindent\textup{(ii)} Along a null characteristic of the actual normalized surface,
\[
ZJ=(BA+T_z)J+vDJ,\qquad
v=\frac{\mathcal L_B(ZJ)}{\Delta(J)}.
\]

\noindent\textup{(iii)} On the null slice $[xz^2]J=0$ the last identity simplifies to $v=Z([yz^3]J)/\Delta(J)$.
\end{lemma}
\begin{proof}
Direct coefficient extraction gives $[yz^3](BAJ)=B[xz^3]J$, $[xz^2](T_zJ)=3[xz^3]J$, and the other contributions to their difference vanish by the degree-four restriction. Also $[yz^3](DJ)=2[y^2z^2]J$ and $[xz^2](DJ)=[xyz]J$. The normalized frame along $Z$ has body velocity $BA+vD+T_z$, giving the stated phase equation. Substitution of $a(x-r y^2/2-B yz)^2+\ell$ gives $\Delta=(8/3)aB^2$.
\end{proof}

Apply Lemma~\partIref{geo:characteristic-preparation} to $Z=\partial_z+d\partial_y$ on each regular branch, retaining its nonzero pivots and branch boundaries in the stopping relation. It gives bounded total variation for the bounded algebraic scalars used below.

\begin{theorem}[Uniform characteristic transport]\label{low:4:sf:thm:transport}

\noindent\textup{(i)} On each connected active null characteristic, with the normalized phase in the stated compact neighborhood and a consistently chosen null branch,
\[
\int |v|\,dz\lesssim_P1,\qquad
\int |u|\,dz\lesssim_P1.
\]

\noindent\textup{(ii)} Consequently the change of the frame shear along the characteristic is bounded, the transverse flow Jacobian is uniformly comparable to one, and the accumulated connection $\int |rv/(2B)|\,dz$ is bounded independently of the length of the characteristic.
\end{theorem}
\begin{proof}
Every normalized coefficient of $J$ is a bounded algebraic function of the original surface point, with fixed format. Lemmas~\ref{low:4:sf:lem:connection} and~\partIref{geo:characteristic-preparation} therefore bound $\int|v|$. The same variation argument applies to $B$. The identity $ZB=-2Bu+rv$ gives
\[
u=-\frac12 Z\log|B|+\frac r{2B}v.
\]
Since $B,B^{-1},r$ are bounded, integration bounds $\int|u|$. For a flow $y=\Psi(s,z)$ from a horizontal transversal, $\partial_z\partial_s\Psi=u\partial_s\Psi$. The signed identity gives
\[
 \partial_s\Psi(s,z)=\left|\frac{B(s,z_0)}{B(s,z)}\right|^{1/2}\exp\left(\frac12\int_{z_0}^z\frac rBv\,dt\right).
\]
Choose the null cutoff $\sup|r|\le P^{-C}$ after the bound for $\int|v|$ so that the exponent has absolute value at most one. Both the flow Jacobian and its inverse are then bounded by powers of $P$. This proves the asserted quantitative comparison; exponentiating only the unsigned $P^C$ estimate for $\int|u|$ would not give it. Finally $v=Zd$ gives the bounded shear variation.
\end{proof}

\paragraph{A geometric assignment for the remaining null region}
The null geometry can be prepared without treating characteristic flow maps as semialgebraic. The construction below is combined with the complementary rescaled counts in Section~\ref{low:4:sf:sec:newrescaling}.

\begin{definition}[Width function]\label{def:width-function}
Let $E$ be an outer admissible set on a fixed regular critical and selected null branch, including the full normalized phase and physical restrictions, with its closure in a slightly larger regular domain. For $p=(y_0,z_0)$ define the fixed-shear square
\[
Q_p(t)=\{(y_0+s+d(p)w,z_0+w):|s|,|w|<t\},
\]
and, for $p\in E$, set $R_E(p)=\sup\{t\ge0:Q_p(t)\subset E\}$. Put $M_p(t)=\max_{q\in\overline{Q_p(t)}}|d(q)-d(p)|$ and $\rho(p)=\max_{0\le t\le R_E(p)}t/(1+M_p(t))$. Closed squares may be taken in the slightly enlarged regular domain. For an actual fixed physical frame the admissibility is bounded by the physical-center restrictions; alternatively, every construction can first be restricted to a finite inner window.
The function $\rho$ is the width surrogate used in this construction. A balanced box at width $W$ has comparable horizontal and null-coordinate widths, up to powers of $P$, and bounded shear oscillation as in the containing sets below.
\end{definition}
\begin{proposition}[Common ribbons under characteristic transport]\label{low:4:sf:common-ribbons}\label{low:5:qng:ribbons}
Consider a regular critical/null admissible set of fixed algebraic format, with conditioned complete phase and curvature windows, a consistently chosen null branch $Z=\partial_z+d\partial_y$, at most $P^C$ normalized frames above each tangential center, and the local projected packing bound $P^C\Lambda$, meaning that each square $Q_p(1)$ with bounded relative normalized frames contains at most $P^C\Lambda$ original labels, counted with multiplicity. Suppose every connected active characteristic has shear variation at most $P^C$, and its transverse flow Jacobian and inverse are bounded by powers of $P$. Retain the nested outer domain and the inner buffer of Lemma~\partIref{geo:core}.

\noindent\textup{(i)} The inner labels admit an assignment to pieces with a contained algebraic ruled ribbon and
\[
\begin{gathered}
L\ge W,\qquad S=LW,\qquad K=L/W,\\
\text{assigned-label count}\le P^C\Lambda S,\qquad \text{ribbon area}\ge P^{-C}S.
\end{gathered}
\]

\noindent\textup{(ii)} Relative to every assigned anchor, the endpoint ranges are
\[
|d|\le P^CK,\qquad |e|\le P^CLK,\qquad |f|\le P^CL.
\]

\noindent\textup{(iii)} The ribbon shear differs from its centerline shear by at most $P^C$, and its coframe and inverse are bounded by powers of $P$. There are boundedly many assigned components and oriented arcs at each width.

\noindent\textup{(iv)} For every prescribed polynomial lower width $M$, a sufficiently large inner buffer and then a sufficiently small coefficient defect give $W\ge M$.
\end{proposition}
\begin{proof}
Use the outer set and width function of Definition~\ref{def:width-function}. Compact maximization and quantifier elimination show that $\rho$ is a continuous semialgebraic function of fixed format.

Write $D_E(p)$ for intrinsic distance to the boundary in the metric $\theta^2+dz^2$. While a path stays in $Q_p(t)$, one has $|dy-d(p)dz|\le |\theta|+M_p(t)|dz|$. A path that exits the square therefore has intrinsic length at least $c\,t/(1+M_p(t))$. This proves $\rho(p)\lesssim D_E(p)$.

For the converse, put $D=D_E(p)$. Start null characteristics at a horizontal $Y$ interval of length $c_PD$ through $p$, and run them for $|z-z_0|\le c_PD$. Such paths first travel along $Y$ and then along $Z$, so they stay inside the intrinsic $D$ ball. The assumed signed Jacobian bound makes their transverse separation comparable, by powers of $P$, to their initial separation. Their shear changes individually by at most $P^C$. For two such characteristics, integrate $y'=d$ from $z_0$ to $z$ to obtain
\[
y_2(z)-y_1(z)=(y_2(z_0)-y_1(z_0))+(z-z_0)(d_2(z_0)-d_1(z_0))+O_P(|z-z_0|).
\]
At a time $|z-z_0|\asymp_PD$, both the initial and final separations are $O_P(D)$; hence the initial shear difference is $O_P(1)$. Thus the entire characteristic rectangle has bounded shear oscillation. Its transverse Jacobian bounds imply that its central portion contains $Q_p(c_PD)$. Therefore $D_E(p)\le P^C\rho(p)$. All constants remain polynomial because the flow estimate used the signed compatibility formula. The inner-buffer lemma makes $\rho$ at selected states at least an arbitrarily large prescribed polynomial constant, after increasing its buffer by a fixed polynomial factor.

Assign each state its dyadic band $W\le\rho<2W$ and its component $C$ of the enlarged open band $W/4<\rho<4W$. Fixed semialgebraic format bounds the number of these components at each $W$. Since $D_E\asymp_P\rho$, every point of $C$ has a path of length $P^CW$ toward the outer boundary; stop it at its first exit from $C$. Before that exit $\rho\ge W/4$, so the balanced-box comparison and signed transport put this access path in a containing set of the form
\[
\rho\ge P^{-C}W,\quad |z-z_0|\le P^CW,\quad |y-y_0-d_0(z-z_0)|\le P^CW,\quad |d-d_0|\le P^C.
\]
Its endpoints are therefore connected in the relevant component of this fixed-format semialgebraic set. Lemma~\partIref{geo:paths} replaces the access path by a piecewise algebraic path there. Coordinate-monotone subdivision and the displayed ranges give intrinsic length $P^CW$, with the same buffer and bounded shear oscillation; no semialgebraicity of intrinsic distance is asserted.

Split the boundedly many boundary arcs and vertices at singularities and tangent-chart changes. For $p=(y_0,z_0)\in C$, put $d_0=d(p)$ and let $S_p^+\subset E$ be the strict enlargement of the preceding containing set:
\[
S_p^+=\left\{(y,z)\in E:
\begin{gathered}
\rho(y,z)>\tfrac12P^{-C}W,\quad |z-z_0|<2P^CW,\\
|y-y_0-d_0(z-z_0)|<2P^CW,\quad |d(y,z)-d_0|<2P^C
\end{gathered}\right\}.
\]
The constants are those of the containing bounds, fixed before this enlargement. For each boundary arc or vertex $\Gamma$, define
\[
U_\Gamma=\{p\in C:\text{some }b\in\Gamma\text{ lies in the same connected component of }S_p^+\text{ as }p\}.
\]
The access paths above have a strict margin in $S_p^+$, so the sets $U_\Gamma$ cover $C$. They are relatively open: a compact path from $p$ to such a point $b$ retains its strict inequalities when $p$ is perturbed slightly, and a short segment joins the new point to $p$ in the same regular branch. Component membership in this fixed-format semialgebraic family is semialgebraic, by the parameter-compatible trivialization used in Lemma~\partIref{geo:conditioning}. Hence the $U_\Gamma$ and their connected components have uniformly bounded format and number. For any $p\in U_\Gamma$, Lemma~\partIref{geo:paths} gives a fixed-format path to $\Gamma$ inside $S_p^+$. Subdivide it into monotone pieces in $z$ and $y-y_0-d_0(z-z_0)$. The displayed bounds then give coframe length at most $P^CW$.

The overlap graph of these connected relative-open sets is connected because they cover the connected set $C$. Its number of vertices is bounded independently of $L/W$. Choose a spanning tree. At a point in an overlap, concatenate the two access paths to the corresponding boundary arcs or vertices; this gives an algebraic bridge of coframe length at most $P^CW$. Thus only boundedly many bridges are needed. The balanced boxes covering a long boundary arc are not vertices of this finite graph.

The resulting common center curve has boundedly many full oriented algebraic arcs, including the bridges. Split them into the $Y$ and $Z$ tangent charts, and split the $Y$ chart at the fixed threshold $|dz/ds|=c_P$, where $s$ is its $\theta$-flux parameter. For the quartic application choose $c_P=P^{-C}$, decreasing it polynomially if reconstruction requires, and then choose $\delta_{\rm null}$ so that $\delta_{\rm null}/c_P$ is less than a fixed small multiple of the lower bound for $|B|$. Assign each state to the first full oriented arc reached by a buffered bridge; endpoints use balanced caps, namely endpoint patches at the local width $W$. Give the full arc its dyadic $L$, comparable to the larger of its intrinsic length and $W$. A later retained subarc is reached from every assigned fixed frame by its initial bridge and then along this full arc, at cost $P^CL$. Derivative preparation therefore changes neither the assignment nor its original $L$.

Chaining balanced boxes along the arc and its bridge gives, relative to every assigned anchor, at most $P^C L/W$ boxes. Shear oscillation is $O_P(1)$ in each box, while the $z$ increment is $O_P(W)$ and the $\theta$ increment is $O_P(W)$. Consequently $|d|\le P^CK$, $|f|\le P^CL$, and $|e|\le P^CLK$. The same cover has area $O_P(LW)$. Each balanced box is covered by $O_P(W^2)$ intrinsic unit boxes with bounded relative normalized frames, so the projected-square count assumed in the proposition bounds its original-label count, with multiplicity, by $P^C\Lambda W^2$. Summation gives $P^C\Lambda LW$. This also follows directly from projected-area packing, with the assumed finite normal-fiber bound.

For a full arc of length at least $W$, use finite derivative preparation to retain a subarc of comparable length, in the respective $Y$- or $Z$-oriented parameter
\[
\theta(\gamma')=1,\ z'=m,\ |m|\le1,\qquad\text{or}\qquad z'=1,\ \theta(\gamma')=m,\ |m|\le1.
\]
Every fixed derivative in either parameter is an algebraic function of the original algebraic curve parameter divided by a nonzero algebraic flux. Thus the same finite-difference preparation gives $|m^{(j)}|\le P^CL^{-j}$ and $|d_\gamma^{(j)}|\le P^CKL^{-j}$ to any prescribed fixed order, including every order needed for the coefficient maps. The two corresponding ruled ribbons are
\[
\Psi(s,t)=(\gamma_y(s)+d_\gamma(s)t,\gamma_z(s)+t),\qquad\text{or}\qquad\Psi(s,t)=(\gamma_y(s)+t,\gamma_z(s)),\quad |t|<c_PW.
\]
The interior margin of the admissible set and bounded shear oscillation in balanced boxes keep these ribbons inside $E$. The first Jacobian is $1+t d_\gamma'$, bounded away from zero after fixing $c_P$; the second Jacobian has absolute value one. In the original algebraic curve parameter, both maps have fixed semialgebraic format and bounded multiplicity. Their images therefore have area at least $P^{-C}LW$. Arcs of length below $W$ use their balanced caps with $L=W$.

For the first ribbon let $\Delta=d(\Psi(s,t))-d_\gamma(s)$, which is $O_P(1)$ by the balanced-box construction. Its coframe matrix is
\[
\begin{pmatrix}1+t d_\gamma'-m\Delta&-\Delta\\m&1\end{pmatrix},\qquad \det=1+t d_\gamma'.
\]
The entries and inverse are bounded. For the second ribbon,
\[
\begin{pmatrix}\Psi^*\theta\\\Psi^*dz\end{pmatrix}
=\begin{pmatrix}m-\Delta&1\\1&0\end{pmatrix}\begin{pmatrix}ds\\dt\end{pmatrix},
\qquad \det=-1.
\]
The entries and inverse are bounded by powers of $P$. Hence the actual critical value has bounded gradient in ribbon coordinates, and the ordinary scalar derivative preparation applies there. All statements are relative to every assigned fixed frame because its initial bridge was included before the full-arc assignment. The construction depends on actual phases and fixed branch/admissibility relations and precedes the choice of anchor.
\end{proof}

Theorem~\ref{low:4:sf:thm:transport}, the normalized-frame packing of Section~\ref{quartic:normalization}, and Lemma~\partIref{geo:core} verify these hypotheses for the quartic null admissibility.

The endpoint range in Proposition~\ref{low:4:sf:common-ribbons} also gives $|\eta|=|e-df|\lesssim_PLK$. Consequently its full possible $(\eta,f)$ domain has area $O_P(L^2K)$, while the actual ribbon has area $\gtrsim_PLW$. In Lemma~\ref{low:4:sf:null-incidence}, integrating over the possible point labels and dividing by the good area costs
\begin{equation}\label{low:4:sf:null-tonelli}
 \frac{L^2K}{LW}=K^2.
\end{equation}
The $30$-output map in the next proof uses the ordered derivatives in Definition~\partIref{def:coefficient-atlas} and the following transformation rule: $Y$, the chosen $Z$, $B$, and hence $X=|B|^{-1/2}Y$ transform under a fixed upper-affine map as the corresponding objects for the transformed phase. With the external seed jet fixed, every resulting endpoint jet depends only on $J$. Thus the $H_0$ conditional integration is legitimate for that map.

\paragraph{Direction parameters when curvature times aspect is small}
\begin{lemma}[Finite seed jets at small curvature--aspect]\label{low:surface:seed-rounding}
Use a common ruled ribbon from Proposition~\ref{low:4:sf:common-ribbons}, with scales $L,W$ and $K=L/W$. Put $\alpha=\sqrt{|B|}\theta$, let $X,Z$ be dual to $\alpha,dz$, and suppose $d\alpha=-\mathfrak a\alpha\wedge dz$, where $\mathfrak a=rv/(2B)$. Fix $k_0\in\{4,5,6\}$ and a sufficiently high fixed preparation order.

\noindent\textup{(i)} If $|r|\le\epsilon_r$ and $\epsilon_rK\le H_0$, with $H_0\ge1$, the $k_0$ external seed derivatives have a net of size $P^C(1+H_0)^CK^{k_0}$. The external coefficients are held fixed in the equation $X_J+n(u)Z_J$, or $Z_J+b(u)X_J$ in the other angular chart; the vector fields remain phase dependent. For outputs with longitudinal index at most $k_0$, this net preserves the prepared bounded-gradient widths
\[
M_fL^{1-i}\quad(j=0,\ i\ge1),\qquad M_fL^{-i}W^{1-j}\quad(j\ge1),
\]
and the bounded-scalar widths $M_gL^{-i}W^{-j}$, up to polynomial constants. Here $M_f$ bounds the ribbon gradient of each scalar $f$, and $M_g$ bounds each scalar $g$; all derivatives used by the finite chain rules are prepared simultaneously.

\noindent\textup{(ii)} The same conclusions hold, with an additional factor $c^{-C}$, if $|r|K\le H_0$ only on a fixed-format subset of area at least $cLW$ of the assigned ribbon.
\end{lemma}
\begin{proof}
In the $Y$-oriented ribbon write
\[
\alpha=A\,ds+C\,dt,\qquad dz=m(s)\,ds+dt,\qquad D=A-mC.
\]
Prepare the bounded coframe entries and inverse combinations, the bounded shear difference $\Delta=d\circ\Psi-d_\gamma$, and the required scalar list. The seed bounds $|d_\gamma^{(j)}|\le P^CKL^{-j}$ and the prepared bounds for $\Delta$ give $|v|\le P^CW^{-1}$ at a common good set. Small-value preparation of $\mathfrak a$ gives
\[
|\partial_s^i\partial_t^j\mathfrak a|\le P^C\epsilon_rL^{-i}W^{-j-1}.
\]
The exterior identity is $A_t-C_s=\mathfrak aD$. In the scaled variables $\sigma=s/L$, $T=(t-t_0)/W$, the common prepared points therefore satisfy
\[
|\partial_\sigma^i A|\le P^C,\qquad |\partial_\sigma^i\partial_T^j A|\le P^C(K^{-1}+\epsilon_r)\quad(j\ge1)
\]
through the prescribed finite derivative orders.

Where $|A|\ge P^{-C}$, parameterize the reference seed by its $\alpha$ flux $u$. Its external slope is $n_0=m/A$, with $|n_0^{(j)}|\le P^CL^{-j}$. Put $U=u/L$, write $\sigma=\sigma_{\rm ref}(U)+K^{-1}\chi(U)$, $T=\tau(U)$, and $n=n_0+K^{-1}\delta n$. The equations for the differences are
\[
\tau'=K\frac{-m+nA}{D},\qquad
\chi'=K\bigl(A(\sigma,T)^{-1}-A(\sigma_{\rm ref},0)^{-1}\bigr)-\frac CA\tau'.
\]
The reference numerator vanishes. Since $m$ has no transverse dependence, every difference in a numerator has a factor $K^{-1}$ from a long displacement or a factor $K^{-1}+\epsilon_r$ from a transverse derivative of $A$. After multiplication by $K$, finite differentiation therefore gives recurrence coefficients and their derivatives in the external jet bounded by $P^C(1+H_0)^C$. The denominators are powers of the conditioned base-point pivots. This is a finite Taylor recurrence at one point.

In the $z$-parameter chart write $\alpha=A\,ds+C\,dt$, $dz=ds$, with $C$ conditioned away from zero. The exterior identity becomes $A_t-C_s=-\mathfrak a C$, and the seed equation $Z+b(u)X$ gives $T'=K(b-A)/C$. It has the same cancellation. If the first chart has a small $A$ pivot, $D=A-mC$ bounds $m$ and $C$ away from zero; taking $z$ as the long coordinate gives this second chart, with scales comparable by powers of $P$.

In either chart round the external derivatives of orders $0\le j<k_0$ at mesh
\[
c_P(1+H_0)^{-C}\frac{W}{L^{j+1}}.
\]
Each range $P^CL^{-j}$ costs $P^C(1+H_0)^CK$ choices. The finite recurrence controls the rounded displacement jets at scale $W/L^j$. For the subsequent transverse $Z$ flow in the first chart, with its parameter divided by $W$, the scaled long component is $-K^{-1}C/D$ and the transverse component is $A/D$. Every fixed jet of its long component retains $K^{-1}$. In the second chart the transverse $X$ flow has no long component. Hence a purely longitudinal derivative entering a mixed row carries the required factor $W/L$. The finite chain rule preserves the displayed bounded-gradient widths and, separately, the bounded-scalar widths. The scalar amplitudes are never combined. No higher external seed derivative is needed for a row whose longitudinal index is at most $k_0$.

For the subset assertion first make the exceptional fraction for the bounded scalars smaller than $c/2$. On the retained part, $|\mathfrak a|\le P^CH_0/(KW)$. The subset derivative estimate in Lemma~\partIref{geo:good-points-strips} gives its finite jets with a factor $c^{-C}$ and retains area bounded below by a fixed power of $c$ times $LW$. Apply the same recurrence with $1+H_0$ replaced by $P^Cc^{-C}(1+H_0)$. The preparation uses algebraic derivative expressions in the original seed parameter; no algebraicity of a flux primitive or of a characteristic flow is required.
\end{proof}

\begin{theorem}[Small curvature--aspect branch]\label{low:4:sf:thm:smallaspect}
On the common ribbons above, suppose $|r|\le\varepsilon_r$ throughout the retained outer piece and $\varepsilon_r K\le c_P$.

\noindent\textup{(i)} Then the external direction jets admit a $K^4$ net preserving the 24 critical-output widths.

\noindent\textup{(ii)} Their common coefficient-volume bound at the fixed frame is $O_P(S^{-38}K^{-2})$, with the original physical coefficient factor.
\end{theorem}
\begin{proof}
Apply Lemma~\ref{low:surface:seed-rounding} with $k_0=4$, preparing the complete phase, coframe, shear difference and inverse combinations through all required orders. The identity $d\alpha=-\mathfrak a\alpha\wedge dz$ follows from the quartic connection identities. The resulting four-jet net costs $P^CK^4$. In the intrinsic coefficient map the fields $X_J,Z_J$ remain phase dependent; their variations contribute only terms three critical degrees below the model restriction blocks. Thus the six-normal/24-critical atlas retains its nonzero determinant. Its 24 widths have product $S^{-38}K^{-8}$, and the joint $\eta$ integration justified above costs $K^2$. Thus $\operatorname{vol}E_j\lesssim_P(S^{-38}K^{-8})K^4K^2=S^{-38}K^{-2}$.
\end{proof}

For $\varepsilon_rK\le H_0$, Lemma~\ref{low:surface:seed-rounding} gives the same estimate with the factor $P^C(1+H_0)^C$; in particular $H_0=P^C$ still gives polynomial constants.

\begin{corollary}[A version on subsets of comparable area]
Fix $c>0$. The same bound holds if $\{|r|\le c_P/K\}$ occupies at least a fraction $c$ of the existing common ribbon, with a sufficiently small $c_P$ depending also on $c$. No bound for $r$ on the whole outer piece is required.
\end{corollary}
\begin{proof}
Use the subset assertion of Lemma~\ref{low:surface:seed-rounding} with $k_0=4$. The quartic atlas and the point comparison in Theorem~\ref{low:4:sf:thm:smallaspect} are unchanged; the additional constants depend polynomially on $c^{-1}$.
\end{proof}

\paragraph{Bounded total shear on center curves with a substantial null direction}
The preceding curvature--aspect restriction is unnecessary when the assigned long center curve has a substantial $Z$ component. This is an actual varying-generator statement.

\begin{theorem}[Oblique bounded-shear branch]\label{low:4:sf:thm:oblique}
Consider an oriented center curve arc in the assignment of Proposition~\ref{low:4:sf:common-ribbons} whose tangent satisfies $|\theta(\gamma')/z'|\le C_P$, and put $L=\max(W,\operatorname{length}(\gamma))$ for its piece.

\noindent\textup{(i)} Relative to every assigned fixed frame, the shear on its actual ribbon is uniformly bounded.

\noindent\textup{(ii)} Its coefficient-volume bound at the fixed frame is $O_P(S^{-38}K^{-3})$ in the grid chart and $O_P(S^{-38}K^{-15})$ in the weighted chart, with $S=LW$ and $K=L/W$. No curvature--aspect restriction is imposed.
\end{theorem}
\begin{proof}
Parameterize the arc by $z$ and put $\beta=\theta(\gamma')$. The exact moving-frame equation is
\[
J'=V_\beta J+d'DJ,\qquad
V_\beta=(B+\beta r)A+\beta BC+\beta T_y+T_z.
\]
For $q=x-r y^2/2-Byz$, the coefficient functional below satisfies $L_\beta(V_\beta T)=0$ for every polynomial $T$ of degree at most four, and $L_\beta(Dq^2)=1$. For a reproducible explicit formula, its seven nonzero coefficient weights are
\[
\begin{array}{c|c@{\qquad}c|c}
x^2&\dfrac{\beta(4B+5r\beta)}{8(B+r\beta)}
 &xz^2&\dfrac{-2B^2+3r^2\beta^2}{16B^2(B+r\beta)}\\[7pt]
xyz&-\dfrac{\beta(2B+3r\beta)}{16B(B+r\beta)}
 &xy^2&-\dfrac{\beta^2}{8(B+r\beta)}\\[7pt]
z^4&\dfrac{3r(-4B^2-4Br\beta+r^2\beta^2)}{16B^4(B+r\beta)}
 &yz^3&\dfrac{3(2B^2+2Br\beta-r^2\beta^2)}{16B^3(B+r\beta)}\\[7pt]
y^3z&\dfrac{3\beta^2}{8B(B+r\beta)}&&
\end{array}
\]
Applying these weights to all monomials of degree at most four verifies the two identities. Bounded $\beta$, controlled nonzero $B$, and the small null cutoff make these coefficients bounded. On the compact null phase neighborhood, $L_\beta(DJ)$ is bounded away from zero. Thus
\[
|d'|\lesssim_P\sum_{\nu}|J_\nu'|,
\qquad \operatorname{Var}_\gamma(d)\lesssim_P1.
\]
The second assertion uses bounded variation of the bounded, fixed-format semialgebraic coefficient functions along the original arc parameter. Varying $\beta$ causes no derivative of $L_\beta$ in this pointwise identity.

Use the oriented-arc refinement already made in Proposition~\ref{low:4:sf:common-ribbons}. Each assigned anchor reaches its full arc by a buffered $P^CW$ bridge inside a bounded-shear containing set. The exact one-form estimate applies along the entire original arc. Thus an anchor far from that retained subarc still has bounded relative shear to every point of it. The actual ribbon differs in shear by $O_P(1)$ from its seed. This proves the every-anchor assertion without changing the assignment or its $L$. Arcs shorter than $W$ use their balanced caps and $L=W$.

Remove the fixed anchor shear. Ordinary $(y,z)$ coordinates are now uniformly comparable to the intrinsic coframe on the whole actual ribbon. Its critical value has bounded ordinary gradient. The ordinary semialgebraic curved-strip good-point argument and the fixed centerline-jet net apply, without a phase-dependent seed ODE. At a coefficient point, condition on the actual relative $d$ first. The critical value in these fixed coordinates is that of $J$ under the bounded substitution $(y,z)\mapsto(y-dz,z)$, modulo a constant. Its variation is just the envelope restriction $T|_{F_x=0}$; no variation of the conditional $d$ is taken. The six-normal/24-critical coefficient atlas is uniformly nonsingular over this compact family of bounded substitutions. The width products are $S^{-38}K^{-8}$ and $S^{-38}K^{-24}$, and the ordinary curve nets cost $K^4$ and $K^8$.

Finally use the full joint incidence density in actual centers,
\[
dF\,de\,df=C|\lambda^4B|\,dJ_{30}\,da\,db\,dc\,dd\,de\,df.
\]
For fixed $J,d,e,f$, the normal subgroup $H_3=\{d=e=f=0\}$ gives a uniformly bounded conditional integral in $(a,b,c)$ under the same-phase anchor window. The relative $d$ interval has bounded length. Candidate points lie in an anchor-relative ordinary disk of area $O_P(L^2)$ and every phase supplies good area $\gtrsim_PLW$, giving the remaining factor $K$. The grid and weighted products, respectively, are
\[
(S^{-38}K^{-8})K^4K=S^{-38}K^{-3},\qquad (S^{-38}K^{-24})K^8K=S^{-38}K^{-15}.
\]
These prove both tails. The original fixed-frame weights are unchanged.
\end{proof}

\subsection{Quartic coefficient estimates}\label{quartic:null-counts}\label{low:4:sf:sec:newrescaling}
We retain the original ribbon assignment, the complete physical augmentation, and all $34$ coefficients. Write
\[
S=LW,\quad K=L/W,\quad Y=\partial_y,\quad Z=\partial_z+d\partial_y,
 \quad r=c_{yy},\quad B=c_{yz}+dr,
\]
and $P_1=Y\phi$, $Q_1=Z\phi$. The following propositions give the two large-parameter coefficient estimates. Their input is a positive-area certificate inside the assigned ribbon, with its selection costs included.

For the square-quadric model, apply Lemmas~\partIref{geo:model-box} and~\partIref{geo:core} with normal pivot $\lambda$ and null-root pivot $2B$. Choose the model-box radius first and the complete coefficient defect last. Every selected inner state then has the prescribed intrinsic boundary margin $M$ in the full unbanded outer sheet. Enlarge the physical window using $p_{\mathrm{phys}}+B_{\mathrm{phys}}h(0)$ and $B_{\mathrm{phys}}Dh$, preserving the physical scales up to powers of $P$.

Use Lemma~\partIref{geo:relative-preparation} on the original affine transversals, simultaneously for the scalar list below. Its fixed fractional loss and relative derivative bounds precede all amplitude and area-fraction selections; zero and constant branches are kept separately.

\begin{lemma}[Transverse derivative bounds before interval selection]\label{low:5:qsa:cellwise-preparation}
Let an original ruled ribbon have long length $L$ and transverse width $W$.

\noindent\textup{(i)} There is a family of prepared cells on its affine transverse lines, independent of a later dyadic amplitude $h$ or selected longitudinal fraction $q$, with a fixed fractional loss and boundedly many width-$c_PW$ cell charts per line. On each cell $r$ is either zero or comparable to its representative value, and its relative transverse jets satisfy Lemma~\partIref{geo:relative-preparation}. For each of the finitely many bounded scalars $f$ needed in reconstruction, with $|f|\le M_f$, the bound
\[
|\partial_t^jf|\le P^CM_fW^{-j}\qquad(0\le j\le N)
\]
holds throughout every retained cell, with constants independent of $h$ and $q$. The scalars include $B,\lambda,P_1,Q_1,u,v$, the bounded coframe coefficients and inverse combinations, and the bounded shear difference $\Delta$.

\noindent\textup{(ii)} For $0<q\le1$, a fixed-format certificate of area at least $qLW$, consisting of whole prepared cells at one dyadic amplitude $h$ with $|r|\lesssim_Ph$, contains a curved subribbon of width comparable to $W$ and long length $L_0\ge P^{-C}q^CL$. At a further fixed positive fraction of this subribbon, the simultaneous mixed-jet estimates can be written
\[
\begin{aligned}
|\partial_s^i\partial_t^jf|&\le P^CM_fL_0^{-i}W^{-j},\\
|\partial_s^i\partial_t^j(r/h)|&\le P^CL_0^{-i}W^{-j}.
\end{aligned}
\]
The pure transverse case $i=0$ retains its original constants and contains no factor $q^{-C}$. If a coarser longitudinal selection is used, its extra polynomial losses may instead be attached only to rows with $i\ge1$.
\end{lemma}
\begin{proof}
Apply Lemma~\partIref{geo:relative-preparation} to $r$ on each original affine transversal, fixing the exceptional fraction before amplitude pigeonholing. This leaves boundedly many intervals of length at least $c_PW$ with comparable amplitude. For each additional bounded scalar $f$, apply the one-dimensional absolute derivative-good-point argument on these intervals with the same exceptional fraction. Fixed-format algebraicity bounds the number of bad components. Discard short good components and trim endpoints; a fixed positive fraction of each cell remains in boundedly many intervals of width at least $c_PW$, with the pure transverse bounds throughout. Prepare the finite scalar list and derivative orders simultaneously. Fubini gives a fixed area loss independent of later amplitudes or selected fractions. Parameter-uniform semialgebraic selection gives boundedly many cell charts over boundedly many long-parameter intervals.

Now choose $h$ using an amplitude of a whole prepared cell and give the entire cell that label. To make the branch split exhaustive, put $R_{\rm cell}=\sup_{\rm cell}|r|$, a fixed-format semialgebraic function of the long parameter. If $R_{\rm cell}\le\epsilon_0|n|$, use the slope alternative and choose $h\asymp|n|$. Otherwise value comparability gives $|r|\ge c_P\epsilon_0|n|$ throughout the cell; use the curvature alternative with $h\asymp R_{\rm cell}\asymp_P|r|$. Thus no thin pointwise subset of a cell is used for either branch. This changes only fixed chart and overlap constants in the thresholds. Since $n$ depends only on the original long parameter, its amplitude band introduces no transverse restriction. Division of the relative derivative estimates for $r$ by the selected $h$ now gives uniform bounds for $\partial_t^j(r/h)$ on the entire cell, even when $h$ is extremely small.

For fixed $h$, a selected set of area $qLW$ is a union of these whole cell fibers. Their widths are bounded above and below by constant multiples of $W$ and their number per line is bounded. Its long projection therefore has length at least $c_PqL$. Fixed format bounds the number of cell charts and connected intervals in that projection. Select one chart and one interval with length at least a fixed power of $q$ times $L$. Its cell midpoint $\tau(s)$ is algebraic, has range $O_P(W)$, and has a margin $c_PW$ to both cell endpoints. One-dimensional derivative preparation of $\tau$, followed by selection of a component of its good set, retains length $L_0\ge P^{-C}q^CL$ with $|\tau^{(j)}|\le P^CW L_0^{-j}$. Thus $t=\tau(s)+t_0$, $|t_0|\le c_PW$, is a contained subribbon. All previously established pure transverse estimates hold throughout it; they are not re-estimated from its area fraction.

In these shifted coordinates, each scalar $\partial_{t_0}^jf$ is already bounded by $P^CM_fW^{-j}$ throughout the contained rectangle of dimensions $L_0,W$. Apply the one-dimensional argument of Lemma~\partIref{geo:preparation} only in $s$ to each such scalar, at fixed $t_0$, and then use Fubini and a fixed union bound. A common fixed positive fraction has the mixed estimates displayed in the statement. The same argument applies to $r/h$, whose pure transverse bounds were established before choosing the selected interval, and to $n/h$, which has bounded values and no original transverse dependence. This proves the claimed separation: longitudinal loss changes $L$ to $L_0$, whereas the transverse scale remains $W$. 
\end{proof}

\begin{definition}[Preparation for rescaling]\label{def:prepared-rescaling}
Use the cellwise order of Lemma~\ref{low:5:qsa:cellwise-preparation}: transverse preparation at the original width $W$ precedes amplitude selection, and only the stated longitudinal bounds incur the subsequent selection loss. For rescaling, the prepared list consists of $n/h,r/h$, the bounded geometric scalars and inverse combinations, $u,v,\lambda,P_1,Q_1$, to the sufficiently high fixed order required below. The slope and curvature cases are the cellwise alternatives in that lemma.
\end{definition}
Use the fixed global scaling $F\mapsto\widehat F$, $\widehat F(X,Y,z)=h^2F(X/h,Y/h,z)$. Its graph and null fields are $\widehat c(Y,z)=h\,c(Y/h,z)$, $\widehat d(Y,z)=h\,d(Y/h,z)$, $\widehat Y=\partial_Y$, and $\widehat Z=\partial_z+\widehat d\,\partial_Y$; also put $\widehat r=\widehat c_{YY}$ and $\widehat\theta=dY-\widehat d\,dz$. Compact-model closeness below concerns its jet in the actual endpoint frame, at a point satisfying the simultaneous bounds of Definition~\ref{def:prepared-rescaling}. Translate this point and normalize the normal root and tangential directions, with $\widehat d$ set to zero there. The global scaling is used only for the coefficient and Haar Jacobians.

\begin{lemma}[Reconstruction after a fixed dyadic scaling]\label{low:4:sf:lem:newreconstruction}
Suppose a prepared part of an actual null ribbon has slope $n=dz/ds$ and a fixed dyadic $h$ such that $|n|\lesssim h$, $|r|\lesssim h$, $hL\ge M$, and $W\ge M$. Assume that the good-point preparation includes the relative finite jets of $n/h,r/h$ and the bounded normalized geometric scalars.

\noindent\textup{(i)} Then on a positive-area good set
\[
 \widehat F(X,Y,z)=h^2F(X/h,Y/h,z)
\]
in its actual endpoint-normalized frame belongs to a fixed small compact neighborhood of square-quadric models with mixed coefficient bounded away from zero. The set has fixed format in the original algebraic seed parameter and affine transversal.

\noindent\textup{(ii)} If $|r|\ll h$, this is a null neighborhood.

\noindent\textup{(iii)} If $|r|\asymp h$, its first curvature is bounded away from zero.
\end{lemma}
\begin{proof}
Write $\theta=A\,ds+C\,dt$, $dz=n\,ds+dt$, with bounded coframe and inverse. In $U=hs,T=t$,
\[
 \widehat\theta=A\,dU+hC\,dT,\qquad dz=(n/h)\,dU+dT.
\]
Both scales $hL,W$ are large. Set $u=Yd$, $v=Zd$, which are small in the original null model neighborhood. The identities are
\[
\begin{aligned}
 \widehat d&=hd,\qquad \widehat Y\widehat d=u,\qquad \widehat Z\widehat d=hv,\\
 \widehat c_{YY}&=r/h,\qquad \widehat c_{Yz}=B-\widehat d(r/h),\\
 \widehat c_{zz}&=-2\widehat dB+\widehat d^2(r/h).
\end{aligned}
\]
Normalize $\widehat d=0$ at the good point. Prepared inverse-coframe jets and these identities bound every required ordinary graph jet near
\[
Q_0(Y,z)=\tfrac12(r/h)Y^2+BYz.
\]
No division by $r$ is involved. Normal-derivative jets equal those of $\lambda$. For $\widehat\phi(Y,z)=h^2\phi(Y/h,z)$, the identities $\widehat Y\widehat\phi=hP_1$ and $\widehat Z\widehat\phi=h^2Q_1$ hold with the original scalars evaluated at $(Y/h,z)$. The nonconstant critical-value jets are consequently $O_P(h)$.

For the arbitrary cubic $G=\widehat F_X$, use its restriction jets to $X=\widehat c$ through order six and the four graph-normal jets of $G_X$ with indices $0,Y,z,Yz$. At $Q_0$, the restriction kernel is
\[
G=(X-Q_0)(\alpha+\beta X+\gamma Y+\eta z).
\]
The four normal jets recover its four coefficients because $(Q_0)_{Yz}=B\ne0$. For an explicit square minor, take the restriction Taylor rows with $(i,j)$ in
\[
 \{(i,j):i+j\le3\}\cup\{(1,3),(2,2),(3,1),(2,3),(3,2),(3,3)\},
\]
and the four graph-normal rows of $G_X$ at $(i,j)=(0,0),(1,0),(0,1),(1,1)$. On the twenty monomials of total degree at most three, elementary elimination gives determinant $B^{11}$, independently of $r/h$. To verify nonsingularity without expanding the determinant, the sixteen restriction rows annihilate a cubic precisely when it is $(X-Q_0)$ times an affine-linear polynomial; the four normal rows then successively recover its constant, $Y$, $z$ and $X$ coefficients, the last with pivot $B$. It remains invertible for the prepared graph perturbation. Thus $G$ is close to $\lambda_0(X-Q_0)$. Integrating in $X$, the remaining $14$ nonconstant binary coefficients are recovered from the critical-value jets through order four. This proves the claimed compactness without an a priori bound for the rescaled coefficients.
\end{proof}

A positive-area fixed-format good set contains a curved substrip of comparable dimensions: select one of its finitely many vertical interval fibers of width $cW$, then prepare its bounded midpoint function. In an intrinsic long parameter, all derivatives are still algebraic functions of the original parameter divided by its nonzero algebraic flux. Any fractional-area losses remain part of the derivative preparation. The contained strip lies entirely in the reconstructed compact set.

\paragraph{Scalar rows for the rescaled counts}

Use a prepared rescaling from Lemma~\ref{low:4:sf:lem:newreconstruction}. Write $a=hP_1$ and $b=h^2Q_1$, with $h=m$ in the slope case. The external rounded long direction is $\widehat Y+\beta\widehat Z$; its purely longitudinal rows use the frozen scalar $a+\beta b$, and its mixed rows use $b$. At a model with affine critical part $\widehat\ell$, write $\widehat\ell_Y=\partial_Y\widehat\ell$ and, for a polynomial variation $T$, put $T_c=T|_{X=Q_0}$. Then
\[
  \delta a=\widehat YT_c,\qquad
  \delta b=\widehat ZT_c+\widehat\ell_Y\,\delta\widehat d.
\]
The last term lowers critical degree by three.

\begin{lemma}[The two scalar coefficient determinants]\label{low:4:sf:scalar-determinants}

\noindent\textup{(i)} The scalar rows in Propositions~\ref{low:4:sf:prop:newslopecount} and~\ref{low:4:sf:prop:newcurvaturecount} have the nonzero determinants and output volumes given below. For the weighted rows put $q_*=2B\beta+\widehat r$. For the grid rows use the mixed coefficient $B'$ after the bounded tangent normalization. In the row convention where an $(i,j)$ purely longitudinal row extracts $[u^{i-1}]\partial_uT_c$, and a row with $j\ge1$ extracts $[u^iv^{j-1}]\partial_vT_c$, their ordinary-degree block determinants are
\[
\begin{array}{c|r|r}
w&\text{weighted scalar block}&\text{grid scalar block}\\ \hline
 1&-1&-1\\
 2&-4&-4\\
 3&18&18\\
 4&96&96\\
 5&15q_*^4/8&24(B')^4\\
 6&-3q_*^6/16&-24(B')^6\\
 7&-7q_*^6/64&-12(B')^6\\
 8&q_*^4/2&4(B')^4
\end{array}
\]

\noindent\textup{(ii)} Their products are $(8505/64)q_*^{20}$ and $191102976(B')^{20}$.
\end{lemma}
\begin{proof}
For $w\le4$, use the homogeneous binary columns $u^{w-j}v^j$, $0\le j\le w$. For $w=4+k>4$, use $Q^k u^{4-k-j}v^j$, $0\le j\le4-k$, with $Q=(\widehat r/2)u^2+Bu(\beta u+v)$ in the weighted case and $Q=A'u^2+B'uv$ in the grid case. The indicated row extracts are obtained by differentiating these explicit polynomials once. Binomial expansion followed by triangular elimination gives the displayed blocks; this is the same restriction calculation above with the appropriate derivative multiplier on each row. In the weighted case all eight purely longitudinal rows use the long scalar $a+\beta b$, so their principal variation is $\partial_uT_c$. Replacing it by $a$ alone would change the last block and can make the determinant vanish. Every other row uses $b$, whose principal variation is $\partial_vT_c$. The null-root correction in $\delta b$ has degree three below that principal variation, so it does not change any diagonal block. Variations of bounded direction coefficients are removed by first retaining the six normal outputs and the two zeroth scalar outputs.

The weighted index sums are $\sum i=70$, $\sum j=30$, with eight purely longitudinal and sixteen mixed rows. Their widths therefore have product $h^8h^{32}(hL)^{-62}W^{-14}=h^{40}(hL)^{-62}W^{-14}$. The grid sums are $50,50$, with four purely longitudinal and twenty mixed rows, giving $h^4h^{40}(hL)^{-46}W^{-30}=h^{44}(hL)^{-46}W^{-30}$. These scalar amplitudes give the factors used in the two volume estimates.
\end{proof}


\begin{proposition}[The slope case]\label{low:4:sf:prop:newslopecount}
On a selected certificate with $|n|\asymp m$, $|r|\ll m$ and $mK\gtrsim_PM$, the unchanged original assignment has the coefficient-volume bound at the fixed frame
\[
 \boxed{\quad S^{-38}K^{-4}(mK)^{-9}\quad}
\]
up to the stated derivative preparation losses.
\end{proposition}
\begin{proof}
Apply Lemma~\ref{low:4:sf:lem:newreconstruction} with $h=m$. The selected compact substrip has dimensions $mL,W$ and uniformly oblique tangent. The oblique quartic one-form of Theorem~\ref{low:4:sf:thm:oblique} bounds the total variation of $\widehat d$ along its centerline by the variation of its bounded algebraic normalized phase coefficients. Original transverse shear variation is $O_P(1)$, hence rescaled variation is $O_P(m)$. The whole substrip therefore has bounded relative shear in ordinary coordinates.

Use the weighted $24$-row restriction atlas and the scalar convention of Lemma~\ref{low:4:sf:scalar-determinants}. Prepare the coordinate derivatives of $a$ and $b$ separately. The leading restriction determinant is exactly
\[
 \frac{8505}{64}(2B\beta+r/m)^{20},
\]
in the row convention of Lemma~\ref{low:4:sf:scalar-determinants}. Six normal outputs complete the uniformly invertible $30$-coordinate map. The separate scalar amplitudes give output volume
\[
m^{40}(mL)^{-62}W^{-14};
\]
the ordinary curve-jet net costs $(mK)^8$. No phase-dependent seed is rounded.

Hold the endpoint shear fixed when applying the coefficient map, and integrate it afterwards over the interval stated below. The exact full coefficient scaling is $dF=m^2d\widehat F$. The original compact anchor bounds the $H_3=\{d=e=f=0\}$ Haar fiber; its rescaled coordinates are $(a,mb,mc)$, giving another $m^2$. This does not require a compact rescaled anchor. Integrate $\widehat d$ over its original geometric range $O_P(mK)$. The generic original point box $|e|\lesssim LK$, $|f|\lesssim L$ has rescaled area $O_P(mL^2K)$; good area is $\gtrsim_PmLW$, so the candidate area divided by the good area is at most $P^CK^2$. Therefore
\[
m^4\,m^{40}(mL)^{-62}W^{-14}(mK)^8(mK)K^2
 =S^{-38}K^{-4}(mK)^{-9}.
\]
The selected slope band need not contain the whole assigned arc.
\end{proof}

\begin{proposition}[The curvature case]\label{low:4:sf:prop:newcurvaturecount}
On a selected certificate with $|r|\asymp h$, $|n|\lesssim h$ and $hK\gtrsim_PM$, the coefficient-volume bound at the fixed frame is
\[
 \boxed{\quad h^7S^{-38}K^{-2}\quad}
\]
with the same preparation convention. The fixed frame may have zero first curvature.
\end{proposition}
\begin{proof}
The reconstructed compact substrip has $|\widehat r|\asymp1$. Its Gauss coordinates $p=\widehat c_Y,z$ satisfy
\[
dp=\widehat r\widehat\theta+B\,dz;
\]
their metric, area, and strip dimensions are comparable to $hL,W$. Four purely longitudinal grid rows come from the frozen long scalar $a+\beta b$, and twenty transverse rows from $b$. The long direction remains uniformly transverse to the null direction even if its quadratic coefficient vanishes. Retain the six normal outputs and the two zeroth scalar outputs. The direction-normalization coefficients are functions of these retained values. By the chain rule, their variations in the other outputs are linear combinations of the retained output differentials. Subtract those combinations from the remaining Jacobian rows; this preserves the determinant and holds these direction coefficients fixed in those rows. The variations of the phase-dependent null fields remain in the lower-degree correction described above. The exact grid determinant is $191102976(B')^{20}$, with $|B'|\asymp_P1$. The same degree-three correction gives a uniform actual map, output volume $h^{44}(hL)^{-46}W^{-30}$, and net $(hK)^4$.

Here a four-dimensional conditional fiber is legitimate. For every upper affine $g$, critical Gauss labels satisfy $p_F=a_g+p_J$, $z_F=f_g+z_J$. The scalars $P_1,Q_1$ and their Gauss derivatives are covariant. Thus the fixed rounded output map depends only on $J$, including when the remaining frame shear varies. 

On the endpoint Gauss slice $J_x=J_{xy}=J_{xz}=J_{xyz}=0$, the four physical-center body columns are $(A,C,T_x,D-a_gC)$; their slice determinant is $\pm\lambda^4\widehat r$. Set $\widehat\eta=\widehat e-d_gf$ on the \emph{right side} of this joint incidence identity. At fixed $(\widehat\eta,f)$ write $g=k g_0$, with $k\in H_0=\{e=f=0\}$ and $g_0$ fixed by those two labels. Conditioning on the $30$ outputs therefore leaves the direct $H_0$ action on the fixed phase $J\circ g_0^{-1}$. Its original compact-anchor pivot is $\lambda^4B$, including at $r=0$. The scaling sends its coordinates to $(a_g,hb_g,hc_g,hd_g)$, so its Haar volume contributes $h^3$, and $dF=h^2d\widehat F$ gives total $h^5$. No fixed-phase inverse of an $\eta$ point map is used.

The original null-to-Gauss endpoint shear changes $\eta$ by $O_P(L/h)$. Hence $|\widehat\eta|\lesssim_P hLK+L\lesssim_P hLK$, and $|f|\lesssim_PL$. Candidate area divided by good area again costs $K^2$. The full product of the volume factors is
\[
h^5h^{44}(hL)^{-46}W^{-30}(hK)^4K^2
 =h^7S^{-38}K^{-2}.
\]
\end{proof}

\begin{lemma}[Selection of an amplitude certificate]\label{low:surface:amplitude-selection}
Let an assigned $Y$-oriented null ribbon have scales $L\ge W\ge M$, $S=LW$, $K=L/W$, seed slope $n=n(s)$, and first curvature $r$. Suppose $n,r$ have fixed upper bounds and $|r|K>H_0$ on a fixed positive area fraction. Fix a sufficiently small slope threshold $\epsilon_0>0$. Fix a finite preparation list.

\noindent\textup{(i)} There are transverse cells satisfying Lemma~\ref{low:5:qsa:cellwise-preparation}, chosen independently of the amplitude $h$ and fraction $q$ below, with the original width and transverse constants.

\noindent\textup{(ii)} For sufficiently small fixed $\epsilon>0$, there is a dyadic $h$, with $x=hK\ge c_PH_0$, and a certificate of area at least $qS$, where $q\ge c_{P,\epsilon}x^{-\epsilon}$. Its whole transverse cells satisfy one of
\[
|n|\asymp_Ph,\quad |r|\lesssim\epsilon_0h,
\qquad\text{or}\qquad
|r|\asymp_Ph,\quad |n|\lesssim_Ph.
\]

\noindent\textup{(iii)} The certificate contains a prepared subribbon with
\[
L_0\ge P^{-C}q^CL,\qquad W_0\asymp_PW,\qquad
|\tau^{(j)}|\le P^CW L_0^{-j},\qquad
\frac{hL_0}{W}\ge c_{P,\epsilon}x^{1-C\epsilon},
\]
where $t=\tau(s)$ is its transverse center. Pure transverse estimates retain their original width and constants. If $C\epsilon<1/2$ and $H_0$ is sufficiently large, recentering preserves the slope alternatives and every fixed rescaled-length threshold. Signs, root sheets, the original assignment and its anchors are retained.

\noindent\textup{(iv)} If complete reconstruction holds on a fixed positive fraction of this subribbon and its admissibility has fixed format in the original algebraic parameters, one can select a contained compact substrip at a further fixed fractional cost.
\end{lemma}
\begin{proof}
Choose the exceptional fraction in Lemma~\ref{low:5:qsa:cellwise-preparation} smaller than the stated large-curvature fraction. With $R_{\rm cell}=\sup_{\rm cell}|r|$, assign each whole prepared cell to $R_{\rm cell}\le\epsilon_0|n|$ or its complement. In the first case use $h\asymp|n|$; in the second use $h\asymp R_{\rm cell}$. Relative value comparability gives the displayed alternatives on the entire cell, and $n$ depends only on the long parameter. Adjusting the fixed threshold absorbs the dyadic comparison constants. Cells meeting the retained large-curvature set have $hK\ge c_PH_0$.

The dyadic weights $x^{-\epsilon}$ have a bounded sum for $x\ge c_PH_0$. Weighted pigeonholing supplies the stated area fraction. The fixed-format interval and midpoint selection in Lemma~\ref{low:5:qsa:cellwise-preparation} gives $L_0,W_0,\tau$ with the displayed bounds. It changes the seed slope by $\tau'=O_P(W/L_0)$, so the relative change is $O_P((hL_0/W)^{-1})$. Since $hL_0/W\gtrsim_Pq^Cx\gtrsim_{P,\epsilon}x^{1-C\epsilon}$, the threshold choices preserve the required slopes and lengths. The pure transverse bounds were obtained before $h$ and $q$ were chosen.

For the last assertion, Fubini gives transverse good length at least $c_PW$ on a fixed fraction of the retained long interval. Fixed format bounds the number of fiber intervals and selection charts. Choose one interval of length at least $c_PW$ and its midpoint over an interval of comparable long length; prepare that midpoint and trim a fixed width about it. This gives a contained compact substrip with the same dimensions up to powers of $P$. Make the selections in the original algebraic curve parameter and affine transversal. A monotone flux change preserves interval counts, and its differentiated expressions satisfy Lemma~\partIref{geo:preparation}; algebraicity of its primitive is unnecessary.
\end{proof}

\begin{theorem}[Common exhaustion of null ribbons]\label{low:4:sf:thm:newexhaustion}

\noindent\textup{(i)} The original common null radius-band/center curve assignment has count $O_P(\Lambda S)$. Its certificates include vanishing first curvature and vanishing critical gradients.

\noindent\textup{(ii)} Its coefficient-volume bound at the fixed frame is
\[
 \boxed{\qquad C_P S^{-38}K^{-3/2}.\qquad}
\]

\noindent\textup{(iii)} The null-surface summation consequently allows $p<39$ without changing the physical coefficient character.
\end{theorem}
\begin{proof}
The $Z$-oriented arcs and the $Y$-oriented arcs with $|n|\ge c_P$ are already covered by Theorem~\ref{low:4:sf:thm:oblique}; in the latter case its $z$ parameter has $|\beta|=|\theta(\gamma')/z'|\le c_P^{-1}$. The threshold split was made before assignment. On the remaining arcs $|n|\le c_P$, and the null window has $|r|\le\delta_{\rm null}\ll c_P$. Choose this cutoff below the polynomial reconstruction tolerances, so the small rescaled critical-gradient terms are within the required compact neighborhood. Fix a large constant $H_0$. If $|r|K\le H_0$ holds on a fixed positive area fraction, the positive-fraction version of Theorem~\ref{low:4:sf:thm:smallaspect} gives $C_PS^{-38}K^{-2}$; the fixed finite-jet recurrence costs only a fixed polynomial in $H_0$.

Otherwise apply Lemma~\ref{low:surface:amplitude-selection} with the scalar list of Lemma~\ref{low:4:sf:lem:newreconstruction}. Write $x=hK$ and retain its certificate fraction $q\ge c_{P,\epsilon}x^{-\epsilon}$. Reconstruction gives actual compact rescaled phases on a fixed good fraction; the last part of the selection lemma gives a contained compact substrip. The slope case uses its actual centerline in the oblique one-form estimate, while the curvature case uses its uniformly regular Gauss image.

Choose one fixed $C$ covering all finite derivative, area, and counting losses. Propositions~\ref{low:4:sf:prop:newslopecount} and~\ref{low:4:sf:prop:newcurvaturecount} give, respectively,
\[
C_Pq^{-C}S^{-38}K^{-4}x^{-9},\qquad
C_Pq^{-C}h^7S^{-38}K^{-2}.
\]
Choose $\epsilon C<1/2$, then $H_0$ large enough for the seed and long-scale requirements, and finally the fixed buffer and inner defect in Lemma~\partIref{geo:core}. The slope sum is $\lesssim_PS^{-38}K^{-4}\sum_xx^{-8}$. Since $h$ is bounded above, the curvature sum is
\[
\lesssim_PS^{-38}K^{-2+C\epsilon}\sum_h h^{7+C\epsilon}
 \lesssim_PS^{-38}K^{-3/2}.
\]
All choices are certificates inside the original piece. This proves the exhaustion and the summable aspect factor. The fixed-catalogue transfer and the global shape summation retain the original factor $v_j^{p-35}$.
\end{proof}

\begin{theorem}[The complete surface branch]\label{low:4:sf:surface-completion}

\noindent\textup{(i)} All frames from the fixed family assigned to the sufficiently small fifth-singular-value neighborhood of Corollary~\partIref{geo:rank-four-entry-low} admit common actual-surface pieces. If $S=LW$ is a piece's area scale, its count is at most $P^C\Lambda S$.

\noindent\textup{(ii)} Its anchored coefficient set has tail at most $P^CS^{-38}$. The null pieces have the stronger summable aspect factor $K^{-3/2}$.

\noindent\textup{(iii)} Consequently the common summation over surface pieces is valid whenever $q=(1+\varepsilon)(p-1)<38$, with the weight of the original localization piece $v_j^{p-35}$.
\end{theorem}
\begin{proof}
Normalize and select the first valid actual root and curvature chart as in Section~\ref{quartic:normalization}, retaining the complete coefficient and physical windows. The nonnull assignment and tail are Theorem~\ref{low:4:sf:thm:gauss}; the null assignment and tail are Proposition~\ref{low:4:sf:common-ribbons} and Theorem~\ref{low:4:sf:thm:newexhaustion}. Their outer windows retain each original state after bounded normalization.

Boundedly many pieces at each width, $S\ge W^2$, and the summable aspect factor give Proposition~\partIref{an:common-incidence} for $q<38$, with the stated weight. The nonnull aspect factors are stronger; subunit scales cost powers of $P$.
\end{proof}

\paragraph{The quartic incidence input}\label{quartic:assembly}\label{low:4:ga:section}
The regular exponent $36$, cutoff $H\le P^C/d_2$, and Theorem~\ref{low:4:sf:surface-completion} supply the quartic row of Proposition~\ref{an:unified-incidence}, with loss $d_2^{-(q-36)}$.

\paragraph{Exact row lists and determinant data}
The companion supplement gives all seven regular-chart row lists in Section~\ref{ext-supp:low:4:rd:rows}, and the ten residual row lists, model conventions, exact determinants and rational verification in Section~\ref{ext-supp:low:4:rr:certificates} and Lemma~\ref{ext-supp:low:4:rr:minor_lemma}. Every list retains all $34$ input coefficients. The regular determinant table and the chart-specific dilation and additional-sample net cardinalities above remain the inputs to the coefficient estimates.


\subsection{Quintic coordinates and characteristic transport}\label{low:degree5}\label{quintic:coordinates}\label{low:5:qsf:normalization}\label{low:5:qsf:sec:packing}\label{low:5:qsf:preparation}\label{low:5:qsf:sec:gauss}\label{low:5:qsf:null-incidence-section}\label{low:5:qsf:sec:affinearea}\label{low:5:qsc:atlas}\label{low:5:qsc:nonnull}\label{low:5:qsc:reconstruction}\label{quintic:transport}\label{low:5:qnt:section}\label{low:5:qng:section}

\par\smallskip\noindent\textit{Step 1. Normalization and conditional frame measures.}\par\nobreak\smallskip

In Sections~\ref{quintic:coordinates} and~\ref{quintic:counts}, $k=5$. Use the physical weighted mass bound, coefficient framework, rank models and regular estimates of Section~\ref{reg:finite5}, with the same fixed frames. The transport results concern complete quintic phases under their stated curvature and rate cutoffs; Section~\ref{low:6:s36q:module} explains their use on a sextic derivative phase.

Use the conditioned square-quadric family of Corollary~\partIref{geo:rank-four-entry-low} and the ordinary normalization of Section~\ref{quartic:normalization}, with $k=5$ and closeness in all $55$ coefficients. Completing the square includes a term $b(x-Q)$ by shifting the constant term of $Q$. The correction in $\mathcal U_0$ preserves $(e,f)$; the normal-fiber bound is twelve, and the pivots, fiber-volume bound, and projected packing are~\eqref{low:4:sf:eq:normalized}--\eqref{low:4:sf:eq:area}.

Lemma~\partIref{geo:good-points-strips} gives~\eqref{low:4:sf:critical-widths}. The centerline nets cost $P^CK^{10}$ or $P^CK^5$, with every chain-rule row through ordinary order ten prepared before rounding. The Gauss intrinsic-point root bound is $4\cdot3=12$. Retain the common all-root assignment, switching sets, artificial-cut coverage, and original physical restrictions; coefficient inversion uses Lemma~\partIref{geo:coefficient-map}.

\paragraph{Joint coefficient and frame integration}
The conditional $H_0,H_3,H$ fiber bounds are Lemma~\partIref{surf:haar}, applied at the fixed frame with the corresponding phase-only pivot. Bounded right displacement changes the pivot only by powers of $P$.

The Gauss disintegration is \partIeqref{surf:gauss-density} with $m=2$ and slice dimension $N_k-4=51$. In chart II impose $J_x=J_{xy}=J_{xz}=J_{xzz}=0$ at the origin and write $\lambda=J_{xx}(0)$, $r=(c_*)_{yy}(0)$ and $B=(c_*)_{yz}(0)$. The selected null window conditions $\lambda$ and $B$ and their inverses; $r$ may vanish. Lemma~\partIref{surf:null-density} gives the joint coefficient--point density \partIeqref{surf:null-density-eq}. A $J$-only full slice map uses $H_0$ conditioning at fixed $(\eta,f)$; an external endpoint shear uses $H_3$ after fixing $(d,e,f)$. For physical-center Gauss incidence use Lemma~\partIref{surf:physical-gauss-density}, retaining the pivot of the fixed frame independently of the endpoint curvature.

The selected curvature and physical windows give $|\det D^2c|\ge P^{-C}$ and a physical radius at most $P^C$. Theorem~\partIref{atlas:affine-area}, with $\deg F_x\le4$, and Proposition~\partIref{atlas:projected-packing} yield $N\lesssim_P\Lambda v^{-1/2}$. Since $dp\,dz=|r|\,dy\,dz$ and $|r|\asymp_P1$ in the Gauss chart, a contained Gauss strip satisfies $LW\lesssim_Pv^{-1/2}$. A null strip has the same bound once its signed Jacobian is proved comparable to tangential area.

The square-quadric model is $J_0=a(x-Q(y,z))^2+\ell_y y+\ell_z z$, where $a\ne0$ and $Q$ is a nonsingular quadratic form. The actual normal critical graph is defined by $J_x=0$.

\paragraph{Step 2. The $16+35$ splitting}
Let $\mathcal C$ be the following set of $35$ input monomials:
\[
 \mathcal C=\{y^iz^j:1\le i+j\le5\}\ \cup\ \bigcup_{a=1}^{5}\{x^ay^iz^j:i+j=5-a\}.
\]
The first set has $20$ elements and the second has $5+4+3+2+1=15$. Its complement in the $55$ nonconstant quintic monomials is the set of $20$ monomials containing $x$ and having total degree at most four. In the Gauss chart the four slice rows are the coefficients of $x,xy,xz,xyz$. In the null chart they are the coefficients of $x,xy,xz,xz^2$. In either chart, take as normal outputs the other $16$ coefficients from this $20$-element set. Their coefficient differentials are coordinate rows. The slice and normal rows eliminate exactly these $20$ input directions, leaving $\mathcal C$.

For $Q(u,v)=Au^2+Buv+Cv^2$, restrictions of the columns in $\mathcal C$ are
\[
 \Pi_{\le5}(u,v)/\mathbb R,\quad Q\Pi_4^{\rm hom},\quad Q^2\Pi_3^{\rm hom},\quad Q^3\Pi_2^{\rm hom},\quad Q^4\Pi_1^{\rm hom},\quad Q^5.
\]
Their dimensions by tangential degree are $2,3,4,5,6,5,4,3,2,1$. This also proves that the image of restriction has dimension $35$: division by the monic polynomial $x-Q$ shows that its kernel on the quintic space is $(x-Q)\mathbb R[x,y,z]_{\le3}$, of dimension $20$.

\paragraph{Step 3. Two exact restriction determinants}
Use the index sets
\[
 \mathcal I_{\rm w}=\{(i,j)\in\mathbb{Z}_{\ge0}^2:0<i+2j\le10\},\qquad
 \mathcal I_{\rm g}=\{(i,j):0\le i,j\le5,\ (i,j)\ne(0,0)\}.
\]
Both sets have $35$ elements. Write $\mathcal R_{\rm w}$ and $\mathcal R_{\rm g}$ for the matrices taking the listed restriction columns to their Taylor coefficients on the respective row sets. The restriction calculation in the proof of Lemma~\partIref{surf:normal-filtration}, with degree five and the constant row and column omitted, gives
\[
\det\mathcal R_{\rm w}=A^{35},\qquad
\det\mathcal R_{\rm g}\big|_{A=0}=B^{35},\qquad
\det\mathcal R_{\rm g}\big|_{C=0}=B^{35},
\]
up to row-ordering sign. The nontrivial ordinary-degree block exponents are $5,8,9,8,5$. For a nonsingular conditioned quadratic form, fix a sufficiently small cutoff $\delta_{\rm q}=P^{-C}$ from the current conditioning bounds. Use the weighted rows when $|A|\ge\delta_{\rm q}$; otherwise $|B|\ge P^{-C}$ and the grid determinant remains bounded away from zero by continuity from $A=0$. The weighted row space is preserved triangularly by a prescribed curve substitution with identity linear part. For the grid use the actual higher-jet bounds $|\zeta^{(j)}|\le P^CL^{1-j}$ and compact continuity for $L\ge P^C$; bounded $L$ has bounded count, as in the quartic argument.

\paragraph{Gauss coordinates}
Use the Gauss coordinates of Section~\ref{quartic:coordinates}, with $r$ bounded away from zero. Formula~\eqref{low:4:sf:eq:gaussvariation} applies to every remaining quintic input column. At the model $\lambda=2a$, and its correction lowers tangential degree by three. Hence it leaves the ordinary-degree restriction determinants unchanged. The $16$ normal outputs and $35$ Gauss outputs form an invertible $51$-coordinate map, including bounded affine tangential terms. The implicit jets have only the retained normal and curvature pivots in their denominators, so compact continuity gives the actual nearby map.

\paragraph{Null coordinates}\label{low:5:qsc:null-atlas}
On the slice $J_x=J_{xy}=J_{xz}=J_{xzz}=0$ at the origin, write $Q=Ay^2+Byz$ with $B\ne0$. Let $d(y,z)$ be the chosen root of
\[
c_{zz}+2dc_{yz}+d^2c_{yy}=0,\qquad d(0)=0,
\]
Use the fields of Definition~\partIref{def:null-coordinates} and write $\phi=J|_c$. The $16$ normal outputs and the $35$ ordered jets $(Y^iZ^j\phi)(0)/(i!j!)$, $(i,j)\in\mathcal I_{\rm g}$, form a nonsingular $51$-coordinate map near the compact null family, including $A=0$.

The variation calculation in the proof of Theorem~\ref{low:4:sf:thm:nulljets} applies to these $35$ columns: the implicit denominators are $\lambda$ and $2B$, and the null-field correction lowers critical degree by three. Thus it preserves the determinant $B^{35}$ without division by $A$. The normal outputs are coordinate pivots, and compact continuity proves the nearby assertion. Ordered-jet bounds on ribbons are established by the transport argument below.

\par\smallskip\noindent\textit{Step 4. Nonnull coefficient-volume factors.}\par\nobreak\smallskip

For the nonnull count, apply the proof of Theorem~\ref{low:4:sf:thm:gauss} with the present $51$-coordinate map. The common Gauss construction gives count $O_P(\Lambda LW)$, good area $\gtrsim_PLW$, and candidate endpoint area $O_P(L^2)$. Lemma~\partIref{surf:null-density} gives $dF=C|\lambda|^4|r|^2\,dJ_{51}\,dh_H$, and Lemma~\partIref{surf:haar} bounds the conditional $H$ integral at the fixed frame. Preparing every ordinary derivative through order ten validates both centerline nets. The numerical factors are
\[
\begin{array}{c|c|c|c}
\text{rows}&\text{critical width product}&\text{jet net}&\text{coefficient bound}\\\hline
\mathcal I_{\rm w}&L^{-115}W^{-30}&K^{10}&S^{-145/2}K^{-63/2}\\
\mathcal I_{\rm g}&L^{-85}W^{-60}&K^5&S^{-145/2}K^{-13/2}
\end{array}
\]
Here $S=LW$, $K=L/W$, and each bound includes the point factor $K$. The weighted index sums are $125,55$, with ten purely longitudinal rows; the grid sums are $90,90$, with five. All $16$ normal windows stay bounded. Proposition~\partIref{an:common-incidence} gives summation for $q_p=(1+\eta)(p-1)<145/2$, in particular for $73<p<147/2$ and sufficiently small $\eta$.

\par\smallskip\noindent\textit{Step 5. Reconstruction in all 55 coefficients.}\par\nobreak\smallskip

The null rescaling also requires the following algebraic calculation. If $G=F_X$ is quartic and $q=X-Q(Y,z)$, then $G|_Q=0$ gives $G=qH$ with $\deg H\le2$. The graph values of $G_X$ determine $H$ only modulo a multiple of $q$. Therefore graph and normal-derivative jets alone do not reconstruct a quintic: they miss $G=cq^2$. The critical value recovers exactly this missing direction.

More precisely, fix $Q=AY^2+BYz$, $B\ne0$. On the $55$-dimensional space of quintics modulo constants take the following $55$ linear outputs, with the graph held fixed:
\[
 \begin{array}{ll}
 {}[Y^iz^j]F_X(Q,Y,z),&0\le i,j\le4\quad(25\text{ rows}),\\
 {}[Y^iz^j]F_{XX}(Q,Y,z),&0\le i,j\le2\quad(9\text{ rows}),\\
 {}[Y^iz^j]F(Q,Y,z),&1\le i+j\le5\quad(20\text{ rows}),\\
 {}[Y^3z^3]F(Q,Y,z),&\quad(1\text{ row}).
 \end{array}
\]
This map is invertible. If its first $25$ rows vanish, the degree-four restriction determinant gives $G=qH$, $\deg H\le2$. The next nine rows are the degree-two restriction of $H$, and their vanishing gives $H=cq$. Integration in $X$ yields
\[
F=c\bigl(X^3/3-X^2Q+XQ^2\bigr)+R(Y,z),\qquad \deg R\le5.
\]
Its critical value is $cQ^3/3+R$. The $(3,3)$ coefficient is $cB^3/3$, so it forces $c=0$. The remaining twenty rows force $R$ to be constant. This proves injectivity and hence invertibility. In the coordinate system using the coefficients of $G$ and the twenty binary coefficients, the reduced $35$ by $35$ determinant is exactly $B^{27}/3$ for the specified row order. The change from nonbinary $F$ coefficients to $G$ coefficients adds only fixed positive integer factors.

Uniformity under sufficiently small finite-jet perturbations of $Q$ follows by compact continuity. Thus, once a proposed null rescaling has supplied graph jets close to $Q$, graph jets of $F_{XX}$ close to a nonzero constant, and small nonconstant critical-value jets through the indicated orders, the complete rescaled quintic is close to the corresponding square-quadric. The extra sixth-order critical row is essential; omitting it leaves the explicit nonzero kernel above.

\paragraph{Characteristic transport and common ribbons}

Let $F$ have total degree at most five. Work on a fixed critical branch $x=c(y,z)$ with the null coordinates and rates of Definition~\partIref{def:null-coordinates}, and write $\phi=F|_c$. Suppose the actual null-normalized phases belong to a sufficiently small fixed neighborhood of the compact square-quadric family. In particular, all complete normalized coefficients have bounds $P^C$, $|a|^{-1}+|B|^{-1}\le P^C$, where $2a=F_{xx}|_{\mathrm{graph}}$, and $|v|\le\nu_0$ for a sufficiently small inverse power $\nu_0=P^{-C_0}$. The last bound follows from the normalized coefficient identity $v=3[xz^3]J/(2aB)$ and the vanishing of $[xz^3]J$ at every null-normalized square-quadric model.

\begin{definition}[Frozen null derivatives]\label{def:frozen-null}
The frozen null derivative of order $j$ at $p_*=(y_*,z_*)$ is
\[
\left.\frac{d^j}{dt^j}\phi\bigl(y_*+d(p_*)t,z_*+t\bigr)\right|_{t=0}.
\]
It is the ordinary $\partial_z^j\phi$ in the affine null-normalized coordinates at $p_*$; the direction is held fixed as $t$ varies. It is generally different from the iterated variable-field derivative $Z^j\phi$.
\end{definition}
\begin{lemma}[Exact sixth critical derivative]\label{low:5:qnt:sixth}
At every null-normalized point, the sixth ordinary derivative of the critical value in the frozen null direction satisfies
\[
\partial_z^6\phi=-80aB^2v^2.
\]
\end{lemma}
\begin{proof}
Write the normalized phase as $J$, with $J_x=J_{xy}=J_{xz}=J_{xzz}=0$ at the origin, and set $b=[xz^3]J$. Its nondegenerate critical graph satisfies $c(0,z)=-b z^3/(2a)+O(z^4)$. The coefficient of $z^6$ in $J(c(0,z),0,z)$ receives contributions exactly from $a c(0,z)^2$ and $b c(0,z)z^3$; there is no pure $z^6$ term because $\deg J\le5$. Hence $[z^6]\phi=-b^2/(4a)$. Differentiating $c_{zz}+2dc_{yz}+d^2c_{yy}=0$ by $Z$ at the origin gives $c_{zzz}+2Bv=0$, while the critical equation gives $c_{zzz}=-3b/a$. Thus $b=2aBv/3$. Multiplication by $6!$ proves the formula. All other terms containing $x$ have order at least seven after this substitution, so no higher normal coefficient has been omitted.
\end{proof}

Lemma~\partIref{geo:characteristic-preparation} applies to the algebraic scalars and their finite derivative lists below on the retained regular $Z$-branch. It supplies simultaneous good points and boundedly many connected intervals in each $v$ band; branch boundaries and nonzero pivots remain in the stopping relation.

\begin{lemma}[The moving-curve identity]\label{low:5:qnt:chain}
Along a characteristic put $D=d/dt$, $P_1=\phi_y\circ\gamma$, $R_2=\phi_{yy}\circ\gamma$, and $T_3=\phi_{yyy}\circ\gamma$. At each point, the sixth derivative in the frozen null direction is
\begin{align*}
\partial_z^6\phi={}&D^6(\phi\circ\gamma)-15vD^4P_1-20(Dv)D^3P_1-15(D^2v)D^2P_1-6(D^3v)DP_1-(D^4v)P_1\\
&+45v^2D^2R_2+60v(Dv)DR_2+\bigl(15vD^2v+10(Dv)^2\bigr)R_2-15v^3T_3.
\end{align*}
\end{lemma}
\begin{proof}
Fix a base point $t_0$ and use its affine null coordinates. The transverse displacement between the actual curve and the frozen null line at time $t_0+s$ is
\[
\delta(s)=\frac{v}{2}s^2+\frac{Dv}{6}s^3+\frac{D^2v}{24}s^4+\frac{D^3v}{120}s^5+\frac{D^4v}{720}s^6+O(s^7),
\]
with all displayed derivatives evaluated at $t_0$. Taylor expansion in the transverse variable gives, through order six,
\[
\phi(\text{frozen line at }t_0+s)=\phi(\gamma(t_0+s))-\delta(s)P_1(t_0+s)+\frac{\delta(s)^2}{2}R_2(t_0+s)-\frac{\delta(s)^3}{6}T_3(t_0+s)+O(s^7).
\]
Extract the coefficient of $s^6$ and multiply by $6!$. This gives the displayed formula, including every derivative of the varying null direction. In particular, an affine tangential summand $\lambda y$ is correctly canceled by these correction terms.
\end{proof}

\begin{theorem}[Uniform quintic characteristic transport]\label{low:5:qnt:transport}

\noindent\textup{(i)} On every connected active null characteristic under the normalized coefficient and pivot bounds above,
\[
\int |v|\,dt\le P^C\nu_0^{3/5},\qquad \int |u|\,dt\le P^C.
\]
Consequently the accumulated shear is bounded independently of characteristic length.

\noindent\textup{(ii)} After choosing the null curvature cutoff $\sup|r|$ sufficiently small by an inverse power of $P$, the transverse flow Jacobian and its inverse are bounded by powers of $P$.
\end{theorem}
\begin{proof}
Fix one connected interval of length $H$ on which $\nu/2<|v|\le\nu$, where $0<\nu\le\nu_0$. The scalars $P_1,R_2,T_3$ and $Q_1=Z\phi$ are uniformly $P^C$-bounded. Indeed, at the current null-normalized origin, $(P_1,R_2,T_3,Q_1)=(J_y,J_{yy},J_{yyy},J_z)(0)$; the critical equation and the three first normalizations remove every additional term in these identities. Apply Lemma~\partIref{geo:characteristic-preparation} simultaneously to $v$ with bound $\nu$, to $P_1,R_2$ with their fixed bounds, and to $Q_1$ through derivative order five. At a common good point one has
\[
\begin{aligned}
|D^jv|&\le P^C\nu H^{-j} &&(1\le j\le4),\\
|D^jP_1|&\le P^CH^{-j} &&(0\le j\le4),\\
|D^jR_2|&\le P^CH^{-j} &&(0\le j\le2),\\
|D^6(\phi\circ\gamma)|&=|D^5Q_1|\le P^CH^{-5}.&
\end{aligned}
\]
The identities in Lemmas~\ref{low:5:qnt:sixth} and~\ref{low:5:qnt:chain}, together with $|aB^2|\ge P^{-C}$ and $|v|>\nu/2$, imply
\[
\nu^2\le P^C\bigl(H^{-5}+\nu H^{-4}+\nu^2H^{-2}+\nu^3\bigr).
\]
Choose the fixed upper cutoff $\nu_0$ small enough to absorb the last term. For $H\ge P^{C_1}$, enlarge $C_1$ to absorb the third term. Young's inequality absorbs $P^C\nu H^{-4}$ into a fixed small multiple of $\nu^2$ plus $P^CH^{-8}$, which is at most $P^CH^{-5}$ when $H\ge1$. Thus $H\le P^C\nu^{-2/5}$. If $H<P^{C_1}$, the same inequality holds after enlarging $C$, because $\nu\le1$. Each such interval therefore satisfies
\[
\int_I|v|\,dt\le\nu H\le P^C\nu^{3/5}.
\]
Use the half-open dyadic bands $\nu/2<|v|\le\nu$, where $\nu=2^{-j}\nu_0$ and $j\ge0$. Lemma~\partIref{geo:characteristic-preparation} bounds the number of interval components in each band by a fixed constant. Summing the preceding estimate gives a convergent geometric series. A constant nonzero rate is included in its half-open band and obeys the same interval estimate; the zero set contributes nothing. Apply the length estimate to arbitrary finite subintervals to include infinite intervals.

The total variation of $\log|B|$ is at most $P^C$, since $B$ has bounded range and inverse and uniformly bounded level crossings. Apply the signed compatibility and flow-Jacobian calculation in the proof of Theorem~\ref{low:4:sf:thm:transport}. The bound $\int|v|\le P^C\nu_0^{3/5}$ just proved gives $\int|u|\le P^C$. Choose $\sup|r|\le P^{-C_2}$ after these polynomial constants are fixed, so that $\tfrac12\left|\int(r/B)v\,dt\right|\le1$. The polynomial bounds for $|B|$ and $|B|^{-1}$ then give polynomial upper and lower bounds for the transverse flow Jacobian.
\end{proof}

The small complete phase neighborhood is essential. For example, $q_\gamma=x-yz+\gamma z^3/3$ and $t_\gamma=y-\gamma z^2/2$ give the degree-four invariant $q_\gamma^2+\frac{8}{9\gamma}t_\gamma^3=x^2-2xyz+\frac{2\gamma}{3}xz^3-\frac13y^2z^2+\frac{8}{9\gamma}y^3$, which has unbounded shear along the flow of $y\partial_x+\gamma z\partial_y+\partial_z$. Its mixed-curvature square coefficient is not close to the square-quadric model and its $y^3$ coefficient is large when $\gamma$ is small. In the moving-curve identity its cubic transverse term has exactly the size needed to prevent absorption. The theorem applies in the normalized neighborhood supplied by the rank-four reduction.

All geometry in the quintic subsections concerns the actual degree-five phase. The normal equation is the quartic equation $F_x=0$, and the chosen null line is retained as a branch label. At a tangential point $(e,f)$, there are at most four regular normal roots and at most two selected null directions at each root. Thus the null normalization has at most eight frames over that point; allowing the complementary nonnull normalization gives at most twelve. These bounds replace the quartic bounds of six and nine, respectively. The transverse normalization Jacobian remains $2|\lambda|^4|B|$, since its computation uses only the four slice equations and not the total degree.

Choose the prescribed margin $M$ before the inner defect in Lemmas~\partIref{geo:model-box}--\partIref{geo:core}, retaining the full outer critical/null sheet and physical window. Theorem~\ref{low:5:qnt:transport} supplies the two transport hypotheses of Proposition~\ref{low:5:qng:ribbons}; the finite normal fibers and projected packing above supply its packing hypothesis. It gives the common quintic assignment, contained ribbons, prepared centerlines and fixed-frame endpoint ranges, with $W\ge M$.

The endpoint ranges imply $|\eta|=|e-df|\le P^CLK$. Thus either $(e,f)$ or $(\eta,f)$ ranges over an anchor-relative set of area at most $P^CL^2K$. Dividing by the good ribbon area $P^{-C}LW$ gives the common coefficient--point factor $P^CK^2$. If an ordinary-coordinate coefficient map retains the actual endpoint shear as an external parameter, its conditional shear integration has range $P^CK$ and the remaining $\operatorname{span}(A,C,T_x)$ Haar fiber is bounded by the fixed frame's normal pivot. If an output map depends only on the intrinsic endpoint phase, the four-dimensional null-fiber disintegration applies instead. These are distinct conditional integrations; they are not independent volume estimates multiplied for different phases.

\subsection{Quintic coefficient estimates and exhaustion}\label{quintic:counts}\label{low:5:qsa:section}\label{low:5:qsa:reconstruction}\label{low:5:qsc:interfaces}\label{quintic:assembly}\label{low:5:qsc:exhaustion-section}\label{low:5:qga:actual-section}

Use the ribbon assignment and its scales. In the following small curvature--aspect count, the critical coefficient map has $35$ grid outputs and $16$ normal outputs. The grid width product is $S^{-145/2}K^{-25/2}$, and the full null coefficient--point incidence costs $K^2$ when its conditioned output map depends only on the intrinsic endpoint phase. These facts concern all $55$ nonconstant quintic coefficients.

The buffer $M$ is fixed from compact-model, determinant and metric conditioning before the inner defect; later enlargement of the master calculation parameter to include reciprocal defects does not change it.

\begin{theorem}[A net for the direction derivatives]\label{low:5:qsa:smallaspect}
Suppose $|r|\le\epsilon_r$ on a retained outer ribbon and $\epsilon_rK\le H_0$, where $H_0\ge1$ is fixed.

\noindent\textup{(i)} The external direction jets have a net of size at most $P^C(1+H_0)^CK^5$ that preserves the $35$ critical-output widths.

\noindent\textup{(ii)} The corresponding common coefficient-volume bound at the fixed frame is
\[
P^C(1+H_0)^CS^{-145/2}K^{-11/2}.
\]

\noindent\textup{(iii)} The same conclusion, with an additional factor $c^{-C}$, holds if $|r|K\le H_0$ only on a subset of area at least $cS$ of the already assigned common ribbon.
\end{theorem}
\begin{proof}
Apply Lemma~\ref{low:surface:seed-rounding} with $k_0=5$ and $\mathfrak a=rv/(2B)$. The quintic connection identity gives its exterior equation. Prepare the normalized phase, coframe, inverse combinations and bounded shear difference through a sufficiently high fixed order, for example $24$. The five external derivatives of the seed slope have ranges $P^CL^{-j}$ and the meshes in that lemma. Their net costs $P^C(1+H_0)^CK^5$ and preserves the critical-output widths.

The $16$ normal outputs and $35$ critical outputs form the uniformly invertible quintic intrinsic atlas of Section~\ref{low:5:qsc:null-atlas}. At a model, variation of the phase-dependent null direction contributes only terms three critical degrees below the restriction blocks; the additional normal rows retain the other coefficient directions. Fixed higher external seed jets are small at the large buffer scale, and compact continuity gives the actual nearby atlas. Its output depends only on the normalized endpoint phase and the fixed external jet. The four-dimensional conditional null Haar integral is therefore legitimate under the fixed frame condition, and no separate endpoint shear range is introduced. The five seed choices cost $P^C(1+H_0)^CK^5$. Multiplying by the grid volume and the common point factor gives $S^{-145/2}K^{-25/2}K^5K^2=S^{-145/2}K^{-11/2}$.

The positive-area assertion follows from the subset part of Lemma~\ref{low:surface:seed-rounding}, including its polynomial area loss.
\end{proof}

\begin{proposition}[Prepared dyadic rescaling]\label{low:5:qsa:preparedscaling}
Consider a prepared $Y$-oriented part of an actual common null ribbon, with $\theta=A\,ds+C\,dt$, $dz=n\,ds+dt$, and bounded coframe and inverse. Fix a dyadic $h\le c_P$ such that $|n|\le P^Ch$, $|r|\le P^Ch$, $hL\ge M$, and $W\ge M$. Assume the preparation of Definition~\ref{def:prepared-rescaling}. Here $L$ denotes the retained long length; if it arose from a fraction $q$ of a longer ribbon, assume also $hL/W\ge M$. The original pure transverse bounds are retained and the longitudinal loss has already been recorded.

\noindent\textup{(i)} Then, for $M$ a sufficiently large power of the earlier compact-model conditioning and the original null neighborhood sufficiently small, the polynomial
\[
\widehat F(X,Y,z)=h^2F(X/h,Y/h,z)
\]
in its actual endpoint-normalized frame belongs to a fixed small compact neighborhood of square-quadric models.

\noindent\textup{(ii)} Its mixed curvature is bounded away from zero; its first curvature is $r/h$. This conclusion controls all $55$ coefficients.
\end{proposition}
\begin{proof}
Use the coframe and graph identities in the proof of Lemma~\ref{low:4:sf:lem:newreconstruction}, with $U=hs,T=t$. The coframe determinant is the original $A-nC$, and its entries and inverse are bounded. Prepared finite jets have long scale $hL$ and transverse scale $W$. More explicitly, the pure transverse derivatives of the bounded rescaled geometric and normal scalars have bounds $P^CW^{-j}$ independent of a selected fraction $q$, while their mixed derivatives with at least one long component have bounds $P^C(hL)^{-i}W^{-j}$, $i\ge1$, at the selected positive-area good points. This follows by applying the preceding cellwise lemma before the change $U=hs$. If an original width-$W$ cell was recentered on a retained interval of length $L_0$, the new slope differs by $\tau'$, with $|\tau^{(j)}|\le P^CW L_0^{-j}$. The condition $hL_0/W\ge c_Px^{1-C\epsilon}$, where $x=hK$ is the selected amplitude-aspect scale, makes $\tau'/h$ and its higher rescaled jets small and preserves the slope bounds. Any extra polynomial longitudinal losses are absorbed by the same positive power of $x$. There is no corresponding requirement that $q^{-C}/W$ be small.

Normalize $\widehat d=0$ and the first graph jet to zero at the evaluation point. The second graph jet is then $Q_0(Y,z)=\tfrac12(r/h)Y^2+BYz$. The original small null window makes $u,v$ small, and the prepared relative jets make all derivatives of $r/h,B$ and all higher derivatives of $\widehat d$ small by powers of $(hL)^{-1}$ and $W^{-1}$. Differentiating those graph identities therefore makes the finite graph jets, through every order needed below, arbitrarily close to those of $Q_0$ after choosing $M$ large and the original defect small. No division by $r$ is used.

The rescaled graph value of $\widehat F_{XX}$ is exactly the original $\lambda$. Its nonconstant graph jets are small by the same preparation and the large scales. For $\widehat\phi(Y,z)=h^2\phi(Y/h,z)$, $\widehat Y\widehat\phi=hP_1$ and $\widehat Z\widehat\phi=h^2Q_1$, with the original scalars evaluated at $(Y/h,z)$. The prepared inverse coframe consequently bounds its nonconstant ordinary critical jets through order six, with all derivatives beyond its affine jet small. These finite-jet bounds hold at the selected good point.

To recover the complete phase, hold these actual graph jets fixed and use the following $55$ linear output rows: the $25$ grid rows of $\widehat F_X|_{\widehat c}$ with $0\le i,j\le4$; the nine grid rows of $\widehat F_{XX}|_{\widehat c}$ with $0\le i,j\le2$; the twenty nonconstant binary rows of $\widehat F|_{\widehat c}$ with $1\le i+j\le5$; and its additional $(3,3)$ row. The exact quintic reconstruction determinant in Section~\ref{low:5:qsc:reconstruction} proves that this map is uniformly invertible near $Q_0$. The first $25$ outputs vanish because this is the actual critical graph, the next nine are close to those of the constant $\lambda_0$, and the remaining critical outputs are close to those of the retained affine critical value. Applying the bounded inverse compares $\widehat F$ to $(\lambda_0/2)(X-Q_0)^2$ plus that affine value. This proves closeness of all $55$ coefficients without presupposing it.

The $55$-output reconstruction above retains the extra $(3,3)$ critical-value row. Fixed-format selection of a contained curved substrip from the positive-area good set with its selection costs then supplies a contained rescaled piece on which the normalized coefficient vectors lie in the stated fixed compact neighborhood of square-quadric models. All selections use the original algebraic curve parameter and affine transversal.
\end{proof}

\paragraph{The oblique count}\label{low:5:qsc:paid-oblique-section}
This argument uses the common actual ribbons and the every-anchor ranges proved in the quintic common-geometry section. Let a $Z$-oriented arc have its own length $L\ge W$, with $S=LW$ and $K=L/W$. In an anchor-relative tangential frame it is parameterized by $z=s$ and has an actual ruled ribbon
\[
 \Psi(s,t)=(\gamma_y(s)+t,s),\qquad |t|<c_PW,
\]
on a parameter interval of length comparable to $L$. The prepared slope $\beta=\gamma_y'-d_\gamma$ is bounded, the endpoint shear satisfies $|d|\le P^CK$, and the transverse shear difference is $|d(\Psi)-d_\gamma|\le P^C$. Consequently the pullback of the actual critical value has bounded gradient on the rectangle. Preparation gives $|d_\gamma^{(j)}|\le P^CKL^{-j}$ and $|\beta^{(j)}|\le P^CL^{-j}$. 

\begin{theorem}[Oblique quintic count]\label{low:5:qsc:paid-oblique}

\noindent\textup{(i)} These common oblique pieces have coefficient-volume bound at the fixed frame
\[
 \boxed{\qquad O_P\bigl(S^{-145/2}K^{-1/2}\bigr).\qquad}
\]
The weighted restriction chart has the stronger tail $O_P(S^{-145/2}K^{-41/2})$.

\noindent\textup{(ii)} The common count is $O_P(\Lambda S)$.
\end{theorem}
\begin{proof}
Retain a large-fraction simultaneous good set for all actual critical derivatives through order ten and all needed prepared curve derivatives. At a fixed evaluation point $(e,f)$, hold its null-normalized endpoint shear $d$ fixed before selecting an output map. If the point corresponds to a transverse parameter $t_0$, the centerline through it is $y=\gamma_y(s)+t_0$, whose nonconstant jets equal those of $\gamma_y$. Its first slope relative to the endpoint frame is $\gamma_y'(f)-d=\beta(f)+d_\gamma(f)-d$, so its range is bounded independently of $K$. Its derivatives of orders $j\ge2$ have ranges $O_P(KL^{1-j})$. No zeroth curve coefficient is discretized, because the evaluation point is integrated.

Round each order-$j$ curve derivative with mesh $c_PW/L^j$. The first derivative costs $O_P(K)$ possibilities. Each later derivative costs $O_P(K^2)$ possibilities. The grid net has size $\lesssim_P K(K^2)^4=K^9$, and the weighted net has size $\lesssim_P K(K^2)^9=K^{19}$. The first mesh is centered relative to the shear $d$ already held fixed; replacing its bounded range by the full global slope range would introduce an unnecessary additional factor $K$ and would lose the grid margin.

Here is the complete width-preservation argument. Write the true critical function in rectangle coordinates as $f(s,t)=\phi(\gamma_y(s)+t,s)$. A rounded curve gives, at the selected point, the finite-jet substitution $f(s,t+\varepsilon(s))$, where $\varepsilon(0)=0$ and $|\varepsilon^{(j)}(0)|\le c_PW/L^j$. In any mixed output of index $(i,j)$ with $j\ge1$, the chain rule uses only derivatives of $f$ with at least one transverse derivative. Every inserted derivative of $\varepsilon$ supplies the precise factor $WL^{-a}$ needed to retain the bound $L^{-i}W^{1-j}$. In a purely longitudinal output of order $i$, the additional terms have bounds $WL^{-i}\le L^{1-i}$. The direct purely longitudinal term already has the required bound. Thus all $35$ widths are preserved with a fixed constant. 

After subtracting the endpoint shear held fixed, the effective linear tangent coordinates are $(y,z)=(t+\beta_0s,s)$ with $|\beta_0|\le P^C$. At a null model $Q=(r/2)y^2+Byz$,
\[
Q(t+\beta_0s,s)=\left(B\beta_0+\frac r2\beta_0^2\right)s^2+(B+r\beta_0)st+\frac r2t^2.
\]
The null cutoff makes the mixed coefficient $B+r\beta_0$ bounded away from zero. The grid determinant at $r=0$ is therefore $(B+r\beta_0)^{35}$ evaluated at $r=0$, and compact continuity preserves its lower bound for the actual small null neighborhood. For the nonlinear curve jets, the bound $KL^{1-j}=L^{2-j}/W$ is $O_P(W^{-1})$ for every $j\ge2$, since $L\ge W\ge1$. The common inner buffer can be chosen so that $W\ge M=P^C$ with $M$ larger than the reciprocal determinant tolerance. Thus all these higher jets, and their rounded values, are sufficiently small for the grid coefficient atlas. Alternatively the weighted chart is available whenever its long quadratic coefficient is bounded below.

The coefficient map consists of the $16$ normal rows and the $35$ critical rows in these fixed rounded ordinary coordinates. With $d$ and the curve net held fixed, its phase derivative is the envelope restriction, so no derivative of the phase-dependent null root or endpoint shear is frozen illicitly. The actual null incidence is
\[
dF\,de\,df=C|\lambda^4B|\,dJ_{51}\,da\,db\,dc\,dd\,de\,df.
\]
For fixed $J,d,e,f$, the three-dimensional subgroup generated by $A,C,T_x$ has a bounded conditional Haar integral under the same fixed frame window. The quantitative three-column pivot is already present at every compact square-quadric anchor and persists in its actual neighborhood. The $d$ range contributes $K$. The every-anchor ranges $|e|\le P^CLK$ and $|f|\le P^CL$ give a candidate point area $O_P(L^2K)$, while each phase supplies good area $\gtrsim_PLW$. The ratio of these areas is at most $P^CK^2$. These integrations use the full joint incidence.

The grid output volume is $S^{-145/2}K^{-25/2}$, so its complete product of the volume factors is
\[
S^{-145/2}K^{-25/2}\,K^9\,K\,K^2=S^{-145/2}K^{-1/2}.
\]
The weighted output volume is $S^{-145/2}K^{-85/2}$, and its product of the volume factors is $S^{-145/2}K^{-85/2}K^{19}KK^2=S^{-145/2}K^{-41/2}$. The original assignment and count give the theorem.
\end{proof}

The following rescaled estimates use the same ribbon assignment. The quintic characteristic transport and common-ribbon theorems supply its count, every-anchor ranges, and actual ruled strips. The prepared-scaling proposition supplies complete rescaled coefficient compactness, using the $55$-output reconstruction map proved above. The following coefficient estimates are applied to contained compact substrips; a pointwise reconstruction statement is not substituted for their existence.

For a fixed dyadic $h$, write $\widehat F(X,Y,z)=h^2F(X/h,Y/h,z)$, $\widehat L=hL$, and $\widehat K=hK$. If $\widehat\phi$ is its actual critical value and $\widehat Y,\widehat Z$ are its actual null fields, put
\[
a_1=\widehat Y\widehat\phi=hP_1,\qquad b_1=\widehat Z\widehat\phi=h^2Q_1.
\]
These equalities already incorporate the coordinate scaling. Thus $a_1$ and $b_1$ are the scalar gradients of the rescaled phase, bounded by $P^Ch$ and $P^Ch^2$. Their coefficient variations therefore carry no extra factors of $h$ in the determinant.

The full coefficient determinant of this fixed scaling is $h^{2\cdot55-2\cdot70}=h^{-30}$, because the sums of the $x$ and $y$ monomial exponents are both $70$. Thus its measure character is exactly $dF=h^{30}d\widehat F$.

\paragraph{A contained ordinary ribbon in the slope case}
\begin{lemma}[An ordinary strip inside a ruled ribbon]\label{low:surface:ordinary-strip}
\textup{(i)} Suppose a contained ruled ribbon is parameterized by
\[
(y,z)=(\eta(s)+d_0(s)t,s+t),\qquad s\in I,\quad |t|<c_PW,
\]
where $|I|=L\ge W$. Put $\beta_0=\eta'-d_0$ and assume
\[
P^{-C}\le|\beta_0|\le P^C,\qquad |d_0'|\le P^C/W.
\]
After reducing the width constant and trimming the long interval, the image contains an ordinary strip $(y,z)=(\eta(u)+v,u)$ of dimensions comparable by powers of $P$ to $L,W$. If the actual shear differs from $d_0(s)$ by at most $\delta$ on the ruled ribbon, its difference from $d_0(u)$ on the ordinary strip is at most $\delta+O_P(c_P)$. In particular sufficiently small $\delta$ and width constant preserve a lower bound for $|\eta'(u)-d(y,z)|$.

\textup{(ii)} For a prepared $Y$-oriented slope certificate from Proposition~\ref{low:4:sf:common-ribbons}, retain a selected length $L_0$ and assume
\[
0<h\le c_P,\qquad |n|\asymp_Ph,\qquad |\tau^{(j)}|\le P^CW L_0^{-j},\qquad hL_0/W\text{ sufficiently large}.
\]
Suppose $d=d_{\rm seed}+\Delta$, with $|\Delta|\le P^C$, has the prepared bounds along the recentered seed
\[
|d_{\rm seed}'|\le P^C/W,\qquad |\Delta_s|\le P^C/L_0,\qquad |\Delta_t|\le P^C/W.
\]
Then the same conclusion holds with long length $hL_0$ and width $W$ after rescaling the shear by $h$.
\end{lemma}
\begin{proof}
At fixed $z=u$, substitute $t=u-s$. The derivative of the first coordinate is
\[
\frac{d}{ds}\bigl(\eta(s)+d_0(s)(u-s)\bigr)=\beta_0(s)+d_0'(s)(u-s).
\]
For $|u-s|<c_PW$, a sufficiently small width constant keeps its sign and a polynomial lower bound. Its image contains $|y-\eta(u)|<c'_PW$ for $u$ in the trimmed central interval. If $L\asymp_PW$, a contained central cap gives the same dimensions up to powers of $P$. The same derivative gives injectivity on each retained parameter rectangle. The corresponding seed parameters satisfy $|u-s|\le c_PW$, so $|d_0(s)-d_0(u)|=O_P(c_P)$; this proves the shear assertion.

After recentering on $t=\tau(s)$, the new $z$ derivative is $n+\tau'$, comparable to $h$. The rescaled shear derivative in that parameter is
\[
\frac{h(d_{\rm seed}'+\Delta_s+\tau'\Delta_t)}{n+\tau'}=O_P(W^{-1}).
\]
Indeed $W/(hL_0)$ is small and the numerator divided by $h$ is $O_P(W^{-1}+L_0^{-1})$. The original width was chosen so that $|t d_{\rm seed}^\prime|\le1/4$. The recentered seed has $\theta$-flux $1+\tau d_{\rm seed}^\prime-(n+\tau^\prime)\Delta$, which stays bounded away from zero when $h\le c_P$ and the width constant is small. Its rescaled relative slope is this flux times $h/(n+\tau^\prime)$, hence is conditioned above and away from zero. The first part applies at the retained scales with the original transverse width.
\end{proof}

Suppose a prepared quintic slope certificate has $|n|\asymp_Ph$, $|r|\le\epsilon_0h$, and $\widehat K\ge M$. Complete rescaling gives a compact ruled strip whose seed can be parameterized by $z$, with $\beta_0=\eta'-d_0$ conditioned above and away from zero and $|d_0'|\le P^C/W$. Apply Lemma~\ref{low:surface:ordinary-strip}, using its selected-interval version when necessary. The rescaled transverse shear difference is $O_P(h)$, so the contained ordinary strip has $|\eta'-\widehat d|\asymp_P1$ after choosing the width and $h$ small. Both scalar bounds $|a_1|\le P^Ch$ and $|b_1|\le P^Ch^2$ hold throughout it. For $f(u,v)=\widehat\phi(\eta(u)+v,u)$,
\[
\begin{aligned}
\partial_vf&=a_1(\eta(u)+v,u),\\
\partial_uf&=b_1(\eta(u)+v,u)+\bigl(\eta'(u)-\widehat d(\eta(u)+v,u)\bigr)a_1(\eta(u)+v,u).
\end{aligned}
\]
Hence $|\nabla f|\le P^Ch$ on this strip.

For determinant stability use compactness of the actual rescaled phase: the intrinsic jets of $\widehat Y\widehat d$ and $\widehat Z\widehat d$ vanish at the square-quadric model and are small in its prescribed neighborhood. Along the compact centerline, $d_0'=\widehat Z\widehat d+\beta_0\widehat Y\widehat d$ is small. Preparing this scalar and $\beta_0$ at the retained scale $\widehat L_0=hL_0$ makes the higher ordinary seed jets small enough for the grid determinant. Their net cardinality retains the stated longitudinal selection losses.

\paragraph{The oblique scalar grid determinant}
In ordinary strip coordinates $(u,v)$, remove the conditioned endpoint shear and let $\beta=\eta'(0)-\widehat d(0)$, so $|\beta|\asymp_P1$. At the limiting null model with zero first curvature, the graph is
\[
Q=B u(\beta u+v),\qquad \widehat Y=\partial_v,\qquad \widehat Z=\partial_u-\beta\partial_v.
\]
Keep the $16$ normal coefficient outputs. For $1\le i\le5$ use the purely longitudinal scalar rows $[u^{i-1}](\beta_*a_1+b_1)$, where $\beta_*$ is the fixed external rounded tangent, and for $0\le i\le5$, $1\le j\le5$ use the mixed scalar rows $[u^iv^{j-1}]b_1$. The notation denotes Taylor coefficients after composition with the fixed external curve jet. These are $5+30=35$ rows.

At $\beta_* =\beta$, their principal coefficient variations are obtained by applying $\partial_u$ and $\partial_u-\beta\partial_v$, respectively, to the $35$ restriction columns. Their ordinary-degree block determinants, in increasing degrees one through ten, are
\[
 \begin{array}{c|rrrrr}
 \text{degree}&1&2&3&4&5\\ \hline
 \det&-\beta&-4\beta^2&18\beta^3&96\beta^4&-600\beta^5
 \end{array}
\]
and
\[
 \begin{array}{c|rrrrr}
 \text{degree}&6&7&8&9&10\\ \hline
 \det&120B^5\beta^5&600B^8\beta^4&-600B^9\beta^3&-200B^8\beta^2&25B^5\beta
 \end{array}
\]
Their product is $-895795200000000000B^{35}\beta^{30}$. These formulas follow by differentiating the explicit columns $u^{m-j}v^j$ in degrees $m\le5$ and $Q^a u^{5-a-j}v^j$ in degrees $5+a$ and then taking the displayed coefficient rows. The accompanying exact certificate expands every block symbolically. The normal rows are coordinate pivots.

Let $\ell_Y$ denote the $Y$ derivative of the affine critical part of the model. For an actual model perturbation $R$, one has $\delta a_1=\widehat Y(R|_Q)$ and $\delta b_1=\widehat Z(R|_Q)+\ell_Y\delta\widehat d$. The null-root correction enters three critical degrees below the restriction column, as in the intrinsic atlas above, and hence does not change the diagonal blocks. Finite implicit jets have controlled normal and mixed-curvature denominators. Compact continuity consequently proves a uniform $51$-coordinate atlas for small $|r|/h$, small higher curve jets, actual quintic perturbations, and sufficiently accurate rounding of $\beta_*$. 

\begin{proposition}[Quintic slope estimate]\label{low:5:qsc:slope}

\noindent\textup{(i)} On the prepared slope certificate just described, the unchanged original piece has coefficient-volume bound at the fixed frame
\[
 \boxed{\qquad P^C h^{22}S^{-145/2}K^{-1/2}.\qquad}
\]

\noindent\textup{(ii)} If the certificate occupies a fraction $q$ of the original ribbon, the right side is multiplied by a fixed power $q^{-C}$, with all its scale and reconstruction requirements retained in that certificate.
\end{proposition}
\begin{proof}
Use the contained ordinary strip above and prepare the two bounded scalars $a_1,b_1$ separately through total derivative order ten. The purely longitudinal outputs have widths $P^Ch\widehat L^{1-i}$ and the mixed outputs have widths $P^Ch^2\widehat L^{-i}W^{1-j}$. Since the grid index sums are $90,90$, their product is $P^Ch^{65}\widehat L^{-85}W^{-60}$. The smaller mixed amplitude comes from the independently bounded scalar $b_1$. 

At a fixed point, hold its endpoint shear $\widehat d$ fixed and round the five ordinary centerline derivatives. The relative first slope has bounded range and costs $\widehat K$. Derivatives of orders two through five have ranges $P^C\widehat K\widehat L^{1-j}$ and mesh $c_PW/\widehat L^j$, costing $\widehat K^2$ each. Thus the net has size $\lesssim_P\widehat K(\widehat K^2)^4=\widehat K^9$. Higher seed jets are $O_P(W^{-1})$ or smaller; the large fixed buffer supplies the compact determinant tolerance.

Round the external curve by a finite transverse shift, separately in the scalar $\beta_*a_1+b_1$ of amplitude $h$ and in $b_1$ of amplitude $h^2$. Its derivatives have size $P^CW\widehat L^{-j}$, so the chain rule preserves each scalar's prepared widths with its own amplitude. The phase-dependent null field stays in the coefficient map with the lower-degree correction computed above.

The complete fixed scaling has $dF=h^{30}d\widehat F$. With the endpoint shear and center fixed, the original compact anchor bounds the $H_3=\{d=e=f=0\}$ fiber. Its rescaled coordinates are $(a_g,hb_g,hc_g)$, so the corresponding Haar volume is $h^2$ times the original bounded fiber. This uses the fixed frame, whose rescaled coefficient vector need not be bounded. The endpoint shear interval has length $P^ChK=P^C\widehat K$. The rescaled candidate center range has area $O_P(hL^2K)$, while the contained strip has area $\gtrsim_PhLW$. Thus the joint incidence and Tonelli step cost $K^2$. The full product of the volume factors is
\[
h^{32}h^{65}(hL)^{-85}W^{-60}(hK)^9(hK)K^2
 =h^{22}S^{-145/2}K^{-1/2}.
\]
All coefficient integrations are performed in the same original anchored integral. A fraction $q$ changes the finite preparation constants, available long length, and good-area denominator only by fixed powers; these are the stated $q^{-C}$ loss.
\end{proof}

\begin{proposition}[Quintic curvature estimate]\label{low:5:qsc:curvature}

\noindent\textup{(i)} On a prepared certificate with $|r|\asymp_Ph$, $|n|\le P^Ch$, and $hK\ge M$, the corresponding coefficient-volume bound at the fixed frame is
\[
 \boxed{\qquad P^Ch^{18}S^{-145/2}K^{-11/2}.\qquad}
\]

\noindent\textup{(ii)} A fraction $q$ again introduces only a fixed factor $q^{-C}$. The fixed frame is allowed to have zero first curvature.
\end{proposition}
\begin{proof}
Proposition~\ref{low:5:qsa:preparedscaling} bounds every coefficient of the normalized rescaled phase, with $|\widehat r|\asymp_P1$ on this branch. Its $Y$-oriented seed has $\widehat Y$ component bounded away from zero. The Gauss-strip calculation in the proof of Proposition~\ref{low:4:sf:prop:newcurvaturecount} therefore gives a contained prepared strip of long length $\widehat L=hL$ and width $W$, with comparable metric and area and fixed format.

Prepare the actual scalar gradients $a_1,b_1$ separately. At the evaluation point, retain the $16$ normal outputs determining curvature and the two zeroth scalar outputs among the $35$ rows before differentiating the direction-normalized outputs. The Gauss null direction $(B,1)$ and a transverse long direction form a uniformly invertible coordinate system. At the model write $y=u$, $z=\beta u+v$, so $Q=A'u^2+B'uv$ with $|B'|\asymp_P1$. Purely longitudinal rows of the frozen long scalar have amplitude $h$ and mixed rows of $b_1$ have amplitude $h^2$, with principal variations $\partial_u(R|_Q)$ and $\partial_v(R|_Q)$. The restriction determinant and derivative multipliers give $358318080000000(B')^{35}$. Null-root variation lowers critical degree by three. Apply the chain-rule row subtraction in the proof of Proposition~\ref{low:4:sf:prop:newcurvaturecount}, using these retained normal and scalar outputs to determine the direction coefficients. The null-field variation is the lower-degree correction just described. Continuity from the model determinant then bounds the absolute Jacobian of the $51$-output map below by $P^{-C}$ on the prescribed neighborhood.

The width product is $h^{65}(hL)^{-85}W^{-60}$. The Gauss centerline derivatives have ranges $P^C(hL)^{1-j}$; the five required meshes cost $(hK)^5$. The output map uses Gauss derivatives of the normalized endpoint phase, with the external coordinate jet held fixed. Indeed Gauss labels transform by translations under every fixed upper-affine map, and the transformed scalars $a_1,b_1$ and their Gauss derivatives equal those computed from the transformed phase. Thus this map uses the four-dimensional conditional incidence.

Apply Lemma~\partIref{surf:physical-gauss-density} to the intrinsic $51$-output map. Its endpoint density has pivot $|\lambda|^4|\widehat r|$; the fixed-frame $H_0$ bound uses the original mixed curvature $B$ and remains valid at zero first curvature. The substitution $\widehat\eta=\widehat e-d_gf$ is made on the right of the joint incidence identity, where fixing $(\widehat\eta,f)$ leaves one $H_0$ orbit. As in Proposition~\ref{low:4:sf:prop:newcurvaturecount}, its coordinates scale to $(a_g,hb_g,hc_g,hd_g)$, giving $h^3$; together with $dF=h^{30}d\widehat F$ this gives $h^{33}$.

The endpoint-range comparison in the proof of Proposition~\ref{low:4:sf:prop:newcurvaturecount} applies since $hK\ge M$ and gives candidate area divided by good area at most $P^CK^2$. The full product of the volume factors is
\[
h^{33}h^{65}(hL)^{-85}W^{-60}(hK)^5K^2
 =h^{18}S^{-145/2}K^{-11/2},
\]
as required. The same finite preparation and area accounting give a fixed power of $q^{-1}$ when a fractional certificate is used.
\end{proof}

\paragraph{Exhaustion and the surface estimate}
\begin{theorem}[Common exhaustion of actual quintic null pieces]\label{low:5:qsc:null-exhaustion}

\noindent\textup{(i)} Every common null piece has count at most $P^C\Lambda S$. The classification includes zero first curvature and zero critical gradients. It preserves the original assignment, its anchors, and its physical coefficient character.

\noindent\textup{(ii)} Its coefficient-volume bound at the fixed frame is
\[
 \boxed{\qquad P^CS^{-145/2}K^{-1/4}.\qquad}
\]
\end{theorem}
\begin{proof}
Theorem~\ref{low:5:qsc:paid-oblique} covers the $Z$-oriented arcs and the $Y$-oriented arcs with $|n|\ge c_P$, using the subdivision made before assigning anchors. Choose the null curvature cutoff below this angle threshold and the reconstruction tolerances. If $|r|K\le H_0$ on a fixed positive area fraction, Theorem~\ref{low:5:qsa:smallaspect} gives $P^CS^{-145/2}K^{-11/2}$.

On the remaining positive fraction, apply Lemma~\ref{low:surface:amplitude-selection} with the scalar list of Definition~\ref{def:prepared-rescaling}. It supplies $q\ge c_{P,\epsilon}(hK)^{-\epsilon}$, full transverse width, and the rescaled-length and slope bounds. Proposition~\ref{low:5:qsa:preparedscaling} gives complete compact phases on a fixed good fraction; the last part of the selection lemma supplies a contained compact substrip inside the original ribbon.

Let $C$ dominate the preparation and counting losses. Apply the dyadic summation in the proof of Theorem~\ref{low:4:sf:thm:newexhaustion} to Propositions~\ref{low:5:qsc:slope} and~\ref{low:5:qsc:curvature}, using $q^{-C}\le P^C(hK)^{C\epsilon}$. Their $(h,K)$ powers are $(22,-1/2)$ and $(18,-11/2)$, so the sums are bounded by $P^CS^{-145/2}K^{-1/2+C\epsilon}$ and $P^CS^{-145/2}K^{-11/2+C\epsilon}$. The dyadic $h$ powers $22+C\epsilon$ and $18+C\epsilon$ are positive and $h$ has a fixed upper bound. Choose $C\epsilon<1/4$, then $H_0$, the common buffer and the inner defect in that order. These sums, the small-aspect branch and the oblique branch prove the claimed tail.
\end{proof}

\begin{theorem}[The complete quintic surface branch]\label{low:5:qsc:surface-count}

\noindent\textup{(i)} Every fixed-catalogue state in the sufficiently small fifth-singular-value neighborhood has a common actual-surface assignment. A piece of scales $S=LW$, $K=L/W$ has count at most $P^C\Lambda S$.

\noindent\textup{(ii)} Its coefficient-volume bound at the fixed frame is at most
\[
P^CS^{-145/2}K^{-1/4}.
\]

\noindent\textup{(iii)} The summation over surface pieces is valid for $q_p=(1+\eta)(p-1)<145/2$, with the original weight of the fixed localization piece $v_j^{p-70}$.
\end{theorem}
\begin{proof}
The rank-four neighborhood and normalization give at most twelve actual frames per tangential center. Select the first valid chart as in Section~\ref{quintic:coordinates}, retaining the complete coefficient and physical restrictions. The nonnull Gauss tail proved there is stronger than the asserted bound; the common null assignment and Theorem~\ref{low:5:qsc:null-exhaustion} give the null tail. The prescribed outer windows retain each original state after normalization.

The same width multiplicity and $S\ge W^2$ verify Proposition~\partIref{an:common-incidence} for $q_p<145/2$, since the aspect exponent is $1/4$. This gives the stated weight; subunit scales cost powers of $P$.
\end{proof}


The nine complete residual row lists are printed in the certificate supplement, Section~\ref{ext-supp:supp:quintic-residual-rows}; the determinant values, exponent sums, grades, and all case estimates remain in the main argument.


\section{Sextic surface estimates}\label{low:degree6}

Throughout this section $k=6$. We use the physical weighted mass bound of Paper I, Section~\partIref{an:section}, and the regular estimates of Section~\ref{reg:finite6}, retaining full admissibility and all original coefficients. The inflection branch uses the graph $F_{xx}=0$ and quintic transport for $G=F_x$; its coefficient count and the nonnull estimate are in Section~\ref{low:6:inflection-counts}. Sections~\ref{low:6:ordinary-connection} and~\ref{low:6:ordinary-counts} treat critical null graphs. Theorem~\ref{low:6:as36:all-surface} separates their positive rate bands from the identically zero branch. Section~\ref{completion:physical-II} completes the physical sum.

\begin{theorem}\label{low:6:s36:main}
Every $a\in C_c^\infty(\mathbb R^3)$ satisfies $E_a\in L^p(V_6)$ whenever $p>129$.
\end{theorem}

The proof of Theorem~\ref{low:6:s36:main} is completed in Section~\ref{completion:physical-II}, using the coefficient estimates below.

\par\smallskip
\noindent\textit{Proof stages for Theorem~\ref{low:6:s36:main}.}\par\nobreak
\noindent\textup{Step 1.} Rank-four entry and the two normal-root assignments.\par\nobreak
\noindent\textup{Step 2.} Inflection geometry and its complete coefficient estimates.\par\nobreak
\noindent\textup{Step 3.} Critical-graph transport and buffered cover.\par\nobreak
\noindent\textup{Step 4.} The three critical coefficient counts.\par\nobreak
\noindent\textup{Step 5.} Branch assembly and physical completion.\par\nobreak
\smallskip

\subsection{Rank-four models and normal-root assignment}\label{low:6:s36m:section}\label{low:6:s36a:section}

\paragraph{The exact stopped rank-four sextic models}

Use the fixed upper-affine flag and the augmented states of Definition~\partIref{def:augmented-state}, without an ordered moment-balance restriction. Write $\mathcal V_d=V_d$ for $d=5,6$, and put $\ell(u)=p+Bu$ and $\widehat J=([J],\ell)$. Assume
\begin{equation}\label{low:6:s36m:stop}
\|\widehat J\|\le P,\qquad \|[h\cdot\nabla J]\|^2+|Bh|^2\ge\kappa^2|h|^2\quad(h\in\mathbb R^3),\qquad P\ge2,\quad\kappa\ge P^{-1}.
\end{equation}
The physical matrix $B$ may be singular. Denote the orbit matrix in the ordered generator basis of Definition~\partIref{def:augmented-state} by $R_{\widehat J}$. Constants are removed only from the phase output, so a kernel field may differentiate $J$ to a nonzero constant.

\paragraph{Sextic model profiles}

Put $q_0=x-yz$ and $q_\gamma=x-yz+\gamma z^3/3$. The sextic $q_\gamma^2$ has critical graph $x=yz-\gamma z^3/3$ and identically zero critical value. Its null field on this graph is $Z=\partial_z+\gamma z\partial_y$, and its shear satisfies $Zd=\gamma$ for $d=\gamma z$. Thus no bound for null shear can follow from derivatives of the critical value alone in the sextic model family. The phase $q_0^3$ has $F_x=F_{xx}=0$ and $F_{xxx}=6$ on $q_0=0$, so a lower bound for $F_{xx}$ is also unavailable. Finally, $F=q_0^2+\varepsilon q_0^3$ satisfies
\[
\left.(F,F_x,F_{xx},F_{xxx})\right|_{q_0=0}=(0,0,2,6\varepsilon).
\]
The normal third derivative therefore supplies an additional coefficient direction. The full $83$-row orbit matrices of these three examples have rank four and their phase translation matrices have rank three, as can also be checked directly from the two kernel fields in Theorem~\ref{low:6:s36m:classification} below.

Lemma~\partIref{geo:normal-pivot} applies to~\eqref{low:6:s36m:stop}; it supplies the normal-column pivot and $\sigma_4(R_{\widehat J})\ge P^{-C}$.

\paragraph{A parametrization that includes both branches and their intersection}

\begin{theorem}[Exact stopped rank-four sextics]\label{low:6:s36m:classification}

\noindent\textup{(i)} Under \eqref{low:6:s36m:stop}, the equality $\operatorname{rank}R_{\widehat J}=4$ holds only if $B=0$ and
\begin{equation}\label{low:6:s36m:model}
\begin{gathered}
J=f(q)+\lambda\bigl(y-\gamma z^2/2\bigr)+\mu z\pmod{\mathbb R},\qquad f(q)=a_3q^3+a_2q^2+a_1q,\\
q=x-\frac r2y^2-syz-ty+\frac{t\gamma-v}{2}z^2+\frac{s\gamma}{3}z^3-wz,\\
r\gamma=0,\qquad a_3\gamma=0,\qquad \Delta:=rv-s^2\ne0,\qquad (a_2,a_3)\ne(0,0).
\end{gathered}
\end{equation}

\noindent\textup{(ii)} Conversely every phase in \eqref{low:6:s36m:model} has injective phase translation differential and, with $B=0$, upper-affine orbit rank four.

\noindent\textup{(iii)} On the bounded stopped stratum its displayed parameters satisfy
\begin{equation}\label{low:6:s36m:parameter-bounds}
|r|+|s|+|t|+|v|+|w|+|\gamma|+|a_1|+|a_2|+|a_3|+|\lambda|+|\mu|+|\Delta|^{-1}+\bigl(a_2^2+a_3^2\bigr)^{-1/2}\le P^C.
\end{equation}
There is no uniform lower bound for $|\gamma|$ or $|a_3|$ separately.
\end{theorem}
\begin{proof}
Theorem~\partIref{geo:classification} with $k=6$ gives the displayed invariant and $r\gamma=0$. When $\gamma=0$, its profile has degree at most three; when $\gamma\ne0$, it has degree at most two. Thus $a_3\gamma=0$, while nonlinearity gives $(a_2,a_3)\ne(0,0)$. The theorem gives $B=0$, the converse and phase translation injectivity. Its compact conditioned parameter bounds give~\eqref{low:6:s36m:parameter-bounds}, including the inverse norm of the nonlinear coefficient vector, without dividing by $\gamma$ or $a_3$ separately.
\end{proof}

The quadratic-invariant branch $\gamma=0$ has five graph coefficients, three coefficients of $f$, and two tangential linear coefficients. Its dimension is ten where $a_3\ne0$. The cubic-invariant branch $\gamma\ne0$ has $r=a_3=0$ and dimension nine. Their closures meet where $r=a_3=\gamma=0$. The coefficient of $z^2$ in the tangential summand is $-\lambda\gamma/2$ in the cubic branch; deleting it would change the exact model family and its kernel. The parameter bound does not require division by $\gamma$ and therefore does not fail as the two branches meet.

\begin{corollary}[Quantitative entry into the enlarged model family]\label{low:6:s36m:entry}
For each fixed $L\ge1$, a sufficiently large $C_L$ has the following property. If \eqref{low:6:s36m:stop} holds and $\sigma_5(R_{\widehat J})\le P^{-C_L}$, then $\|B\|\le\kappa/2$ and $J$ is within $P^{-L}$ in all $83$ coefficients of a model \eqref{low:6:s36m:model} satisfying \eqref{low:6:s36m:parameter-bounds} with a fixed enlarged exponent. In particular
\[
\|[h\cdot\nabla J]\|\ge (\sqrt3/2)\kappa|h|.
\]
\end{corollary}
\begin{proof}
Apply Lemma~\partIref{geo:rank-entry} with $k=6$ and complete defect $P^{-L}$ to the model family in Theorem~\ref{low:6:s36m:classification}. Choose the small-fifth-pivot threshold to make $\|B\|\le\kappa/2$ as well. Subtraction of $|Bh|^2\le\kappa^2|h|^2/4$ from the augmented stop proves the final phase-only estimate.
\end{proof}

The proximity is in the original upper-affine flag; the graph is constructed from the actual sextic.


\paragraph{Critical and inflection graphs}

This reduction uses the model family and parameter bounds of Theorem~\ref{low:6:s36m:classification}. For each label $j$, replace its representative $g_j$ by the normalized comparison frame $g_j\circ k_j$ specified below; coefficient integration remains in the original frame.

\begin{lemma}[Two normal-root alternatives]\label{low:6:s36a:two-types}
Fix the compact model conditioning in \eqref{low:6:s36m:parameter-bounds}. There are powers $\eta=P^{-C_1}$ and $\delta=P^{-C_2}$ with the following property.

\noindent\textup{(i)} Every actual sextic within coefficient distance $\delta$ of that family admits a normalized frame $g\circ k$, where $g$ is its original frame, $k\in\mathcal U_0=\{e=f=0\}$, and $\|k\|+\|k^{-1}\|\le P^C$ in homogeneous affine matrices, satisfying at least one of the following alternatives:
\begin{align*}
\mathrm I:&\quad F_{xx}=F_{xxy}=F_{xxz}=0,\qquad |F_{xxx}|\ge P^{-C},\\
\mathrm C:&\quad F_x=F_{xy}=F_{xz}=0,\qquad |F_{xx}|\ge P^{-C}.
\end{align*}
Type I is an inflection sheet and type C a critical sheet in Definition~\partIref{def:normal-sheet}; these labels are independent of the tangential shear charts I and II. All displayed values are at the origin of the corrected frame.

\noindent\textup{(ii)} In either type, the tangential Hessian of the defining normal equation is nonsingular with a polynomial inverse bound and can be put in the usual nonnull or null shear chart.

\noindent\textup{(iii)} Type I has a complete derivative phase neighborhood of a conditioned square-quadric model. Type C includes all cubic-invariant models; its exact critical graph may have degree three.
\end{lemma}
\begin{proof}
Use a model representation \eqref{low:6:s36m:model} only to prove existence, and write its invariant as $q=x-Q(y,z)$, where $Q$ may initially have degree three. If $|a_3|\ge\eta$, then $\gamma=0$, $Q$ is quadratic, and the inflection graph is $x=Q-a_2/(3a_3)$. Its shift and all model parameters are polynomially bounded. Its normal pivot is $6a_3$, and the tangential Hessian of $F_{xx}$ after tangent normalization is $-6a_3D^2Q$. Choose $\delta$ small enough for the actual inflection normalization to persist. The actual derivative phase is close in its complete quintic coefficient vector to $3a_3(x-Q+a_2/(3a_3))^2$ modulo its constant. This proves type I.

Suppose $|a_3|<\eta$. Choose $\eta$ below a sufficiently small power of the lower bound for $(a_2^2+a_3^2)^{1/2}$. Then $|a_2|\ge P^{-C}$. Set $q_\ell=-a_1/(2a_2)$, so $|q_\ell|\le P^C$. On $|q-q_\ell|\le1$, the normal second derivative $2a_2+6a_3q$ has the constant sign of $a_2$ and absolute value at least $|a_2|$, provided $\eta$ is chosen small enough. At the center the normal first derivative is $3a_3q_\ell^2$, whose size is arbitrarily small relative to $|a_2|$ with the same choice. The endpoint signs therefore give a unique simple model critical root in this interval. Coefficient proximity on the polynomially bounded corresponding spatial strip preserves the endpoint signs, derivative lower bound, and root. For the exact model its critical graph is the graph $q=\text{constant}$, and its tangent-normalized Hessian is $-f''(q)D^2Q$. In the cubic-invariant branch the Hessian determinant of $Q$ is $-s^2$ at every point, while in the quadratic branch it is $rv-s^2$; both have polynomial inverse bounds. Thus the Hessian and its selected shear pivot persist for the actual root. This proves type C. The successive choices of $\eta$, then the finite normalization bounds, then $\delta$ involve only powers of the original compact conditioning parameter.
\end{proof}

\begin{proposition}[An assignment retaining original labels]\label{low:6:s36a:assignment}
Let $g_j$ be the fixed upper-affine family and let the selected original labels satisfy the small-fifth-pivot model neighborhood.

\noindent\textup{(i)} The labels admit a fixed-format measurable assignment, chosen for all labels of each phase before a label is fixed for coefficient integration, to a type I or type C normalized frame $h_j=g_j\circ k_j$ with $k_j,k_j^{-1}$ bounded by $P^C$. The tangential center is unchanged. Original coefficients, physical restrictions, weights, and labels remain attached to this assignment.

\noindent\textup{(ii)} At each prescribed tangential center, there are at most twelve normalized type I frames and at most fifteen normalized type C frames satisfying the nonzero pivot conditions.
\end{proposition}
\begin{proof}
Define validity by the actual normal-root equations, the actual first tangential derivative equations, the selected nonnull or null shear equation, polynomial coefficient and pivot bounds, and the bounded-right-comparison inequality. These are fixed-format semialgebraic conditions. The model representation in Lemma~\ref{low:6:s36a:two-types} proves that at least one candidate is valid; a model need not be selected as part of the assignment. Select the first valid type in a fixed order, then the first valid real normal root and shear root in their ordinary ordered lists. The normal translation and the two normal shears are determined by those roots. This gives a measurable fixed-format map in the phase and the original label.

For type I, the normal equation is $F_{xx}(c,e,f)=0$, of degree at most four in $c$. At a root with $F_{xxx}\ne0$, its first tangential derivative equations determine the two normal shears as in the quintic normalization calculation. There is at most one nonnull shear and at most two null shears. Hence there are at most twelve normalized frames. For type C, the equation $F_x(c,e,f)=0$ has degree at most five, giving at most fifteen frames by the same calculation. These are bounds on all actual normalized frames satisfying the pivots, not just the near-model roots used to establish existence.

Since $k_j\in\mathcal U_0$, the composition $g_j\circ k_j$ preserves the tangential center. The bounded norms of $k_j$ and $k_j^{-1}$ give $\sigma_i(B_jDk_j)\asymp_P\sigma_i(B_j)$, while the new physical center is $p_j+B_jk_j(0)$. This comparison with the original physical data is polynomially controlled. Two original labels whose normalized frames lie in one bounded right neighborhood have original relative frames in a bounded right neighborhood as well. Thus the established local packing bound transfers with only its polynomial enlargement.
\end{proof}

For a common component construction, use the actual normalized root sheets of each type with nested outer coefficient, pivot, and transported physical windows. Original priority choices determine which labels are assigned to each type, but do not restrict a continuation path to its original priority region. In the inflection type the outer bounds on all nonconstant coefficients of the original sextic $F$ and its physical data may be retained in the derivative phase common-ribbon construction because normalized tangential translations of the compact cubic-normal model preserve those windows after a polynomial enlargement, and the inner defect is chosen last. In type C, Proposition~\ref{low:6:s36a:assignment} supplies the finite actual normalization and packing comparison. The nonnull geometry is treated in Proposition~\ref{low:6:s36i:nonnull-count}; the null geometry uses the sextic connection identity, buffered relocation and the corridor assignment of Theorem~\ref{low:6:as36:all-surface}. No quintic critical-value transport estimate is used for this type.


\subsection{Normalization, geometry of the derivative phase, and coefficient maps}\label{low:6:inflection-geometry}\label{low:6:qsf:normalization}\label{low:6:qsf:sec:packing}\label{low:6:qsf:preparation}\label{low:6:qsf:sec:gauss}\label{low:6:qsf:null-incidence-section}\label{low:6:s36q:module}\label{low:6:sai:section}\label{low:6:inflection-atlas}\label{low:6:nc36:section}

\subsubsection{Normalization and rescaling}
Fix $D_{\mathrm{ph}}\in\{5,6\}$ and apply the ordinary square-quadric normalization of Section~\ref{quartic:normalization} in all $N_{D_{\mathrm{ph}}}$ coefficients, with its conditioned model family and nested physical and stopping windows. Completing the square includes $b(x-Q)$. The regular sextic graph has $D_{\mathrm{ph}}=6$; the inflection derivative phase has $D_{\mathrm{ph}}=5$.

For $G=F_x$, the covariance and normalized equations are~\eqref{low:4:sf:eq:covariance}--\eqref{low:4:sf:eq:normalized}, with the same $\mathcal U_0$ correction, two curvature charts, and retained inverse pivots. Lemma~\partIref{atlas:normal-fibers} gives at most $3(D_{\mathrm{ph}}-1)$ frames per center, Jacobians~\eqref{low:4:sf:eq:jacI}--\eqref{low:4:sf:eq:jacII}, and $a=a_0$, $b=b_0+da_0$ as in~\partIeqref{atlas:normal-ab}. The fiber and projected-packing bounds are~\eqref{low:4:sf:eq:fibervol}--\eqref{low:4:sf:eq:area}. Lemma~\partIref{geo:good-points-strips} gives~\eqref{low:4:sf:critical-widths}; rounding through order $N$ costs $P^CK^N$, with every chain-rule row prepared and the evaluation point integrated.

The Gauss intrinsic-point root bound is $(D_{\mathrm{ph}}-1)(D_{\mathrm{ph}}-2)$. Theorem~\partIref{geo:gauss} retains all regular branches, artificial cuts, and fixed-frame restrictions; coefficient inversion uses Lemma~\partIref{geo:coefficient-map}.

\paragraph{Joint coefficient and frame integration}
Use Lemma~\partIref{surf:haar} for the conditional $H_0,H_3,H$ fibers, with their fixed-frame phase-only pivots and the prescribed polynomially bounded right displacements.

The Gauss and null coefficient disintegrations are \partIeqref{surf:gauss-density} and \partIeqref{surf:null-density-eq} with $m=2$, phase degree $D_{\mathrm{ph}}$, and $N_{D_{\mathrm{ph}}}-4$ free slice coordinates. In the null chart use $J_x=J_{xy}=J_{xz}=J_{xzz}=0$ at the origin, $\lambda=J_{xx}(0)$ and $B=(c_*)_{yz}(0)$ with their conditioned inverses; $r=(c_*)_{yy}(0)$ may vanish. Apply Lemma~\partIref{surf:null-density} to a full intrinsic output map, using $H_0$ or, for an external endpoint shear, $H_3$ as specified there. Physical-center Gauss incidence is Lemma~\partIref{surf:physical-gauss-density}; its conditional Haar bound retains the fixed-frame pivot independently of the endpoint's $r$.

\paragraph{Application to a quintic derivative phase}
The reconstruction, characteristic transport and common-ribbon results of Sections~\ref{quintic:coordinates} and~\ref{quintic:counts} will be applied to the actual derivative phase $G=F_x$. Their hypotheses are a complete degree-five square-quadric neighborhood for $G$, the regular root $G_x=F_{xx}=0$, a consistently selected null root, bounds $P^C$ for the complete normalized $G$ coefficients, $|G_{xx}|^{-1}$ and $|B|^{-1}$, and $|v|\le\nu_0=P^{-C_0}$ with the small constant of Theorem~\ref{low:5:qnt:transport}. For the polynomial transverse-Jacobian and common-ribbon conclusions, choose the null curvature cutoff $\sup|r|\le P^{-C_2}$ after the preceding transport constants, as in Theorem~\ref{low:5:qnt:transport}; states above this cutoff use the nonnull chart. Their sixth critical derivative is the derivative of $G$'s critical value; it is not an identity for the sextic critical value.

The complete quintic reconstruction retains $25$ graph-equation rows, nine second-normal rows, twenty nonconstant binary critical-value rows and the extra $(3,3)$ critical row. The grid restriction in Lemma~\partIref{surf:normal-filtration} supplies the degree-four and degree-two blocks used in that proof.

Theorem~\ref{low:5:qnt:transport} gives polynomial signed characteristic transport. Proposition~\ref{low:5:qng:ribbons} gives pieces $L\ge W\ge M$, $S=LW$, $K=L/W$, count $P^C\Lambda S$, contained ribbon area $P^{-C}S$, and fixed-frame endpoint bounds $|d|\le P^CK$, $|e|\le P^CLK$, $|f|\le P^CL$. Its center curve uses the finite adjacency graph of buffered boundary-arc neighborhoods and bounded access bridges; the graph is not indexed by all balanced boxes. The initial access bridge and the full assigned arc remain available after selecting a prepared subarc, so neither the assignment nor $L$ changes.

Use Lemmas~\partIref{geo:relative-preparation} and~\ref{low:5:qsa:cellwise-preparation} before amplitude selection. Pure transverse estimates keep their original width $W$; later longitudinal selection replaces $L$ by $L_0\ge P^{-C}q^CL$ without imposing $q^{-C}/W$. For Proposition~\ref{low:5:qsa:preparedscaling}, prepare $n/h,r/h$, $B$, the normal scalar, $P_1,Q_1$, $u=Yd$, $v=Zd$, the coframe and its inverse combinations, and the bounded shear difference to the stated finite order. Retain $|n|,|r|\le P^Ch$, $hL_0\ge M$, $W\ge M$, and $hL_0/W\ge M$ on the contained ribbon. The sextic application below verifies these hypotheses with its additional scalar fields.

\paragraph{Normal multiplicity and derivative phase geometry}

Work in $V_6$ with its coefficient measure. The derivative phase results give cubic-normal geometry. The full $83$-coefficient incidence estimate combines this geometry with the reconstruction of the original phase below. Here $P$ bounds the fixed number of coefficients, inverse pivots and description complexities.


\paragraph{The derivative phase supplies the cubic-normal geometry}

Let $Q$ be a quadratic polynomial in $y,z$ with nonsingular Hessian, put $q=x-Q(y,z)$, and consider a conditioned cubic-normal model
\begin{equation}\label{low:6:sai:cubic-model}
F_0=a_3q^3+a_2q^2+a_1q+\ell(y,z),\qquad |a_3|\ge P^{-C},\qquad |\det D^2Q|\ge P^{-C},
\end{equation}
where $\ell$ is affine and all displayed coefficients are bounded by powers of $P$. The coefficient space is taken modulo additive constants; the coefficient $a_1$ of the nonconstant polynomial $q$ is retained. The polynomially sized outer normal interval contains $q=-a_2/(3a_3)$. The model parameter bounds make this enlargement polynomial in $P$.

\begin{proposition}[Derivative phase reduction]\label{low:6:sai:derivative-reduction}

\noindent\textup{(i)} If $F\in\mathcal V_6$ is sufficiently close to \eqref{low:6:sai:cubic-model} in its complete coefficient vector, then $G=F_x$, modulo its constant, lies in a conditioned square-quadric neighborhood in $\mathcal V_5$. Its actual critical equation is $G_x=F_{xx}=0$, and its simple-root derivative is $G_{xx}=F_{xxx}$.

\noindent\textup{(ii)} The normal shift, all inverse pivots, and the complete derivative phase translation stop are bounded by powers of $P$.
\end{proposition}
\begin{proof}
The exact derivative of the model is
\[
(F_0)_x=3a_3\left(q+\frac{a_2}{3a_3}\right)^2+a_1-\frac{a_2^2}{3a_3}.
\]
Its graph is the quadratic graph $x=Q-a_2/(3a_3)$, and its normal second derivative is $6a_3$. Normal translation to this graph is bounded by a power of $P$. Pullback under that fixed translation is a bounded linear map on the full finite-dimensional coefficient space, so a sufficiently small complete sextic perturbation gives a sufficiently small complete quintic derivative perturbation.

The derivative phase has a quantitative translation stop without using the physical entries. Let $\widetilde Q=Q-a_2/(3a_3)$, $\widetilde q=x-\widetilde Q$, and let $h=(h_x,h_T)$. The nonconstant part of the directional derivative of $3a_3\widetilde q^2$ is
\[
6a_3\widetilde q\bigl(h_x-\nabla\widetilde Q\cdot h_T\bigr).
\]
Writing $D_h=h\cdot\nabla(3a_3\widetilde q^2)$, we have
\[
\begin{pmatrix}[xy]D_h\\{[xz]D_h}\end{pmatrix}=-6a_3D^2Q\,h_T,\qquad [x]D_h=6a_3\bigl(h_x-\nabla Q(0)\cdot h_T\bigr).
\]
Consequently $\|[D_h]\|\ge P^{-C}|h|$. A sufficiently small complete perturbation preserves this inequality. The same polynomial bounds give persistence of the simple normal root and of the nonsingular tangential Hessian by quantitative implicit inversion.
\end{proof}

Write the upper-affine group as
\[
g(x,y,z)=(x+ay+bz+c_0,\ y+dz+e,\ z+f).
\]
The first column of $Dg$ is $e_x$, so
\begin{equation}\label{low:6:sai:derivative-covariance}
\partial_x(F\circ g)=(F_x)\circ g,\qquad \partial_x^2(F\circ g)=(F_{xx})\circ g.
\end{equation}
This covariance permits use of the derivative phase in the original sextic frames, with coefficient integration over the full sextic space.

\begin{lemma}[Inflection normalization]\label{low:6:sai:inflection-normalization}
Suppose the phase expressed in a frame $g$ lies in a sufficiently small coefficient neighborhood of \eqref{low:6:sai:cubic-model}.

\noindent\textup{(i)} There is $k\in\mathcal U_0$ such that the affine matrices of $k,k^{-1}$ have norms at most $P^C$ and, in the frame $g\circ k$ with the same tangential center,
\[
F_{xx}(0)=F_{xxy}(0)=F_{xxz}(0)=0,
\]
and the tangential Hessian of $F_{xx}$ is in either the nonnull chart with zero mixed entry or the null chart with zero second diagonal entry. The pivots $\mu=F_{xxx}(0)$ and the relevant Hessian entry have polynomial lower bounds.

\noindent\textup{(ii)} At a fixed tangential center there are at most twelve such normalized frames.

\noindent\textup{(iii)} The transverse four-variable normalization Jacobian is $|\mu|^3|H_{11}|$ in the nonnull chart and $2|\mu|^3|H_{12}|$ in the null chart, where $H$ is the Hessian of $F_{xx}$ after the tangent normalization.
\end{lemma}
\begin{proof}
At the model $F_{xx}=6a_3(q+a_2/(3a_3))$, the absolute normal pivot is at least $P^{-C}$, and the tangent Hessian and its inverse have norms at most $P^C$. The implicit function theorem with these bounds and Lemma~\partIref{geo:two-chart-shear} therefore give the bounded root, tangent and shear corrections for a sufficiently small complete perturbation. Apply Lemma~\partIref{atlas:normal-fibers} to the normal polynomial $F_{xx}$, of degree at most four, with pivot $\mu=F_{xxx}$: it gives the twelve-frame bound and both displayed Jacobians.
\end{proof}

Let $x=c_F(y,z)$ be the selected actual inflection graph, and define two scalars
\begin{equation}\label{low:6:sai:two-scalars}
\phi(y,z)=F(c_F(y,z),y,z),\qquad \psi(y,z)=F_x(c_F(y,z),y,z).
\end{equation}
The constant in $\phi$ may be suppressed. The constant in $\psi$ is an ordinary sextic coefficient output and is retained. Since the graph is defined by $F_{xx}=0$, both quantities transform as scalars under the induced tangential affine action. In particular, if $T_g(y,z)=(y+dz+e,z+f)$, then the transformed graph is $c_F\circ T_g-ay-bz-c_0$, and
\begin{equation}\label{low:6:sai:scalar-covariance}
\phi_{F\circ g}=\phi_F\circ T_g\pmod{\mathbb R},\qquad \psi_{F\circ g}=\psi_F\circ T_g.
\end{equation}

\begin{lemma}[Bounds for the two inflection scalars]\label{low:6:sai:scalar-compatibility}
Retain the normalized sextic coefficient window on a common inflection sheet.

\noindent\textup{(i)} At the origin of each normalized tangent frame, $|\nabla\phi(0)|+|\nabla\psi(0)|+|\psi(0)|\le P^C$.

\noindent\textup{(ii)} These bounds also hold in the intrinsic nonnull Gauss coordinates and in every bounded-coframe null ribbon supplied by the quintic derivative phase geometry.
\end{lemma}
\begin{proof}
Write $p=c_y$, $q_1=c_z$, and use tangent vectors $\partial_y+p\partial_x$ and $\partial_z+q_1\partial_x$. Then
\[
\phi_y=F_y+pF_x,\qquad \phi_z=F_z+q_1F_x,
\]
which are exactly the tangential linear coefficients of the tangent-normalized sextic phase. It is not necessary that $F_x$ vanish. Since $F_{xx}=0$ on the graph,
\[
\psi_y=F_{xy},\qquad \psi_z=F_{xz}.
\]
In a frame with tangential shear $d$, the corresponding bounded pairs are $(\phi_y,\phi_z+d\phi_y)=(J_y,J_z)$ and $(\psi_y,\psi_z+d\psi_y)=(J_{xy},J_{xz})$. Also $\psi=J_x$. All are bounded by the sextic admissibility conditions.

In the nonnull Gauss chart $p=c_y$, put $r=c_{yy}$ and solve $y=Y(p,z)$. Since $Y_p=1/r$ and $Y_z=-c_{yz}/r=d$, the two gradient pairs become $(J_y/r,J_z)$ and $(J_{xy}/r,J_{xz})$. The inverse curvature bound controls $1/r$. In a null ribbon, the coframe $(dy-d\,dz,dz)$ and its inverse have bounded coefficients in the ribbon parameters. The preceding pairs therefore control both ordinary ribbon gradients. 
\end{proof}

\begin{proposition}[Common geometry with the full sextic admissibility retained]\label{low:6:sai:common-inflection-geometry}

\noindent\textup{(i)} On a common inflection sheet with the full sextic admissibility conditions retained, the derivative phase $G=F_x$ in a conditioned square-quadric neighborhood satisfies the characteristic transport bounds of Theorem~\ref{low:5:qnt:transport}.

\noindent\textup{(ii)} The null states admit an assignment before anchoring to pieces of scales $L\ge W\ge M$, $S=LW$, $K=L/W$, with count $P^C\Lambda S$, contained ribbon area at least $P^{-C}S$, and anchor-relative endpoint ranges
\[
|d|\le P^CK,\qquad |e|\le P^CLK,\qquad |f|\le P^CL.
\]
The common coefficient--point factor is $P^CK^2$.

\noindent\textup{(iii)} The nonnull Gauss assignment and finite-fiber packing have the same statement as for the derivative phase. 
\end{proposition}
\begin{proof}
Apply Proposition~\ref{low:6:sai:derivative-reduction} to $G=F_x$. Its degree is at most five, its critical graph is exactly $F_{xx}=0$, and every normalized upper-affine frame is simultaneously the corresponding derivative phase frame by \eqref{low:6:sai:derivative-covariance}. Its normalized coefficients and pivots satisfy the square-quadric hypotheses. The finite root and null-direction multiplicities and the transverse Jacobians are precisely those in Lemma~\ref{low:6:sai:inflection-normalization}.

Assign the original labels to their priority branches. Use Lemma~\partIref{geo:model-box} with $k=6,m=3$ and the core construction of Lemma~\partIref{geo:core}, retaining the full sextic coefficient, root-pivot and transported physical restrictions in the outer domain. Fix the polynomial buffer $M$ before the complete inner defect. Polynomial model pullback and the conditioned root corrections then preserve these outer restrictions throughout the prescribed box. Continue on this full admissible domain, with priority and later amplitude labels attached to the original states.

Apply Theorem~\ref{low:5:qnt:transport} on each connected part of the nested outer characteristic for $G$, then Proposition~\ref{low:5:qng:ribbons} on that outer domain, retaining the full sextic bounds. The width surrogate $\rho$, balanced boxes, selected arcs and good-point restrictions have fixed format, with boundedly changed complexity. Assign the original labels using the buffered ribbons as certificates. They have the same count, contained area, endpoint ranges and bounded coframe. Lemma~\ref{low:6:sai:scalar-compatibility} gives bounded scalar data throughout. The anchor-relative endpoint area is $O_P(L^2K)$ and the good area is $\gtrsim_PS=LW$, giving ratio $O_P(K^2)$. This disintegration uses the original fixed-frame sextic coefficients.
\end{proof}

The null transport invoked here is a theorem about the quintic polynomial $G$, whose critical value is $\psi$. In particular, the exact sixth frozen derivative identity acts on $\psi$ in this branch. The obstruction $F=q_\gamma^2$ does not satisfy this reduction: then $F_x=2q_\gamma$ has constant normal derivative and is not a stopped square-quadric derivative phase. That model belongs to the regular-root branch treated below.
Use the buffer fixed in Proposition~\ref{low:6:sai:common-inflection-geometry}. Apply the full-width transverse preparation of Lemma~\ref{low:5:qsa:cellwise-preparation}, including the additional sextic scalars listed below, before selecting $h$ or an area fraction.
\begin{proposition}[Prepared inflection rescaling in all $83$ coefficients]\label{low:6:sai:inflection-rescaling}
Consider a prepared $Y$-oriented inflection ribbon from Proposition~\ref{low:6:sai:common-inflection-geometry}, and a dyadic $h\le c_P$ for which $|n|,|r|\le P^Ch$, where $n$ is the long $z$-slope and $r=c_{yy}$. Assume that the already prepared full-width transverse cells and the retained long length $L_0$ satisfy the derivative phase conditions of Proposition~\ref{low:5:qsa:preparedscaling}, including $W\ge M$ and $hL_0/W\ge M$. Assume that the finite preparation list, fixed before amplitude selection, also includes the full-sextic scalars $P_F=Y\phi$, $Q_F=Z\phi$, and the scalar value $\psi$.

\noindent\textup{(i)} Then
\[
\widehat F(X,Y,z)=h^3F(X/h,Y/h,z)
\]
in its actual endpoint-normalized frame has all $83$ coefficients bounded and lies in a prescribed small neighborhood of a conditioned cubic-normal model
\[
\frac{\mu_0}{6}(X-Q_0(Y,z))^3+b_0(X-Q_0(Y,z))+\ell_0(Y,z),\qquad \mu_0=F_{xxx}|_{\rm endpoint},
\]

\noindent\textup{(ii)} Here $\widehat G=\partial_X\widehat F$, $Q_0=\tfrac12(r/h)Y^2+BYz$, $|\mu_0|,|B|$ have polynomial lower bounds, $b_0=\widehat G|_{\rm graph,endpoint}=h^2\psi|_{\rm endpoint}$, and $\ell_0$ is affine.
\end{proposition}
\begin{proof}
Set $G=F_x$. Differentiating the rescaling gives the identity
\[
\widehat G:=\partial_X\widehat F=h^2G(X/h,Y/h,z).
\]
Thus $\widehat G$ is precisely the usual degree-five square-quadric rescaling of the derivative phase. Apply Proposition~\ref{low:5:qsa:preparedscaling}, including the complete reconstruction in Section~\ref{low:5:qsc:reconstruction}, with normal scalar $\mu=G_{xx}|_{\rm graph}=F_{xxx}|_{\rm graph}$ and critical value $\psi$. The needed normalized derivative phase scalars are derivatives of the full sextic and were retained in the preparation. This controls all $55$ nonconstant coefficients of $\widehat G$ and the finite graph jet $\widehat c$ near $Q_0$, with an arbitrarily small error after the stated choices. The affine tangential part retained by the quintic theorem has derivative sizes $O_P(h)$ and $O_P(h^2)$ here, so it is included in that small error by fixing $c_P$ small. Its omitted constant is not free: on the graph its value is $h^2\psi$. The full sextic window bounds $\psi$, so adding this constant gives closeness of the entire $56$-coefficient polynomial $\widehat G$ to
\[
G_0=\frac{\mu_0}{2}(X-Q_0)^2+b_0.
\]
This also identifies why the exponent $h^3$, rather than $h^2$, is required for the sextic phase: its third normal derivative remains $\mu$.

The geometric rescaled coframe is the same one used for $\widehat G$. The two intrinsic components of the rescaled full critical differential are
\[
\widehat Y\widehat\phi=h^2P_F,\qquad \widehat Z\widehat\phi=h^3Q_F.
\]
Prepare $P_F,Q_F$ on the original full-width cells before amplitude selection, exactly as the derivative phase scalars. At selected good points, the inverse coframe and their mixed jets have long scale $hL_0$ and transverse scale $W$. Repeated finite differentiation therefore controls the ordinary nonconstant jets of $\widehat\phi$ through order six needed for the reconstruction below; all terms beyond its affine jet are arbitrarily small when $M$ is large and the earlier geometric defect is small. For the later coefficient atlas, carry out this same preparation through at least total order twelve, retaining all $\phi$ rows $0\le i,j\le6$ and all $\psi$ rows $0\le i,j\le4$, together with every geometric and normal jet occurring in their chain rules. If a later longitudinal certificate had fraction $q$, its loss has already shortened $L_0$; the pure transverse bounds remain $P^CW^{-j}$. 

To recover the rescaled sextic, let $\mathcal P H$ be the monomial antiderivative in $X$ vanishing at $X=0$, for a polynomial $H(X,Y,z)$ of degree at most five including its constant. This bounded linear map takes values in the sextics divisible by $X$. The identity
\[
\widehat F=\mathcal P\widehat G+K(Y,z)
\]
holds with a binary polynomial $K$ of degree at most six, modulo its constant. It controls the $56$ coefficients of monomials containing $X$ by the just proved bound for $\widehat G$. Evaluating on the actual graph gives
\[
K=\widehat\phi-(\mathcal P\widehat G)(\widehat c(Y,z),Y,z).
\]
The $27$ nonconstant ordinary binary derivatives through total order six at the endpoint determine every coefficient of $K$ exactly. The bounds for $\widehat\phi$, the full polynomial $\widehat G$, and the finite graph jet therefore control all these remaining coefficients. At the limiting model, $\mathcal P G_0$ evaluated at $X=Q_0$ equals $\mu_0Q_0^3/6+b_0Q_0$. Hence the recovered polynomial is $\mu_0(X-Q_0)^3/6+b_0(X-Q_0)+\ell_0$, as claimed. Continuity of these finitely many explicit operations gives the prescribed small complete-coefficient neighborhood. This is a pointwise reconstruction argument with the actual graph jet held fixed, not a change of coefficient variables in a measure calculation.
\end{proof}

\subsubsection{Complete coefficient maps and measure factors}
We construct maps on all $83$ nonconstant sextic coefficients and prove their local volume estimates under explicit geometric hypotheses. The subsequent sections verify these hypotheses for the common assignments. In this atlas construction, $q=x-Q(y,z)$ with $Q$ a conditioned nonsingular homogeneous quadratic form; $P$ bounds all model parameters and inverse pivots. Outputs are Taylor coefficients unless derivatives are specified.

\paragraph{The normal coefficient data}
The normal-multiplicity examples in Subsection~\ref{low:6:s36m:section} show that a full sextic coefficient map must retain a third normal output. For a regular critical graph selected by $F_x(c(y,z),y,z)=0$, graph motion is given by $\delta c=-R_x|_c/F_{xx}|_c$; its coordinates are counted among the $83$ outputs.

\paragraph{The 48+25+9+1 coefficient decomposition}
Use the monomial basis of $\mathcal C_d$ specified in the proof of Lemma~\partIref{surf:normal-filtration} and the index sets $\mathcal I_d^{\mathrm g}=I_d$, $\mathcal I_d^{\mathrm w}=I_d^{\mathrm w}$ of Definition~\partIref{def:jet-indices}. In the output lists below, undecorated $\mathcal I_d$ and $\mathcal E$ mean the index sets for the selected grid or weighted chart. Remove the constant only from the phase-value block.

Lemma~\partIref{surf:normal-filtration} with $k=6$, omitting the constant input and output, gives the following $83$ coordinates:
\[
 \begin{array}{c|c|c}
 \text{field}&\text{indices}&\text{number of outputs}\\ \hline
F|_Q&\mathcal I_6\setminus\{(0,0)\}&48\\
F_x|_Q&\mathcal I_4&25\\
F_{xx}|_Q&\mathcal I_2&9\\
F_{xxx}|_Q&\{(0,0)\}&1.
 \end{array}
\]
Its absolute determinant in the original monomial coefficients is $3072|B|^{80}$ in the grid chart and $3072|A|^{80}$ in the weighted chart.
The block sizes are $49-1,25,9,1$; their factorial factor is $2^9\cdot6=3072$ and their exponent is $\binom83+\binom63+\binom43=80$.

\paragraph{Actual regular and inflection graphs}
First take an exact regular-root model
\[
F_0=a_2q^2+a_3q^3+\ell_y y+\ell_z z,\qquad \lambda=2a_2\ne0,\qquad \mu=6a_3.
\]
For a nearby actual phase, let $c$ be the selected graph $F_x(c,y,z)=0$, and put $\phi=F|_c$, $\lambda_F=F_{xx}|_c$, and $\mu_F=F_{xxx}|_c$. At the model, a perturbation $R$ gives
\begin{equation}\label{low:6:nc36:regular-variation}
 \delta c=-\frac{R_x|_Q}{\lambda},\qquad
 \delta\phi=R|_Q,\qquad
 \delta\lambda_F=R_{xx}|_Q+\mu\delta c,\qquad
 \delta\mu_F=R_{xxx}|_Q.
\end{equation}
Consequently the outputs $(\phi_{\mathcal I_6\setminus0},c_{\mathcal I_4},(\lambda_F)_{\mathcal I_2},\mu_F(0))$ give a full $83$-coordinate map. The $\mu\delta c$ term in the third block is eliminated by graph rows already present. The determinant is the determinant in Lemma~\partIref{surf:normal-filtration} times $(-1/\lambda)^{25}$.

For the normal cubic model, including its lower normal parameter, write
\[
F_0=a_3q^3+a_1q+\ell_y y+\ell_z z,\qquad \mu=6a_3\ne0.
\]
Use the actual inflection graph $F_{xx}(c,y,z)=0$, and set $\phi=F|_c$, $\psi=F_x|_c$, and $\mu_F=F_{xxx}|_c$. Its variations are
\begin{equation}\label{low:6:nc36:inflection-variation}
 \delta c=-\frac{R_{xx}|_Q}{\mu},\qquad
 \delta\phi=R|_Q+a_1\delta c,\qquad
 \delta\psi=R_x|_Q,\qquad
 \delta\mu_F=R_{xxx}|_Q.
\end{equation}
The $a_1\delta c$ term is the contribution of graph motion when $a_1\ne0$.

\begin{proposition}[Two coefficient maps with varying graphs]\label{low:6:nc36:moving-atlases}

\noindent\textup{(i)} The regular-root outputs in the preceding paragraph and the inflection outputs
\[
 \phi_{\mathcal I_6\setminus0},\qquad
 \psi_{\mathcal I_4},\qquad
c_{\mathcal I_2},\qquad
 \mu_F(0)
\]
have Jacobians with absolute determinant at least $P^{-C}$ on sufficiently small coefficient neighborhoods of their respective model families, with the stated model bounds retained. Their derivative matrices and inverse derivative matrices have norms at most $P^C$.

\noindent\textup{(ii)} For ordinary jets the inflection determinant at a model has absolute value $3072|B|^{80}|\mu|^{-9}$ in the grid chart, and $3072|A|^{80}|\mu|^{-9}$ in the weighted chart. No inverse normal root separation and no inverse $F_{xx}$ enter the inflection atlas.
\end{proposition}
\begin{proof}
The regular statement follows from \eqref{low:6:nc36:regular-variation}. For the inflection statement, assign weight $2$ to $x$ and weight $1$ to $y,z$. The principal fixed-graph map in Lemma~\partIref{surf:normal-filtration} is block diagonal by this weight when output weights include the normal derivative order. The correction $a_1\delta c$ in the phase block lowers the weight by four. It is strictly triangular and does not change the determinant. The graph block contributes $(-1/\mu)^9$. All finite graph jets are algebraic functions of the coefficients, with denominators powers of the selected normal pivot. Quantitative continuity gives uniform neighborhoods and inverse bounds depending polynomially on $P$.
\end{proof}

The exact numerical constants above refer to ordinary scalar and graph jets. The actual Gauss and ordered-null scalar outputs used in the quadratic-graph construction satisfy the same quantitative invertibility bounds, with their coordinate factors retained. At a model the Gauss correction in $\delta\phi$ is $-\ell_y\partial_y(\delta c)/r$, where $r=Q_{yy}\ne0$. In the null chart, with $Q=Ay^2+Byz$, one has $\delta d=-(\delta c)_{zz}/(2B)$ and the correction in $\delta(Y^iZ^j\phi)$ is $\ell_y\partial_y^i\partial_z^{j-1}(\delta d)$ for $j\ge1$. These corrections lower the input weight by three in the regular-root case and by five in the inflection case. The model scalar fields $\lambda_F$, $\psi$, and $\mu_F$ are constant, so changing their differentiation operators contributes no further model variation. These lower-weight corrections preserve the determinant after the scalar coordinate change has been accounted for. In the Gauss chart the actual coordinate is $p=c_y$; at $Q=Ay^2+Byz+Cz^2$ it is $p=2Ay+Bz$. The substitution $y=(p-Bz)/(2A)$ is triangular on the weighted output sets, with diagonal $(2A)^{-i}$ on index $(i,j)$. Thus the Gauss atlas carries the corresponding nonzero conditioned coordinate factor, and its determinant lower bound and inverse bound remain powers of $P$. The ordered-null field is already the ordinary coordinate field at the reference model and does not acquire this Gauss factor. These bounds apply on the stated compact neighborhoods. The rescaled estimates below use their stated preparation hypotheses.

\paragraph{Additional coordinates from \texorpdfstring{$F_{xx}$}{Fxx}}
Additional coordinates from the bounded second normal scalar give the following atlas. Define
\[
 \mathcal E^{\mathrm g}=\{(i,j):i+j\le4\}\cup\{(3,2),(2,3),(3,3)\}
 =\mathcal I_3^{\mathrm g}\cup\{(4,0),(0,4)\},
\]
and
\[
 \mathcal E^{\mathrm w}=\{(i,j):i+j\le4\}\cup\{(5,0),(4,1),(6,0)\}.
\]
Both sets have $18$ elements and contain the respective set $\mathcal I_3$.

\begin{proposition}[Sixteen second-normal outputs]\label{low:6:nc36:strong-atlas}

\noindent\textup{(i)} In the nondegenerate critical-graph setting, the outputs
\[
 \phi_{\mathcal I_6\setminus0},\qquad
 (\lambda_F)_{\mathcal I_3},\qquad
c_{\mathcal E},\qquad
 \mu_F(0)
\]
give a square coefficient map with $48+16+18+1=83$ outputs. It gives local coordinates where its Jacobian is nonzero.

\noindent\textup{(ii)} For ordinary jets its model absolute determinant is $56623104|B|^{73}|\lambda|^{-18}$ in the grid chart and $56623104|A|^{73}|\lambda|^{-18}$ in the weighted chart, independent of the cubic profile coefficient $a_3$. For actual Gauss scalar jets with ordinary graph coordinates, the weighted determinant has the additional factor $|2A|^{-237}$.

\noindent\textup{(iii)} On sufficiently small complete-coefficient neighborhoods of the stated conditioned models, the Gauss atlas and the separately normalized ordered-null atlas retain determinant lower bounds $P^{-C}$ and inverse bounds $P^C$.
\end{proposition}
\begin{proof}
Apply Lemma~\partIref{surf:strong-atlas} with $D=6$. Its grid and weighted output lists are exactly those stated here, since $\mathcal E^{\mathrm g}=E_6$ and $\mathcal E^{\mathrm w}=E_6^{\mathrm w}$. It gives a nonzero model Jacobian in each selected chart and local coordinates wherever the corresponding actual Jacobian is nonzero, together with the factor $|\lambda|^{-18}$, and the quadratic exponent $t_6=126-3\cdot16-2-3=73$.

At $Q=yz$, exact integer elimination in the displayed sextic output convention gives raw absolute determinant $56623104=2^{21}3^3$; the weighted model $Q=y^2$ has the same constant. The accompanying exact matrix verification checks both values. Lemma~\partIref{surf:strong-atlas} gives independence from the cubic profile coefficient and the higher graph and fixed external jets. The actual Gauss coordinate change contributes $|2A|^{-237}$, since the longitudinal index sums are $203$ in the critical block and $34$ in the second-normal block. The lower-weight moving-coordinate corrections and compact continuity give the stated Gauss and ordered-null atlas bounds.
\end{proof}

\begin{lemma}[Ordinary-jet elimination with large normal derivatives]\label{low:6:nc36:large-normal-ordinary}
Fix the evaluation coordinates, allowing a fixed externally rounded curve substitution. Suppose $F_x(c,y,z)=0$ and $\lambda_F(0)\ne0$.

\noindent\textup{(i)} For the ordinary-jet version of Proposition~\ref{low:6:nc36:strong-atlas}, the coefficient Jacobian depends, apart from the factor $\lambda_F(0)^{-18}$, only on the finite jets of $c$ and the fixed external coordinate substitution. It has no dependence on the size of $\mu_F$ or its tangential derivatives.

\noindent\textup{(ii)} Thus the same lower bound persists when the graph jets remain sufficiently close to the model, even if $\mu_F$ is large.
\end{lemma}
\begin{proof}
Use the moving-graph row elimination in Lemma~\partIref{surf:strong-atlas} with $D=6$, keeping the external tangential germ fixed. Both graph-index sets $\mathcal E$ are downward closed and contain $\mathcal I_3$. Hence every term from differentiating $\mu_F\delta c$ is removed by graph rows, as is the term $F_{xxxx}(0)\delta c(0)$ in the last output. The eighteen graph rows give the factor $\lambda_F(0)^{-18}$. The remaining determinant is the fixed-graph determinant at the actual finite graph jet and fixed external substitution, independently of the higher normal derivatives. Continuity in those graph data proves the stated lower bound.
\end{proof}

The qualification ``ordinary-jet'' in Lemma~\ref{low:6:nc36:large-normal-ordinary} is necessary for the stronger $18$-graph-output atlas. A phase-dependent Gauss or null operator has its own variation, involving graph jets that need not lie in $\mathcal E$.


\paragraph{Step 1. The full coefficient and frame measure}
Normalize the chosen graph by its value, two first derivatives, and one selected mixed or null second derivative at the origin. In the regular case these four graph equations are equivalent to the corresponding four linear conditions on $F_x$; in the inflection case they are the corresponding four linear conditions on $F_{xx}$. The four chosen graph indices occur in all the atlases above. Removing them leaves $79$ coefficient outputs on the $79$-dimensional normalized slice; no coefficient direction has been discarded.

The complete coefficient action has determinant one on the upper-unipotent group. The usual four transverse orbit columns give, at a fixed Gauss point,
\begin{equation}\label{low:6:nc36:gauss-density}
dF=C|\kappa|^4|r|^2\,dJ_{79}\,dh_H,
\end{equation}
where $\kappa=F_{xx}|_c$ in the regular-root chart and $\kappa=F_{xxx}|_c$ in the inflection chart. At fixed physical center in the null chart, the full coefficient--point density is
\begin{equation}\label{low:6:nc36:null-density}
dF\,de\,df=C'2|\kappa|^4|B|\,dJ_{79}\,dh_{H_0}\,d\eta\,df.
\end{equation}
These identities retain all $83$ coefficients. Differentiate the four constraints in the orbit columns as in Lemma~\partIref{surf:null-density}. Replacing $F_x$ by $F_{xx}$ changes the normal pivot to $\kappa$ and leaves the group coordinates fixed; the other $79$ coefficients retain their identity block. Lemma~\partIref{geo:original-subgroup-pivots} gives the phase-only four-column bound at the fixed frame. Fixing all $79$ outputs leaves finitely many normalized phases, uniformly, and one conditional $H$ or $H_0$ fiber. Lemma~\partIref{surf:haar} bounds its measure by $O_P(1)$; in the $H_0$ case, the factorization $h=k g_0$ of Lemma~\partIref{surf:null-density} places the restriction on that subgroup. If an ordinary output map depends on the external endpoint shear, condition on it and integrate only the $H_3$ fiber as specified there. Its shear interval remains a separate cost.

\paragraph{Step 2. Volume of an assigned surface certificate}\label{low:6:nc36:certificate-section}
Suppose a commonly assigned surface piece has length $L$, width $W$, count $O_P(\Lambda LW)$, and a positive good proportion on which the required derivatives of its bounded scalar fields satisfy the usual anisotropic bounds. Put $S=LW$ and $K=L/W$. The graph output intervals and the undifferentiated third-normal output interval are kept bounded. These hypotheses are imposed on the actual piece.

The $48$ critical-value grid outputs have product
\[
L^{-141}W^{-105}=S^{-123}K^{-18}.
\]
For a bounded scalar, the full $\mathcal I_d^{\mathrm g}$ output grid instead has widths $L^{-i}W^{-j}$, including its bounded constant row. Consequently the sixteen second-normal outputs contribute $L^{-24}W^{-24}=S^{-24}$, and the twenty-five inflection-slope outputs contribute $L^{-50}W^{-50}=S^{-50}$. The raw full-coefficient products are therefore
\begin{equation}\label{low:6:nc36:raw-products}
 \begin{array}{c|c|c}
 \text{atlas}&\text{grid product}&\text{weighted product}\\ \hline
 \text{regular root, strong}&S^{-147}K^{-18}&S^{-147}K^{-78}\\
 \text{inflection}&S^{-173}K^{-18}&S^{-173}K^{-88}.
 \end{array}
\end{equation}
For comparison, the weaker regular atlas gives $S^{-132}K^{-18}$ in the grid chart. 

\par\smallskip\noindent\textit{Step 3. Geometric costs and paid products.}\par\nobreak\smallskip

In a Gauss certificate with the usual fixed-format curved-strip preparation, the grid and weighted curve nets cost $K^6$ and $K^{12}$, respectively, and the endpoint-area comparison costs $K$. The products from \eqref{low:6:nc36:raw-products} are
\[
\begin{aligned}
\text{regular grid:}\quad&S^{-147}K^{-18}K^6K=S^{-147}K^{-11},\\
\text{inflection grid:}\quad&S^{-173}K^{-18}K^6K=S^{-173}K^{-11},\\
\text{regular weighted:}\quad&S^{-147}K^{-78}K^{12}K=S^{-147}K^{-65},\\
\text{inflection weighted:}\quad&S^{-173}K^{-88}K^{12}K=S^{-173}K^{-75}.
\end{aligned}
\]
These bounds use \eqref{low:6:nc36:gauss-density} and include the fixed-frame Haar restriction.

For an ordinary oblique certificate, assume that the relative first curve slope has bounded range, that the $j$th external centerline derivative has range $O_P(KL^{1-j})$ for $2\le j\le6$, and that the endpoint shear has range $O_P(K)$. Rounding the first jet with mesh $W/L$ costs $K$; rounding each higher jet with mesh $W/L^j$ costs $K^2$. The external net therefore has size $\lesssim_P K(K^2)^5=K^{11}$. The endpoint shear interval costs $K$, and the common point comparison costs $K^2$. With the $H_3$ fiber kept conditional as above, the resulting full tails are
\begin{equation}\label{low:6:nc36:paid-products}
\begin{aligned}
(S^{-147}K^{-18})K^{11}KK^2&=S^{-147}K^{-4},\\
(S^{-173}K^{-18})K^{11}KK^2&=S^{-173}K^{-4}.
\end{aligned}
\end{equation}
These bounds apply to geometric certificates with the derivative ranges and endpoint shear range assumed in this paragraph. 

\paragraph{Step 4. The small scalar output map in an oblique inflection ribbon}
Let $u$ be the long coordinate and $v$ the transverse coordinate. At a frozen model write $Q=B u(\beta u+v)$, with $B\beta\ne0$, and $a=\partial_v\phi$, $b=(\partial_u-\beta\partial_v)\phi$. Use six purely longitudinal rows
\[
 \partial_u^{i-1}(\beta_*a+b),\qquad 1\le i\le6,
\]
with the value $\beta_*=\beta$ frozen at the evaluation point, together with forty-two mixed rows
\[
 \partial_u^i\partial_v^{j-1}b,\qquad 0\le i\le6,\quad 1\le j\le6.
\]
On the nonconstant monomial basis of $\mathcal C_6$ their determinant, in the derivative convention just displayed, has absolute value
\begin{equation}\label{low:6:nc36:scalar-det}
 2^{172}3^{74}5^{31}|\beta|^{42}|B|^{56}.
\end{equation}
To see the nonvanishing and the characters directly, group by tangential degree. In degrees $m\le6$, the pure row fixes the purely longitudinal coefficient and the remaining triangular scalar block contributes $\beta^m$ times nonzero integer factors. In degree $6+a$, there are $7-a$ columns $Q^a u^{6-a-j}v^j$, and the scalar block contributes $\beta^{7-a}B^{a(7-a)}$ times nonzero integer factors. Summing gives $21+21=42$ for $\beta$ and $56$ for $B$. Multiplying the derivative factorials gives the constant in \eqref{low:6:nc36:scalar-det}.

Each scalar is rounded separately. In particular one does not replace the smaller scalar $b$ by a difference of two rounded quantities of the larger size of $a$. Six centerline jets suffice: the mixed critical rows contain at most six transverse-to-long ordered derivatives, and the additional bounded normal scalar has order at most four. All extra coordinate derivatives entering the chain rule are prepared before rounding; rounding a degree-$j$ centerline jet in mesh $W/L^j$ then preserves each scalar's own derivative width.

\paragraph{Inflection rescaling}\label{low:6:nc36:inflection-scaling}

The two rescaled inflection estimates are proved with their contained certificates in Propositions~\ref{low:6:s36i:slope} and~\ref{low:6:s36i:curvature}. Their coefficient factors, distinct scalar amplitudes, fixed-frame Haar measures, seed nets, and point comparisons are included in \eqref{low:6:nc36:inflection-paid-h} and \eqref{low:6:nc36:inflection-intrinsic-h}. The exhaustive amplitude selection and uniform sum are proved in Theorem~\ref{low:6:s36i:null-count}.

\paragraph{Quadratic-graph coefficient exponents}
The conditioned regular-root quadratic atlases have raw area exponent $147$, and the conditioned inflection quadratic atlases have raw area exponent $173$, both exceeding the endpoint requirement $128$. These are raw coefficient exponents before geometric costs. The inflection graph is treated through $G=F_x$ and Proposition~\ref{low:6:sai:common-inflection-geometry}, retaining the full original $F$ admissibility condition and the extra $F$ scalars. The subsequent critical assignment retains the corresponding fixed frames and supplies the actual certificates used by its coefficient counts.

The critical count below uses algebraic corridors and the ordinary coefficient map specified in Theorem~\ref{low:6:so36:exact-atlas}. That theorem proves, by triangular elimination with respect to the stated weights, that this map's Jacobian depends on the normal scalar and the quadratic graph coefficients and is independent of the higher graph jets.

\subsection{Inflection and nonnull coefficient estimates}\label{low:6:inflection-counts}\label{low:6:s36i:small-aspect-section}\label{low:6:s36i:nonnull-section}\label{low:6:inflection-exhaustion}\label{low:6:s36i:exhaustion-section}
The graph for this count is the actual inflection graph $F_{xx}=0$. We retain $\phi=F|_c$, $\psi=F_x|_c$, and $\mu_F=F_{xxx}|_c$ as in \eqref{low:6:sai:two-scalars}. Its geometry is supplied by applying the quintic transport theorem to $G=F_x$, with the full sextic outer admissibility retained. As required by the examples in Subsection~\ref{low:6:s36m:section}, all coefficient maps retain $\mu_F$ and use its conditioned pivot; they do not divide by the vanishing second normal derivative.

Consider a $Y$-oriented common ribbon from Proposition~\ref{low:6:sai:common-inflection-geometry}, with $L\ge W\ge M$, $S=LW$, and $K=L/W$. Its seed has $\theta(\gamma')=1$ and $z'=m$, where $\theta=dy-d\,dz$. The bounded-obliqueness cases can be counted in the ordinary ribbon chart; here we suppose $|m|\le c_P$ for a sufficiently small fixed constant. The original ribbon assignment, fixed frames, and count $O_P(\Lambda S)$ remain fixed.

\begin{proposition}[Six external seed jets at small curvature--aspect]\label{low:6:s36i:small-aspect}
Let $\mathcal H\ge1$ be fixed. Suppose $|r|\le\epsilon_r$ on the outer ribbon and $\epsilon_rK\le\mathcal H$.

\noindent\textup{(i)} There is a finite external seed net of size at most $P^C(1+\mathcal H)^CK^6$ preserving the $48$ critical-value widths and the $25$ inflection-slope widths of the grid atlas.

\noindent\textup{(ii)} The anchored full-sextic coefficient-volume bound is
\[
\operatorname{vol}E_j(L,W)\lesssim_P (1+\mathcal H)^C S^{-173}K^{-10}.
\]

\noindent\textup{(iii)} The same conclusion, multiplied by $c^{-C}$, holds when $|r|K\le\mathcal H$ is imposed only on a subset of area at least $cS$ of the already assigned ribbon. 
\end{proposition}
\begin{proof}
The compatibility identity $ZB=-2Bu+rv$, with $u=Yd$ and $v=Zd$, gives $d\alpha=-\mathfrak a\alpha\wedge dz$ for $\alpha=\sqrt{|B|}\theta$ and $\mathfrak a=rv/(2B)$. In the ruled coordinates, $A=\sqrt{|B|}(1+t d_\gamma'-m\Delta)$, where $\Delta=d\circ\Psi-d_\gamma$ is bounded. The width is chosen so that $|t d_\gamma'|\le1/4$; decreasing $c_P$ in the present $|m|\le c_P$ branch gives $|A|\ge P^{-C}$.

Apply Lemma~\ref{low:surface:seed-rounding} with $k_0=6$ and threshold $\mathcal H$. Prepare the coframe, inverse combinations, shear difference, graph and normal scalars, and both full-sextic scalars through every required order, for example $40$. The lemma preserves both the $48$ critical-value widths and the $25$ bounded-scalar widths of $\psi$. The six external coefficients of $n$ are fixed when the actual fields $X_J,Z_J$ vary with $J$. Its subset assertion supplies the stated positive-area version.

Use the inflection atlas of Proposition~\ref{low:6:nc36:moving-atlases}: $48$ critical rows, $25$ rows of $\psi$, nine graph rows, and the value $\mu_F$. Four graph normalizations leave $79$ free slice outputs. At the model, variations of the actual null operators are lower triangular in the stated weight, and the model $\psi$ and $\mu_F$ are constant. The determinant is therefore the nonzero inflection determinant from that proposition, up to bounded coframe factors. The six fixed seed jets have small higher coefficients at the chosen large buffer scale, so compact continuity supplies the actual atlas. This uses the full normalized sextic neighborhood.

The output map depends only on $J$ and the fixed external jet. The fixed-frame $H_0$ disintegration is therefore the one in \eqref{low:6:nc36:null-density}; no separate endpoint shear range is introduced. The raw output product is $S^{-173}K^{-18}$. The seed net costs $K^6$, and the unchanged fixed-frame point comparison costs $K^2$. Thus $(S^{-173}K^{-18})K^6K^2=S^{-173}K^{-10}$.
\end{proof}

\paragraph{The nonnull surface branches}
\begin{lemma}[A lower bound for the first curvature excludes the cubic graph model]\label{low:6:s36i:nonnull-model}
Fix a first-curvature threshold $\rho_*>0$ that is an inverse power of the conditioning parameter. In the critical type-C neighborhood from Proposition~\ref{low:6:s36a:assignment}, choose the complete coefficient defect sufficiently small relative to $\rho_*$.

\noindent\textup{(i)} Any selected critical graph with $|c_{yy}|\ge\rho_*$ is then in a conditioned neighborhood of the quadratic-graph model family. Its nonnull Gauss normalization and both ordinary or intrinsic coefficient atlases are the quadratic-graph ones.

\noindent\textup{(ii)} The same assertion holds for the inflection type.
\end{lemma}
\begin{proof}
The exact type-C model has invariant $q=x-Q(y,z)$, with $Q_{yy}=r$ and $r\gamma=0$, in the notation of \eqref{low:6:s36m:model}. In particular $\gamma=0$ makes the complete $Q$ quadratic, including its possible $z^2$ and affine terms. The selected ordinary normal root has a conditioned nonzero second normal derivative. At that root the finite graph jet is a rational continuous function of the complete phase coefficients and this inverse normal pivot. Consequently an actual phase within coefficient distance $\delta$ of its conditioned model has first graph curvature within $P^C\delta$ of the model's $r$, after the bounded root and tangent normalization. Tangential shearing and subtracting an affine graph function do not change $c_{yy}$.

Choose $P^C\delta<\rho_*/2$. Then $|c_{yy}|\ge\rho_*$ implies $|r|\ge\rho_*/2$ for the model, whence $\gamma=0$. Its graph is quadratic. The nonnull shear $d=-c_{yz}/c_{yy}$ and its inverse coordinate correction have bounds polynomial in $P,\rho_*^{-1}$. Shrinking the original defect again if necessary retains the complete conditioned coefficient neighborhood after this bounded normalization. The available normal profile can contain its retained cubic term; the critical-graph coefficient map already includes that third-normal coefficient. In the inflection type, $|F_{xxx}|\ge P^{-C}$; choosing the model defect after this pivot bound gives a nonzero model coefficient $a_3$. Then $a_3\gamma=0$ forces $\gamma=0$, without a first-curvature exclusion.

These alternatives cover the normalization charts. Indeed choose $\rho_*$ sufficiently small relative to the already fixed upper bound for $D^2c$ and lower bound for $|\det D^2c|$. If $|c_{yy}|<\rho_*$, then $c_{yz}^2-c_{yy}c_{zz}$ is positive with a polynomial lower bound. The stable root in Lemma~\partIref{geo:two-chart-shear} therefore gives the null chart with its mixed pivot bounded away from zero. If the Hessian is definite, this small-first-curvature alternative is impossible. Thus the small-curvature states pass to the null counts, while the present lemma covers the nonnull states.
\end{proof}

\begin{proposition}[The full nonnull counts]\label{low:6:s36i:nonnull-count}

\noindent\textup{(i)} The actual nonnull Gauss branches admit the assignment of Section~\ref{low:6:qsf:sec:gauss}, with count $O_P(\Lambda S)$ and scales $S=LW$, $K=L/W$.

\noindent\textup{(ii)} Their full fixed-frame coefficient-volume bounds are
\[
\begin{aligned}
&O_P(S^{-147}K^{-11})&&\hbox{for the ordinary-critical branch},\\
&O_P(S^{-173}K^{-11})&&\hbox{for the inflection branch}.
\end{aligned}
\]
The constants may include the fixed polynomial cost of the threshold $\rho_*$. 
\end{proposition}
\begin{proof}
Lemma~\ref{low:6:s36i:nonnull-model} places the actual phases in the compact quadratic-graph neighborhoods required by the coefficient atlases. The intrinsic Gauss coordinates are $p=c_y,z$. Their translation covariance and finite-fiber packing, proved in Section~\ref{low:6:qsf:sec:gauss}, apply to the actual graph. On the inflection branch the same geometry is that of $G=F_x$, while Lemma~\ref{low:6:sai:scalar-compatibility} retains the full sextic scalars throughout the outer admissibility. In either case the common contained strip has dimensions $L,W$ and its centerline derivatives satisfy $|\zeta^{(j)}|\lesssim_PL^{1-j}$. The anchor-relative evaluation region has area $O_P(L^2)$, whereas the simultaneous good set has area $\gtrsim_PLW$. The quotient is $O_P(K)$.

At a point in this common set, where the scalar bounds in \eqref{low:6:qsf:critical-widths} and the required centerline bounds hold simultaneously, round each centerline derivative of order $j$ with mesh $c_PW/L^j$. The grid atlas uses six such derivatives, costing $O_P(K^6)$; the weighted atlas uses twelve, costing $O_P(K^{12})$. This transverse-shift rounding preserves both bounded-gradient and bounded-scalar widths by the chain rule described after \eqref{low:6:qsf:critical-widths}. The linear tangent transformation is bounded. Write the transformed quadratic form as $A_0u^2+B_0uv+C_0v^2$ and fix a sufficiently small cutoff $\delta_{\rm q}=P^{-C}$ from its conditioning bounds. Use the weighted atlas if $|A_0|\ge\delta_{\rm q}$. If $|A_0|<\delta_{\rm q}$, then $|B_0|\ge P^{-C}$, and use the grid atlas near $A_0=0$. Choose the coefficient defect after this cutoff. The higher centerline jets are small at the large strip scale; bounded strip scales have local count and are included by enlarging the fixed constant.

The full Gauss incidence density is \eqref{low:6:nc36:gauss-density}. Fixing the intrinsic outputs leaves a bounded fixed-frame $H$ Haar fiber, with no independently multiplied geometric window. The four Gauss products following \eqref{low:6:nc36:raw-products} now apply. Their grid bounds prove the proposition, and their weighted bounds are stronger. 
\end{proof}

\paragraph{Exhaustion of the inflection branch}

The critical value $\phi$ and inflection slope $\psi$ may both vanish and are not divided by. In the family $q^2+\epsilon q^3$ from Subsection~\ref{low:6:s36m:section}, the regular graph retains its second normal pivot as $\epsilon\to0$, and the regular atlas retains its third normal coefficient. Inflection normalization uses a separately conditioned third normal pivot; no inverse bound for it is uniform in $\epsilon$.

Use the common inflection assignment of Proposition~\ref{low:6:sai:common-inflection-geometry}. Its scales are $L\ge W\ge M$, $S=LW$, and $K=L/W$, and its count is $O_P(\Lambda S)$. The fields $\phi=F|_c$ and $\psi=F_x|_c$ are those of \eqref{low:6:sai:two-scalars}, with $F_{xx}|_c=0$. All coefficient maps retain the $48$ nonconstant critical rows, $25$ inflection-slope rows, nine graph rows, and the third normal value. After the four normalizations they have $79$ slice outputs.

Every incidence calculation below is in these $83$ coefficients of $F$. For the inflection graph the normalization constraints are $(F_{xx},F_{xxy},F_{xxz},F_{xxzz})$ in the null chart and $(F_{xx},F_{xxy},F_{xxz},F_{xxyz})$ in the Gauss chart. The four-column calculation in \eqref{low:6:nc36:null-density} has normal pivot $\kappa=\mu_F$: its $H_0$ determinant is $2|\mu_F|^4|B|$, and its three-column $H_3$ determinant is $|\mu_F|^3$. At fixed physical center in the Gauss chart the determinant is $|\mu_F|^4|r|$; at fixed Gauss point it is \eqref{low:6:nc36:gauss-density}. These follow from the same transverse rows with their normal derivative index increased by one; all $79$ unconstrained original coefficient directions retain the identity block. The degree-stable Haar proofs therefore apply on $F$, without treating $F\mapsto F_x$ as a coefficient substitution.

\paragraph{The unscaled oblique charts}

\begin{lemma}[Unscaled inflection oblique count]\label{low:6:s36i:unscaled}
The $Z$-oriented ribbons, and the $Y$-oriented ribbons whose seed slope $n=dz/ds$ satisfies $|n|\ge c_P>0$, have fixed-frame coefficient-volume bound
\[
P^CS^{-173}K^{-4}.
\]
The constants are uniform over the original pieces.
\end{lemma}
\begin{proof}
A $Z$-oriented ribbon is already an ordinary strip with bounded relative slope. In the $Y$-oriented case $|n|\ge c_P$, so parameterizing the seed by $z$ gives $\beta_0=\eta'-d_0=1/n$ conditioned above and away from zero and $|d_0'|\le P^C/W$. Lemma~\ref{low:surface:ordinary-strip} gives a contained ordinary strip of comparable dimensions, including the balanced-cap case. The bounded transverse shear difference from Proposition~\ref{low:5:qng:ribbons} keeps the relative first slope and the ordinary scalar gradients bounded. The strip remains inside the full inflection admissibility.

In both orientations the relative first curve derivative has bounded range, and the derivatives of orders $j=2,\ldots,6$ have ranges $P^CKL^{1-j}$. The first mesh $c_PW/L$ costs $K$, and the other five meshes $c_PW/L^j$ cost $K^2$ each. The complete six-jet net has size $\lesssim_P K(K^2)^5=K^{11}$. The null curvature cutoff is chosen sufficiently small relative to the fixed angle bounds. In endpoint-normalized coordinates the limiting graph after the external first shear has quadratic part
\[
Q=\tfrac r2(v+\beta u)^2+B(v+\beta u)u.
\]
At $r=0$ its grid restriction pivot is the conditioned mixed coefficient $B$. Uniform compact continuity gives the actual inflection grid atlas of Proposition~\ref{low:6:nc36:moving-atlases}, for sufficiently small $r$ and higher external curve jets. The latter jets are small after the large buffer is fixed, since $KL^{1-j}\le P^C/W$ for $j\ge2$. The normalized phase remains in its prescribed inflection neighborhood throughout the contained strip.

Prepare the two scalars separately. The $48$ critical rows have product $L^{-141}W^{-105}$, and the $25$ bounded $\psi$ rows have product $L^{-50}W^{-50}$. The graph and third normal intervals are bounded, so the raw product is $S^{-173}K^{-18}$. The endpoint shear is held fixed in the coefficient map and then integrated over an interval of length at most $P^CK$, and the original $H_3$ integral is bounded by the inflection incidence calculation above and Lemma~\partIref{surf:haar}. The candidate center area divided by the good area is at most $P^CK^2$. These are simultaneous restrictions in the same original coefficient--point integral. Multiplication gives $S^{-173}K^{-18}K^{11}KK^2=S^{-173}K^{-4}$.
\end{proof}

\paragraph{Two rescaled coefficient estimates}

For the remaining $Y$-oriented ribbons, take $|n|<c_P$ and the small null-curvature window. Let $0<h\le c_P$ be dyadic and write
\[
\widehat F(X,Y,z)=h^3F(X/h,Y/h,z),\qquad \widehat L=hL,\qquad \widehat K=hK.
\]
The coefficient and scalar identities are
\[
dF=h^3d\widehat F,\qquad \widehat\psi=h^2\psi,\qquad a=\widehat Y\widehat\phi=h^2Y\phi,\qquad b=\widehat Z\widehat\phi=h^3Z\phi.
\]
The third normal pivot is unchanged. Proposition~\ref{low:6:sai:inflection-rescaling}, using the full-width preparation before amplitude selection, supplies complete compact rescaled phases and contained compact substrips. The following statements include the precise two alternatives needed in the exhaustion.

\begin{proposition}[Rescaled inflection slope estimate]\label{low:6:s36i:slope}
Suppose the prepared certificate has $|n|\asymp_Ph$, $|r|\le\epsilon_0h$, and $hL/W\ge M$, with $\epsilon_0$ sufficiently small.

\noindent\textup{(i)} Its original piece has coefficient-volume bound
\[
P^Ch^{14}S^{-173}K^{-4}.
\]

\noindent\textup{(ii)} If obtaining the certificate used a fraction $q$ of the original ribbon, the bound is multiplied by $q^{-C}$ for a fixed exponent $C$.
\end{proposition}
\begin{proof}
All normalized coefficients of the rescaled phase lie in the fixed compact model neighborhood of Proposition~\ref{low:6:sai:inflection-rescaling}, its seed has $z$ slope bounded above and away from zero, and $|\widehat r|\le\epsilon_0$. Apply the selected-interval part of Lemma~\ref{low:surface:ordinary-strip}, with the full-width preparation and $hL_0/W\ge M$ in the certificate. The rescaled transverse shear difference is $O_P(h)$, so this gives a contained ordinary strip of dimensions $hL_0,W$ with relative slope bounded above and away from zero. The rescaling and finite-jet bounds in that proposition make the ordered null jets of $\widehat Y\widehat d,\widehat Z\widehat d$ small; preparation along its centerline at scale $hL_0$ gives the higher external-jet tolerance for the inflection atlas. No bound for the total endpoint shear is needed.

At a normalized endpoint the limiting graph has $Q=Bu(\beta u+v)$. Use the six purely longitudinal rows of the frozen scalar $\beta_*a+b$ and the forty-two mixed rows of $b$ in \eqref{low:6:nc36:scalar-det}. Keep all twenty-five $\widehat\psi$ outputs, all nine graph coordinates, and the third normal value. In the fixed-graph model, the $48$-row scalar block has the nonzero determinant given in \eqref{low:6:nc36:scalar-det}. The remaining normal restrictions are exactly the inflection atlas of Proposition~\ref{low:6:nc36:moving-atlases}; the retained graph rows remove the lower-normal contribution to phase variation. Changes of the null fields contribute strictly lower-weight rows, as in that proposition. The full map is therefore uniformly invertible on the compact rescaled neighborhood for sufficiently small $\epsilon_0$ and sufficiently accurate first-slope rounding.

The purely longitudinal rows have amplitude $h^2$, the mixed rows amplitude $h^3$, and the bounded scalar $\widehat\psi$ amplitude $h^2$. Their width product is
\[
h^{12+126}(hL)^{-141}W^{-105}\cdot h^{50}(hL)^{-50}W^{-50}
=h^{188}(hL)^{-191}W^{-155}.
\]
Prepare and round the two gradient scalars separately; the chain rule for each finite transverse shift preserves its own amplitude. Six ordinary curve jets cost $(hK)^{11}$, the endpoint shear costs $hK$, and the fixed-frame $H_3$ Haar coordinates scale by $h^2$. The ratio of center areas remains $K^2$: both candidate area and good area acquire the same factor $h$. Including $dF=h^3d\widehat F$ gives
\begin{equation}\label{low:6:nc36:inflection-paid-h}
h^3h^2h^{188}(hL)^{-191}W^{-155}(hK)^{11}(hK)K^2
=h^{14}S^{-173}K^{-4}.
\end{equation}
A longitudinal fraction changes the retained length, derivative constants, finite nets, and area denominator by fixed powers of $q^{-1}$. All these factors are included in the stated $q^{-C}$, while full-width preparation preserves the transverse atlas tolerance.
\end{proof}

\begin{proposition}[Rescaled inflection curvature estimate]\label{low:6:s36i:curvature}
Suppose the prepared certificate has $|r|\asymp_Ph$, $|n|\le P^Ch$, and $hL/W\ge M$.

\noindent\textup{(i)} Its original piece has coefficient-volume bound
\[
P^Ch^9S^{-173}K^{-10}.
\]

\noindent\textup{(ii)} A fraction $q$ again costs a fixed power $q^{-C}$. The fixed frame is allowed to have zero first curvature.
\end{proposition}
\begin{proof}
Proposition~\ref{low:6:sai:inflection-rescaling} gives the fixed compact model neighborhood and $|\widehat r|\asymp_P1$. The long seed has $\widehat Y$ component bounded away from zero. Apply the degree-independent Gauss-strip calculation in the proof of Proposition~\ref{low:4:sf:prop:newcurvaturecount} to obtain a contained prepared strip with fixed format and metric, area, and dimensions comparable to $hL,W$.

At the evaluation point retain the nine graph coordinates determining curvature and the two zeroth critical-gradient outputs before differentiating the outputs in the long and null directions. These directions form a uniformly conditioned coordinate pair. At the limiting model, $y=u$, $z=\beta u+v$, and $Q=A'u^2+B'uv$ with $|B'|\asymp_P1$. The six purely longitudinal rows have principal variation $\partial_u(R|_Q)$ and the forty-two mixed rows have $\partial_v(R|_Q)$; their determinant is the nonzero grid restriction determinant times fixed derivative factors. Take the twenty-five $\widehat\psi$ rows and nine graph rows in the same coordinates and retain the third normal output. All blocks then use the same $Q$ and axes. The four original graph-normalization equations remain as linear combinations of low-order graph rows, through a bounded invertible row change. These rows complete the inflection map. The lower-weight Gauss/null corrections preserve its determinant by Proposition~\ref{low:6:nc36:moving-atlases}. The direction coefficients are determined by the retained graph and zeroth scalar outputs. Apply the chain-rule row subtraction in the proof of Proposition~\ref{low:4:sf:prop:newcurvaturecount} to hold those coefficients fixed in the remaining Jacobian rows. This is a determinant-preserving row operation; the coefficient variables and the lower-weight field corrections are unchanged.

The scalar width product is again $h^{188}(hL)^{-191}W^{-155}$. The Gauss centerline derivatives of orders one through six have ranges $P^C(hL)^{1-j}$ and meshes $c_PW/(hL)^j$, so the net costs $(hK)^6$. This output map depends only on the normalized endpoint phase and the fixed seed jet. The physical-center Gauss incidence with pivot $|\mu_F|^4|\widehat r|$ described above and the original inflection null-anchor $H_0$ pivot therefore apply. The latter pivot uses the nonzero mixed curvature and does not require nonzero first curvature at the anchor.

The endpoint-range and $H_0$ scaling calculation in Proposition~\ref{low:5:qsc:curvature} applies with the same $h,L,W,K$ and $hK\ge M$. It gives the candidate-area quotient $K^2$ and Haar factor $h^3$. As there, introduce $\widehat\eta$ only on the right side of the full coefficient--point identity; no fixed-phase inverse point map is used. Together with the sextic coefficient rescaling $dF=h^3d\widehat F$, the product is
\begin{equation}\label{low:6:nc36:inflection-intrinsic-h}
h^3h^3h^{188}(hL)^{-191}W^{-155}(hK)^6K^2
=h^9S^{-173}K^{-10}.
\end{equation}
The same finite preparations and area comparisons give the stated fixed power of $q^{-1}$ for a fractional certificate.
\end{proof}

\paragraph{Exhaustion and the uniform amplitude sum}

\begin{theorem}[The complete inflection null count]\label{low:6:s36i:null-count}

\noindent\textup{(i)} Every original inflection null piece has count at most $P^C\Lambda S$. The count includes zero first curvature, zero seed slope, and zero critical gradients, and preserves the original assignment and anchors.

\noindent\textup{(ii)} Its coefficient-volume bound at the fixed frame is
\[
\operatorname{vol}_{83}E_j(S,K)\le P^CS^{-173}K^{-3}.
\]
\end{theorem}
\begin{proof}
Lemma~\ref{low:6:s36i:unscaled} covers the $Z$-oriented arcs and the $Y$-oriented arcs with $|n|\ge c_P$. Choose the remaining null-curvature cutoff below the angle and reconstruction tolerances. If $|r|K\le H_0$ on a fixed positive area fraction, Proposition~\ref{low:6:s36i:small-aspect} gives $P^CS^{-173}K^{-10}$, including identically zero curvature.

On the remaining positive fraction, apply Lemma~\ref{low:surface:amplitude-selection}, including every full-sextic scalar and graph jet required by the $83$-output maps. Its cell alternatives retain signs and both root-sheet labels and give $q\ge c_{P,\epsilon}(hK)^{-\epsilon}$. Choose $\epsilon$ so that the preparation exponent satisfies $C\epsilon<1/2$, then $H_0$ large enough for the reconstruction thresholds. Proposition~\ref{low:6:sai:inflection-rescaling} reconstructs complete compact phases on a fixed good fraction. The last part of the selection lemma gives a contained compact substrip in the original algebraic parameters, preserving the original transverse width and the assignment.

Let $C_*$ dominate the preparation and counting losses. Cover the anchored coefficient set by its possible dyadic certificates. Apply the summation calculation in Theorem~\ref{low:5:qsc:null-exhaustion} to the bounds in Propositions~\ref{low:6:s36i:slope} and~\ref{low:6:s36i:curvature}, replacing their powers of $(h,K)$ by $(14,-4)$ and $(9,-10)$ and their area factor by $S^{-173}$. The inequality $q^{-C_*}\le P^C(hK)^{C_*\epsilon}$ gives convergent dyadic sums in $h$ and the bounds $P^CS^{-173}K^{-4+C_*\epsilon}$ and $P^CS^{-173}K^{-10+C_*\epsilon}$. Choose also $C_*\epsilon<1$. Since $K\ge1$, both are bounded by $P^CS^{-173}K^{-3}$. The choices remain in order: compact and angle tolerances, $\epsilon$, $H_0$, the common buffer, and the inner coefficient defect. The physical weighted mass bound is unchanged. Together with the unscaled and small-aspect branches, this exhausts the original inflection null pieces.
\end{proof}


\subsection{Critical graphs: transport and assignment}\label{low:degree6:ordinary}\label{low:6:ordinary-connection}\label{low:6:s36cg:section}\label{low:6:ordinary-cover}\label{low:6:as36:relocation-section}\label{low:6:r3:band-summation}\label{low:6:og:cover}

Use the full sextic coefficient space, admissibility and physical family of Sections~\ref{reg:five-six} and~\ref{low:6:inflection-geometry}. On the selected critical graph, $F_x=0$ and $|F_{xx}|\ge P^{-C}$; Section~\ref{low:6:inflection-counts} treats the inflection branch. The connection estimates below apply on positive rate bands. Corollary~\ref{low:6:r1:zero-common} supplies the separate zero-rate input. The cover of Theorem~\ref{low:6:og:common-cover} gives the certificates counted in Section~\ref{low:6:ordinary-counts}, and Theorem~\ref{low:6:as36:all-surface} combines their critical null bounds.

\subsubsection{Connection and distortion in one rate band}

We prove the connection identities and bandwise transverse distortion for the cubic-invariant branch. The buffered-corridor construction below gives the assignment; its coefficient integrals remain in $\mathcal V_6$ at the fixed frames.

\paragraph{Step 1. The original affine flows}

Put $q_\gamma=x-yz+\gamma z^3/3$ and $U_\gamma=y-\gamma z^2/2$. The fields $Y_0=z\partial_x+\partial_y$ and $Z_0=y\partial_x+\gamma z\partial_y+\partial_z$ satisfy $[Y_0,Z_0]=0$, $Y_0q_\gamma=Z_0q_\gamma=0$, $Y_0U_\gamma=1$, and $Z_0U_\gamma=0$. Hence the phase $q_\gamma^2$ has critical value zero and null shear $d=\gamma z$, with $Z d=\gamma$. 

The coordinates $(q,u,z)=(q_\gamma,U_\gamma,z)$ have inverse
\[
x=q+uz+\frac{\gamma}{6}z^3,\qquad y=u+\frac{\gamma}{2}z^2,
\]
and Jacobian one. In these coordinates the two original affine flows are translations in $u,z$. Explicitly, their joint action is
\[
g_{s,t}(x,y,z)=\left(x+yt+\frac{\gamma}{2}zt^2+\frac{\gamma}{6}t^3+s(z+t),\ y+\gamma zt+\frac{\gamma}{2}t^2+s,\ z+t\right).
\]
For each fixed $(s,t)$ this is an affine map with determinant one. Thus its complete action on $\mathcal V_6$ has coefficient Jacobian one. The polynomial coordinate substitution is useful for computing this fixed $83$-dimensional representation, but does not create an unrestricted space of polynomials of degree eighteen. Its image consists of the span of the $83$ explicitly substituted original monomials.

\paragraph{An exact connection identity for cubic implicit surfaces}

Let $G_0$ be a real polynomial of degree at most three, and let $G_0(c(y,z),y,z)=0$ be a regular graph with $\rho=(G_0)_x\ne0$. Let $d$ be a regular null root, and set
\[
Y=\partial_y,\quad Z=\partial_z+d\partial_y,\quad r=c_{yy},\quad B=c_{yz}+dr,\quad v=Zd,\quad u=Yd.
\]
Assume $Bv\ne0$ on the branch. Then
\begin{equation}\label{low:6:s36cg:cubic-connection}
u=\frac43 Z\log|\rho|+\frac13Z\log|B|+\frac13Z\log|v|,
\qquad
\frac{rv}{B}=\frac83Z\log|\rho|+\frac53Z\log|B|+\frac23Z\log|v|.
\end{equation}

To prove the identity, lift $Z$ and $Y$ to tangent vectors in the ambient space:
\[
X=(c_z+dc_y,d,1),\qquad U=(c_y,1,0).
\]
Along a null trajectory, the null equation gives $X'=vU$ and $U'=Be_x$. Write $H=D^2G_0$ and $T=D^3G_0[X,X,X]$. Differentiating $G_0=0$ on the graph gives $H[U,X]=-\rho B$ and $H[X,X]=0$. Three derivatives along the trajectory give $T=2\rho Bv$. On the other hand $YX=Be_x+uU$, so differentiation of $H[X,X]=0$ by $Y$ gives
\[
0=D^3G_0[U,X,X]+2B Z\rho-2u\rho B.
\]
The cubic tensor $D^3G_0$ is constant, hence $ZT=3vD^3G_0[U,X,X]$. Eliminate the mixed third derivative and use $T=2\rho Bv$ to obtain the first formula. The exact graph compatibility identity $ZB=-2Bu+rv$ then gives the second formula.

Suppose $\rho,B$ and their inverses are bounded and $\nu/2\le|v|\le\nu$ on a connected characteristic segment. Its transverse flow Jacobian $J(t,t_0)$ satisfies $J(t_0,t_0)=1$ and $\partial_t\log J=u$. Integrating the first formula gives
\[
J(t,t_0)=\left|\frac{\rho(t)}{\rho(t_0)}\right|^{4/3}\left|\frac{B(t)}{B(t_0)}\right|^{1/3}\left|\frac{v(t)}{v(t_0)}\right|^{1/3}.
\]
The ratio bounds give two-sided uniform bounds for $J$ on this segment. This conclusion permits shear drift of order $\nu L$. It concerns a single $v$ band and does not imply a uniform Jacobian when $v$ traverses arbitrarily many dyadic bands.

\paragraph{Step 2. The sextic square-root identity}

For an arbitrary sextic $F$ with regular normal critical graph and $\lambda=F_{xx}\ne0$, write locally
\[
F(x,y,z)=\phi(y,z)+\sigma Q(x,y,z)^2,
\quad \sigma=\operatorname{sgn}\lambda,
\quad Q(c(y,z),y,z)=0,
\quad \rho=Q_x>0.
\]
This factorization is analytic and algebraic on a selected branch: Taylor's integral formula gives $F-\phi=(x-c)^2a(x,y,z)$ with $\sigma a>0$, and one takes $Q=(x-c)\sqrt{\sigma a}$. In particular $\lambda=2\sigma\rho^2$. The preceding differential calculation applies to any regular defining function $Q$, except that $Z(D^3Q[X^3])$ also contains $D^4Q[X^4]$. It gives
\[
u=\frac43Z\log|\rho|+\frac13Z\log|B|+\frac13Z\log|v|-\frac{D^4Q[X^4]}{6\rho Bv}.
\]

At the base point, derivatives in this formula are taken along the fixed ambient vector $X$. Because $Q=Q_X=Q_{XX}=0$ there and $\deg F\le6$, seven derivatives of $F=\phi+\sigma Q^2$ give
\[
0=70\sigma Q_{XXX}Q_{XXXX}+\phi_{\mathrm{fr},7},
\qquad Q_{XXX}=2\rho Bv,
\]
where $\phi_{\mathrm{fr},7}$ is the seventh derivative of $\phi$ in the frozen tangential direction $(d,1)$. Consequently
\begin{equation}\label{low:6:s36cg:sextic-connection}
\begin{aligned}
u&=\frac23Z\log|\lambda|+\frac13Z\log|B|+\frac13Z\log|v|+\frac{\phi_{\mathrm{fr},7}}{420\lambda B^2v^2},\\
\frac{rv}{B}&=\frac43Z\log|\lambda|+\frac53Z\log|B|+\frac23Z\log|v|+\frac{\phi_{\mathrm{fr},7}}{210\lambda B^2v^2}.
\end{aligned}
\end{equation}

The connection cannot be expressed solely in critical-value derivatives. Set $q=q_\gamma$, $U=U_\gamma$, and
\[
G_0=q+\varepsilon\left(zq-\frac{4}{3\gamma}U^2\right)
=q+\varepsilon\left(zx+\frac13yz^2-\frac{4}{3\gamma}y^2\right).
\]
The quartic terms cancel, so $G_0$ is cubic and $F=G_0^2$ is sextic. Its critical value is identically zero. Along $q=U=0$, one has
\[
v=\gamma,\quad B=1,\quad (G_0)_x=1+\varepsilon z,\quad r=\frac{8\varepsilon}{3\gamma(1+\varepsilon z)},\quad u=\frac{4\varepsilon}{3(1+\varepsilon z)}.
\]
Thus $rv\ne0$ while every critical-value derivative vanishes. The connection is controlled by the normal factor $(G_0)_x$, in agreement with the identity.

\paragraph{Step 3. The frozen seventh derivative}\label{low:6:s36cg:frozen-seven}

Along a null trajectory let $D=d/dt$, $P_1=\phi_y$, $R_2=\phi_{yy}$, and $T_3=\phi_{yyy}$. At any point the frozen derivative of Definition~\ref{def:frozen-null} satisfies
\begin{align*}
\phi_{\mathrm{fr},7}={}&D^7\phi-21vD^5P_1-35(Dv)D^4P_1-35(D^2v)D^3P_1\\
&-21(D^3v)D^2P_1-7(D^4v)DP_1-(D^5v)P_1\\
&+105v^2D^3R_2+210v(Dv)D^2R_2+\bigl(105vD^2v+70(Dv)^2\bigr)DR_2\\
&+\bigl(21vD^3v+35(Dv)D^2v\bigr)R_2\\
&-105v^3DT_3-105v^2(Dv)T_3.
\end{align*}
Indeed, write the difference between the actual trajectory and its tangent line as $\delta(s)=\sum_{j\ge2}D^{j-2}v\,s^j/j!$ and expand $\phi(\text{tangent line})=\phi(\text{trajectory})-\delta P_1+\delta^2R_2/2-\delta^3T_3/6+O(s^8)$. This proves the displayed identity without replacing a frozen derivative by an iterated vector-field derivative.

\paragraph{Step 4. Polynomial transverse distortion on a single \texorpdfstring{$v$}{v} band}\label{low:6:s36cg:polynomial-distortion}

Assume that the actual normalized sextics remain in a sufficiently small neighborhood of the compact regular-root rank-four family. Every exact null-normalized model in that family has $rv=0$ and $Y^3\phi=0$: the quadratic-invariant branch has $v=0$, and the cubic-invariant branch has $r=0$ and critical value affine in $U_\gamma,z$. Therefore, after choosing the neighborhood sufficiently small, one may require
\[
\min(|r|,|v|)\le\varepsilon,\qquad |Y^3\phi|\le\varepsilon,
\]
where $\varepsilon=P^{-M}$ will be chosen after the fixed exponent below. The other scalars, inverse pivots, and $v$ are bounded by $P^{C_0}$.

Let $I$ be any connected trajectory interval contained in $\nu/2<|v|<\nu$. The transverse flow Jacobian between any two points of $I$ is bounded above and below by powers of $P$, uniformly in $\nu$ and the length of $I$.

To prove this assertion, split $I$ into a bounded number of intervals on which the finite derivatives needed in the seventh-derivative expansion are monotone. Such a subdivision follows from Lemma~\partIref{geo:characteristic-preparation}, applied to the finitely many higher derivatives; all functions are algebraic scalars evaluated along the same null trajectory. On each resulting interval, boundedness and the divided-difference formula \cite[equation (44)]{deBoor} give pointwise estimates $|D^jf(t)|\lesssim M_f\operatorname{dist}(t,\partial I)^{-j}$ for the required derivatives of $f=P_1,R_2,Z\phi,v$, with amplitudes $P^{C_0}$ or $\nu$ as appropriate. Write $\operatorname{TV}(f)$ for the total variation of $f$ on the trajectory interval under consideration. The total variation of $T_3=Y^3\phi$ is $O(\varepsilon)$, and the total variation of $v$ is $O(\nu)$, with a fixed-format constant.

Choose a sufficiently large fixed exponent $A$ and put $T=P^A\nu^{-1/2}$. Remove end intervals of length $T$ from every member of this finite subdivision; an interval shorter than $2T$ is removed entirely. On a removed interval, the signed compatibility formula gives
\[
\int u\,dt=-\frac12\Delta\log|B|+\frac12\int\frac{rv}{B}\,dt.
\]
The second term has absolute value at most $P^{A+C}\sqrt{\varepsilon}$ if $\nu\le2\varepsilon$, since $r$ remains $P$-bounded. If $\nu>2\varepsilon$, the neighborhood condition forces $|r|\le\varepsilon$ throughout the band, and that term is at most $P^{A+C}\varepsilon$. Here and below the upper bound $\nu\le P^{C_0}$ is included in $C$. Choose $M$ sufficiently large after $A$ to make the sum of all these errors bounded by an absolute constant. The logarithmic factor is a power of $P$.

On the remaining cores integrate the seventh-derivative expansion, dividing by $\lambda B^2v^2$. Every term except the last two is bounded in absolute integral by
\[
P^C\left(\nu^{-2}T^{-5}+\nu^{-1}T^{-4}+T^{-2}\right)
=P^C\left(P^{-5A}\nu^{1/2}+P^{-4A}\nu+P^{-2A}\nu\right).
\]
The two last terms are bounded by $P^C\nu\operatorname{TV}(T_3)+P^C\sup|T_3|\operatorname{TV}(v)\le P^C\varepsilon\nu$. Thus first choosing $A$ sufficiently large and then $M$ sufficiently large makes the total core defect integral bounded by an absolute constant. No factor proportional to the core length occurs.

On one retained core, put $\mathcal D=\phi_{\mathrm{fr},7}/(420\lambda B^2v^2)$. The transverse flow Jacobian satisfies $J(t_0,t_0)=1$ and $\partial_t\log J=u$, so \eqref{low:6:s36cg:sextic-connection} gives
\[
J(t,t_0)=\left|\frac{\lambda(t)}{\lambda(t_0)}\right|^{2/3}\left|\frac{B(t)}{B(t_0)}\right|^{1/3}\left|\frac{v(t)}{v(t_0)}\right|^{1/3}\exp\left(\int_{t_0}^{t}\mathcal D(\tau)\,d\tau\right).
\]
The ratios of $\lambda,B$ are bounded by powers of $P$, the ratio of $|v|$ is at most two, and $\int|\mathcal D|\lesssim1$. Hence $P^{-C}\le J(t,t_0)\le P^C$. Combine the bounded number of core and trimmed intervals. Applying the same proof to an arbitrary subinterval of $I$ proves the two-sided Jacobian bound between arbitrary points. The argument applies to a single $v$ band.

\paragraph{Step 5. Passage to the geometric assignment}

The preceding distortion estimate applies on a single buffered $v$ band. The buffered-core relocation and Theorem~\ref{low:6:og:common-cover} supply the subsequent assignment without treating an artificial band boundary as a physical boundary. The thin and high-residual cases use the fixed-width count of Theorem~\ref{low:6:sp36:tail}; the moderate cases use the ordinary coefficient count, with its horizontal and wide orientations treated separately. Those counts use the coefficient atlas in Taylor coordinates and keep all seed, endpoint, and fixed-frame incidence factors. 

\subsubsection{Interior representatives and the rate output}

We retain the original indices, coefficient anchors, and weights. On a nondegenerate critical graph write $F_x(c,y,z)=0$, $\phi=F|_c$, $\lambda=F_{xx}|_c$, $Y=\partial_y$, $Z=\partial_z+d\partial_y$, $r=c_{yy}$, $B=c_{yz}+dr$, $u=Yd$, $v=Zd$, and $\theta=dy-d\,dz$. The radius of a prescribed model box is chosen before the inner coefficient defect.

The thin-band example $F=(x-yz)^2+\epsilon xyz^3$ has critical graph $c=yz-\epsilon yz^3/2$, $r=0$, $B=1-3\epsilon z^2/2$, $d=3\epsilon yz/(2B)$, and $v(y,0)=3\epsilon y/2$. Consequently the pointwise band $\nu/2<v<\nu$ at $z=0$ has width $\nu/(3\epsilon)$, which tends to zero with $\nu$. We therefore relocate the geometric representative before selecting its band, keeping the original coefficient frame.

For $q_\gamma=x-yz+\gamma z^3/3$ and $U_\gamma=y-\gamma z^2/2$, the phase $q_\gamma^2+\epsilon(U_\gamma-hz)^2$ has $r=0$, $B=1$, $d=\gamma z$, $v=\gamma$, null coframe $(dU_\gamma,dz)$ and critical value $\epsilon(U_\gamma-hz)^2$. For $h=1$, the seed $U=z=s$ and its null extension $z=s+t$ have coframe $ds,ds+dt$ and area comparable to seed length times width; endpoint shear can be $\gamma L$ and horizontal displacement $O(L+\gamma L^2)$. For general $h$, the seed residual is $h$. The cover below therefore uses both corridor orientations and intrinsic seed length. The full coefficient maps retain the third-normal output, and $q^3$ uses its separate inflection normalization.

\paragraph{Relative norming}

\begin{lemma}[Relative norming on a complex box]\label{low:6:r3:norming}
Write $\mathbb D(0,r)=\{z\in\mathbb C:|z|<r\}$. Let $f$ be holomorphic on the complex polydisc $\mathbb D(0,8R)^n$ and satisfy a nonzero polynomial equation $A(z,f(z))=0$ of total degree at most $D$. Here $R>0$, $n$ and $D$ are fixed, and the coefficients of $A$ are arbitrary.

\noindent\textup{(i)} For every fixed integer $N\ge0$ there is $C=C(n,D,N)$ such that
\[
\sup_{z\in\mathbb D(0,2R)^n}|\partial^\alpha f(z)|\le C R^{-|\alpha|}\sup_{t\in[-R,R]^n}|f(t)|\qquad(|\alpha|\le N).
\]

\noindent\textup{(ii)} There is also a constant $c=c(n,D)>0$ such that
\[
\sup_{t\in[-R/2,R/2]^n}|f(t)|\ge c\sup_{t\in[-R,R]^n}|f(t)|.
\]
No constant depends on the amplitude of $f$, and the assertion includes families for which $f$ tends identically to zero.

\noindent\textup{(iii)} Put $\nu_* = \sup_{[-R,R]^n}|f|$. Either $f$ vanishes identically, or there is $t_*\in[-R/2,R/2]^n$ and a real ball of radius $c_{n,D}R$ about $t_*$ on which
\[
c_{n,D}\nu_*\le |f(t)|\le C_{n,D}\nu_*.
\]
One can choose $t_*$ from a fixed finite grid in $[-R/2,R/2]^n$, after decreasing the constants. The number of grid points depends only on $n,D$.
\end{lemma}

\begin{proof}
In one variable, we claim that
\[
\sup_{|z|\le3R}|f(z)|\lesssim_D \sup_{t\in[-R/2,R/2]}|f(t)|.
\]
For a nonconstant branch, choose an irreducible relation whose coefficients, viewed as polynomials in $z$, have degree at most $D$. The structural Bernstein inequality \cite[Corollary~3.3.2 and Proposition~2.1.5]{RoytwarfYomdin}, applied inside the disk of radius $4R$, gives the claim. Constant branches are immediate.

For several variables, apply the one-variable estimate successively while fixing all other coordinates. Generic slices satisfy a nonzero polynomial equation of degree at most $D$; exceptional slices follow by continuity. This successively bounds the complex supremum in each coordinate by the supremum on its real interval, producing a constant depending only on $n,D$. The same iteration with the smaller real intervals proves the second assertion. Finally apply the multivariable Cauchy estimate to obtain every prescribed derivative.

The second conclusion of the lemma gives a point in the smaller box where $|f|\ge c\nu_*$. Its first derivative bound is $C\nu_*/R$ throughout a fixed enlargement of that box. Therefore the value remains at least $c\nu_*/2$ on a ball of radius $cR/(2C)$. Moving its center to the nearest point of a fixed grid of sufficiently small spacing preserves the same assertion with smaller constants. Its upper bound follows from the relative complex norming estimate. 
\end{proof}

\paragraph{A uniform complex model box}

Use the rank-four model $F_0$ and polynomial affine flow $h_0(s,t)$ from the proof of Lemma~\partIref{geo:model-box} in degree six with a simple critical root. Writing the tangential coefficients as $\ell,m$, the model is $F_0=f(q)+\ell(y-\gamma z^2/2)+mz$, with $f'(0)=0$ and $|f''(0)|,|2B|\ge P^{-C}$. Here $f$ has degree at most three, and its cubic coefficient vanishes when $\gamma\ne0$. The model orbit has $y=s+\gamma t^2/2$, $z=t$, and null coframe $ds,dt$.

\begin{lemma}[Complex continuation on a prescribed model box]\label{low:6:r1:complex-box}
Fix a polynomially large $M=P^A$, with $A$ selected from the preceding conditioning requirements.

\noindent\textup{(i)} There is a sufficiently small inner defect $\delta=P^{-C_A}$ such that, if the complete $83$-coefficient vector of $F$ is within $\delta$ of the model $F_0$, the conclusions of Lemma~\partIref{geo:model-box} hold in degree six and normal order two. In particular, the actual nondegenerate critical graph and selected null root are holomorphic on the complex $(s,t)$ polydisc of radius $20M$, with graph and null-frame corrections and any prescribed finite number of their derivatives bounded by
\[
P^C(1+M)^C\delta.
\]

\noindent\textup{(ii)} On the real polydisc of radius $10M$, the normalized frames satisfy the outer bounds on all nonconstant coefficients of $F$, the normal and null pivots, and the physical data in Definition~\partIref{def:admissibility}. Their relative affine maps and inverses have polynomial bounds in $P,M$. The matrix expressing the null coframe $(\theta,dz)$ in $(ds,dt)$ and its inverse are uniformly bounded.
\end{lemma}
\begin{proof}
Apply Lemma~\partIref{geo:model-box} with $k=6$ and normal order two. Its pivots are $f''(0),2B$, and its coframe formula gives an arbitrarily small perturbation of $ds,dt$ after the defect is chosen. The outer full-sextic coefficient and physical windows are fixed first; the common lemma supplies the holomorphic roots, derivative bounds and relative affine comparisons on the stated boxes.
\end{proof}

The scalar $v=Zd$ is algebraic of uniformly bounded degree on this selected branch. Its expression in the model variables $s,t$ is still algebraic of uniformly bounded degree, because the map from these variables to $(y,z)$ is polynomial of degree at most two. The preceding lemma supplies its single holomorphic branch on the entire complex model box. 

\begin{theorem}[Buffered relocation before amplitude assignment]\label{low:6:r1:relocation}

\noindent\textup{(i)} For each original selected inner label in the critical sextic null branch, either $v$ is identically zero on its connected analytic branch, or one can select a representative $\widehat g$ with the following properties. If the original frame is $g_j$, then $\widehat g=g_j\circ k$, and the affine matrices of $k,k^{-1}$ have norms at most $P^C$. For some dyadic $\nu>0$, the representative has a neighborhood of radius $cM$ in the coframe metric of Definition~\partIref{def:null-coordinates} contained in the enlarged outer admissibility and in the band
\[
c_0\nu\le |v|\le C_0\nu.
\]
Here $c_0,C_0,c>0$ are fixed conditioning constants, independent of $\nu$.

\noindent\textup{(ii)} The selection is Borel measurable in the phase coefficients, has fixed algebraic format after $\nu$ is specified, and is chosen for all original labels before one is fixed for coefficient integration. The fixed coefficient frame and weight are retained.

\noindent\textup{(iii)} Local packing of the relocated labels costs only a polynomial factor: if original relative unit neighborhoods contain at most $\Lambda$ labels, then relocated relative unit neighborhoods contain at most $P^C\Lambda$ labels, counting multiplicity.
\end{theorem}
\begin{proof}
First choose a conditioned rank-four reference model for the original inner normalized label. The compact model parameter set and coefficient-distance test form a fixed-format semialgebraic relation. A fixed lexicographic selection after compact minimization supplies a measurable reference; no differentiable model retraction is required. Fix $M$ before choosing the inner defect in Lemma~\ref{low:6:r1:complex-box}.

Put $\nu_0=\sup_{[-M,M]^2}|v|$. Apply Lemma~\ref{low:6:r3:norming} with $R=M$ to the holomorphic branch on the radius-$20M$ box supplied by Lemma~\ref{low:6:r1:complex-box}. On the radius-$2M$ complex polydisc it gives $|v|\le C\nu_0$ and $|\partial_sv|+|\partial_tv|\le C\nu_0/M$, with $C$ depending only on the fixed algebraic degree. If $\nu_0=0$, the identity theorem gives vanishing throughout the connected analytic branch. Otherwise select a real maximizer $(s_*,t_*)$, with a fixed tie rule. Supremum and maximizer selection are semialgebraic because the branch, square, and coefficient parameters have fixed format. On the real square of radius $cM$ about this maximizer, the relative derivative bound gives
\[
\tfrac12\nu_0\le |v|\le\tfrac32\nu_0
\]
when $c$ is sufficiently small. This square lies in the real outer model box. Its coframe and inverse are bounded by Lemma~\ref{low:6:r1:complex-box}, so it contains an actual intrinsic ball of radius $c'M$ about its center. Normalize the actual frame there and denote it by $\widehat g$. Choose the dyadic $\nu$ with $\nu\le\nu_0<2\nu$; the displayed neighborhood then lies in $\nu/2\le|v|\le3\nu$, and fixed strict outer factors such as $\nu/4<|v|<4\nu$ give the band.

The reference model flow to $(s_*,t_*)$ is polynomially bounded in $P,M$, and the actual normal/null correction and inverse are bounded. Therefore $\widehat g=gk$ with $k,k^{-1}$ bounded by a fixed power of $P$. This is a relocation of the label's geometric representative. The coefficient integral, normalization determinant, and assigned weight still refer to the original fixed family frame $g$; no coefficient Jacobian for the selected model or maximizer is inserted.

For the packing assertion, if a relocated representative lies in $h\mathcal B$, its original label lies in $h\mathcal B\mathcal K_P^{-1}$, where $\mathcal B$ is the fixed relative unit neighborhood and $\mathcal K_P$ is a polynomially bounded subset of the six-dimensional upper-affine group. That set is covered by $P^C$ relative unit neighborhoods: the group multiplication and inverse are polynomials of fixed degree, so an ordinary polynomially fine parameter mesh suffices. Each such neighborhood contains at most $\Lambda$ original labels. The desired bound follows even when many original labels choose the same representative. All selections precede the later common-piece assignment and anchoring.
\end{proof}

\begin{definition}[Relocated representatives and bands]\label{def:relocated-band}
For an original frame $g_j$, its relocated representative is $\widehat g=g_jk(F,j)$ from Theorem~\ref{low:6:r1:relocation}, with $k,k^{-1}$ polynomially bounded. The original coefficient normalization, determinant, Haar restriction and coefficient-independent weight remain attached to $g_j$. On the nonzero branch the dyadic label $\nu$ is comparable to the relocated local supremum; it need not be comparable to $|v(g_j)|$, which may vanish. A banded outer domain for its certificate retains $c\nu\le|v|\le C\nu$ on the endpoint region used in the coefficient count.
\end{definition}
The relocation supplies a buffered representative in one actual rate band. The assignment in Theorem~\ref{low:6:og:common-cover} is made after this relocation and retains all original labels and weights. Its transverse Jacobian input is the single-band transport estimate in Section~\ref{low:6:s36cg:polynomial-distortion}; no product over successive rate bands is taken.

\paragraph{The coordinate \texorpdfstring{$v$}{v}}

Let $x=c(y,z)$ be the actual nondegenerate critical graph, and write $Y=\partial_y$, $Z=\partial_z+d\partial_y$, $B=c_{yz}+dc_{yy}$, and $v=Zd$. The selected null equation is
\[
c_{zz}+2dc_{yz}+d^2c_{yy}=0.
\]
At an actual endpoint normalized so that $d=c_y=c_z=0$, differentiate this identity by $Z$ to obtain the formula
\begin{equation}\label{low:6:r3:v-graph-row}
c_{zzz}+2Bv=0,\qquad B=c_{yz}.
\end{equation}
The fourth graph normalization gives $c_{zz}=0$ there; the row $c_{zzz}$ remains free.

\begin{lemma}[A band interval in the coefficient map]\label{low:6:r3:v-output}
Suppose a complete regular-root coefficient atlas retains the ordinary third graph jet, including the rows $B=c_{yz}$ and $c_{zzz}$, and its conditioning gives $P^{-C}\le |B|\le P^C$.

\noindent\textup{(i)} Replacing the derivative output $c_{zzz}$ by $v$ preserves the finite-fiber bound and changes the inverse output density by the factor $2|B|$.

\noindent\textup{(ii)} On the coefficient set $|v|\asymp\nu$, the new output interval has length $O(\nu)$. 
\end{lemma}
\begin{proof}
Hold all the other output coordinates fixed. Formula~\eqref{low:6:r3:v-graph-row} is the invertible triangular change $c_{zzz}=-2Bv$. Its inverse row Jacobian is $2|B|$, since $B$ is among the retained other coordinates. If $\mathcal T_{\mathrm{old}}$ records $c_{zzz}$ and $\mathcal T_{\mathrm{new}}$ records $v$ in the same row, then $\det D\mathcal T_{\mathrm{new}}=-(2B)^{-1}\det D\mathcal T_{\mathrm{old}}$. Their fibers have exactly the same cardinalities. The coefficient set $c_0\nu\le |v|\le C_0\nu$ is the union of two intervals of total length $O(\nu)$. Applying the full output map and integrating this row gives the asserted gain. The same argument applies after deleting the four graph-normalization rows, because $c_{zzz}$ was not one of them. All four coefficient directions represented by the transverse orbit columns remain present in the complete coefficient--point density.
\end{proof}

Both the strong $48+16+18+1$ ordinary critical-graph coefficient map and the weaker $48+25+9+1$ atlas retain the complete ordinary third graph jet. In a Taylor-coefficient convention the output is $[z^3]c=c_{zzz}/6$, so its replacement by $v$ has inverse row Jacobian $|B|/3$, rather than $2|B|$. This fixed factorial has no effect on the volume exponent. The ordinary graph rows must be retained even when the critical-value or normal-scalar outputs use intrinsic differentiation.

\begin{lemma}[An externally conditioned curved seed]\label{low:6:r3:v-seed}
Suppose the output graph is expressed in externally fixed seed coordinates
\[
y=t+\zeta(s),\qquad z=s,\qquad C(s,t)=c(t+\zeta(s),s),
\]
with $\zeta(0)=0$, and put $\beta=\zeta'(0)$ and $a=\zeta''(0)$.

\noindent\textup{(i)} At an actual endpoint normalized by $d=c_y=c_z=0$, one has
\begin{align}
B&=C_{st}-\beta C_{tt},\label{low:6:r3:seed-B}\\
 2Bv&=3aB-\bigl(C_{sss}-3\beta C_{sst}+3\beta^2C_{stt}-\beta^3C_{ttt}\bigr).\label{low:6:r3:seed-v}
\end{align}

\noindent\textup{(ii)} Let $\mathcal T_{\mathrm{old}}$ record $C_{sss}$ and let $\mathcal T_{\mathrm{new}}$ replace that row by $v$. With the external seed and all other graph outputs fixed, $\det D\mathcal T_{\mathrm{new}}=-(2B)^{-1}\det D\mathcal T_{\mathrm{old}}$, so the inverse derivative-row Jacobian is $2|B|$. Its Taylor-coefficient version has inverse row Jacobian $|B|/3$.
\end{lemma}
\begin{proof}
The first identity is the second-order chain rule. Put $D=\partial_s-\beta\partial_t$, with $\beta$ evaluated at the endpoint and held constant during these spatial derivatives. The third-order chain rule gives
\[
D^3C=c_{zzz}+3\zeta''(0)c_{yz}+\zeta'''(0)c_y.
\]
The last term vanishes by endpoint normalization. Substituting~\eqref{low:6:r3:v-graph-row} proves~\eqref{low:6:r3:seed-v}. With the seed and other graph outputs fixed, the coefficient of $C_{sss}$ is $-1$ and $B$ is fixed, so Lemma~\ref{low:6:r3:v-output} gives the same triangular Jacobian factors. The seed shifts the old graph-row interval by terms depending on other outputs, without removing the factor $\nu$ from its length.
\end{proof}

Under an additional affine tangential change held fixed during coefficient differentiation, with matrix and inverse bounded by $P^C$, the retained complete third graph jet permits a finite choice of a pure third derivative with a polynomially conditioned nonzero coefficient in the frozen null-direction cube. A fixed nonlinear seed contributes only lower graph derivatives to this third-order identity. During the output change the seed parameters and any external endpoint shear are fixed. If they vary in the coefficient count, they are integrated separately over their stated ranges, with their stated net cardinalities. 

\begin{corollary}[Band-refined anchored count]\label{low:6:r3:band-count}
Fix the certificate parameters, including a positive rate scale $\nu$, and fix tail exponents $b,\tau$. Suppose the complete coefficient incidence at the original frame has the output coordinates of Lemma~\ref{low:6:r3:v-output}, with the graph row $c_{zzz}$ in an interval of length at most $P^C$. Assume that the inverse-Jacobian and fiber-multiplicity bounds are uniform over this interval. After integration in the conditional fixed-frame fibers, put $r_3=c_{zzz}$, let $\mathbf w$ collect the remaining output and endpoint variables, and let $\omega\in\mathcal N$ index the fixed external net choices. Assume that the resulting densities, including inverse Jacobians and fiber multiplicities, satisfy
\[
0\le\rho_\omega(r_3,\mathbf w)\le\mathbf1_{I_\omega(\mathbf w)}(r_3)M_\omega(\mathbf w),\qquad |I_\omega(\mathbf w)|\le P^C,
\]
where each $M_\omega$ is nonnegative and independent of $r_3$, and that
\[
\sum_{\omega\in\mathcal N}\int M_\omega(\mathbf w)\,d\mathbf w\le C_{\mathrm{cert}}(\nu)S^{-b}K^{-\tau}.
\]
Here $C_{\mathrm{cert}}(\nu)\ge0$ is independent of the phase, the original index $j$, and $S,K$; it includes all dependence of these remaining bounds on the fixed certificate parameters and on $\nu$. If the prepared endpoints also satisfy $|v|\asymp\nu$, then
\begin{equation}\label{low:6:r3:band-tail}
 |E_j(\nu,S,K)|\le P^CC_{\mathrm{cert}}(\nu)\,\nu S^{-b}K^{-\tau}.
\end{equation}
The same conclusion holds for the externally fixed seed coordinates of Lemma~\ref{low:6:r3:v-seed}, with the majorant independent of the replaced row $C_{sss}$. Proposition~\partIref{low:6:r3:weighted-holder} applies when $C_{\mathrm{cert}}(\nu)$ has a uniform bound over the rates being summed. In that proposition, $\mathcal A$ is the resulting constant independent of $\nu$.
\end{corollary}
\begin{proof}
First integrate in the conditional fixed-frame fibers. Hold the other outputs and external parameters fixed. Lemma~\ref{low:6:r3:v-output} replaces $c_{zzz}$ by $v$ with inverse row Jacobian $2|B|\le P^C$ and unchanged fiber multiplicity. The assumed majorant is independent of the replaced row, so integration over $|v|\asymp\nu$ costs at most $P^C\nu$. Its remaining integral and net sum have the stated bound. Tonelli gives \eqref{low:6:r3:band-tail}. In the coordinates of Lemma~\ref{low:6:r3:v-seed}, the old row is shifted by terms involving only the other outputs and the fixed seed; the same row-independent majorant and inverse row Jacobian give the identical bound.
\end{proof}

\paragraph{Rate labels at certificate endpoints}

The rate factor in~\eqref{low:6:r3:band-tail} uses the actual certificate endpoints, as in Definition~\ref{def:relocated-band}. Apply the conditional Haar estimate in the original fixed-frame window with its established pivot; no conjugation invariance of the null subgroup under relocation is assumed. The polynomial endpoint-range and packing costs were included in Theorem~\ref{low:6:r1:relocation}.

\subsubsection{Corridor cover and zero rate}

The interior set is supplied by Theorem~\ref{low:6:r1:relocation} and Lemma~\ref{low:6:r1:complex-box}; Section~\ref{low:6:s36cg:polynomial-distortion} supplies the transverse Jacobian bound. The construction below gives common labels, count, certificates, and endpoint ranges for every assigned frame. The vertical, parabolic and wide-corridor counts in Section~\ref{low:6:ordinary-counts} then provide coefficient bounds for these certificates. The component count uses algebraic corridors and does not require definability of the full null-flow family.

\paragraph{The outer domain and its interior set}

Use the fields and rates of Definition~\partIref{def:null-coordinates}. Fix an outer regular band domain $E$ of fixed algebraic format. Its normalized sextic coefficient, normal pivot, null pivot, and physical bounds are retained. In the application, $E$ includes the outer complete-coefficient distance condition to the conditioned rank-four family. Choose this outer defect after the prescribed fixed collar width so that the small-model estimates in Lemma~\ref{low:6:og:parabolic-collar} and Section~\ref{low:6:s36cg:polynomial-distortion} hold throughout $E$, and choose the original inner defect smaller so that Lemma~\ref{low:6:r1:complex-box} reaches this outer window. These are the nested model windows used to define the inner buffered set $A$ below. Assume $|u|+|v|\le P^C$, and that every connected contained characteristic interval has transverse Jacobian and inverse at most $K_0=P^C$. The fixed-factor outer band supplied by relocation is covered by a fixed number of adjacent dyadic bands. Lemma~\partIref{geo:characteristic-preparation} bounds the number of times a contained trajectory meets their boundaries. Applying the single-band estimate on these finitely many intervals proves the stated Jacobian bound on the whole contained interval, with a power of $P$ independent of $\nu$. Let $T\ge2$ bound $|y|+|z|+|d|$ throughout $E$. The physical framework supplies such a $T$ with $\log T=O_P(1+\log d_1^{-1})$.

Use the quantified model-box construction in the proof of Lemma~\partIref{geo:core}, adding the strict actual-rate bounds throughout each box to the full coefficient, pivot and physical restrictions. The polynomial model flow and selected algebraic root branches keep this test semialgebraic of fixed format. Coframe comparison and fixed extra box margins give a set $A\subset E$ such that every path from $\overline A$ to $E^c$ has length at least $M$ in the coframe metric. Choose the initial box larger and the inner defect smaller so that the relative band neighborhood from Theorem~\ref{low:6:r1:relocation} contains the required core box; then every relocated representative passes the test.

Choose $m=P^{-C_1}M\ge1$ so that the fixed-width flows and collars use less than $M/10$ intrinsic distance. Here $C_1$ depends only on preceding conditioning constants; choose $M$ sufficiently large before the inner defect. Geometric widths may differ from $m$ by fixed powers of $P$.

\paragraph{Two elementary geometric tools}
Let $\Phi(z,z_0,a)$ denote the horizontal coordinate of the null characteristic solving $dy/dz=d(y,z)$ with $y(z_0)=a$, on its interval of continuation in $E$.

\begin{lemma}[A corridor comparison without exit-time definability]\label{low:6:og:corridor}
Suppose an algebraic midline $y=\gamma(z)$, $z\in I$, has a contained corridor $|y-\gamma(z)|<w(z)$ in $E$, where $w_0\le w(z)\le2w_0$. Let $\alpha$ be its seed residual from Definition~\partIref{def:seeds-collars}.

\noindent\textup{(i)} If $|\alpha|\le h$, then flows from the midline and initial horizontal offsets at most $c_Pw_0$ remain in the corridor for time at most $c_P\min(|I|,w_0/h)$ away from its vertical endpoints. On these time intervals,
\[
 |\Phi(z,z_0,\gamma(z_0)) -\gamma(z)|\le K_0\int_{z_0}^z|\alpha(\tau)|\,d\tau.
\]

\noindent\textup{(ii)} If $\alpha$ has one sign and $h\le|\alpha|\le2h$, the projection from the seed to the horizontal section has derivative $J\alpha$, with $K_0^{-1}\le J\le K_0$. A tube of flow-time width $c_Pw_0/h$ and seed $z$-length comparable to $H$ therefore has area comparable, by powers of $P$, to $Hw_0$, provided $hH\ge C_Pw_0$.
\end{lemma}
\begin{proof}
Use the nonlinear variation-of-constants identity
\[
 \gamma(z)-\Phi(z,z_0,\gamma(z_0))=\int_{z_0}^z\partial_a\Phi(z,\tau,\gamma(\tau))\alpha(\tau)\,d\tau.
\]
Apply a first-exit argument simultaneously to all intermediate trajectories in this formula and the allowed initial offsets. Before any exit, the $K_0$ bound makes each of the two displacement contributions at most $w_0/4$ for sufficiently small $c_P$, contradicting exit from the corridor. Thus the formula and containment hold throughout the stated interval.

For a fixed final horizontal section, differentiate the flow with respect to its initial time and initial point on the midline. The exact derivative is $\partial_a\Phi\,(\gamma'-d)=J\alpha$. For the map $\Psi(z,t)=(\Phi(z+t,z,\gamma(z)),z+t)$ one consequently has
\[
 \Psi^*\theta=J\alpha\,dz,\qquad \Psi^*(dz)=dz+dt,\qquad |\det D\Psi|=J|\alpha|.
\]
A characteristic crosses a seed graph with residual of one strict sign at most once: two crossings with the same orientation require an intervening crossing of the opposite orientation. Thus the tube has multiplicity one on the chosen graph arc. Integrating the displayed determinant proves the area assertion.
\end{proof}

\begin{lemma}[Projection caps]\label{low:6:og:caps}
Let $C\subset A$ be a connected corridor whose $z$ projection is an interval of length less than $m$.

\noindent\textup{(i)} The null-characteristic projection of $C$ to a fixed horizontal section meeting that interval has an interval as its image; every image point retains intrinsic buffer at least $M-m$, and the construction is measurable in all parameters.

\noindent\textup{(ii)} If the image has length $\ell$, then one common cap with scales $L=\max(m,\ell)$ and $W=m$, up to powers of $P$, contains all projected representatives, provides a contained certificate of area comparable to $Lm$, and has count $O_P(\Lambda Lm)$. Original labels retain their anchors and weights.
\end{lemma}
\begin{proof}
Use the projection construction in the proof of Lemma~\partIref{geo:caps}, with the present time cutoff $m$. The core model-box test and Lemma~\ref{low:6:r1:complex-box} retain the flows in $E$ with polynomial relative-frame bounds. Their transverse Jacobians are controlled by the bandwise hypothesis. The image is an interval by connectedness and continuity, with buffer at least $M-m$; extending it to length at least $m$ uses at most another $m$. Its null flows for time $c_Pm$ give a contained rectangle of bounded coframe and area comparable to $Lm$.

For the count, the actual flow tube for $|t|\le2m$ contains all assigned points and is covered by $O_P(Lm)$ parameter unit boxes. The model-box frame comparison and original family packing give $O_P(\Lambda Lm)$ labels. The Borel endpoint construction in the proof of Lemma~\partIref{geo:caps} applies to these same continuous projections. The original anchors and weights remain fixed.
\end{proof}

For a corridor with $|\alpha|\le2h$, the cap length obeys the more explicit bound $\ell\le P^C(w_0+hH)$, where $H$ is its $z$ length. This follows by applying the same variation formula to the midline and the $K_0$ Lipschitz estimate to each horizontal section. 

\paragraph{The finite allocation}

\begin{theorem}[Common buffered corridor cover]\label{low:6:og:common-cover}

\noindent\textup{(i)} Under the stated hypotheses, all assigned labels in $A$ have a coefficient-measurable assignment to at most
\[
P^C(1+\log T)^C
\]
pieces. Each piece has scales $L\ge W\ge P^{-C}M\ge1$, count at most $P^C\Lambda LW$, and a contained geometric certificate of area at least $P^{-C}LW$.

\noindent\textup{(ii)} Every fixed frame assigned to that piece can be joined to its certificate through $E$ by a path of intrinsic length at most $P^CL$. In its original normalized coordinates the certificate therefore has endpoint ranges
\[
 |d|\le P^CL,\qquad |z|\le P^CL,\qquad |y|+|y-dz|\le P^CL^2.
\]
The certificates are characteristic rectangles, characteristic tubes of algebraic seed arcs admitting parametrizations with coframe and inverse bounded by powers of $P$ on the stated scales, and fixed-width algebraic collars if the additional small-model hypotheses of the collar lemma hold.

\noindent\textup{(iii)} Writing $\mathscr P$ for the assigned pieces, for every fixed $\eta>0$ one has $\sum_{\Pi\in\mathscr P}(L_\Pi W_\Pi)^{-\eta}\le P^{C_\eta}(1+\log T)^{C_\eta}$. 
\end{theorem}
\begin{proof}
Apply Lemma~\partIref{geo:bands-packing} to the stated buffered core with \(\tau=m\) and \(w_*=P^{-C_2}m\); write \(\gamma\) for its midpoint \(\Gamma\). The core model boxes give its local geometric hypotheses, and original family packing with the already established relocation multiplicity gives its label hypothesis, uniformly in the specified rate band. Retain the full outer admissibility and fixed root branches. Lower-dimensional arcs use the small-width constructions below, and isolated points use caps. Resolve all ties before choosing a coefficient-integration frame.

On a small-width component with \(|\alpha|\le1\), the seed's intrinsic length is comparable to its \(z\) length. Its horizontal ruled collar of width \(c_Pm\) stays in \(E\) by the original buffer and has coframe
\[
 (\alpha+d(\gamma,z)-d(\gamma+t,z))\,dz+dt,\qquad dz.
\]
The entries and inverse are bounded by powers of $P$, since $|Yd|\le P^C$ and $m$ is a fixed polynomial in $P$. On a small-width component with signed \(|\alpha|>1\), parameterize by the seed flux $ds=\alpha dz$, with orientation adjusted for its sign. Then $|dz/ds|\le1$, and the null collar of time width $c_Pm$ stays in $E$ and has coframe $J\,ds,dz/ds\,ds+dt$. Its area is comparable to $m$ times the seed's intrinsic length. A seed shorter than $m$ uses one balanced cap. Assign the whole original narrow corridor to its seed collar; labels need not lie inside the certificate. Their horizontal bridge to the seed has length at most $2w_*$, and subsequent seed travel has its stated intrinsic length. The area and boundary-length estimates below give the label count for this assignment. Only boundedly many such pieces occur.

Use the width and signed residual components of Lemma~\partIref{geo:bands-packing}, including its single low range and its two high tails.

For every component whose $z$ length $H$ is less than $m$, use Lemma~\ref{low:6:og:caps}. This also assigns its vertical subdivision edges and any artificial thin band pieces. This construction uses no interior margin for a newly selected width or residual band. The buffer used by a cap is the original buffer of every point of $A$.

For a high-residual component with $H\ge m$, write $L_\gamma=\int_I|\alpha|\,dz$. Its seed is $Y$-oriented, and
\[
L_\gamma\ge C_Pw_0H/m\ge C_Pw_0,\qquad Hw_0\le c_PmL_\gamma.
\]
Assign the entire corridor to the fixed-width null collar of its entire midline, with scales $L\asymp_PL_\gamma$ and $W\asymp_Pm$. Its area is comparable to $mL_\gamma$. Every fixed frame reaches the midline by a horizontal path of length $O(w_0)\le O_P(L_\gamma)$, then travels along it and along a collar characteristic. The assignment includes anchors outside the collar image, with the path bounds just established. The fixed-width collar can be replaced by the algebraic parabolic collar proved below.

It remains to treat a moderate-residual component with $H\ge m$. For a signed band write $h\le|\alpha|\le2h$; the small-residual case uses its uniform upper bound instead. If $hH\le C_Pw_0$, a sufficiently short fixed fraction of the $z$ interval and a horizontal seed of width $c_Pw_0$ give a contained characteristic rectangle by Lemma~\ref{low:6:og:corridor}. Its area is at least $P^{-C}Hw_0$. Use $L\asymp_P\max(H,w_0)$ and $W\asymp_P\min(H,w_0)$. The entire original corridor, not only this retained certificate, is assigned to that rectangle.

If $hH>C_Pw_0$, retain a central seed subinterval of $z$ length comparable to $H$ and use the contained null tube of time width $c_Pw_0/h$. For $h\ge1$, the flux parameter gives length $L\asymp_PhH$, width $W\asymp_Pw_0/h$, and coframe $J\,ds,m(s)ds+dt$, where $|m(s)|\le1$. For $h<1$, put $\zeta=z+t$ and $\eta=ht$ in the tube parametrization. A central rectangle in these variables has dimensions $H,w_0$, and its coframe is
\[
J\alpha\,d\zeta-J(\alpha/h)\,d\eta,\qquad d\zeta.
\]
Both this matrix and its inverse are bounded by powers of $P$, since $\alpha/h$ is between fixed positive constants in absolute value. Use $L\asymp_PH$ and $W\asymp_Pw_0$. In both cases $LW\asymp_PHw_0$. The upper moderate cutoff gives $w_0/h\ge P^{-C}m$ when $h\ge1$, and every other width is at least $P^{-C}m$. Thus every transverse width used in the assignment is at least $P^{-C}m$.

Lemma~\partIref{geo:bands-packing} gives the label count from the corridor area and intrinsic boundary length. In a moderate component, its bound is \(P^C(H+hH+w_0+Hw_0)\le P^CHw_0\), using \(H\ge m\), \(w_0\ge P^{-C}m\), and \(h\le P^C\max(1,w_0/m)\). In a high-residual component, \(L_\gamma\ge C_Pw_0H/m\) and \(|\alpha|\ge2\) bound it by \(P^CmL_\gamma\). In a small-width component, put \(\ell_{\mathrm{int}}=H+\int_I|\alpha|\,dz\); the bound is \(P^Cm\max(m,\ell_{\mathrm{int}})\), including when \(\alpha\equiv0\). These are the respective certificate areas up to powers of \(P\). The arc estimate handles lower-dimensional labels. Projection caps retain their enclosing-flow-rectangle count from Lemma~\ref{low:6:og:caps}. Thus every piece has at most \(P^C\Lambda LW\) original labels.

Every path from the fixed frame to the certificate used above has intrinsic length at most $P^CL$: use its horizontal bridge, the appropriate part of its common seed, and its final contained transverse path. For a projection cap use first its flow of length less than $m$, then the horizontal cap. Along any such path $dd=u\theta+v\,dz$, so $|\Delta d|\le P^CL$. Normalize $d=0$ at the fixed frame. Integrating $dy=\theta+d\,dz$ gives $|\Delta z|\le P^CL$ and $|\Delta y|+|\Delta y-d\Delta z|\le P^CL^2$. This proves the endpoint assertions for every assigned fixed frame.

Lemma~\partIref{geo:bands-packing} gives the total piece count and polynomial upper bounds in \(T\) for the seed lengths. The retained geometric constructions therefore have the permitted physical logarithms. Their flux integrals, cap endpoints and flows are Borel by the arguments above; the fixed tie rules give a common coefficient-measurable assignment. Since every area is at least \(P^{-C}M^2\ge1\), the piece count implies the asserted inverse-area bound.
\end{proof}

\paragraph{Algebraic collars of fixed polynomial width}\label{low:6:og:fixed-collars}

\begin{lemma}[Parabolic replacement of a fixed-width null collar]\label{low:6:og:parabolic-collar}
Let an algebraic seed be parameterized by its null flux, so that $\theta(\gamma_s)=1$ and $z_s=m(s)$ with $|m|\le1$. Suppose every seed point has the prescribed outer interior margin of the admissible set, and suppose the complete sextic phase at all points of the relevant outer model boxes is sufficiently close to the conditioned rank-four family.

\noindent\textup{(i)} For a fixed polynomial width $m_0=P^A$, chosen before this coefficient defect, the map
\[
 \Psi(s,t)=\left(\gamma_y(s)+d_\gamma(s)t+\frac12v_\gamma(s)t^2,\ \gamma_z(s)+t\right),\qquad |t|<c_Pm_0,
\]
is a contained algebraic collar of bounded multiplicity. In the original algebraic seed parameter it has fixed semialgebraic format.

\noindent\textup{(ii)} Its coframe and inverse are bounded by powers of $P$, its area is comparable to the seed's intrinsic length times $m_0$, and it has the same polynomial anchor-relative endpoint ranges as the null collar.

\noindent\textup{(iii)} The same statement applies to horizontal seed intervals and hence to projection caps.
\end{lemma}
\begin{proof}
At an exact conditioned rank-four null model, $d=\gamma z$, hence $u=Yv=Zv=0$. Quantitative continuity in the complete sextic coefficients, with the normal and null pivots retained, gives \partIeqref{geo:small-rates} on the designated outer model boxes. Choose this defect after $m_0$, with the prescribed buffer larger than $4m_0$.

Apply Lemma~\partIref{geo:collar} with seed data $(a,b,D,V)=(\gamma_y,\gamma_z,d_\gamma,v_\gamma)$ and $b'=m$. Its Jacobian \partIeqref{geo:collar-jacobian}, containment and coframe estimates give the stated algebraic collar and bounded multiplicity. Since $\theta(\gamma_s)=1$ and $|m|\le1$, intrinsic seed length and flux length are comparable, giving the area claim. The collar remains at fixed polynomial intrinsic distance from the seed; concatenating with the anchor-relative paths constructed in the proof of Theorem~\ref{low:6:og:common-cover} preserves the anchor-relative endpoint ranges. Horizontal seeds have $m=0$ and satisfy the same argument.
\end{proof}

For a $Z$-oriented algebraic seed with $|\gamma'-d_\gamma|\le1$, the calculation in the proof of Lemma~\partIref{geo:narrow-null} gives the horizontal collar of width $c_Pm_0$, with unit area Jacobian and polynomial coframe bounds. Ordinary finite-format preparation applies to both types of algebraic collar. Their external seed grids and complete fixed-frame coefficient counts are proved in Section~\ref{low:6:sp36:section}.


The coefficient proofs below use the ordinary horizontal coordinates on these algebraic corridors and keep their actual amplitudes.

\paragraph{The identically zero rate}

\begin{corollary}[Assignment when the rate vanishes identically]\label{low:6:r1:zero-common}
Suppose $v$ vanishes identically on the connected analytic critical/null branch selected in Theorem~\ref{low:6:r1:relocation}. Then the same corridor construction gives its assignment, with the same count, multiplicity, geometric certificates, and endpoint ranges for every assigned frame. No positive rate label is assigned to these pieces.
\end{corollary}
\begin{proof}
The identity $v\equiv0$ on the local model box continues along the connected analytic branch, so it is not a thin additional constraint on its outer admissible set. Lemma~\ref{low:6:r1:complex-box} supplies the prescribed interior set, with the rate bounds omitted from its semialgebraic model-box test. The compatibility equation $ZB=-2Bu+rv$ becomes $u=-\tfrac12Z\log|B|$. Hence the transverse flow Jacobian is exactly
\[
J(t)=J(t_0)\left|\frac{B(t_0)}{B(t)}\right|^{1/2}.
\]
Both $J$ and its inverse are bounded by powers of $P$ on every connected contained characteristic interval. All the hypotheses used in Theorem~\ref{low:6:og:common-cover} therefore hold on this domain. Its proof uses no positive lower bound for $|v|$ except through the already established Jacobian bound, so it applies without an amplitude subdivision. Lemma~\ref{low:6:og:parabolic-collar} uses the same full-coefficient model buffer and also applies. The same ordinary coefficient counts hold with the rate output given its full bounded interval, rather than an interval of length $O(\nu)$; thus their bounds have the factor $\nu$ omitted. This is also the form used in the final summation.
\end{proof}

\subsection{Critical graphs: coefficient bounds and completion}\label{low:6:ordinary-counts}\label{low:6:so36:section}\label{low:6:sp36:section}\label{low:6:wc36:section}\label{low:6:ordinary-assembly}\label{low:6:as36:assembly}\label{low:6:s36:global}

\subsubsection{Ordinary horizontal ribbons}

We first count a supplied algebraic ribbon whose $z$ length dominates its horizontal width, as provided by Theorem~\ref{low:6:og:common-cover}. Here $L,W$ are these ordinary coordinate dimensions and $K_0=L/W$. The coefficient substitution and conditional Haar restriction use the original fixed frame. The application to all assigned anchors follows after the ribbon calculation. The assigned geometric aspect $K$ can differ from $K_0$; equation~\eqref{low:6:as36:zlong-assigned} gives the required conversion.

\par\smallskip
\noindent\textit{Proof stages for Theorem~\ref{low:6:so36:count}.}\par\nobreak
\noindent\textup{Step 1.} Scalar preparation and endpoint ranges.\par\nobreak
\noindent\textup{Step 2.} Coefficient determinant and slope charts.\par\nobreak
\noindent\textup{Step 3.} Joint graph windows and fixed seed grids.\par\nobreak
\noindent\textup{Step 4.} Fixed-frame integration and the two products.\par\nobreak
\smallskip

\paragraph{The certificate and its scalar preparation}

In the original tangential coordinates $(y,z)$, the actual nondegenerate critical graph satisfies
\[
F_x(c(y,z),y,z)=0.
\]
Use Definition~\partIref{def:null-coordinates}, writing $\phi=F|_c$, $\lambda=F_{xx}|_c$, $u_0=u$ and $v_0=v$. Throughout the ribbon retain the normalized sextic coefficient window, the inverse normal pivot, and the inverse selected null pivot. In particular the full admissibility gives
\[
|Y\phi|+|Z\phi|+|\lambda|+|u_0|+|v_0|\le P^C,\qquad |\lambda|,|B|\ge P^{-C}.
\]
The derivative bounds for $u_0,v_0$ of order zero follow directly from the actual finite graph jets in the normalized endpoint frame; for example $u_0=-c_{yzz}/(2B)$ and $v_0=-c_{zzz}/(2B)$ when its selected null slope is zero. No derivative phase admissibility substitutes for the full phase window.

Assume the certificate is a contained horizontal ribbon
\begin{equation}\label{low:6:so36:ribbon}
(y,z)=(\Gamma(\xi)+\tau,\xi),\qquad |\xi|<L,\quad |\tau|<W,\qquad L\ge W\ge1,
\end{equation}
where $\Gamma$ is algebraic of fixed description complexity. Put $S=LW$, $K_0=L/W$, and assume only the following coframe bound:
\begin{equation}\label{low:6:so36:defect}
e(\xi,\tau):=\Gamma'(\xi)-d(\Gamma(\xi)+\tau,\xi),\qquad |e(\xi,\tau)|\le P^CL.
\end{equation}
The actual area element is $d\tau\,d\xi$. If a rate band is retained, suppose $|v_0|\le P^C\nu$, with $0<\nu\le1$ after a polynomial normalization. A lower band bound is needed by the geometric construction, but the coefficient estimate uses only this upper bound.

\begin{lemma}[The cost occurs only in the purely longitudinal rows]\label{low:6:so36:scalar-widths}

\noindent\textup{(i)} For any fixed differentiation order, a common subset of area at least $P^{-C}LW$ in \eqref{low:6:so36:ribbon} satisfies
\[
|\partial_\xi^i\partial_\tau^j\phi|\le P^C L^{-i}W^{1-j}\quad(j\ge1),\qquad |\partial_\xi^i\phi|\le P^C L\,L^{1-i}\quad(i\ge1),
\]
and
\[
|\partial_\xi^i\partial_\tau^j\lambda|\le P^CL^{-i}W^{-j}.
\]

\noindent\textup{(ii)} More precisely, $|\partial_\xi^i\phi|\le P^C(1+\sup|e|)L^{1-i}$ for $i\ge1$ through the prescribed order on the same common good set, where the supremum is over the full ribbon. In the $48$-row grid this factor occurs six times, while in the $48$-row weighted set it occurs twelve times. The $42$ mixed grid rows do not acquire it.
\end{lemma}
\begin{proof}
In ribbon coordinates, $\partial_\tau\phi=Y\phi$ and $\partial_\xi\phi=Z\phi+eY\phi$, with bounds $P^C$ and $P^C(1+\sup|e|)$ respectively. Both are algebraic of fixed format on the full rectangle. Apply Lemma~\partIref{geo:preparation} to all required mixed derivatives of the first scalar, the purely longitudinal derivatives of the second, and the bounded scalar $\lambda$, choosing simultaneous exceptional proportions sufficiently small. This gives the mixed widths $W^{1-j}$ and purely longitudinal widths $L^{1-i}$ on a positive-area set, while keeping the two gradient amplitudes distinct.
\end{proof}

\begin{lemma}[Endpoint ranges and prepared centerline jets]\label{low:6:so36:endpoints}
Suppose the fixed frame has a polynomially bounded right comparison either to a point of the ribbon or to a point joined to the ribbon by a horizontal segment of length at most $P^CW$, with the full admissibility bounds holding on that segment.

\noindent\textup{(i)} Relative to that fixed frame, its endpoint coordinates and null slopes satisfy
\[
|f|\le P^CL,\qquad |d|\le P^CL^2,\qquad |y|\le P^CL^3,\qquad |\eta|=|y-df|\le P^CL^3.
\]
Thus the fixed-frame $(\eta,f)$ candidate region has area at most $P^CL^4$.

\noindent\textup{(ii)} On a common positive-area good subset, in addition to Lemma~\ref{low:6:so36:scalar-widths},
\[
|\Gamma^{(j)}|\le P^CL^{3-j}\qquad(2\le j\le12).
\]

\noindent\textup{(iii)} At an endpoint normalized so that its selected null slope vanishes, the relative first seed derivative has the stronger range $|e|\le P^CL$.
\end{lemma}
\begin{proof}
If a horizontal bridge is present, its slope change is at most $P^CW$ by $|Yd|\le P^C$, and its horizontal displacement is at most $P^CW$. Since $W\le L$, these changes fit the asserted endpoint ranges. Begin the reference curve at its landing point in the ribbon. Use the transverse level passing through that point as reference centerline; translating its transverse parameter does not change $\Gamma'$ or its higher derivatives. Along this curve, differentiation of the actual slope gives
\[
\frac{d}{d\xi}d(\Gamma(\xi)+\tau_0,\xi)=v_0+e u_0.
\]
Its absolute value is at most $P^CL$. Starting from the initial slope of size at most $P^C(1+W)$, integration gives $|d|\le P^CL^2$ along the reference curve. Across any horizontal section, the bound $|Yd|\le P^C$ changes $d$ by at most $P^CW$, so the same bound holds throughout the ribbon. Since $\Gamma'=d+e$, one has $|\Gamma'|\le P^CL^2$. Integrating this bound gives $|y|\le P^CL^3$, and $|f|\le P^CL$ is immediate. Their combination bounds $\eta=y-df$. A bounded right comparison to the fixed frame adds only polynomially bounded translations and linear shears, which are absorbed in these ranges. 

The algebraic scalar $\Gamma'$ has amplitude $P^CL^2$ on its full interval. Its one-dimensional derivative preparation gives $|\Gamma^{(j)}|\le P^CL^{3-j}$ on all but a prescribed small longitudinal proportion. Intersect its product with the full transverse interval with the good set from Lemma~\ref{low:6:so36:scalar-widths}. No transverse width is lost. At a normalized endpoint the local seed is $\Gamma(\xi_0+\xi)-\Gamma(\xi_0)-d_0\xi$, so its first derivative is exactly $e$ and all its higher derivatives are unchanged.
\end{proof}

\paragraph{The coefficient Jacobian}

At the normalized endpoint let $c(0)=c_y(0)=c_z(0)=c_{zz}(0)=0$, $r=c_{yy}(0)$, and $B=c_{yz}(0)$. Let a fixed external seed have Taylor expansion $\gamma(\xi)=e\xi+\sum_{j\ge2}\gamma_j\xi^j$, and form ordinary output jets using $(y,z)=(\tau+\gamma(\xi),\xi)$. Use $\mathcal I_d^{\rm g},\mathcal I_d^{\rm w},\mathcal E^{\rm g},\mathcal E^{\rm w}$ from Proposition~\ref{low:6:nc36:strong-atlas} and its preceding definitions.
The first index is longitudinal. Use $48$ ordinary critical-value jets indexed by $\mathcal I_6\setminus\{0\}$, $16$ ordinary second-normal jets indexed by $\mathcal I_3$, $18$ graph jets indexed by $\mathcal E$, and the third-normal value $\mu=F_{xxx}|_c(0)$. These are all $83$ outputs before deleting the four graph normalization rows.

\begin{theorem}[Exact quadratic-jet determinant]\label{low:6:so36:exact-atlas}

\noindent\textup{(i)} The Jacobian of this ordinary output map, at the normalized endpoint, is $\lambda(0)^{-18}$ times the fixed-graph determinant evaluated at
\begin{equation}\label{low:6:so36:quadratic}
Q(\xi,\tau)=A\xi^2+b\xi\tau+C\tau^2,\qquad A=eB+\tfrac12re^2,\quad b=B+re,\quad C=\tfrac12r.
\end{equation}
It has no dependence on any higher graph jet, any higher external seed jet, or any higher normal derivative.

\noindent\textup{(ii)} In the weighted chart its absolute value is exactly
\[
56623104\,|A|^{73}|\lambda(0)|^{-18}.
\]

\noindent\textup{(iii)} In the grid chart the fixed-graph determinant has the form
\[
G(A,b,C)=b^{73}g(AC/b^2),\qquad |g(0)|=56623104,
\]
where $g$ is a fixed real polynomial. Consequently there is a fixed $\delta_*>0$ such that $|AC|\le\delta_*b^2$ implies $|G(A,b,C)|\gtrsim|b|^{73}$. The formula at $b=0$ is interpreted through the homogeneous polynomial $G$.
\end{theorem}
\begin{proof}
Fix the external seed during coefficient differentiation. The two output lists are precisely the $D=6$ lists in Lemma~\partIref{surf:strong-atlas}; their graph index sets are downward closed and contain the respective set $\mathcal I_3$. Lemma~\ref{low:6:nc36:large-normal-ordinary} therefore reduces the Jacobian to the fixed-graph determinant with factor $\lambda(0)^{-18}$.

The weight-triangular argument in the proof of Lemma~\partIref{surf:strong-atlas} removes all higher graph and external seed terms from the determinant. The linear substitution $(y,z)=(\tau+e\xi,\xi)$ has coefficient determinant of absolute value one by Lemma~\partIref{an:character} and sends the quadratic graph to \eqref{low:6:so36:quadratic}.

By Proposition~\ref{low:6:nc36:strong-atlas}, the weighted fixed-graph determinant is $56623104|A|^{73}$ in absolute value. Lemma~\partIref{surf:strong-atlas} and \partIeqref{surf:tD} give the grid form $b^{73}g(AC/b^2)$ with $t_6=73$, and the mixed quadratic calculation in Proposition~\ref{low:6:nc36:strong-atlas} gives $|g(0)|=56623104$. Continuity of this fixed polynomial near zero gives the lower bound, while its homogeneous form supplies the polynomial interpretation at $b=0$.
\end{proof}

\begin{lemma}[Two quadratic charts cover every first slope]\label{low:6:so36:two-charts}
Suppose $|r|,|B|\le P^C$ and $|B|\ge P^{-C}$.

\noindent\textup{(i)} There is $\sigma_P\ge P^{-C'}$ such that, writing $D=1+|e|$, every endpoint belongs to one of the following two charts:
\[
\text{weighted: }|A|\ge\sigma_PD,\qquad
\text{grid: }|AC|\le\delta_*b^2\text{ and }|b|\ge P^{-C'}.
\]

\noindent\textup{(ii)} The choice is made by the weighted inequality first, with the grid chart used on its complement. It introduces only two coefficient sets and no slope-dependent number of charts.
\end{lemma}
\begin{proof}
The identity is $b^2-4AC=B^2$. If $|e|\le1$ and $|A|<\sigma_P(1+|e|)$, then $|AC|\le P^C\sigma_P$. If $|e|\ge1$, eliminate $r$ to get $AC=(A/e)(A/e-B)$, and hence $|AC|\le2\sigma_P(2\sigma_P+|B|)$. Choose $\sigma_P$ sufficiently small as a fixed inverse power of $P$. In either case $|AC|$ is a sufficiently small multiple of the lower bound for $B^2$. The identity then gives $|b|\ge P^{-C'}$ and $|AC|\le\delta_*b^2$.
\end{proof}

\paragraph{Graph derivative bounds and nets for relative seed parameters}

The four graph normalization rows may be changed back to their original null-coordinate versions by a unit-determinant triangular row change and then deleted. For example the second-order condition is $c_{\xi\xi}-2e c_{\xi\tau}+e^2c_{\tau\tau}=0$, since the first graph derivatives vanish. The two remaining second-order variables are an area-preserving shear of $(r,B)$. All complete graph jets of ordinary total order three and four also transform with determinant of absolute value one. Let $\mathcal W_{\le4}$ be the joint window of the eleven remaining graph coordinates of order at most four. Then $|\mathcal W_{\le4}|\le P^C$ and $|\mathcal W_{\le4}\cap\{|v_0|\le P^C\nu\}|\le P^C\nu$. For the second bound, the original coordinate $c_{zzz}=-2Bv_0$ is among this complete third-order jet and may be replaced by $v_0$ with a conditioned triangular Jacobian, keeping $B$ as another output. These are joint volumes; individual sheared coordinates may have large ranges.

\begin{lemma}[The three additional graph variables]\label{low:6:so36:graph-windows}
At a prepared endpoint set $D=1+|e|$.

\noindent\textup{(i)} The weighted graph-coordinate set, after the four normalizations are removed, has volume at most $P^C\nu D^{15}$. In the weighted chart, the graph-coordinate volume multiplied by the supremum of the inverse absolute coefficient Jacobian determinant is at most $P^C\nu D^{-58}\le P^C\nu$.

\noindent\textup{(ii)} In the grid chart it has volume at most
\[
P^C\nu\bigl(D^8+|\gamma_2|D^6\bigr)\le P^C\nu L^8,
\]
where $\gamma_2$ is the quadratic Taylor coefficient of the rounded seed.
\end{lemma}
\begin{proof}
The original endpoint graph jets through order six are bounded by powers of $P$, since the full normalized sextic and $\lambda^{-1}$ are bounded. First remove the nonlinear seed part formally. If $c_e(\xi,\tau)=c(\tau+e\xi,\xi)$ and $\zeta(\xi)=\gamma(\xi)-e\xi$, the curved graph is $c_e(\xi,\tau+\zeta(\xi))$. Its jet transformation is triangular by total degree. The weighted graph set is closed under every substitution index $(i,j)\mapsto(i-k,j+1)$ for $k\ge2$. The resulting transformation on its eighteen coordinates has determinant one, and the four normalization rows remain compatible with it. Thus the higher seed derivatives cause no weighted graph-window loss at all. Conditional on the full lower-order graph jet, the three remaining linear-seed rows $(5,0),(4,1),(6,0)$ have ranges at most $P^CD^5,P^CD^4,P^CD^6$. Their product is $P^CD^{15}$.

For the grid graph set the only missing substitution index is $(1,4)$ arising from $(3,3)$ and $k=2$. The two fifth-order outputs have ranges $P^CD^3$ and $P^CD^2$. After the known lower-order terms are removed, the sixth-order output has range $P^C(D^3+|\gamma_2|D)$. The three remaining graph ranges have product $\lesssim_P D^3D^2(D^3+|\gamma_2|D)=D^8+|\gamma_2|D^6$. The preparation gives $|\gamma_2|\le P^CL$ and the relative first-slope range gives $D\le P^CL$, so this is at most $P^CL^8$. The lower-order joint volume, including its factor $\nu$, was retained throughout. Finally the weighted Jacobian lower bound is $P^{-C}D^{73}$, which absorbs $D^{15}$ exactly as asserted.
\end{proof}

\begin{lemma}[Fixed grids for relative endpoint seeds]\label{low:6:so36:seed-grids}

\noindent\textup{(i)} For the grid chart six external seed derivatives suffice, and for the weighted chart twelve suffice. They can be rounded on grids fixed in advance from $(P,L,W)$, with mesh $c_PW/L^j$ for derivative order $j$. The respective numbers of choices are at most
\[
P^C L^{11}K_0^6\quad\hbox{and}\quad P^C L^{23}K_0^{12}.
\]

\noindent\textup{(ii)} For each fixed rounded seed, the coefficient output map depends only on the normalized endpoint phase. There is no additional endpoint shear interval to integrate.
\end{lemma}
\begin{proof}
The constant seed coefficient is zero at the endpoint and requires no grid. Its first coefficient is the relative slope $e=\Gamma'-d_0$, in the fixed interval of size $P^CL$; its mesh $W/L$ costs $P^CLK_0$ choices. For every $j\ge2$, the prepared interval has size $P^CL^{3-j}$ and its mesh is $W/L^j$, so each such coefficient costs $P^CL^2K_0$ choices. Thus $m$ seed derivatives cost $P^C(LK_0)(L^2K_0)^{m-1}$, giving the displayed counts at $m=6,12$. The intervals are centered at zero and fixed before the coefficients are selected. No coefficient-dependent grid center is used.

Rounding changes the ribbon parameter relation by a formal shear $\tau\mapsto\tau+\delta\gamma(\xi)$ whose scaled derivatives satisfy $|\delta\gamma^{(j)}|L^j/W\le c_P$. The chain rule preserves all prepared scalar widths after a fixed enlargement. A mixed critical row still contains a transverse derivative, so it cannot acquire a purely longitudinal term of the larger amplitude. A purely longitudinal row receives transverse corrections with amplitude at most $W$, which is bounded by $(1+\sup|e|)L$. The bounded scalar $\lambda$ is handled by the same scaled chain rule. Prepare every fixed extra scalar derivative needed for these comparisons before rounding.

The seed is expressed in the actual normalized endpoint coordinates. Once its rounded coefficients are fixed, the output map is an ordinary map of $J$ alone. The first seed grid includes the full relative range of the endpoint shear. The graph-window bounds remain valid after rounding, since the mesh changes the first derivative by at most a constant and the second derivative by at most $W/L^2\le1$.
\end{proof}

\paragraph{The coefficient count in the original frame}

\begin{theorem}[Ordinary horizontal-ribbon coefficient estimate]\label{low:6:so36:count}
Fix an original fixed frame, with scales $S=LW$ and $K_0=L/W$ satisfying $1\le K_0\le S$. Suppose a measurable set of coefficient vectors supplies an already assigned ribbon satisfying \eqref{low:6:so36:ribbon} and \eqref{low:6:so36:defect}, the full regular-null admissibility bounds above, and a polynomially bounded comparison between that anchor and one ribbon point, allowing also the full-activity horizontal bridge of Lemma~\ref{low:6:so36:endpoints}.

\noindent\textup{(i)} Then its complete $83$-coefficient volume obeys
\begin{equation}\label{low:6:so36:two-tails}
\operatorname{vol}E_j(L,W,\nu)\le P^C\nu\left(S^{-267/2}K_0^{5/2}+S^{-257/2}K_0^{-93/2}\right).
\end{equation}

\noindent\textup{(ii)} The coefficient set also satisfies
\begin{equation}\label{low:6:so36:common-tail}
\operatorname{vol}E_j(L,W,\nu)\le P^C\nu S^{-257/2}K_0^{-1}.
\end{equation}
\end{theorem}
\begin{proof}
Use the common good set of area at least $P^{-C}S$. Apply the scalar width products in~\eqref{low:6:nc36:raw-products} in the present coordinate scales, the purely longitudinal amplitudes of Lemma~\ref{low:6:so36:scalar-widths}, the joint graph-coordinate volumes of Lemma~\ref{low:6:so36:graph-windows}, and the seed grids of Lemma~\ref{low:6:so36:seed-grids}. The third-normal output retains its bounded interval. Table~\ref{app:table-oblique-factors} collects these factors and the point factor justified below.

For each fixed rounded seed, the complete $79$-output slice map has the determinant bound of Theorem~\ref{low:6:so36:exact-atlas} and bounded inverse multiplicity by Lemma~\partIref{geo:coefficient-map}. It depends only on the endpoint phase, so Lemma~\partIref{surf:null-density} gives the original conditional $H_0$ Haar bound. The candidate $(\eta,f)$ area is $P^CL^4$ by Lemma~\ref{low:6:so36:endpoints}; division by the good area gives $P^CL^4/(LW)=P^CL^2K_0$. The finite grids account for all seed choices.

Multiplication gives
\[
\nu S^{-147}K_0^{-18}L^6L^8(L^{11}K_0^6)(L^2K_0)=\nu S^{-267/2}K_0^{5/2},
\]
\[
\nu S^{-147}K_0^{-78}L^{12}(L^{23}K_0^{12})(L^2K_0)=\nu S^{-257/2}K_0^{-93/2}.
\]
There are only two charts. Since $1\le K_0\le S$, the first term is at most $\nu S^{-130}K_0^{-1}$ and the second at most $\nu S^{-257/2}K_0^{-1}$. This proves both stated bounds, with the single rate factor supplied by the retained graph row.
\end{proof}

\paragraph{Application to every vertically long assigned corridor}

\begin{proposition}[All moderate corridors with $H\ge w_0$ fit the ordinary horizontal count]\label{low:6:oi:z-long}
Consider a moderate corridor from Theorem~\ref{low:6:og:common-cover} with $H\ge w_0$.

\noindent\textup{(i)} It admits an ordinary algebraic horizontal-ribbon coefficient certificate of area $\asymp_P Hw_0$, with $|\gamma'(z)-d(\gamma(z)+\tau,z)|\le P^CH$ throughout the ribbon and original-label count $O_P(\Lambda Hw_0)$.

\noindent\textup{(ii)} The coefficient-volume bound in Theorem~\ref{low:6:so36:count} applies with coordinate aspect $K_0=H/w_0$.
\end{proposition}
\begin{proof}
The full cell contains the ordinary ribbon $y=\gamma(z)+\tau$, $|\tau|<c_Pw_0$, over a fixed central fraction of the same $z$ interval. No characteristic parametrization is needed for this certificate. On a signed moderate band the upper cutoff gives $h\le P^C\max(1,w_0/m)\le P^CH$; the merged small-residual band satisfies the same bound. At every ribbon point,
\[
 |\gamma'(z)-d(\gamma(z)+\tau,z)|
 \le |\alpha(z)|+\left|\int_0^\tau u(\gamma(z)+s,z)\,ds\right|
 \le P^C(h+w_0)\le P^CH.
\]
Its dimensions therefore satisfy all scalar and atlas hypotheses of the horizontal count, with $L\asymp_PH$, $W\asymp_Pw_0$. The full cell's boundary-perimeter calculation gives the count of the original assignment.

An original representative outside the central ribbon first follows its horizontal bridge of length $O_P(w_0)$ to the midline. This bridge lies in the full admissible set and changes the slope by at most $P^Cw_0$. The subsequent seed interval has $z$ length $O(H)$ and residual $O_P(H)$. Thus the proof of Lemma~\ref{low:6:so36:endpoints} gives the same ranges $|d|\lesssim_PH^2$, $|y|+|y-dz|\lesssim_PH^3$, and $|z|\lesssim_PH$. The original bounded relocation only changes their constants. Hence the same fixed-frame coefficient--point incidence and seed nets apply to every assigned anchor. Applying Theorem~\ref{low:6:so36:count} proves the assertion.
\end{proof}

The coordinate aspect in this proposition is $K_0=H/w_0$. The assigned geometric aspect $K=L_{\mathrm{geom}}/W_{\mathrm{geom}}$ may instead equal $h^2H/w_0$ up to powers of $P$. Since Theorem~\ref{low:6:og:common-cover} gives $W_{\mathrm{geom}}\ge1$, one has $K\lesssim_PS$. Dropping the coordinate factor $K_0^{-1}$ in Theorem~\ref{low:6:so36:count} and using $S^{-257/2}\le P^CS^{-513/4}K^{-1/4}$ gives the estimate in the assigned variables
\begin{equation}\label{low:6:as36:zlong-assigned}
\operatorname{vol}E_j\lesssim_P\nu S^{-513/4}K^{-1/4}.
\end{equation}
This conversion uses $K^{1/4}\lesssim_PS^{1/4}$.

\subsubsection{Fixed-width parabolic collars}

The small-width and high-residual pieces and projection caps of Theorem~\ref{low:6:og:common-cover} have transverse width a fixed power of $P$. Lemma~\ref{low:6:og:parabolic-collar} supplies algebraic parabolic certificates. We prove their fixed-frame count, including the external seed grid.

\par\smallskip
\noindent\textit{Proof stages for Theorem~\ref{low:6:sp36:tail}.}\par\nobreak
\noindent\textup{Step 1.} The actual parabolic ribbon.\par\nobreak
\noindent\textup{Step 2.} Preparation and fixed-frame endpoint ranges.\par\nobreak
\noindent\textup{Step 3.} Atlas validity and the fixed seed grid.\par\nobreak
\noindent\textup{Step 4.} Conditional integration and the final product.\par\nobreak
\smallskip

\paragraph{The parabolic ribbon}

Let an algebraic seed $\gamma(s)=(a(s),b(s))$ on the actual critical sheet be parameterized by $\theta(\gamma_s)=1$, where $\theta=dy-d\,dz$, and write $m=b'(s)$, with $|m|\le1$. Let $D(s)=d(\gamma(s))$ and $V(s)=v_0(\gamma(s))$. Then
\[
a'=1+D m,\qquad D'=u_0+Vm.
\]
For a prescribed fixed width $M_0=P^A$, use
\begin{equation}\label{low:6:sp36:ribbon}
\Psi(s,t)=\left(a(s)+D(s)t+\tfrac12V(s)t^2,\ b(s)+t\right),\qquad |s|<L,\quad |t|<M_0.
\end{equation}
Direct differentiation gives the exact tangential Jacobian
\begin{equation}\label{low:6:sp36:jacobian}
\det D\Psi=1+u_0(s)t+\tfrac12V'(s)t^2,\qquad V'=Yv_0+mZv_0.
\end{equation}
Assume the geometric construction supplies a contained full-activity ribbon of this form with polynomially bounded coframe and inverse. In the application Lemma~\ref{low:6:og:parabolic-collar} supplies it after $M_0$ is fixed and the inner defect is chosen. Retain $|v_0|\le P^C\nu$ and the full regular-null pivots. Its area is comparable to $L$, up to powers of $P$.

\begin{lemma}[Preparation and fixed-frame endpoint ranges]\label{low:6:sp36:preparation}

\noindent\textup{(i)} At a common set of evaluation points of area at least $P^{-C}L$, the scalar outputs satisfy
\[
|\partial_s^i\partial_t^j\phi|\le P^CL^{-i}\quad(j\ge1),\qquad |\partial_s^i\phi|\le P^CL^{1-i}\quad(i\ge1),\qquad |\partial_s^i\partial_t^j\lambda|\le P^CL^{-i}
\]
through any prescribed fixed orders, simultaneously with the seed-jet bounds below. In particular, the common set supports scalar preparation through total order twelve and unrounded seed preparation through order thirteen.

\noindent\textup{(ii)} In actual normalized endpoint coordinates, after centering the parameters at the evaluation point, the tangential germ has the form
\[
Y=A(s)+D_*(s)t+\tfrac12V_*(s)t^2,\qquad Z=B_*(s)+t.
\]
Here $A(0)=B_*(0)=0$, and its finite jets have ranges
\[
|A_1|\le P^C,\quad |A_j|\le P^CL^{2-j}\ (j\ge2),\quad |(B_*)_j|\le P^CL^{1-j}\ (j\ge1),
\]
\[
|(D_*)_0|\le P^C,\quad |(D_*)_j|\le P^CL^{1-j}\ (j\ge1),\qquad |(V_*)_j|\le P^CL^{-j}\ (j\ge0).
\]

\noindent\textup{(iii)} On the same set, the conditioned reciprocal of \eqref{low:6:sp36:jacobian} satisfies $|\partial_s^i\partial_t^j(\det D\Psi)^{-1}|\le P^CL^{-i}$ through any prescribed fixed finite order. The linear part is conditioned.

\noindent\textup{(iv)} Relative to a fixed frame with bounded relocation and a full-activity horizontal bridge of length at most $P^CL$, the possible $(\eta,f)$ positions have area at most $P^CL^3$ and the endpoint shear has range $P^CL$.
\end{lemma}
\begin{proof}
The coframe bounds and complete coefficient admissibility give bounded ordinary ribbon gradients of $\phi$ and bounded $\lambda$. In the original algebraic seed parameter $\zeta$, $\partial_s=\alpha(\zeta)^{-1}\partial_\zeta$ with nonvanishing algebraic $\alpha$, so every prescribed flux derivative remains algebraic in $(\zeta,t)$. Lemma~\partIref{geo:preparation} therefore gives the asserted scalar bounds, including the extra derivatives, on the actual rectangle of dimensions $L,M_0$. A horizontal bridge of length at most $P^CL$ changes the initial slope and horizontal coordinate by at most $P^CL$, since $|Yd|\le P^C$, without changing $z$. Along the seed, boundedness of $D'=u_0+Vm$ keeps $D$ in a range $P^CL$ relative to the fixed frame. Then $a'=1+Dm$ is bounded by $P^CL$, while $b'$ and $V$ are bounded. The one-dimensional part of the same lemma prepares $a,b,D,V$ simultaneously with the scalar bounds. The coframe and its inverse bound $(\det D\Psi)^{-1}$ by $P^C$. Include this algebraic reciprocal and the full scalar and unrounded seed lists in the same preparation. Translation of the parameter origin and the determinant-one endpoint shear preserve these bounds.

Use the centered functions $A,B_*,D_*,V_*$ defined in the proof of Lemma~\ref{low:78:s37g:graph-window}, with transverse shift $t_0$ and the actual endpoint null slope $d_0$. Here $|t_0|\le M_0$ and $|d_0|\le P^CL$, so their higher jets have the displayed ranges. The first entry $A_1$ is bounded by the actual coframe; $D_*(0)$ is the small transverse null-direction error supplied by the parabolic approximation. Thus the linear part is a small conditioned perturbation of $\left(\begin{smallmatrix}1&0\\m&1\end{smallmatrix}\right)$, with its actual endpoint entries retained.

Finally $|b|\le P^CL$, $|D|\le P^CL$, and integration of $a'=1+Dm$ gives $|a|\le P^CL^2$. Fixed transverse width adds polynomially bounded changes. Therefore $|f|\le P^CL$, $|y|\le P^CL^2$, $|d|\le P^CL$, and $|\eta|=|y-df|\le P^CL^2$. The candidate $(\eta,f)$ area is at most $P^CL^3$. These estimates include the bounded relocation from the original frame.
\end{proof}

\paragraph{Atlas validity and the seed grid}

The weight-triangular proof of Theorem~\ref{low:6:so36:exact-atlas} applies to every fixed tangential germ with invertible linear part. Its higher terms strictly raise tangential degree. Only its linear part and the original graph Hessian enter the determinant; its coefficient character is the explicit power $126$ of that linear tangential determinant. Here this determinant and its inverse are bounded. At the exact linear part $\left(\begin{smallmatrix}1&0\\m&1\end{smallmatrix}\right)$, the quadratic form is
\[
Q(s,t)=(r/2+Bm)s^2+Bst,
\]
so the grid determinant is exactly $56623104B^{73}\lambda^{-18}$ in absolute value. At the actual nearby linear part, the finite quadratic polynomial $G$ from Theorem~\ref{low:6:so36:exact-atlas} has the same conditioned lower bound. This uses only the small error in the linear part, not smallness of the complete higher sheared jet. The ranges of the parameter derivatives in Lemma~\ref{low:6:sp36:preparation} also bound all graph-coordinate ranges by $P^C$. Its full third-order subblock retains the original coordinate $c_{zzz}$, so replacing it by $v_0=-c_{zzz}/(2B)$ gives the factor $\nu$. The fixed nonlinear parameter substitution is an invertible triangular change on this full lower-order graph jet; it does not destroy that factor.

\begin{lemma}[A fixed relative seed grid]\label{low:6:sp36:rounding}

\noindent\textup{(i)} Six long derivatives of the four seed functions suffice for the $83$-output grid map. Use the following meshes in their prepared endpoint ranges:
\[
\begin{array}{c|c|c}
\text{parameter}&\text{mesh}&\text{number of choices}\\ \hline
A_1&c_PL^{-1}&P^CL\\
A_j,\ 2\le j\le6&c_PL^{-j}&P^CL^{10}\\
(B_*)_j,\ 1\le j\le6&c_PL^{-j-1}&P^CL^{12}\\
(D_*)_j,\ 1\le j\le6&c_PL^{-j}&P^CL^6\\
(D_*)_0,\ (V_*)_j,\ 0\le j\le6&c_P,\ c_PL^{-j}&P^C.
\end{array}
\]
The total cost is $P^CL^{29}$. These grids are fixed in advance in the relative endpoint parameter coordinates.

\noindent\textup{(ii)} They preserve every critical and second-normal output width. The width preservation holds for independent rounding of the seed functions on these grids.
\end{lemma}
\begin{proof}
The critical grid has longitudinal index at most six, as do the sixteen normal rows and eighteen graph rows. The tangential germ is quadratic in $t$, so these outputs use only the first six longitudinal derivatives of the seed functions, even when the total output order is twelve. The centered constants $A(0)=B_*(0)=0$ are fixed; the full range of $A_1$ is retained.

Use the finite-jet comparison in the proof of Lemma~\partIref{surf:collar-net} with seed order six and functions $A,B_*,D_*,V_*$. Lemma~\ref{low:6:sp36:preparation} supplies its endpoint ranges, scalar bounds through total order twelve, and auxiliary unrounded seed bounds through order thirteen. The same lemma supplies the reciprocal-Jacobian bounds. These are all the inputs to the finite-jet comparison.

The displayed meshes are precisely those of that proof. They give $P^CL^{5\cdot6-1}=P^CL^{29}$ and preserve the critical and second-normal widths. In particular, the finer $B_*$ mesh controls the product of the potentially large $D_*$ coefficient with the rounding error. The auxiliary higher seed jets remain unrounded during comparison and can then be discarded, since no output depends on them.
\end{proof}

\begin{theorem}[Fixed-width coefficient-volume bound at the fixed frame]\label{low:6:sp36:tail}
Suppose the parabolic ribbon \eqref{low:6:sp36:ribbon} lies in the outer admissible set of Definition~\partIref{def:admissibility}. Assume its coframe matrix and inverse have norms at most $P^C$ and $|v_0|\le P^C\nu$ throughout the ribbon. Retain the fixed-frame comparison hypotheses of Lemma~\ref{low:6:sp36:preparation}. Then its complete coefficient set at the original fixed frame has volume at most
\[
P^C\nu L^{-133}.
\]
\end{theorem}
\begin{proof}
Use the common prepared set of Lemma~\ref{low:6:sp36:preparation}. With transverse width a fixed power of $P$, the critical grid product is $P^CL^{-141}$ and the second-normal product is $P^CL^{-24}$. The coefficient atlas in Taylor coordinates, the bounded graph-coordinate ranges, the retained third-normal value, and the rate row therefore give raw coefficient product $P^C\nu L^{-165}$. Lemma~\ref{low:6:sp36:rounding} gives the seed-grid cardinality $P^CL^{29}$.

The endpoint candidate area is $P^CL^3$ and the area where the bounds hold is $P^{-C}L$, so their ratio is $P^CL^2$. Hold the original endpoint shear fixed, integrate it over its interval of length at most $P^CL$, and use the original $H_3$ Haar fiber from Lemma~\partIref{surf:null-density}. The product is $\nu L^{-165}L^{29}L^2L=\nu L^{-133}$. 
\end{proof}

The exponent $133$ suffices for the required incidence exponent approaching $128$.

The small-width pieces whose seeds satisfy $|\alpha|\le1$ have horizontal collars with bounded coframe and fixed transverse width. Theorem~\ref{low:6:so36:count} applies to those collars after bounded-length cases are placed in a cap. All other small-width pieces, high-residual pieces, and projection caps are covered by Theorem~\ref{low:6:sp36:tail}. Their assigned width is a fixed power of $P$, so $S\asymp_PL$ and $K\lesssim_PS$. In particular their bound $P^C\nu S^{-133}$ implies $P^C\nu S^{-513/4}K^{-1/4}$. All these estimates hold with the rate factor omitted on the identically zero branch.

\subsubsection{Wide corridors}

Put $w=w_0$ in the assignment of Theorem~\ref{low:6:og:common-cover}. A moderate wide component has coordinate aspect $K_0=w/H$ and $w>H\ge m\ge1$ and residual bounded by $P^Ch$, where $h\le P^Cw/m$. 

The geometric assignment supplies a contained algebraic ribbon
\begin{equation}\label{low:6:wc36:eq-ribbon}
 (y,z)=(\xi+\Gamma(\tau),\tau),\qquad |\xi|<c_Pw,\quad \tau\in I,\quad |I|\asymp_PH,
\end{equation}
inside the full regular-null admissible set, with $|\Gamma'-d(\Gamma,\tau)|\lesssim_P h$. Its actual area is comparable, by powers of $P$, to $wH$. Every original assigned anchor, after its allowed bounded right comparison, reaches the midline by a contained horizontal bridge of length $O_P(w)$, then travels along the same midline interval and horizontally to the ribbon. This assertion holds for every anchor, including those outside the retained central subribbon. The count of the entire assigned piece is $O_P(\Lambda S)$, where $S\asymp_PwH$. The number of pieces has only the polynomial and logarithmic factors in Theorem~\ref{low:6:og:common-cover}.

Use the same critical graph and scalar fields, with $\mu=F_{xxx}|_c$.
Normalized coefficient admissibility supplies
\begin{equation}\label{low:6:wc36:eq-activity}
 |Y\phi|+|Z\phi|+|r|+|u|+|v|+|\lambda|+|\mu|\lesssim_P1,
 \qquad |\lambda|,|B|\ge P^{-C}.
\end{equation}
It also bounds every fixed graph jet in the endpoint's original normalized frame. The buffered band has $|v|\lesssim_P\nu$, $0<\nu\le1$ after polynomial normalization. The identically zero branch and the inflection branch have already been assigned separately and are not included in this regular band estimate.

For each fixed relative seed, use the original $H_0$ coefficient--point incidence of Lemma~\partIref{surf:null-density}, as in Theorem~\ref{low:6:so36:count}, with endpoint variables $(\eta,f)$ and $\eta=y-df$. 

\begin{theorem}[The wide-corridor estimate]\label{low:6:wc36:thm-main}
Let $E_j(w,H,\nu)$ be the coefficient set for any one original assigned anchor in a moderate wide component satisfying the preceding inputs.

\noindent\textup{(i)} Then
\begin{equation}\label{low:6:wc36:eq-rawfinal}
 \vol E_j(w,H,\nu)\le P^C\nu w^{-129}H^{-161}.
\end{equation}

\noindent\textup{(ii)} Let $L,W$ be the original assigned geometric lengths, $S=LW\asymp_PwH$, and $K=L/W\ge1$. Then
\begin{equation}\label{low:6:wc36:eq-final}
 \vol E_j(w,H,\nu)\le P^C\nu S^{-257/2}K^{-1/2}.
\end{equation}
The constants are uniform in the residual $h$, including the subcase $hH\gg_Pw$. All external seed grids and every required change of the ordinary output map are included in this bound.

\noindent\textup{(iii)} If no rate-band restriction is imposed, the same conclusions hold with $\nu$ replaced by $1$; in particular they apply to the separately assigned identically zero-rate branch.
\end{theorem}

\par\smallskip
\noindent\textit{Proof stages for Theorem~\ref{low:6:wc36:thm-main}.}\par\nobreak
\noindent\textup{Step 1.} A common ribbon, amplitudes and endpoint bounds.\par\nobreak
\noindent\textup{Step 2.} The ordinary atlas and joint graph window.\par\nobreak
\noindent\textup{Step 3.} Fixed seed grids and acceptance after slope selection.\par\nobreak
\noindent\textup{Step 4.} The original-frame point factor.\par\nobreak
\noindent\textup{Step 5.} Raw product and assigned-aspect conversion.\par\nobreak
\smallskip

\paragraph{The ribbon and coefficient map}

Use the ordinary rectangle~\eqref{low:6:wc36:eq-ribbon} and the complete regular-root atlas, including its third-normal output.

Define
\begin{equation}\label{low:6:wc36:eq-D0}
D_0=w+hH\lesssim_P wH,
\end{equation}
using $H\ge1$ and $h\lesssim_P w/m\lesssim_P w$. For the coefficient set in Theorem~\ref{low:6:wc36:thm-main}, fix once and for all the larger bound $D_*=P^CwH\ge D_0$. Its actual seed grids and candidate point boxes use $D_*$ and are independent of $h$. At an arbitrary ribbon point put $e=\Gamma'-d(\xi+\Gamma,\tau)$. The horizontal variation identity and \eqref{low:6:wc36:eq-activity} give $|e|\lesssim_P h+w\lesssim_P w$. Choose the uniform amplitude $E=P^Cw$ large enough that $|Z\phi+eY\phi|\le E$ everywhere in the ribbon.

Apply Lemma~\partIref{geo:preparation} on the actual rectangle of dimensions $w,H$ to the finite list of bounded algebraic scalars used below, with their stated amplitudes. Their selected derivatives have fixed format. Choose the exceptional proportions small enough that the common good set has area at least $P^{-C}wH$, and include every extra mixed derivative needed by the later finite chain rules. Each derivative then has its amplitude times $w^{-i}H^{-j}$ as its bound, up to the permitted polynomial constant.

\begin{lemma}[Scalar widths, graph curvature, and endpoint ranges]\label{low:6:wc36:lem-preparation}

\noindent\textup{(i)} There is a common set of endpoints of area at least $P^{-C}wH$ at which the following inequalities hold to every required fixed order:
\begin{align}
 |\partial_\xi^i\partial_\tau^j\phi|&\lesssim_P w^{1-i}H^{-j} &&(i\ge1),\label{low:6:wc36:eq-mixedwidth}\\
 |\partial_\tau^j\phi|&\lesssim_P E H^{1-j} &&(j\ge1),\label{low:6:wc36:eq-purewidth}\\
 |\partial_\xi^i\partial_\tau^j\lambda|+|\partial_\xi^i\partial_\tau^j r|&\lesssim_P w^{-i}H^{-j}.\label{low:6:wc36:eq-boundedwidth}
\end{align}

\noindent\textup{(ii)} Relative to every original assigned anchor, the endpoint ranges satisfy
\begin{equation}\label{low:6:wc36:eq-endpoint}
 |d|\lesssim_P D_0,\qquad |f|\lesssim_P H,\qquad |y|+|\eta|\lesssim_P D_0H.
\end{equation}

\noindent\textup{(iii)} The relative seed's first derivative in the normalized endpoint frame is $O_P(w)$; its derivatives of orders $j\ge2$ can simultaneously be prepared to be $O_P(D_0H^{1-j})$.
\end{lemma}
\begin{proof}
The actual identities are $\partial_\xi\phi=Y\phi$ and $\partial_\tau\phi=Z\phi+eY\phi$. Prepare these two algebraic scalars, with respective amplitudes $P^C$ and $E$, and also the bounded scalars $\lambda$ and $r$. Differentiating the first identity gives \eqref{low:6:wc36:eq-mixedwidth}; differentiating the second in the short direction gives \eqref{low:6:wc36:eq-purewidth}. In particular the larger amplitude $E$ occurs only in the six purely transverse critical rows used below. Preparing $r$ gives \eqref{low:6:wc36:eq-boundedwidth}; this additional preparation supplies the graph factor used in this orientation.

Normalize the fixed frame's selected slope to zero. Its bounded right comparison changes only powers of $P$. The initial horizontal bridge changes $d$ by at most $P^Cw$. Along the midline,
\[
 \frac{d}{d\tau}d(\Gamma(\tau),\tau)=v+(\Gamma'-d)u,
\]
so its variation over the interval is at most $P^C(H+hH)$. Since $H<w$, this and the final horizontal bridge give $|d|\lesssim_P D_0$. Hence $|\Gamma'|\lesssim_P D_0$, and integration, together with the horizontal bridges, gives $|y|\lesssim_P D_0H$. The bound for $\eta=y-df$ follows immediately. These paths start from each fixed frame.

Prepare the one-dimensional algebraic scalar $\Gamma'$ on the same interval, discarding another small proportion. It yields $|\Gamma^{(j)}|\lesssim_P D_0H^{1-j}$ for $j\ge2$. Passing to the normalized endpoint subtracts its actual slope from the first derivative, leaving $e$, already bounded by $P^Cw$. Higher seed derivatives are unaffected by this linear normalization. Intersecting all the good sets proves the lemma.
\end{proof}

At each prepared endpoint use the same normalized graph, $c(0)=c_y(0)=c_z(0)=c_{zz}(0)=0$ and $d(0)=0$. The quadratic graph is $\tfrac12ry^2+Byz$. Fix an external short-direction seed germ $\gamma(t)$ and use $(y,z)=(s+\gamma(t),t)$. The first index in this wide-certificate calculation is the long horizontal index. Use the grid sets $\mathcal I_d^{\rm g}$ and $\mathcal E^{\rm g}$ from Theorem~\ref{low:6:so36:exact-atlas}: there are $48$ nonconstant critical-value jets, $16$ second-normal jets, $18$ graph jets, and the third-normal value $\mu(0)$. The four original graph normalization rows are restored by a triangular coordinate change before passing to the $79$-dimensional normalized slice.

Denote this full $83$-output coefficient map, before deleting the normalization rows, by $\Phi_{\mathrm{coef}}$. The exact ordinary-atlas theorem is invariant under exchanging the two tangential variables. With the present first-index convention it gives
\begin{equation}\label{low:6:wc36:atlas}
 |\det D\Phi_{\mathrm{coef}}|=|\lambda(0)|^{-18}|G(A,b,C)|,
 \qquad A=\frac r2,\quad b=B+re,\quad C=Be+\frac r2e^2,
\end{equation}
where $e=\gamma'(0)$. The polynomial identity is
\begin{equation}\label{low:6:wc36:G}
G(A,b,C)=b^{73}g(AC/b^2),\qquad \deg g\le36.
\end{equation}
Its meaning at $b=0$ is polynomial continuation. The weight-triangular proof in Theorem~\ref{low:6:so36:exact-atlas} shows that this Jacobian is independent of every higher graph jet and every higher external seed jet. No boundedness of the full phase after a large shear is used. The complete polynomial $g=g_6$ and its slope-root multiplicities are calculated in Section~\partIref{atlas:slopes}.

Apply the degree-six factorization and finite slope estimate in Section~\partIref{atlas:slopes}, with $g=g_6$ and the same Taylor-coordinate convention.
The slope changes have size at most $a_PK_0$, so in the anisotropic variables $(\xi/w,\tau/H)$ their size is at most $a_P$. They therefore preserve the prepared scalar widths and the prepared curvature widths after a fixed enlargement. The final slope is rounded first, and the finite perturbation is chosen after that rounding; no subsequent rounding is used that could remove the determinant lower bound. This is a finite union of fixed external coefficient maps. It neither changes the phase by a coefficient-dependent substitution nor adds an endpoint shear integral.

\paragraph{Joint graph derivative bounds on the longitudinal grid}\label{low:6:wc36:graph-window}

At a prepared endpoint let $\gamma$ first denote the unrounded relative seed germ. Write $c_\gamma(\xi,\tau)=c(\xi+\gamma(\tau),\tau)$ and $r_\gamma(\xi,\tau)=c_{yy}(\xi+\gamma(\tau),\tau)$. For every $i\ge2$,
\[
\partial_\xi^i\partial_\tau^j c_\gamma=\partial_\xi^{i-2}\partial_\tau^j r_\gamma.
\]
The preparation of $r$ on the actual ribbon in Lemma~\ref{low:6:wc36:lem-preparation} gives $|\partial_\xi^i\partial_\tau^jc_\gamma|\le P^Cw^{2-i}H^{-j}$ at each retained graph index with $i\ge2$. Among the complete graph indices of total order at most four, retain the original endpoint-normalized unsheared graph coordinates with $i\le1$, and use the prepared curved coordinates with $i\ge2$. This is a triangular coordinate change of determinant one. Indeed the linear shear substitutes only indices $(i+k,j-k)$ with larger horizontal index, while nonlinear seed terms decrease total derivative order. The normalizing rows remain original coordinates, and the original $c_{zzz}$ remains among the unsheared coordinates. Its replacement by $v=-c_{zzz}/(2B)$ supplies the rate factor $\nu$.

The three additional strong-grid graph rows $(3,2),(2,3),(3,3)$ all satisfy $i\ge2$ and are bounded directly by the prepared scalar identity. No expansion into missing higher unsheared graph indices is required. The controlled nonconstant graph widths are
\[
\begin{array}{c|ccccccccc}
(i,j)&(2,0)&(2,1)&(3,0)&(2,2)&(3,1)&(4,0)&(3,2)&(2,3)&(3,3)\\\hline
\text{width}&1&H^{-1}&w^{-1}&H^{-2}&w^{-1}H^{-1}&w^{-2}&w^{-1}H^{-2}&H^{-3}&w^{-1}H^{-3}.
\end{array}
\]
The five complementary free coordinates have indices $(1,1),(0,3),(1,2),(0,4),(1,3)$ in the original frame. Their joint window has volume at most $P^C\nu$: condition on the other graph coordinates and use $c_{zzz}=-2Bv$, with $|B|\asymp_P1$, for the remaining interval. Without a rate restriction the same volume is at most $P^C$. The finite rounding and slope changes preserve the prepared graph bounds at these same endpoints by Lemma~\ref{low:6:wc36:lem-nets}; they do not require a new ribbon for a perturbed germ. Therefore the joint range of all free graph coordinates has volume at most $P^C\nu w^{-6}H^{-12}$. Combining this with the finite slope perturbation costs at most $H^3$, giving a graph coordinate volume divided by the Jacobian at most $P^C\nu w^{-6}H^{-9}$, uniformly in the relative slope and the curvature $r$.

\begin{lemma}[Six seed grids, followed by finitely many slope choices]\label{low:6:wc36:lem-nets}

\noindent\textup{(i)} With the fixed bound $D_*$, all required short-seed jets can be covered by fixed grids of total size
\begin{equation}\label{low:6:wc36:eq-net}
N_{\rm seed}\lesssim_P H\left(\frac{D_*H}{w}\right)^5\lesssim_P H^{11}.
\end{equation}

\noindent\textup{(ii)} After rounding, the finite slope choices just constructed preserve every prepared scalar and curvature width and cost only a constant number of additional maps.

\noindent\textup{(iii)} Every endpoint in the original common good set is accepted by at least one of these maps, with determinant at least $P^{-C}H^{-3}$.
\end{lemma}
\begin{proof}
At short order $j$, use mesh $c_Pw/H^j$. The first relative derivative has range $P^Cw$, so costs $P^CH$ choices. Each order $2\le j\le6$ has range $P^CD_*H^{1-j}$ and costs at most $P^CD_*H/w$ choices. The quantity $D_*H/w$ is at least one. Although the largest scalar total order is twelve, its short index is at most six. Since the seed depends only on the short variable, its derivatives through order six suffice. This proves \eqref{low:6:wc36:eq-net}.

Fix the grids in normalized endpoint coordinates. The ranges depend only on $P,w,H$. Rounding gives $s\mapsto s+\delta\gamma(t)$, with $\delta\gamma(0)=0$ and $|\delta\gamma^{(j)}|H^j/w\le c_P$. Prepare all extra scalar derivatives to sufficiently high fixed order first. The anisotropic chain rule preserves bounded-scalar widths $w^{-i}H^{-j}$ and mixed critical widths $w^{1-i}H^{-j}$, since a mixed row retains a horizontal derivative. Purely transverse rows acquire terms at most $P^CwH^{-j}$, absorbed by $EH^{1-j}$ since $E=P^Cw$ and $H\ge1$.

The same argument applies to $r$ and to the finite slope changes $\delta e=t_0w/H$: their scaled size is $|\delta e|H/w=|t_0|\le a_P$. Choose a rounded jet first; then use one of the finitely many slope offsets, with no subsequent rounding. The finite slope selection lemma applies because the rounded first slope is still $O_P(w)$. The controlled graph jets are derivatives of $r$, so their bounds persist under these changes. Only the output germ changes; the evaluation point is fixed.

Let $G(F)$ be the common endpoint set in Lemma~\ref{low:6:wc36:lem-preparation}, retaining also the extra derivative bounds used in the preceding chain-rule comparison. Let $G_\omega(F)\subseteq G(F)$ be the endpoints accepted by the fixed seed-grid and slope-offset map $\omega$. At least one map accepts each point of $G(F)$, so
\[
\sum_\omega\int_{G(F)}\mathbf1_{G_\omega(F)}(x)\,dx\ge|G(F)|\ge P^{-C}wH.
\]
Thus the original $P^{-C}wH$ incidence denominator is retained, with the grid cardinality in the numerator.
\end{proof}

\begin{lemma}[The anchored point factor]\label{low:6:wc36:lem-point}
For each fixed rounded external jet and one of its fixed slope offsets, the fixed-frame incidence contributes at most
\begin{equation}\label{low:6:wc36:eq-point}
P^C\frac{D_*H^2}{wH}=P^C\frac{D_*H}{w}\lesssim_P H^2.
\end{equation}
\end{lemma}
\begin{proof}
For each fixed rounded seed and slope offset, the ordinary coefficient map at the selected slope depends only on the endpoint phase and has bounded inverse multiplicity by Lemma~\partIref{geo:coefficient-map}. Apply the original $H_0$ incidence of Lemma~\partIref{surf:null-density}. Equation~\eqref{low:6:wc36:eq-endpoint} gives candidate area $P^CD_*H^2$, and Lemma~\ref{low:6:wc36:lem-preparation} gives good area $P^{-C}wH$. Their ratio proves~\eqref{low:6:wc36:eq-point}; the full relative first-seed range is included in~\eqref{low:6:wc36:eq-net}.
\end{proof}

\begin{proof}[Proof of Theorem~\ref{low:6:wc36:thm-main}]
In the grid $\mathcal I_6^{\rm g}\setminus\{0\}$, there are $42$ rows with positive horizontal index and six purely transverse rows. The horizontal index sum is $147$; the short index sum is also $147$. Thus \eqref{low:6:wc36:eq-mixedwidth}--\eqref{low:6:wc36:eq-purewidth} give the critical width product
\[
P^CE^6w^{42-147}H^{6-147}=P^CE^6w^{-105}H^{-141}.
\]
The sixteen second-normal rows have index sums $24,24$, contributing $P^Cw^{-24}H^{-24}$. The combined scalar product is therefore $P^CE^6w^{-129}H^{-165}$. The retained third-normal value has range $P^C$. The prepared graph output window contributes $P^C\nu w^{-6}H^{-12}$; the inverse determinant at the selected slope costs at most $P^CH^3$, including the conditioned power of $\lambda$.

Use $E=P^Cw$ and the fixed $D_*=P^CwH$ in the seed and endpoint bounds. Equations~\eqref{low:6:wc36:eq-net} and~\eqref{low:6:wc36:eq-point} give the seed cost $P^CH^{11}$ and point factor $P^CH^2$. Combining the scalar, graph, inverse-determinant, seed and point factors gives
\[
\operatorname{vol}E_j\lesssim_P\nu(w^6w^{-129}H^{-165})(w^{-6}H^{-12})H^3H^{11}H^2=\nu w^{-129}H^{-161}.
\]
This proves \eqref{low:6:wc36:eq-rawfinal} for each fixed frame. Table~\ref{app:table-wide-factors} collects these factors.

Express the bound in the aspect used by the assignment. The raw estimate implies
\[
w^{-129}H^{-161}=(wH)^{-129}H^{-32}\le(wH)^{-129}\lesssim_P S^{-129}.
\]
The original assigned short length satisfies $W\ge P^{-C}$, hence $K=S/W^2\le P^CS$. It follows that
\[
S^{-129}\le P^C S^{-257/2}K^{-1/2}.
\]
This proves \eqref{low:6:wc36:eq-final} for the original geometric aspect, whether the original certificate was a horizontally long characteristic rectangle or an oblique tube with $L\asymp_PhH$ and $W\asymp_Pw/h$.
\end{proof}

\begin{corollary}[Summation of moderate wide pieces]\label{low:6:wc36:summation}
The moderate wide pieces of Theorem~\ref{low:6:og:common-cover}, with their assigned count $O_P(\Lambda S)$, satisfy the same rate, area, and aspect summation as the other counted regular pieces.
\end{corollary}
\begin{proof}
Apply the proof of Proposition~\partIref{low:6:r3:weighted-holder} with $b=257/2$, $\tau=1/2$ and $0<a<1$. The assigned bound $K\lesssim_PS$ changes only its permitted conditioning and logarithmic factors. For $\eta=1/2048$ and $129<p<1+(513/4)/(1+\eta)$, one has $(1+\eta)(p-1)<513/4<257/2$. Thus the rate, area and aspect sums converge with the required common exponent and absorb the physical logarithms. Any dyadic organization by $H$ at fixed area adds only a logarithm.
\end{proof}

\subsubsection{Completion of the critical branch}

\begin{theorem}[Complete critical null assignment]\label{low:6:as36:all-surface}
Fix the conditioned critical null branch near the regular-root rank-four family and the original physical family. Choose the prescribed model-box radius before its inner coefficient defect.

\noindent\textup{(i)} All selected inner labels admit an assignment Borel measurable in the phase coefficients, chosen for all original labels before one is fixed for coefficient integration, to pieces with dyadic geometric scales and counts satisfying
\[
L\ge W\ge1,\qquad S=LW,\qquad K=L/W,\qquad \text{assigned-label count}\le P^C\Lambda S.
\]
For each positive dyadic rate $\nu\le P^C$,
\[
\#\{\text{pieces at rate }\nu\}\le P^C(1+\log T)^C.
\]
The identically zero branch has the same piece-count bound as one separate type.

\noindent\textup{(ii)} Each original assigned anchor $j$ satisfies the complete coefficient estimate
\begin{equation}\label{low:6:as36:uniform-tail}
\operatorname{vol}E_j(\nu,S,K)\le P^C\nu S^{-513/4}K^{-1/4}
\end{equation}
on the positive-rate pieces.

\noindent\textup{(iii)} On the identically zero branch the same estimate holds with the factor $\nu$ omitted.
\end{theorem}
\begin{proof}
Theorem~\ref{low:6:r1:relocation} supplies either a buffered representative in a positive-rate band or an identically zero connected branch. Its polynomial right displacement retains the fixed frames and increases local packing only by a power of $P$. The fixed-format core in Section~\ref{low:6:og:cover} is defined using the full model-box test, so it contains these representatives and has the required outer buffer. The single-band transport estimate gives the Jacobian hypothesis on positive bands. Corollary~\ref{low:6:r1:zero-common} supplies that same hypothesis on the identically zero branch. Theorem~\ref{low:6:og:common-cover} now gives the common labels, count, finite multiplicity, and anchor-relative paths.

For a moderate corridor whose $z$ length $H$ is at least its horizontal width $w_0$, Proposition~\ref{low:6:oi:z-long} supplies the ordinary horizontal certificate. Theorem~\ref{low:6:so36:count} gives its bound with coordinate aspect $K_0=H/w_0$. Equation~\eqref{low:6:as36:zlong-assigned} converts it to \eqref{low:6:as36:uniform-tail} in the original assigned aspect $K$.

For a moderate corridor with $w_0>H$, Theorem~\ref{low:6:wc36:thm-main} gives $P^C\nu S^{-257/2}K^{-1/2}$ in the original assigned variables. Since $S,K\ge1$, this is bounded by the right side of \eqref{low:6:as36:uniform-tail}. This case includes both horizontally long characteristic rectangles and the oblique subcase whose geometric length is comparable to $hH$.

The remaining pieces are the small-width or high-residual collars and projection caps. Their algebraic certificates are supplied by Lemma~\ref{low:6:og:parabolic-collar} and its horizontal counterpart. Theorem~\ref{low:6:sp36:tail} gives $P^C\nu S^{-133}$ for the parabolic certificates, while Theorem~\ref{low:6:so36:count} covers the horizontally ruled ones. Their fixed polynomial transverse width implies $K\lesssim_PS$, so these estimates again imply \eqref{low:6:as36:uniform-tail}. Bounded-length caps have all scales bounded by powers of $P$ and the same bounded coefficient map with its rate row gives the stated estimate after increasing $P^C$.

Each coefficient set permits any certificate with the stated ordinary dimensions and bounds. Its seed grids and endpoint boxes are uniform over residual components. For assigned $(S,K)$, let $S_0=Hw_0\asymp_PS$, with polynomially many dyadic choices. Summing the vertical estimate $P^C\nu S_0^{-257/2}K_0^{-1}$ over dyadic $K_0=H/w_0\ge1$ gives $P^C\nu S^{-257/2}$. The wide estimate~\eqref{low:6:wc36:eq-rawfinal}, $P^C\nu S_0^{-129}H^{-32}$, sums over dyadic $H\ge1$ to $P^C\nu S^{-129}$. Its grids use $D_*=P^Cw_0H$, so no residual-$h$ parameter enters the union. Fixed-width certificates have $L\asymp_PS$ and polynomially many dyadic length choices. The aspect conversions therefore give~\eqref{low:6:as36:uniform-tail} for the full coefficient set $E_j(\nu,S,K)$.

Every coefficient count uses the fixed-frame ranges supplied by the common construction and the original conditional Haar incidence. In particular, neither the retained central subribbon nor the slope choice in the wide count changes the assigned anchors. On the identically zero branch each graph rate coordinate is integrated over its full bounded interval. The proofs of all three counts then give exactly the stated estimates without $\nu$. This exhausts the pieces in the cover.
\end{proof}

Section~\ref{an:assembly} combines the sextic branches, using Proposition~\partIref{low:6:r3:weighted-holder} for positive critical bands and treating the identically zero branch as one unbanded type.

The sixteen regular row lists appear in Section~\ref{ext-supp:low:6:s6r:rows}.


\section{Completion of the proof}\label{an:assembly}\label{an:assembly-section}

\subsection{Coefficient incidence}\label{completion:incidence}\label{an:coefficient-assembly}


\paragraph{From coefficient estimates to the physical incidence bound}

For a fixed degree $k$, exponent $p>1$, and $\gamma\ge0$, use the physical incidence input \partIeqref{an:physical-incidence} in each coarse shape $D$, for arbitrary fixed $b_j\ge0$. Its annulus is $\mathcal A_R=\{F:R\le\|F\|\le2R\}$, and its labels, original weights, positive reconstruction and permitted admissibility enlargement are those of Section~\partIref{completion:incidence}. We prove this degree-independent input below with the losses and degree ranges in Table~\ref{an:incidence-inputs}.

{\setlength{\LTcapwidth}{\textwidth}
\begin{longtable}{ccl}
\caption{Surface inputs and incidence ranges for degrees four through thirteen. The tails omit conditioning and permitted logarithmic factors. The regular exponents $b_k$ are recorded in Section~\ref{reg:combined}. The last column specifies $\chi_k+3<p<\chi_k+3+\varepsilon$.}\label{an:incidence-inputs}\\
\toprule
$k$ & Complete surface tail & $\varepsilon$\\
\midrule
$4$ & $S^{-38}K^{-3/2}$ & $1$\\
$5$ & $S^{-145/2}K^{-1/4}$ & $1/2$\\
$6$ & $\nu S^{-513/4}K^{-1/4}$ & $1/16$\\
$7$ & $S^{-429/2}K^{-1/2}$ & $1/16$\\
$8$ & $S^{-673/2}K^{-1/2}$ & $1/16$\\
$9$ & $S^{-1009/2}K^{-1/2}$ & $1/4$\\
$10\le k\le13$ & $S^{-\chi_k-7/2}K^{-1/2}$ & $1/4$\\
\bottomrule
\end{longtable}}

In the sextic row the rate factor belongs to the critical positive bands; their identically zero branch has the same tail with $\nu$ omitted. The inflection and nonnull branches have stronger unbanded tails.

\begin{proposition}[Coefficient incidence]\label{an:unified-incidence}\label{low:4:ga:incidence}\label{low:quartic-incidence}\label{low:5:qga:actual-incidence}\label{low:quintic-incidence}\label{low:6:s36:incidence}\label{low:sextic-incidence}\label{low:78:s78:incidence}\label{low:seven-eight-incidence}\label{low:9:root9:incidence}\label{low:nine-incidence}\label{low:incidence-ten-thirteen}
Fix $4\le k\le13$ and $p$ in the corresponding range of Table~\ref{an:incidence-inputs}. Choose $\eta>0$ so that $q=(1+\eta)(p-1)$ is smaller than the surface area exponent and, when $k\ge5$, the regular exponent $b_k$. Put
\[
\gamma=\begin{cases}q-36,&k=4,\\0,&k\ge5.\end{cases}
\]
Let $\{T_j\}$ be a fixed family in one physical dyadic shape family, with packing multiplicity $\Lambda$, and let $\mathcal A_j(F)$ be its complete stopped admissibility condition, including the physical restrictions and the common assignments proved above. Suppose its coefficient windows and reciprocal pivots are bounded by powers of $P$.

\noindent\textup{(i)} For fixed $b_j\ge0$,
\begin{equation}\label{an:unified-window-incidence}
\int_{V_k}\left(\sum_jb_jv_j\mathbf1_{\mathcal A_j(F)}\right)^p\,dF\lesssim_{p,k} P^{A_{p,k}}\Lambda^{p-1}(1+\log P+\log d_1^{-1})^{B_{p,k}}d_2^{-\gamma}\sum_jb_j^pv_j^{p-\chi_k}.
\end{equation}
\equationalias{low:4:ga:estimate}\equationalias{low:5:qga:actual-bound}\equationalias{low:finite-incidence-equation}

\noindent\textup{(ii)} For the physical family and low-linear selection, this gives~\partIeqref{an:physical-incidence}.

\noindent\textup{(iii)} The same bounds hold for any measurable selected subcollection of the admissibility condition.
\end{proposition}
\begin{proof}
Work in a fixed RQ bin. Lemma~\partIref{geo:normal-pivot} excludes rank below four, and Lemma~\partIref{reg:regular-preparation} assigns the quantitative full-rank region to bounded pieces and the regular region to prepared pieces. The remaining labels have the exhaustive surface assignments: Theorems~\ref{low:4:sf:surface-completion} and~\ref{low:5:qsc:surface-count} for $k=4,5$; Proposition~\ref{low:6:s36a:assignment}, Proposition~\ref{low:6:s36i:nonnull-count}, and Theorems~\ref{low:6:s36i:null-count} and~\ref{low:6:as36:all-surface} for $k=6$; Proposition~\ref{low:78:s78:surface} and~\eqref{low:9:root9:common-surface} for $k=7,8,9$; and Theorem~\partIref{surf:common-tail} for $10\le k\le13$. Their inner and outer windows and priority choices precede coefficient anchoring.

On unbanded pieces set $s_I=1$ on bounded pieces, $s_I=H_I$ on regular curve pieces, and $s_I=S_I$ on surface pieces. The common assignments give count $P^C\Lambda s_I$ and
\[
\sum_I s_I^{-\eta}\le P^C(1+\log P+\log d_1^{-1})^C.
\]
For $k=4,5$, use the bounded number of pieces at each width $W\ge P^{-C}$ and $S\ge W^2$; the other unbanded surface assignments have the stated area multiplicity and logarithmically many aspects. Lemma~\partIref{reg:regular-preparation} supplies regular multiplicity. The resulting dyadic sums converge, with subunit scales costing powers of $P$.

At an original fixed anchor, write the unbanded surface tails in Table~\ref{an:incidence-inputs} as $S^{-\sigma}K^{-\tau}$. Their dyadic sums satisfy Proposition~\partIref{an:common-incidence} since $q<\sigma$ and $\tau>0$; the regular sums for $5\le k\le13$ converge since $q<b_k$. Bounded coefficient sets have complete window volume $P^C$.

In degree four, Lemma~\partIref{reg:middle-cutoff} supplies the cutoff $H\le P^C/d_2$. Its regular sum is
\begin{equation}\label{an:quartic-paid-sum}
\sum_{1\le H\le P^C/d_2}H^{q-36}\lesssim_{p,\eta} P^Cd_2^{-(q-36)}.
\end{equation}
The surface sum still converges because $q<38$. Since $d_2\le1$, the bounded and surface contributions satisfy the same bound with the same factor $d_2^{-(q-36)}$.

For the sextic critical positive bands, apply Proposition~\partIref{low:6:r3:weighted-holder} with $a=1/2$. The complete assignment of Theorem~\ref{low:6:as36:all-surface} gives the convergent sum
\[
\sum_{\nu,S,K}\nu^{1/2}S^{q-513/4}K^{-1/4}<\infty.
\]
Its identically zero branch is one unbanded type. The inflection and nonnull tails are stronger, and their auxiliary amplitude sums have already been included in their complete estimates. Thus the band proposition and Proposition~\partIref{an:common-incidence}, followed by finite H\"older over the piece types, prove~\eqref{an:unified-window-incidence}. At each anchor the fixed coefficient substitution contributes $v_j^{-\chi_k}$; combined with $(b_jv_j)^p$, this gives exactly $b_j^pv_j^{p-\chi_k}$.

For the physical selection, $P\lesssim K_{\mathrm{lin}}^CL_R^C$, $\Lambda\lesssim L_R^3$, and $d_1\gtrsim R^{-1}$. Use only the degree-independent conditioning, logarithm and RQ-bin conversion at the end of the proof of Proposition~\partIref{an:unified-incidence} to obtain~\partIeqref{an:physical-incidence}, retaining $d_2^{-\gamma}$. Positivity gives the assertion for selected subcollections.
\end{proof}

The certificate conventions, arithmetic procedures and fixed row data are specified in supplementary Sections~\ref{ext-supp:low:verification-method}--\ref{ext-supp:low:row-masks}. The quadratic and cubic arguments retain their separate conclusions from Proposition~\ref{root:quadratic-upper} and Theorem~\ref{low:3:thm:main}.

\subsection{Physical summation and completion}\label{completion:physical-II}

Apply Proposition~\partIref{an:physical-assembly} to the incidence estimate just proved. Its hypotheses allow the retained side loss $\gamma$.

For the quartic loss~\eqref{an:quartic-paid-sum}, the three exponents in~\partIeqref{an:physical-positivity} are
\[
\mu_1=p-38,\qquad \mu_2=p-36-(q-36)=1-\eta(p-1),\qquad \mu_3=p-34.
\]
They are all positive by the choices already made. In every degree with $\gamma=0$, positivity follows directly from $p>\chi_k+3$. Thus the concrete incidence estimates establish the smooth-amplitude upper bound on an interval immediately above the asserted threshold in each degree. For each $p>\chi_k+3$, choose $p_0<p$ in that interval and apply the pointwise bound \partIeqref{an:larger-p}.

\begin{proof}[Proof of Theorem~\ref{root:main}]\label{an:main-proof}
The smooth upper bounds for $k=2$ and $k=3$ are Proposition~\ref{root:quadratic-upper} and Theorem~\ref{low:cubic}. For $4\le k\le13$, Proposition~\ref{an:unified-incidence} and Proposition~\partIref{an:physical-assembly} give the smooth upper bound; the quartic middle-side loss is summed above. The pointwise estimate \partIeqref{an:larger-p} extends these bounds to all larger exponents, and Lemma~\partIref{an:cube-reduction} gives the cube upper bound. Proposition~\partIref{an:necessity} proves divergence for $p\le\chi_k+3$, including equality.
\end{proof}
\appendix
\section{Reference tables}\label{app:reference-tables}
\renewcommand{\thetable}{A.\arabic{table}}
\renewcommand{\theHtable}{A.\arabic{table}}
\setcounter{table}{0}


\begin{table}[H]
\centering\small
\caption{Dimensions, characters and physical exponents for degrees four through six, from Lemmas~\partIref{an:character} and~\partIref{an:resource}. The exponents precede any incidence side loss.}\label{app:table-physical}
\[
\begin{array}{c|rr|ccc}
k&N_k&\chi_k&\multicolumn{3}{c}{\text{full-catalogue physical exponents}}\\\hline
4&34&35&p-38&p-36&p-34\\
5&55&70&p-73&p-71&p-69\\
6&83&126&p-129&p-127&p-125
\end{array}
\]
\end{table}

\begin{table}[H]
\centering\small
\caption{Dimensions, characters and initial incidence ranges in degrees seven and eight; see Proposition~\ref{an:unified-incidence}.}\label{app:table-seven-eight}
\[
\begin{array}{c|rr|c}
k&N_k&\chi_k&\text{initial }p\text{ interval}\\\hline
7&119&210&213<p<213+1/16\\
8&164&330&333<p<333+1/16
\end{array}
\]
\end{table}

\begin{table}[H]
\centering\small
\caption{Gauss row counts and fixed-normal powers from Lemmas~\partIref{surf:normal-filtration} and~\partIref{surf:gauss-atlas}.}\label{app:table-gauss-rows}
\[
\begin{array}{c|c|c|c}
k&\text{full row count}&N_k-4&\tau_k\\\hline
7&64+36+16+4-1=119&115&130\\
8&81+49+25+9+1-1=164&160&200\\
9&100+64+36+16+4-1=219&215&295.
\end{array}
\]
\end{table}

\begin{table}[H]
\centering\small
\caption{Raw scalar exponents in the proof of Proposition~\ref{low:9:u9:gauss-tails}, obtained from Proposition~\partIref{surf:gauss-tail}.}\label{app:table-gauss-pairs}
\[
\begin{array}{c|c|rr|rr}
k&m&a_{\mathrm g}&b_{\mathrm g}&a_{\mathrm w}&b_{\mathrm w}\\\hline
7&2&243&194&331&106\\
7&3&309&260&422&147\\
8&2&375&311&511&175\\
8&3&472&408&644&236\\
8&4&513&449&701&261\\
9&2&557&476&759&274\\
9&3&691&610&942&359\\
9&4&757&676&1033&400.
\end{array}
\]
\end{table}

\begin{table}[H]
\centering\small
\caption{Ordinary-surface dimensions and fixed-center normal-frame bounds in Section~\ref{quartic:normalization}.}\label{app:table-normal-frames}
\[
\begin{array}{c|ccc}
k&4&5&6\\ \hline
\deg F_x\text{ at most}&3&4&5\\
\text{normal frames at a center, at most}&9&12&15\\
\text{coefficient slice dimension }N_k-4&30&51&79
\end{array}
\]
\end{table}

\begin{table}[H]
\centering\small
\caption{Factors in the proof of Theorem~\ref{low:6:so36:count}, with $S=LW$ and $K_0=L/W$. Powers of $P$ are suppressed.}\label{app:table-oblique-factors}
\begin{tabular}{lll}
\toprule
Factor & Grid chart & Weighted chart\\
\midrule
Critical and second-normal widths & $S^{-147}K_0^{-18}$ & $S^{-147}K_0^{-78}$\\
Purely longitudinal amplitudes & $L^6$ & $L^{12}$\\
Joint graph volume / determinant & $\nu L^8$ & $\nu$\\
Relative seed grid & $L^{11}K_0^6$ & $L^{23}K_0^{12}$\\
Point comparison & $L^2K_0$ & $L^2K_0$\\
\bottomrule
\end{tabular}
\end{table}

\begin{table}[H]
\centering\small
\caption{Factors in the proof of Theorem~\ref{low:6:wc36:thm-main}. Powers of $P$ are suppressed except in the last row.}\label{app:table-wide-factors}
\[
\begin{array}{lc}
\toprule
\text{Factor}&\text{Contribution}\\
\midrule
\text{Critical and second-normal widths}&w^{-123}H^{-165}\\
\text{Prepared graph window, including rate}&\nu w^{-6}H^{-12}\\
\text{Inverse grid determinant at the selected slope}&H^3\\
\text{Six external seed grids}&H^{11}\\
\text{Original-anchor point comparison}&H^2\\
\text{Third-normal range and finite offsets}&P^C\\
\bottomrule
\end{array}
\]
\end{table}
\begin{thebibliography}{99}
\bibitem{BGZZK} S. Basu, S. Guo, R. Zhang, and P. Zorin-Kranich, \emph{A stationary set method for estimating oscillatory integrals}, J. Eur. Math. Soc. \textbf{28} (2026), 2067--2100. \href{https://doi.org/10.4171/JEMS/1506}{doi:10.4171/JEMS/1506}.
\bibitem{BPR} S. Basu, R. Pollack, and M.-F. Roy, \emph{Algorithms in Real Algebraic Geometry}, second edition, Algorithms and Computation in Mathematics, vol. 10, Springer, Berlin, 2006. \href{https://doi.org/10.1007/3-540-33099-2}{doi:10.1007/3-540-33099-2}.
\bibitem{low:lojasiewicz} S. Basu and A. Mohammad-Nezhad, \emph{Improved effective {\L}ojasiewicz inequality and applications}, Forum Math. Sigma \textbf{12} (2024), e115. \href{https://doi.org/10.1017/fms.2024.66}{doi:10.1017/fms.2024.66}.
\bibitem{low:roots} M. Bafna, M. Sudan, S. Velusamy, and D. Xiang, \emph{Elementary analysis of isolated zeroes of a polynomial system}, arXiv:2102.00602, 2021. \url{https://arxiv.org/abs/2102.00602}.
\bibitem{low:area} R. L. Jerrard, \emph{Area and coarea}, MAT1501 lecture notes, University of Toronto, 7--17 October 2013. \url{https://www.math.toronto.edu/rjerrard/1501/notes/week5-6.pdf}.

\bibitem{Moc96} G. Mockenhaupt, \emph{Bounds in Lebesgue spaces of oscillatory integral operators}, Habilitationsschrift, Universit\"at Siegen, 1996.
\bibitem{LewisSendov} A. S. Lewis and H. S. Sendov, \emph{Twice differentiable spectral functions}, SIAM J. Matrix Anal. Appl. \textbf{23} (2001), 368--386. \href{https://doi.org/10.1137/S089547980036838X}{doi:10.1137/S089547980036838X}.
\bibitem{EtingofSL2} P. Etingof, \emph{Representations of Lie groups and Lie algebras}, Lecture 11, MIT 18.745, Lie Groups and Lie Algebras I, Fall 2020, Theorems 11.16 and 11.18. \href{https://ocw.mit.edu/courses/18-745-lie-groups-and-lie-algebras-i-fall-2020/mit18_745_f20_lec11.pdf}{MIT OpenCourseWare lecture notes}.
\bibitem{deBoor} C. de Boor, \emph{Divided differences}, Surv. Approx. Theory \textbf{1} (2005), 46--69. \url{https://surveys-in-approximation-theory.com/papers/2/2.pdf}.
\bibitem{FriedlandYomdin} O. Friedland and Y. Yomdin, \emph{$(s,p)$-Valent functions}, arXiv:1503.00325v1, 2015. \url{https://arxiv.org/abs/1503.00325v1}. Theorem 3.1 records the nonvanishing multivalent distortion theorem.
\bibitem{GuthArea} L. Guth, \emph{Unexpected applications of polynomials in combinatorics}, expository notes, Proposition 2.5. \url{https://math.mit.edu/~lguth/Exposition/erdossurvey.pdf}.
\bibitem{DLMF} F. W. J. Olver et al., eds., \emph{NIST Digital Library of Mathematical Functions}, equations (5.9.8)--(5.9.9). \url{https://dlmf.nist.gov/5.9}.
\bibitem{BedratyukCovariants} L. Bedratyuk, \emph{Bivariate Poincar\'e series for algebra of covariants of a binary form}, arXiv:1006.1974v1, 2010. \url{https://arxiv.org/abs/1006.1974v1}.
\bibitem{RoytwarfYomdin} N. Roytwarf and Y. Yomdin, \emph{Bernstein classes}, Ann. Inst. Fourier \textbf{47} (1997), no.~3, 825--858. \href{https://doi.org/10.5802/aif.1582}{doi:10.5802/aif.1582}.
\bibitem{PaperI} \emph{Complete affine jets and sharp integrability in three variables I. Degrees at least fourteen}, companion paper, 2026.
\bibitem{AKC80} G. I. Arkhipov, A. A. Karatsuba, and V. N. Chubarikov, \emph{Trigonometric integrals}, Math. USSR-Izv. \textbf{15} (1980), no.~2, 211--239. \href{https://doi.org/10.1070/IM1980v015n02ABEH001212}{doi:10.1070/IM1980v015n02ABEH001212}.
\end{thebibliography}

\end{document}
